% !TeX program = xelatex
% !TeX encoding = UTF-8
\documentclass[UTF8,10pt,a4paper,fontset=fandol]{ctexart}
\usepackage{amsmath,amssymb,graphicx,tabularx,array,multirow,enumitem,titlesec,xcolor,ragged2e,etoolbox}
\usepackage[a4paper,top=1.6cm,bottom=1.6cm,left=1.6cm,right=1.6cm]{geometry}
\setlength{\parindent}{2em}
\setlength{\parskip}{0.15em}
\setlength{\tabcolsep}{4pt}
\renewcommand{\arraystretch}{1.15}
\setlist[enumerate]{leftmargin=2.6em,itemsep=0.1em,topsep=0.2em}
\titleformat{\section}{\large\bfseries}{\thesection}{0.6em}{}
\titleformat{\subsection}{\normalsize\bfseries}{\thesubsection}{0.6em}{}
\titleformat{\subsubsection}{\normalsize\bfseries}{\thesubsubsection}{0.6em}{}
\titlespacing*{\section}{0pt}{0.7em}{0.25em}
\titlespacing*{\subsection}{0pt}{0.55em}{0.2em}
\titlespacing*{\subsubsection}{0pt}{0.4em}{0.15em}
\raggedbottom
\BeforeBeginEnvironment{tabularx}{\par\noindent}
\newcommand{\fieldvalue}[1]{#1}
\newcommand{\handwritten}[1]{#1}
\newcommand{\checkboxfield}[1]{\ifstrequal{#1}{checked}{☑}{\ifstrequal{#1}{unclear}{?}{☐}}}

\title{原生物检验单}
\author{}
\date{}

\begin{document}
\maketitle
\vspace{-2em}

\begin{center}
\small
维生素生物科技（北京）有限公司\\
\textit{PLATINUM LIFE EXCELLENCE BIOTECH}
\hfill 文件编号：REC-YZ-SMP-QC-01-001-d
\end{center}

\vspace{0.4em}

\begin{center}
\begin{tabular}{|p{2.0cm}|p{3.2cm}|p{2.0cm}|p{3.2cm}|}
\hline
\multicolumn{4}{c|}{\textbf{原生物检验单}}\\
\hline
样品名称 &
% #VALUE_ID: LEX-P0001-V0001
% #FIELD_VALUE: 样品名称
% TODO #HANDWRITTEN: 一只羊的肠道菌群; handwriting is partially illegible
\fieldvalue{\handwritten{一只羊的肠道菌群}} &
样品代码 &
% #VALUE_ID: LEX-P0001-V0002
% #FIELD_VALUE: 样品代码
% TODO #HANDWRITTEN: A1b40; handwriting is partially illegible
\fieldvalue{\handwritten{A1b40}}\\
\hline
样品批号 &
% #VALUE_ID: LEX-P0001-V0003
% #FIELD_VALUE: 样品批号
% TODO #HANDWRITTEN: A1b20240308; handwriting is partially illegible
\fieldvalue{\handwritten{A1b20240308}} &
样品规格/数量 &
% #VALUE_ID: LEX-P0001-V0004
% #FIELD_VALUE: 样品规格/数量
% TODO #HANDWRITTEN: 520 mL瓶$\times$6瓶; handwriting is partially illegible
\fieldvalue{\handwritten{520\,mL瓶$\times$6瓶}}\\
\hline
\multirow{2}{*}{样品类别} &
% #VALUE_ID: LEX-P0001-V0005
% #FIELD_VALUE: 样品类别--生产过程中检验样品
% #HANDWRITTEN: ✓
\fieldvalue{\handwritten{\checkboxfield{checked}}}
生产过程中检验样品
&
% #VALUE_ID: LEX-P0001-V0006
% #FIELD_VALUE: 样品类别--待包装产品
\fieldvalue{\checkboxfield{unchecked}}
待包装产品
&
\\[-0.2em]
&
% #VALUE_ID: LEX-P0001-V0007
% #FIELD_VALUE: 样品类别--成品
\fieldvalue{\checkboxfield{unchecked}}
成品
&
% #VALUE_ID: LEX-P0001-V0008
% #FIELD_VALUE: 样品类别--其他
\fieldvalue{\checkboxfield{unchecked}}
其他（\hspace{1.0cm}）
&
\\
\hline
检验类型 &
% #VALUE_ID: LEX-P0001-V0009
% #FIELD_VALUE: 检验类型--初验
% #HANDWRITTEN: ✓
\fieldvalue{\handwritten{\checkboxfield{checked}}}
初验
\quad
% #VALUE_ID: LEX-P0001-V0010
% #FIELD_VALUE: 检验类型--复验
\fieldvalue{\checkboxfield{unchecked}}
复验
&
% #VALUE_ID: LEX-P0001-V0011
% #FIELD_VALUE: 检验类型--退货
\fieldvalue{\checkboxfield{unchecked}}
退货
\quad
% #VALUE_ID: LEX-P0001-V0012
% #FIELD_VALUE: 检验类型--稳定性
\fieldvalue{\checkboxfield{unchecked}}
稳定性
\quad
% #VALUE_ID: LEX-P0001-V0013
% #FIELD_VALUE: 检验类型--留样
\fieldvalue{\checkboxfield{unchecked}}
留样
&
\\
\hline
请验目的 &
% #VALUE_ID: LEX-P0001-V0014
% #FIELD_VALUE: 请验目的--检验
% #HANDWRITTEN: ✓
\fieldvalue{\handwritten{\checkboxfield{checked}}}
检验
\quad
% #VALUE_ID: LEX-P0001-V0015
% #FIELD_VALUE: 请验目的--留样
\fieldvalue{\checkboxfield{unchecked}}
留样
&
% #VALUE_ID: LEX-P0001-V0016
% #FIELD_VALUE: 请验目的--其他
\fieldvalue{\checkboxfield{unchecked}}
其他（\hspace{1.0cm}）
&
\\
\hline
请验部门 &
% #VALUE_ID: LEX-P0001-V0017
% #FIELD_VALUE: 请验部门--生产部
% #HANDWRITTEN: ✓
\fieldvalue{\handwritten{\checkboxfield{checked}}}
生产部
\quad
% #VALUE_ID: LEX-P0001-V0018
% #FIELD_VALUE: 请验部门--物料管理部
\fieldvalue{\checkboxfield{unchecked}}
物料管理部
&
% #VALUE_ID: LEX-P0001-V0019
% #FIELD_VALUE: 请验部门--其他
\fieldvalue{\checkboxfield{unchecked}}
其他（\hspace{1.0cm}）
&
\\
\hline
样品情况 &
% #VALUE_ID: LEX-P0001-V0020
% #FIELD_VALUE: 样品情况--随请验单一起提供
% #HANDWRITTEN: ✓
\fieldvalue{\handwritten{\checkboxfield{checked}}}
随请验单一起提供
&
% #VALUE_ID: LEX-P0001-V0021
% #FIELD_VALUE: 样品情况--未随请验单一起提供
\fieldvalue{\checkboxfield{unchecked}}
未随请验单一起提供
&
\\
\hline
检验人 &
% #VALUE_ID: LEX-P0001-V0022
% #FIELD_VALUE: 检验人
% TODO #HANDWRITTEN: 沙斌; handwriting is difficult to read
\fieldvalue{\handwritten{沙斌}}
&
检验日期 &
% #VALUE_ID: LEX-P0001-V0023
% #FIELD_VALUE: 检验日期
% TODO #HANDWRITTEN: 2024.09.18; handwriting is partially illegible
\fieldvalue{\handwritten{2024.09.18}}
\\
\hline
收验人 &
% #VALUE_ID: LEX-P0001-V0024
% #FIELD_VALUE: 收验人
% TODO #HANDWRITTEN: 陈艳; signature is difficult to read
\fieldvalue{\handwritten{陈艳}}
&
收验日期 &
% #VALUE_ID: LEX-P0001-V0025
% #FIELD_VALUE: 收验日期
% TODO #HANDWRITTEN: 2024.09.15; handwriting is partially illegible
\fieldvalue{\handwritten{2024.09.15}}
\\
\hline
备注 &
% #VALUE_ID: LEX-P0001-V0026
% #FIELD_VALUE: 备注--第一项
% TODO #HANDWRITTEN: 20L; handwriting is partially illegible
\fieldvalue{\handwritten{20L}}
\hspace{2em}
% #VALUE_ID: LEX-P0001-V0027
% #FIELD_VALUE: 备注--第二项
% TODO #HANDWRITTEN: 24h; handwriting is partially illegible
\fieldvalue{\handwritten{24h}}
&
\multicolumn{2}{c|}{}
\\
\hline
\end{tabular}
\end{center}

% LEXOID_PAGE_COMPLETED: 1/70

\newpage

\begin{center}
\textbf{取样记录}
\end{center}

\noindent
文件编码：
% #VALUE_ID: LEX-P0002-V0001
% #FIELD_VALUE: 文件编码
\fieldvalue{REC-YZ-SMP-QC-01-004-b}
\hfill
版本号：
% #VALUE_ID: LEX-P0002-V0002
% #FIELD_VALUE: 版本号
\fieldvalue{1.4}

\vspace{0.5em}

\begin{tabular}{|p{0.15\linewidth}|p{0.35\linewidth}|p{0.15\linewidth}|p{0.25\linewidth}|}
\hline
名称 &
% #VALUE_ID: LEX-P0002-V0003
% #FIELD_VALUE: 名称
\fieldvalue{一级生物反应器细胞液} &
代码 &
% #VALUE_ID: LEX-P0002-V0004
% #FIELD_VALUE: 代码
\fieldvalue{QY0401}
\\ \hline
规格 &
% #VALUE_ID: LEX-P0002-V0005
% #FIELD_VALUE: 规格
\fieldvalue{20ml/袋} &
总数量 &
% #VALUE_ID: LEX-P0002-V0006
% #FIELD_VALUE: 总数量
\fieldvalue{1袋}
\\ \hline
批号 &
% #VALUE_ID: LEX-P0002-V0007
% #FIELD_VALUE: 批号
\fieldvalue{A3122202508006} &
进厂批号 &
% #VALUE_ID: LEX-P0002-V0008
% #FIELD_VALUE: 进厂批号
\fieldvalue{/}
\\ \hline
\multicolumn{1}{|p{0.15\linewidth}|}{供货单位/厂家} &
\multicolumn{3}{p{0.75\linewidth}|}{
% #VALUE_ID: LEX-P0002-V0009
% #FIELD_VALUE: 供货单位/厂家
\fieldvalue{/}}
\\ \hline
取样目的 &
\multicolumn{3}{p{0.75\linewidth}|}{
% #VALUE_ID: LEX-P0002-V0010
% #FIELD_VALUE: 正常检验
% #HANDWRITTEN: checked
\fieldvalue{\handwritten{\checkboxfield{checked}}}\ 正常检验
\quad
% #VALUE_ID: LEX-P0002-V0011
% #FIELD_VALUE: 复检
\fieldvalue{\checkboxfield{unchecked}}\ 复检
\quad
% #VALUE_ID: LEX-P0002-V0012
% #FIELD_VALUE: 稳定性
\fieldvalue{\checkboxfield{unchecked}}\ 稳定性

\par
% #VALUE_ID: LEX-P0002-V0013
% #FIELD_VALUE: 退货检验
\fieldvalue{\checkboxfield{unchecked}}\ 退货检验
\quad
% #VALUE_ID: LEX-P0002-V0014
% #FIELD_VALUE: 留样
\fieldvalue{\checkboxfield{unchecked}}\ 留样
\quad
% #VALUE_ID: LEX-P0002-V0015
% #FIELD_VALUE: 其他
\fieldvalue{\checkboxfield{unchecked}}\ 其他：\rule{2cm}{0.4pt}}
\\ \hline
取样地点 &
\multicolumn{3}{p{0.75\linewidth}|}{
% #VALUE_ID: LEX-P0002-V0016
% #FIELD_VALUE: 产品出口
% #HANDWRITTEN: checked
\fieldvalue{\handwritten{\checkboxfield{checked}}}\ 产品出口
\quad
% #VALUE_ID: LEX-P0002-V0017
% #FIELD_VALUE: 物料库
\fieldvalue{\checkboxfield{unchecked}}\ 物料库
\quad
% #VALUE_ID: LEX-P0002-V0018
% #FIELD_VALUE: 细胞库
\fieldvalue{\checkboxfield{unchecked}}\ 细胞库
\quad
% #VALUE_ID: LEX-P0002-V0019
% #FIELD_VALUE: 包装间
\fieldvalue{\checkboxfield{unchecked}}\ 包装间}
\\ \hline
贮存条件 &
\multicolumn{3}{p{0.75\linewidth}|}{
% #VALUE_ID: LEX-P0002-V0020
% #FIELD_VALUE: 2--8$^\circ$C
% #HANDWRITTEN: checked
\fieldvalue{\handwritten{\checkboxfield{checked}}}\ 2--8$^\circ$C
\quad
% #VALUE_ID: LEX-P0002-V0021
% #FIELD_VALUE: -20$^\circ$C
\fieldvalue{\checkboxfield{unchecked}}\ -20$^\circ$C
\quad
% #VALUE_ID: LEX-P0002-V0022
% #FIELD_VALUE: -80$^\circ$C
\fieldvalue{\checkboxfield{unchecked}}\ -80$^\circ$C
\quad
% #VALUE_ID: LEX-P0002-V0023
% #FIELD_VALUE: 常温
\fieldvalue{\checkboxfield{unchecked}}\ 常温
\quad
% #VALUE_ID: LEX-P0002-V0024
% #FIELD_VALUE: 其他
\fieldvalue{\checkboxfield{unchecked}}\ 其他：\rule{1.5cm}{0.4pt}}
\\ \hline
盛装容器 &
\multicolumn{3}{p{0.75\linewidth}|}{
% #VALUE_ID: LEX-P0002-V0025
% #FIELD_VALUE: 保温箱
% #HANDWRITTEN: checked
\fieldvalue{\handwritten{\checkboxfield{checked}}}\ 保温箱
\quad
% #VALUE_ID: LEX-P0002-V0026
% #FIELD_VALUE: 泡沫箱
\fieldvalue{\checkboxfield{unchecked}}\ 泡沫箱
\quad
% #VALUE_ID: LEX-P0002-V0027
% #FIELD_VALUE: 手提篮
\fieldvalue{\checkboxfield{unchecked}}\ 手提篮
\quad
% #VALUE_ID: LEX-P0002-V0028
% #FIELD_VALUE: 其他
\fieldvalue{\checkboxfield{unchecked}}\ 其他：\rule{1.5cm}{0.4pt}}
\\ \hline
外观复核 &
\multicolumn{3}{p{0.75\linewidth}|}{
% #VALUE_ID: LEX-P0002-V0029
% #FIELD_VALUE: 合格
% #HANDWRITTEN: checked
\fieldvalue{\handwritten{\checkboxfield{checked}}}\ 合格
\quad
% #VALUE_ID: LEX-P0002-V0030
% #FIELD_VALUE: 不合格
\fieldvalue{\checkboxfield{unchecked}}\ 不合格}
\\ \hline
取样量 &
% #VALUE_ID: LEX-P0002-V0031
% #FIELD_VALUE: 取样量
\fieldvalue{1袋} &
检验量 &
% #VALUE_ID: LEX-P0002-V0032
% #FIELD_VALUE: 检验量
\fieldvalue{1袋}
\\ \hline
\multicolumn{1}{|p{0.15\linewidth}|}{留样量} &
\multicolumn{1}{p{0.18\linewidth}|}{
% #VALUE_ID: LEX-P0002-V0033
% #FIELD_VALUE: 留样量
\fieldvalue{/}} &
\multicolumn{2}{p{0.42\linewidth}|}{} \\
\hline
取样人/日期 &
\multicolumn{1}{p{0.35\linewidth}|}{
% #VALUE_ID: LEX-P0002-V0034
% #FIELD_VALUE: 取样人
% #HANDWRITTEN: [签名]
\fieldvalue{\handwritten{[签名]}}\newline
% #VALUE_ID: LEX-P0002-V0035
% #FIELD_VALUE: 取样日期
% #HANDWRITTEN: 2025.8.15
\fieldvalue{\handwritten{2025.8.15}}} &
复核人/日期 &
\multicolumn{1}{p{0.25\linewidth}|}{
% #VALUE_ID: LEX-P0002-V0036
% #FIELD_VALUE: 复核人
% #HANDWRITTEN: [签名]
\fieldvalue{\handwritten{[签名]}}\newline
% #VALUE_ID: LEX-P0002-V0037
% #FIELD_VALUE: 复核日期
% #HANDWRITTEN: 2025.8.15
\fieldvalue{\handwritten{2025.8.15}}}
\\ \hline
备注 &
\multicolumn{3}{p{0.75\linewidth}|}{
% #VALUE_ID: LEX-P0002-V0038
% #FIELD_VALUE: 备注
% #TODO #HANDWRITTEN: T.2.6 ？；字迹不清
\fieldvalue{\handwritten{T.2.6 ？}}}
\\ \hline
\end{tabular}

\begin{center}
第1页共1页
\end{center}
% LEXOID_PAGE_COMPLETED: 2/70

\newpage

\begin{center}
\textbf{生物反应器细胞浓液操作记录}
\end{center}

文件编号：REC-YZ-SOP-QC-99-006-a

\section*{1\quad 仪器设备、试剂耗材：}

\subsection*{1.1\quad 仪器设备}

\begin{tabular}{|p{2.4cm}|p{4.0cm}|p{2.2cm}|p{3.0cm}|p{3.4cm}|}
\hline
名称 & 厂家 & 型号 & 设备编号 & 检查确认 \\
\hline
血糖仪 &
% #VALUE_ID: LEX-P0003-V0005
% #FIELD_VALUE: 厂家
\fieldvalue{三诺生物传感股份有限公司} &
% #VALUE_ID: LEX-P0003-V0006
% #FIELD_VALUE: 型号
\fieldvalue{安稳+} &
% #VALUE_ID: LEX-P0003-V0007
% #FIELD_VALUE: 设备编号
% #TODO #HANDWRITTEN: /; 仅见斜线状手写标记，具体内容不清
\fieldvalue{\handwritten{/}} &
正常\quad
% #VALUE_ID: LEX-P0003-V0008
% #FIELD_VALUE: 血糖仪检查确认—正常
\fieldvalue{\checkboxfield{checked}}
\quad 异常\quad
% #VALUE_ID: LEX-P0003-V0009
% #FIELD_VALUE: 血糖仪检查确认—异常
\fieldvalue{\checkboxfield{unchecked}} \\
\hline
血糖仪 &
% #VALUE_ID: LEX-P0003-V0010
% #FIELD_VALUE: 厂家
\fieldvalue{Eppendorf Corporate} &
\fieldvalue{0.1--2.5\,\textmu l} &
% #VALUE_ID: LEX-P0003-V0011
% #FIELD_VALUE: 设备编号
% #TODO #HANDWRITTEN: 131?3?81; 手写数字辨识不清
\fieldvalue{\handwritten{131?3?81}} &
正常\quad
% #VALUE_ID: LEX-P0003-V0012
% #FIELD_VALUE: 血糖仪检查确认—正常
\fieldvalue{\checkboxfield{checked}}
\quad 异常\quad
% #VALUE_ID: LEX-P0003-V0013
% #FIELD_VALUE: 血糖仪检查确认—异常
\fieldvalue{\checkboxfield{unchecked}} \\
\hline
\end{tabular}

\subsection*{1.2\quad 试剂耗材}

\begin{tabular}{|p{2.8cm}|p{6.2cm}|p{3.0cm}|p{3.2cm}|}
\hline
名称 & 厂家 & 批号 & 有效期至 \\
\hline
血糖测试条 &
% #VALUE_ID: LEX-P0003-V0014
% #FIELD_VALUE: 厂家
\fieldvalue{三诺生物传感股份有限公司} &
% #VALUE_ID: LEX-P0003-V0015
% #FIELD_VALUE: 批号
% #HANDWRITTEN: 444444-06
\fieldvalue{\handwritten{444444-06}} &
% #VALUE_ID: LEX-P0003-V0016
% #FIELD_VALUE: 有效期至
% #HANDWRITTEN: 2027.1.31
\fieldvalue{\handwritten{2027.1.31}} \\
\hline
\end{tabular}

\section*{2\quad 操作步骤：}

\begin{tabular}{|p{8.8cm}|p{5.8cm}|}
\hline
操作过程 & \begin{center}取上清液\end{center} \\
\hline
取出测试条，取样品自然沉降后的上清液0.6ul加入测试条电极反应室，血糖仪自动发出“嘀”提示音，吸入样品完成，血糖仪自动进入倒计时，开始测试； &
取上清液：
% #VALUE_ID: LEX-P0003-V0017
% #FIELD_VALUE: 取上清液体积
% #HANDWRITTEN: 0.4
\fieldvalue{\handwritten{0.4}}\,\textmu l
\\
\hline
重复以上步骤2次，共计数3次。 &
每个样品分别共计数
% #VALUE_ID: LEX-P0003-V0018
% #FIELD_VALUE: 每个样品计数次数
% #HANDWRITTEN: 3
\fieldvalue{\handwritten{3}} 次。 \\
\hline
\multicolumn{2}{|c|}{检测结果} \\
\hline
\multicolumn{2}{|c|}{葡萄糖浓度（mmol/L）} \\
\hline
\multicolumn{2}{|c|}{\begin{tabular}{cccc}
第一次 & 第二次 & 第三次 & 平均值（保留一位小数） \\
\hline
% #VALUE_ID: LEX-P0003-V0019
% #FIELD_VALUE: 第一次检测结果
% #HANDWRITTEN: 18.4
\fieldvalue{\handwritten{18.4}} &
% #VALUE_ID: LEX-P0003-V0020
% #FIELD_VALUE: 第二次检测结果
% #HANDWRITTEN: 19.4
\fieldvalue{\handwritten{19.4}} &
% #VALUE_ID: LEX-P0003-V0021
% #FIELD_VALUE: 第三次检测结果
% #HANDWRITTEN: 20.2
\fieldvalue{\handwritten{20.2}} &
% #VALUE_ID: LEX-P0003-V0022
% #FIELD_VALUE: 平均值
% #HANDWRITTEN: 19.3
\fieldvalue{\handwritten{19.3}}
\end{tabular}} \\
\hline
\end{tabular}

\begin{tabular}{|p{2.5cm}|p{12.1cm}|}
\hline
备注 &
% #VALUE_ID: LEX-P0003-V0023
% #FIELD_VALUE: 备注
% #HANDWRITTEN: ✓
\fieldvalue{\handwritten{\checkmark}} \\
\hline
实验人/日期 &
% #VALUE_ID: LEX-P0003-V0024
% #FIELD_VALUE: 实验人
% #TODO #HANDWRITTEN: 姓名难辨; 手写签名无法准确识别
\fieldvalue{\handwritten{[姓名难辨]}}\quad
% #VALUE_ID: LEX-P0003-V0025
% #FIELD_VALUE: 实验日期
% #HANDWRITTEN: 2024.8.17
\fieldvalue{\handwritten{2024.8.17}} \\
\hline
复核人/日期 &
% #VALUE_ID: LEX-P0003-V0026
% #FIELD_VALUE: 复核人
% #TODO #HANDWRITTEN: 姓名难辨; 手写签名无法准确识别
\fieldvalue{\handwritten{[姓名难辨]}}\quad
% #VALUE_ID: LEX-P0003-V0027
% #FIELD_VALUE: 复核日期
% #HANDWRITTEN: 2025.8.16
\fieldvalue{\handwritten{2025.8.16}} \\
\hline
\end{tabular}
% LEXOID_PAGE_COMPLETED: 3/70

\newpage

\section*{生物反应器细胞液检验操作记录}

\noindent
检验编号：REC-YZ-SOP-QC-99-006-a

\begin{center}
\begin{tabular}{|p{0.16\linewidth}|p{0.30\linewidth}|p{0.16\linewidth}|p{0.25\linewidth}|}
\hline
\multicolumn{4}{c}{密度与活率检测} \\
\hline
样品名称 &
% #VALUE_ID: LEX-P0004-V0001
% #FIELD_VALUE: 样品名称
% #TODO #HANDWRITTEN: 一级生物反应器细胞液; handwriting partly unclear
\fieldvalue{\handwritten{一级生物反应器细胞液}} &
样品规格 &
% #VALUE_ID: LEX-P0004-V0002
% #FIELD_VALUE: 样品规格
% #HANDWRITTEN: 20
\fieldvalue{\handwritten{20}} ml/袋 \\
\hline
样品批号 &
% #VALUE_ID: LEX-P0004-V0003
% #FIELD_VALUE: 样品批号
% #TODO #HANDWRITTEN: A?J?7?7?0?; handwriting unclear
\fieldvalue{\handwritten{A?J?7?7?0?}} &
取样日期 &
% #VALUE_ID: LEX-P0004-V0004
% #FIELD_VALUE: 取样日期
% #HANDWRITTEN: 2024年08月17日
\fieldvalue{\handwritten{2024年08月17日}} \\
\hline
\end{tabular}
\end{center}

\section{仪器设备、试剂耗材}

\subsection{仪器设备}

\begin{center}
\small
\begin{tabular}{|p{0.18\linewidth}|p{0.16\linewidth}|p{0.16\linewidth}|p{0.18\linewidth}|p{0.18\linewidth}|p{0.16\linewidth}|}
\hline
名称 & 厂家 & 型号 & 设备编号 & 校准有效期至 & 检查确认 \\
\hline
% #VALUE_ID: LEX-P0004-V0005
% #FIELD_VALUE: 仪器名称
\fieldvalue{细胞计数仪} &
% #VALUE_ID: LEX-P0004-V0006
% #FIELD_VALUE: 厂家
\fieldvalue{Nexcelom} &
% #VALUE_ID: LEX-P0004-V0007
% #FIELD_VALUE: 型号
\fieldvalue{Cellometer K2} &
% #VALUE_ID: LEX-P0004-V0008
% #FIELD_VALUE: 设备编号
\fieldvalue{1410905} &
% #VALUE_ID: LEX-P0004-V0009
% #FIELD_VALUE: 校准有效期至
% #HANDWRITTEN: 2025.01.10
\fieldvalue{\handwritten{2025.01.10}} &
正常
% #VALUE_ID: LEX-P0004-V0010
% #FIELD_VALUE: 正常
\fieldvalue{\checkboxfield{checked}}
\quad 异常
% #VALUE_ID: LEX-P0004-V0011
% #FIELD_VALUE: 异常
\fieldvalue{\checkboxfield{unchecked}} \\
\hline
% #VALUE_ID: LEX-P0004-V0012
% #FIELD_VALUE: 仪器名称
\fieldvalue{高速冷冻离心机} &
% #VALUE_ID: LEX-P0004-V0013
% #FIELD_VALUE: 厂家
\fieldvalue{Thermo} &
% #VALUE_ID: LEX-P0004-V0014
% #FIELD_VALUE: 型号
\fieldvalue{ST16R} &
% #VALUE_ID: LEX-P0004-V0015
% #FIELD_VALUE: 设备编号
\fieldvalue{1410810} &
% #VALUE_ID: LEX-P0004-V0016
% #FIELD_VALUE: 校准有效期至
% #TODO #HANDWRITTEN: 2025.01.??; final digits unclear
\fieldvalue{\handwritten{2025.01.??}} &
正常
% #VALUE_ID: LEX-P0004-V0017
% #FIELD_VALUE: 正常
\fieldvalue{\checkboxfield{checked}}
\quad 异常
% #VALUE_ID: LEX-P0004-V0018
% #FIELD_VALUE: 异常
\fieldvalue{\checkboxfield{unchecked}} \\
\hline
% #VALUE_ID: LEX-P0004-V0019
% #FIELD_VALUE: 仪器名称
\fieldvalue{电热恒温水槽} &
% #VALUE_ID: LEX-P0004-V0020
% #FIELD_VALUE: 厂家
\fieldvalue{上海一恒} &
% #VALUE_ID: LEX-P0004-V0021
% #FIELD_VALUE: 型号
\fieldvalue{DK-8AXX} &
% #VALUE_ID: LEX-P0004-V0022
% #FIELD_VALUE: 设备编号
\fieldvalue{1411001} &
% #VALUE_ID: LEX-P0004-V0023
% #FIELD_VALUE: 校准有效期至
% #TODO #HANDWRITTEN: 2025.01.??; final digits unclear
\fieldvalue{\handwritten{2025.01.??}} &
正常
% #VALUE_ID: LEX-P0004-V0024
% #FIELD_VALUE: 正常
\fieldvalue{\checkboxfield{checked}}
\quad 异常
% #VALUE_ID: LEX-P0004-V0025
% #FIELD_VALUE: 异常
\fieldvalue{\checkboxfield{unchecked}} \\
\hline
\end{tabular}
\end{center}

\subsection{试剂耗材}

\begin{center}
\small
\begin{tabular}{|p{0.25\linewidth}|p{0.25\linewidth}|p{0.28\linewidth}|p{0.18\linewidth}|}
\hline
名称 & 厂家 & 批号 & 有效期至 \\
\hline
% #VALUE_ID: LEX-P0004-V0026
% #FIELD_VALUE: 名称
\fieldvalue{AO/PI 染液} &
% #VALUE_ID: LEX-P0004-V0027
% #FIELD_VALUE: 厂家
\fieldvalue{Nexcelom} &
% #VALUE_ID: LEX-P0004-V0028
% #FIELD_VALUE: 批号
% #TODO #HANDWRITTEN: 774?49?; handwriting unclear
\fieldvalue{\handwritten{774?49?}} &
% #VALUE_ID: LEX-P0004-V0029
% #FIELD_VALUE: 有效期至
% #HANDWRITTEN: 2025.01.12
\fieldvalue{\handwritten{2025.01.12}} \\
\hline
% #VALUE_ID: LEX-P0004-V0030
% #FIELD_VALUE: 名称
\fieldvalue{0.9\%氯化钠注射液} &
% #VALUE_ID: LEX-P0004-V0031
% #FIELD_VALUE: 厂家
\fieldvalue{石家庄四药} &
% #VALUE_ID: LEX-P0004-V0032
% #FIELD_VALUE: 批号
% #TODO #HANDWRITTEN: 244?3720; handwriting unclear
\fieldvalue{\handwritten{244?3720}} &
% #VALUE_ID: LEX-P0004-V0033
% #FIELD_VALUE: 有效期至
% #HANDWRITTEN: 2025.10.06
\fieldvalue{\handwritten{2025.10.06}} \\
\hline
% #VALUE_ID: LEX-P0004-V0034
% #FIELD_VALUE: 名称
\fieldvalue{20\%人血白蛋白} &
% #VALUE_ID: LEX-P0004-V0035
% #FIELD_VALUE: 厂家
\fieldvalue{Baxter AG} &
% #VALUE_ID: LEX-P0004-V0036
% #FIELD_VALUE: 批号
% #TODO #HANDWRITTEN: A?N?3?8; handwriting unclear
\fieldvalue{\handwritten{A?N?3?8}} &
% #VALUE_ID: LEX-P0004-V0037
% #FIELD_VALUE: 有效期至
% #HANDWRITTEN: 2025.05.24
\fieldvalue{\handwritten{2025.05.24}} \\
\hline
\end{tabular}
\end{center}

\section{操作步骤}

\subsection*{溶液准备}

\noindent
1\%人白蛋白溶液配制：取5 ml 20\%人血白蛋白加入到95 ml的0.9\%氯化钠注射液中，混匀即可。

\noindent
20\%人血白蛋白：\underline{\hspace{1.5cm}} ml

\noindent
0.9\%氯化钠注射液：\underline{\hspace{1.5cm}} ml

\begin{center}
\small
\begin{tabular}{|p{0.25\linewidth}|p{0.35\linewidth}|p{0.28\linewidth}|}
\hline
名称 & 配制批号 & 有效期至 \\
\hline
5$\times$裂解液 &
% #VALUE_ID: LEX-P0004-V0038
% #FIELD_VALUE: 配制批号
% #TODO #HANDWRITTEN: 5?0?0?7?0; handwriting unclear
\fieldvalue{\handwritten{5?0?0?7?0}} &
% #VALUE_ID: LEX-P0004-V0039
% #FIELD_VALUE: 有效期至
% #HANDWRITTEN: 2025.01.07
\fieldvalue{\handwritten{2025.01.07}} \\
\hline
\end{tabular}
\end{center}

\subsection*{操作过程}

\begin{tabular}{|p{0.53\linewidth}|p{0.39\linewidth}|}
\hline
按照50 ml离心管的刻度读取样品体积，往样品管中加入对应体积的5$\times$裂解液（每4 ml样品加入1 ml 5$\times$裂解液，按照“只进不舍”原则，如26 ml样品，加入7 ml 5$\times$裂解液），充分混匀后，分成两份，每份约3 ml，其中一份置于37℃电热恒温水槽中，置于37℃电热恒温水槽中静置10 min得到细胞悬液。取出离心管摇晃混匀，使样品充分裂解；再置于37℃电热恒温水槽中裂解10 min得到细胞悬液。 &
样品体积：
% #VALUE_ID: LEX-P0004-V0040
% #FIELD_VALUE: 样品体积
% #HANDWRITTEN: 12
\fieldvalue{\handwritten{12}} ml

裂解液：
% #VALUE_ID: LEX-P0004-V0041
% #FIELD_VALUE: 裂解液
% #HANDWRITTEN: 2
\fieldvalue{\handwritten{2}} ml

电热恒温水槽：
% #VALUE_ID: LEX-P0004-V0042
% #FIELD_VALUE: 电热恒温水槽温度
% #HANDWRITTEN: 37
\fieldvalue{\handwritten{37}}℃

裂解：
% #VALUE_ID: LEX-P0004-V0043
% #FIELD_VALUE: 裂解时间
% #TODO #HANDWRITTEN: 9? min; handwriting unclear
\fieldvalue{\handwritten{9?}} min --- 
% #VALUE_ID: LEX-P0004-V0044
% #FIELD_VALUE: 裂解结束时间
% #TODO #HANDWRITTEN: 17:??; handwriting unclear
\fieldvalue{\handwritten{17:??}} \\
\hline
取出裂解好的细胞悬液300$\times$g离心5分钟，按照读取样品体积加入对应体积的1\%人白蛋白溶液重悬（每1.2 ml样品体积加入0.1 ml 1\%人白蛋白溶液，按照“只进不舍”原则，如26 ml样品，需加入3 ml 1\%人白蛋白溶液）。 &
% #VALUE_ID: LEX-P0004-V0045
% #FIELD_VALUE: 离心力
% #TODO #HANDWRITTEN: 300; handwriting unclear
\fieldvalue{\handwritten{300}}$\times g$，离心
% #VALUE_ID: LEX-P0004-V0046
% #FIELD_VALUE: 离心时间
% #HANDWRITTEN: 2
\fieldvalue{\handwritten{2}}分钟

加入1 ml人白蛋白溶液总体积：
% #VALUE_ID: LEX-P0004-V0047
% #FIELD_VALUE: 总体积
% #HANDWRITTEN: 1
\fieldvalue{\handwritten{1}} ml

吸取
% #VALUE_ID: LEX-P0004-V0048
% #FIELD_VALUE: 吸取体积
% #HANDWRITTEN: 2
\fieldvalue{\handwritten{2}} $\mu$l \\
\hline
\end{tabular}

\begin{center}
第 2 页 共 5 页
\end{center}
% LEXOID_PAGE_COMPLETED: 4/70

\newpage

\noindent
加人 2.2ml 1\%人台红细胞浓液，反复吹打至单细胞悬液。\\
撕去计数板上下两层保护膜，将细胞悬液混匀，从中吸取 20$\mu$l\\
与 20$\mu$l AO/PI 染料混合均匀。\\
吸取
% #TODO #HANDWRITTEN: 20；蓝色手写修改
\handwritten{20}
$\mu$l 细胞悬液与 AO/PI 染料的混合液，注入到计数板的一\\
个检测室中。\\
将加入混合液的一端插入进样口。\\
选择 Assay 类型，使用 AO/PI 染料进行双荧光活率检测；\\
在主界面左侧 ID 栏中输入样本批号。\\
点击主界面左下角“Preview”预览细胞图片。\\
调整仪器焦距直至观察到荧光通道下的细胞中心透亮，轮廓清\\
晰。\\
点击“count”按钮进行计数并自动拍摄图片。\\
重复以上步骤 2 次，共计数 3 次。打印数据结果。

\medskip
\noindent
文件编码：REC-YZ-SOP-CQ-99-006-a

\medskip
\noindent
与
% #TODO #HANDWRITTEN: 20；蓝色手写修改
\handwritten{20}
$\mu$l AO/PI 染料混合均匀。\\
吸取
% #TODO #HANDWRITTEN: 20；蓝色手写修改
\handwritten{20}
$\mu$l 计数。\\
每个样品共计数
% #HANDWRITTEN: 3
\handwritten{3}
次。

\medskip
\noindent
数据保存路径：
% #VALUE_ID: LEX-P0005-V0001
% #FIELD_VALUE: 数据保存路径
% #TODO #HANDWRITTEN: /1步的图像--第一步中14--牙齿数据--4结存数据--2次；字迹较潦草
\fieldvalue{\handwritten{/1步的图像--第一步中14--牙齿数据--4结存数据--2次}}

\begin{center}
\begin{tabular}{|p{0.17\linewidth}|p{0.20\linewidth}|p{0.16\linewidth}|p{0.13\linewidth}|p{0.12\linewidth}|p{0.12\linewidth}|}
\hline
\multicolumn{6}{|c|}{检测结果}\\ \hline
\multicolumn{6}{|l|}{检测结果：
% #VALUE_ID: LEX-P0005-V0002
% #FIELD_VALUE: 一级反应
\fieldvalue{\checkboxfield{checked}} 一级反应、
% #VALUE_ID: LEX-P0005-V0003
% #FIELD_VALUE: 二级反应
\fieldvalue{\checkboxfield{unchecked}} 二级反应、
% #VALUE_ID: LEX-P0005-V0004
% #FIELD_VALUE: 三级反应
\fieldvalue{\checkboxfield{unchecked}} 三级反应}\\ \hline
检测次数 &
活细胞密度(个细胞/ml) &
样品密度(个细胞/ml) &
活细胞密度 CV 值（\%） &
细胞活率（\%） &
样品活率（\%）\\ \hline
第一次 &
% #VALUE_ID: LEX-P0005-V0005
% #FIELD_VALUE: 第一次活细胞密度
% #TODO #HANDWRITTEN: 47.7×10^?；指数难以辨认
\fieldvalue{\handwritten{47.7$\times 10^{?}$}} &
&
&
% #VALUE_ID: LEX-P0005-V0006
% #FIELD_VALUE: 第一次细胞活率
% #TODO #HANDWRITTEN: 70.0；末位字迹略有不清
\fieldvalue{\handwritten{70.0}} &
\\ \hline
第二次 &
% #VALUE_ID: LEX-P0005-V0007
% #FIELD_VALUE: 第二次活细胞密度
% #TODO #HANDWRITTEN: 47.7×10^?；指数难以辨认
\fieldvalue{\handwritten{47.7$\times 10^{?}$}} &
% #VALUE_ID: LEX-P0008-V0008
% #FIELD_VALUE: 样品密度
% #TODO #HANDWRITTEN: 38.3；数字字迹略有不清
\fieldvalue{\handwritten{38.3}} &
% #VALUE_ID: LEX-P0005-V0009
% #FIELD_VALUE: 活细胞密度 CV 值
% #HANDWRITTEN: 19
\fieldvalue{\handwritten{19}} &
% #VALUE_ID: LEX-P0005-V0010
% #FIELD_VALUE: 第二次细胞活率
% #HANDWRITTEN: 84.4
\fieldvalue{\handwritten{84.4}} &
% #VALUE_ID: LEX-P0005-V0011
% #FIELD_VALUE: 样品活率
% #HANDWRITTEN: 92
\fieldvalue{\handwritten{92}} \\ \hline
第三次 &
% #VALUE_ID: LEX-P0005-V0012
% #FIELD_VALUE: 第三次活细胞密度
% #TODO #HANDWRITTEN: 3.67×10^?；指数难以辨认
\fieldvalue{\handwritten{3.67$\times 10^{?}$}} &
&
&
% #VALUE_ID: LEX-P0005-V0013
% #FIELD_VALUE: 第三次细胞活率
% #TODO #HANDWRITTEN: 90.7；字迹略有不清
\fieldvalue{\handwritten{90.7}} &
\\ \hline
\multicolumn{2}{|l|}{发酵罐细胞体积（ml）} &
\multicolumn{2}{l|}{细胞数量（个细胞/罐）} &
\multicolumn{2}{l|}{}\\[-1.7ex]
\multicolumn{2}{|l|}{} &
\multicolumn{2}{|l|}{} &
\multicolumn{2}{l|}{
% #VALUE_ID: LEX-P0005-V0014
% #FIELD_VALUE: 细胞数量
% #TODO #HANDWRITTEN: 1.78×10^?；指数难以辨认
\fieldvalue{\handwritten{1.78$\times 10^{?}$}}}\\ \hline
\end{tabular}
\end{center}

\noindent
计算要求：样品密度为 3 次活细胞密度结果平均值乘以 1\%人的氯化钠溶液最终体积，再除以生物反应器\\
细胞液体积；细胞数量为样品密度乘以发酵罐体积。

\medskip
\noindent
备注：
% #VALUE_ID: LEX-P0005-V0015
% #FIELD_VALUE: 备注
% #HANDWRITTEN: /
\fieldvalue{\handwritten{/}}

\medskip
\noindent
实验人/日期：
% #VALUE_ID: LEX-P0005-V0016
% #FIELD_VALUE: 实验人
% #TODO #HANDWRITTEN: 娟；签名潦草
\fieldvalue{\handwritten{娟}}
\quad
% #VALUE_ID: LEX-P0005-V0017
% #FIELD_VALUE: 实验日期
% #HANDWRITTEN: 2017.08.15
\fieldvalue{\handwritten{2017.08.15}}

\medskip
\noindent
复核人/日期：
% #VALUE_ID: LEX-P0005-V0018
% #FIELD_VALUE: 复核人
% #TODO #HANDWRITTEN: 姚；签名潦草
\fieldvalue{\handwritten{姚}}
\quad
% #VALUE_ID: LEX-P0005-V0019
% #FIELD_VALUE: 复核日期
% #HANDWRITTEN: 2018.08.15
\fieldvalue{\handwritten{2018.08.15}}

\begin{flushright}
第 3 页 共 5 页
\end{flushright}
% LEXOID_PAGE_COMPLETED: 5/70

\newpage

\begin{center}
华生生物反应器细胞液检测记录
\end{center}

\noindent 文件编码：REC-YZ-SOP-QC-99-006-a

\begin{center}
\begin{tabular}{|p{0.16\linewidth}|p{0.34\linewidth}|p{0.16\linewidth}|p{0.24\linewidth}|}
\hline
\multicolumn{4}{c|}{细胞生长状态检验}\\
\hline
样品名称 &
% #VALUE_ID: LEX-P0006-V0001
% #FIELD_VALUE: 样品名称
\fieldvalue{生物反应器细胞液} &
样品规格 &
% #VALUE_ID: LEX-P0006-V0002
% #FIELD_VALUE: 样品规格
% #HANDWRITTEN: 20
\fieldvalue{\handwritten{20}}\,mL/袋\\
\hline
样品批号 &
% #VALUE_ID: LEX-P0006-V0003
% #FIELD_VALUE: 样品批号
% #TODO #HANDWRITTEN: 难辨；蓝色手写批号无法可靠识别
\fieldvalue{\handwritten{难辨}} &
取样日期 &
% #VALUE_ID: LEX-P0006-V0004
% #FIELD_VALUE: 取样日期
% #HANDWRITTEN: 2024年3月7日
\fieldvalue{\handwritten{2024年3月7日}}\\
\hline
\end{tabular}
\end{center}

\section{仪器设备及试剂}

\subsection{仪器设备}

\begin{center}
\begin{tabular}{|p{0.20\linewidth}|p{0.18\linewidth}|p{0.18\linewidth}|p{0.20\linewidth}|p{0.18\linewidth}|}
\hline
仪器名称 & 厂家 & 型号 & 设备编号 & 检查确认\\
\hline
% #VALUE_ID: LEX-P0006-V0005
% #FIELD_VALUE: 仪器名称
\fieldvalue{荧光显微镜} &
% #VALUE_ID: LEX-P0006-V0006
% #FIELD_VALUE: 厂家
\fieldvalue{徕卡} &
% #VALUE_ID: LEX-P0006-V0007
% #FIELD_VALUE: 型号
\fieldvalue{DmiL LED} &
% #VALUE_ID: LEX-P0006-V0008
% #FIELD_VALUE: 设备编号
\fieldvalue{1431404} &
正常\,
% #VALUE_ID: LEX-P0006-V0009
% #FIELD_VALUE: 正常
\fieldvalue{\checkboxfield{checked}}\quad
另异常\,
% #VALUE_ID: LEX-P0006-V0010
% #FIELD_VALUE: 另异常
\fieldvalue{\checkboxfield{unchecked}}\\
\hline
\end{tabular}
\end{center}

\subsection{试剂耗材}

\begin{center}
\begin{tabular}{|p{0.27\linewidth}|p{0.27\linewidth}|p{0.20\linewidth}|p{0.13\linewidth}|p{0.18\linewidth}|}
\hline
名称 & 厂家 & 货号 & 批号 & 有效期至\\
\hline
% #VALUE_ID: LEX-P0006-V0011
% #FIELD_VALUE: 名称
\fieldvalue{细胞活力检测试剂盒} &
% #VALUE_ID: LEX-P0006-V0012
% #FIELD_VALUE: 厂家
\fieldvalue{江苏凯基生物技术股份有限公司} &
% #VALUE_ID: LEX-P0006-V0013
% #FIELD_VALUE: 货号
\fieldvalue{KGA9501-1000} &
% #VALUE_ID: LEX-P0006-V0014
% #FIELD_VALUE: 批号
% #TODO #HANDWRITTEN: 难辨；蓝色手写批号无法可靠识别
\fieldvalue{\handwritten{难辨}} &
% #VALUE_ID: LEX-P0006-V0015
% #FIELD_VALUE: 有效期至
% #TODO #HANDWRITTEN: 2024.11.07?；末位数字辨识不完全确定
\fieldvalue{\handwritten{2024.11.07?}}\\
\hline
\end{tabular}
\end{center}

\section{操作步骤}

\begin{center}
\begin{tabular}{|p{0.50\linewidth}|p{0.43\linewidth}|}
\hline
\multicolumn{2}{c|}{溶液配制}\\
\hline
取
% #VALUE_ID: LEX-P0006-V0016
% #FIELD_VALUE: 组分A体积
% #HANDWRITTEN: 0.2
\fieldvalue{\handwritten{0.2}}\,\textmu L
组分A 和
% #VALUE_ID: LEX-P0006-V0017
% #FIELD_VALUE: 组分B体积
% #HANDWRITTEN: 0.2
\fieldvalue{\handwritten{0.2}}\,\textmu L
组分B 与
% #VALUE_ID: LEX-P0006-V0018
% #FIELD_VALUE: PBS体积
% #TODO #HANDWRITTEN: 798?；蓝色手写数字辨识不完全确定
\fieldvalue{\handwritten{798?}}\,\textmu L PBS 混合均匀配成20$\times$的荧光染液。 &
\\
\hline
\multicolumn{2}{c|}{操作过程}\\
\hline
取混匀后样品200\textmu L至96孔板中，取2孔； &
样品
% #VALUE_ID: LEX-P0006-V0019
% #FIELD_VALUE: 样品加入量
% #HANDWRITTEN: 200
\fieldvalue{\handwritten{200}}\,\textmu L/孔；\\
&
每个样品取
% #VALUE_ID: LEX-P0006-V0020
% #FIELD_VALUE: 每个样品孔数
% #HANDWRITTEN: 2
\fieldvalue{\handwritten{2}}\,孔\\
\hline
自然沉降后，吸弃上清，注意不要吸到沉淀样品； &
% #VALUE_ID: LEX-P0006-V0021
% #FIELD_VALUE: 吸弃上清
\fieldvalue{\checkboxfield{checked}} 吸弃上清\\
\hline
取20$\times$的荧光染液加入到96孔板中，200\textmu L/孔，避光孵育30--60\,min； &
20$\times$的荧光染液
% #VALUE_ID: LEX-P0006-V0022
% #FIELD_VALUE: 荧光染液加入量
% #HANDWRITTEN: 50
\fieldvalue{\handwritten{50}}\,\textmu L/孔\\
&
避光孵育
% #VALUE_ID: LEX-P0006-V0023
% #FIELD_VALUE: 避光孵育起始时间
% #TODO #HANDWRITTEN: 8?；蓝色手写数字辨识不完全确定
\fieldvalue{\handwritten{8?}}\,min--%
% #VALUE_ID: LEX-P0006-V0024
% #FIELD_VALUE: 避光孵育结束时间
% #TODO #HANDWRITTEN: 12?；蓝色手写数字辨识不完全确定
\fieldvalue{\handwritten{12?}}\,min\\
\hline
染色结束后，吸弃上清，取PBS加入到以96孔板中，200\textmu L/孔，吸弃上清，注意不要吸到沉淀样品，清洗3次，之后每孔加入100\textmu L PBS进行拍照； &
加入PBS
% #VALUE_ID: LEX-P0006-V0025
% #FIELD_VALUE: PBS加入量
% #HANDWRITTEN: 200
\fieldvalue{\handwritten{200}}\,\textmu L/孔\\
&
清洗
% #VALUE_ID: LEX-P0006-V0026
% #FIELD_VALUE: 清洗次数
% #HANDWRITTEN: 3
\fieldvalue{\handwritten{3}}\,次\\
\hline
将荧光显微镜调至4$\times$镜，软件中调节荧光参数，调节到适合的光源和像距，拍摄4张清晰图像。进一步编辑图像，打印并粘贴在实验记录本上。 &
\\
\hline
数据保存路径 &
% #VALUE_ID: LEX-P0006-V0027
% #FIELD_VALUE: 数据保存路径
% #TODO #HANDWRITTEN: 难辨；蓝色手写路径无法可靠识别
\fieldvalue{\handwritten{难辨}}\\
\hline
\multicolumn{2}{c|}{检测结果}\\
\hline
\end{tabular}
\end{center}

% TODO：本页底部“检测结果”表格仅显示标题，具体内容应在后续页面或下方不可见区域。
% LEXOID_PAGE_COMPLETED: 6/70

\newpage

\noindent
\begin{minipage}{\linewidth}
\small
\begin{tabular}{@{}p{0.18\linewidth}p{0.42\linewidth}p{0.34\linewidth}@{}}
\textit{铂生卓越生物科技（北京）有限公司} &
\centering
\begin{tabular}{|c|}
\hline
铂生生物\\
原件\\
\hline
\end{tabular}
&
\raggedleft
文件编码：REC-YZ-SOP-QC-99-006-\hspace{0.5em}
\begin{tabular}{|c|}
\hline
铂生生物\\
版本号：%
% #VALUE_ID: LEX-P0007-V0001
% #FIELD_VALUE: 版本号
\fieldvalue{1.4}\\
\hline
\end{tabular}
\end{tabular}
\end{minipage}

\vspace{0.3em}
\noindent\rule{\linewidth}{0.6pt}

\noindent
\begin{minipage}[t][0.80\textheight][c]{\linewidth}
\centering
% TODO: Embedded microscopic image visible on the page; no source filename is available.
\fbox{\parbox[c][0.18\textheight][c]{0.34\linewidth}{\centering
\textit{显微图像}
}}

\vspace{0.3em}
% #VALUE_ID: LEX-P0007-V0002
% #HANDWRITTEN: （难以辨认）25/08/10；图像下方手写标注
\handwritten{（难以辨认）25/08/10}
\end{minipage}

\noindent
\begin{tabular}{|p{0.16\linewidth}|p{0.78\linewidth}|}
\hline
\centering 备注 &
% #VALUE_ID: LEX-P0007-V0003
% #FIELD_VALUE: 备注
% #HANDWRITTEN: /
\fieldvalue{\handwritten{/}}\\
\hline
\centering 实验人/日期 &
% #VALUE_ID: LEX-P0007-V0004
% #FIELD_VALUE: 实验人
% #HANDWRITTEN: 郭敏
\fieldvalue{\handwritten{郭敏}}\quad
% #VALUE_ID: LEX-P0007-V0005
% #FIELD_VALUE: 实验日期
% #HANDWRITTEN: 2025.08.14
\fieldvalue{\handwritten{2025.08.14}}\\
\hline
\centering 复核人/日期 &
% #VALUE_ID: LEX-P0007-V0006
% #FIELD_VALUE: 复核人
% #HANDWRITTEN: 姚虹
\fieldvalue{\handwritten{姚虹}}\quad
% #VALUE_ID: LEX-P0007-V0007
% #FIELD_VALUE: 复核日期
% #HANDWRITTEN: 2025.08.15
\fieldvalue{\handwritten{2025.08.15}}\\
\hline
\end{tabular}

\vfill
\begin{center}
\scriptsize 第 5 页\quad 共 5 页
\end{center}
% LEXOID_PAGE_COMPLETED: 7/70

\newpage

\noindent
\textbf{Assay:}
% #VALUE_ID: LEX-P0008-V0001
% #FIELD_VALUE: Assay
\fieldvalue{MSC-1}

\medskip
\noindent
\textbf{Cell Type F1:}
% #VALUE_ID: LEX-P0008-V0002
% #FIELD_VALUE: Cell Type F1
\fieldvalue{MSC (AO)}

\noindent
\textbf{Cell Type F2:}
% #VALUE_ID: LEX-P0008-V0003
% #FIELD_VALUE: Cell Type F2
\fieldvalue{MSC (PI)}

\medskip
\noindent
\textbf{Sample ID:}
% #VALUE_ID: LEX-P0008-V0004
% #FIELD_VALUE: Sample ID
\fieldvalue{QC-JS-QY0401-A31Z2202508006-20250815-A01-25H-1}

\noindent
\textbf{Dilution:}
% #VALUE_ID: LEX-P0008-V0005
% #FIELD_VALUE: Dilution
\fieldvalue{2.00}

\bigskip
\noindent\rule{\linewidth}{0.4pt}

\noindent\textbf{Results:}

\medskip
\noindent
\begin{tabular}{@{}p{0.29\linewidth}p{0.34\linewidth}p{0.29\linewidth}@{}}
\textbf{Count} & \textbf{Concentration} & \textbf{Mean Diameter} \\
\underline{\hspace{0.24\linewidth}} &
\underline{\hspace{0.29\linewidth}} &
\underline{\hspace{0.24\linewidth}} \\[2pt]

\textbf{Total:}
% #VALUE_ID: LEX-P0008-V0006
% #FIELD_VALUE: Count -- Total
\fieldvalue{148}
&
\textbf{Total:}
% #VALUE_ID: LEX-P0008-V0010
% #FIELD_VALUE: Concentration -- Total
\fieldvalue{$5.22\times10^{5}$ cells/mL}
&
\textbf{Total:}
% #VALUE_ID: LEX-P0008-V0013
% #FIELD_VALUE: Mean Diameter -- Total
\fieldvalue{14.3 micron}
\\

\textbf{Live:}
% #VALUE_ID: LEX-P0008-V0007
% #FIELD_VALUE: Count -- Live
\fieldvalue{133}
&
\textbf{Live:}
% #VALUE_ID: LEX-P0008-V0011
% #FIELD_VALUE: Concentration -- Live
\fieldvalue{$4.70\times10^{5}$ cells/mL}
&
\textbf{Live:}
% #VALUE_ID: LEX-P0008-V0014
% #FIELD_VALUE: Mean Diameter -- Live
\fieldvalue{14.7 microns}
\\

\textbf{Dead:}
% #VALUE_ID: LEX-P0008-V0008
% #FIELD_VALUE: Count -- Dead
\fieldvalue{15}
&
\textbf{Dead:}
% #VALUE_ID: LEX-P0008-V0012
% #FIELD_VALUE: Concentration -- Dead
\fieldvalue{$5.20\times10^{4}$ cells/mL}
&
\textbf{Dead:}
% #VALUE_ID: LEX-P0008-V0015
% #FIELD_VALUE: Mean Diameter -- Dead
\fieldvalue{11.0 microns}
\\

\textbf{Viability:}
% #VALUE_ID: LEX-P0008-V0009
% #FIELD_VALUE: Viability
\fieldvalue{90.0\%}
&
&
\end{tabular}

\bigskip

\noindent
\begin{tabular}{@{}p{0.46\linewidth}p{0.46\linewidth}@{}}
% TODO: Image panel visible on this page; source filename unavailable.
\fbox{\rule{0pt}{3.0cm}\rule{0.40\linewidth}{0pt}}
&
% TODO: Image panel visible on this page; source filename unavailable.
\fbox{\rule{0pt}{3.0cm}\rule{0.40\linewidth}{0pt}}
\\[8pt]
% TODO: Image panel visible on this page; source filename unavailable.
\fbox{\rule{0pt}{3.0cm}\rule{0.40\linewidth}{0pt}}
&
% TODO: Image panel visible on this page; source filename unavailable.
\fbox{\rule{0pt}{3.0cm}\rule{0.40\linewidth}{0pt}}
\end{tabular}

\vfill

\noindent
操作人/审核人:
% #VALUE_ID: LEX-P0008-V0016
% #FIELD_VALUE: 操作人/审核人
% TODO #HANDWRITTEN: illegible; handwritten signature
\fieldvalue{\handwritten{[illegible signature]}}

\hfill
日期:
% #VALUE_ID: LEX-P0008-V0017
% #FIELD_VALUE: 日期
% #TODO #HANDWRITTEN: 24年08月15; handwritten date is partially unclear
\fieldvalue{\handwritten{24年08月15}}
% LEXOID_PAGE_COMPLETED: 8/70

\newpage

\noindent Assay: MSC-1

\noindent Cell Type F1: MSC (AO)

\noindent Cell Type F2: MSC (PI)

\medskip

\noindent Sample ID: QC-JS-QY0401-A31Z220250806-20250815-A01-25H-2

\noindent Dilution: 2.00

\medskip

\noindent Results:

\medskip

\noindent
\begin{tabular}{@{}p{0.29\linewidth}p{0.33\linewidth}p{0.30\linewidth}@{}}
Count & Concentration & Mean Diameter \\
============ & ============ & ============ \\
Total: 170 & $6.04\times10^{5}$ cells/mL & 16.1 micron \\
Live: 157 & $5.58\times10^{5}$ cells/mL & 16.4 microns \\
Dead: 13 & $4.60\times10^{4}$ cells/mL & 12.8 microns \\
Viability: 92.4\% & & \\
\end{tabular}

\bigskip

\noindent
\begin{tabular}{@{}p{0.46\linewidth}p{0.46\linewidth}@{}}
\centering
% TODO: Image visible on the page; source filename unavailable.
\fbox{\parbox[c][0.22\textheight][c]{0.40\linewidth}{\centering Image placeholder}}
&
\centering
% TODO: Image visible on the page; source filename unavailable.
\fbox{\parbox[c][0.22\textheight][c]{0.40\linewidth}{\centering Image placeholder}}
\tabularnewline[1em]
\centering
% TODO: Image visible on the page; source filename unavailable.
\fbox{\parbox[c][0.22\textheight][c]{0.40\linewidth}{\centering Image placeholder}}
&
\centering
% TODO: Image visible on the page; source filename unavailable.
\fbox{\parbox[c][0.22\textheight][c]{0.40\linewidth}{\centering Image placeholder}}
\tabularnewline
\end{tabular}

\vfill

\noindent 操作人/审核人：
% #VALUE_ID: LEX-P0009-V0001
% #FIELD_VALUE: 操作人/审核人
% #TODO #HANDWRITTEN: illegible signature; handwritten signature is not clearly readable
\fieldvalue{\handwritten{[签名难以辨认]}}

\hfill 日期：
% #VALUE_ID: LEX-P0009-V0002
% #FIELD_VALUE: 日期
% #HANDWRITTEN: 2024.8.14
\fieldvalue{\handwritten{2024.8.14}}
% LEXOID_PAGE_COMPLETED: 9/70

\newpage

\noindent\textbf{Assay:}
% #VALUE_ID: LEX-P0010-V0001
% #FIELD_VALUE: Assay
\fieldvalue{MSC-1}

\medskip
\noindent\textbf{Cell Type F1:}
% #VALUE_ID: LEX-P0010-V0002
% #FIELD_VALUE: Cell Type F1
\fieldvalue{MSC (AO)}

\noindent\textbf{Cell Type F2:}
% #VALUE_ID: LEX-P0010-V0003
% #FIELD_VALUE: Cell Type F2
\fieldvalue{MSC (PI)}

\medskip
\noindent\textbf{Sample ID:}
% #VALUE_ID: LEX-P0010-V0004
% #FIELD_VALUE: Sample ID
\fieldvalue{QC-JS-QY0401-A31Z2202508006-20250815-A01-25H-3}

\noindent\textbf{Dilution:}
% #VALUE_ID: LEX-P0010-V0005
% #FIELD_VALUE: Dilution
\fieldvalue{2.00}

\bigskip
\noindent\rule{\linewidth}{0.4pt}

\noindent Results:

\medskip
\begin{tabular}{@{}p{0.29\linewidth}p{0.34\linewidth}p{0.29\linewidth}@{}}
\textbf{Count} & \textbf{Concentration} & \textbf{Mean Diameter} \\
\underline{\texttt{=============}} &
\underline{\texttt{=============}} &
\underline{\texttt{=============}} \\[2pt]
% #VALUE_ID: LEX-P0010-V0006
% #FIELD_VALUE: Count --- Total
\fieldvalue{Total: 115}
&
% #VALUE_ID: LEX-P0010-V0007
% #FIELD_VALUE: Concentration --- Total
\fieldvalue{$4.08\times10^{5}$ cells/mL}
&
% #VALUE_ID: LEX-P0010-V0008
% #FIELD_VALUE: Mean Diameter --- Total
\fieldvalue{16.1 micron}
\\[2pt]
% #VALUE_ID: LEX-P0010-V0009
% #FIELD_VALUE: Count --- Live
\fieldvalue{Live: 109}
&
% #VALUE_ID: LEX-P0010-V0010
% #FIELD_VALUE: Concentration --- Live
\fieldvalue{$3.87\times10^{5}$ cells/mL}
&
% #VALUE_ID: LEX-P0010-V0011
% #FIELD_VALUE: Mean Diameter --- Live
\fieldvalue{16.2 microns}
\\[2pt]
% #VALUE_ID: LEX-P0010-V0012
% #FIELD_VALUE: Count --- Dead
\fieldvalue{Dead: 6}
&
% #VALUE_ID: LEX-P0010-V0013
% #FIELD_VALUE: Concentration --- Dead
\fieldvalue{$2.10\times10^{4}$ cells/mL}
&
% #VALUE_ID: LEX-P0010-V0014
% #FIELD_VALUE: Mean Diameter --- Dead
\fieldvalue{14.6 micron}
\\[2pt]
% #VALUE_ID: LEX-P0010-V0015
% #FIELD_VALUE: Viability
\fieldvalue{Viability: 94.9\%}
&
&
\end{tabular}

\bigskip

\begin{center}
\begin{tabular}{@{}c@{\hspace{1.2cm}}c@{}}
\fbox{\parbox[c][4.0cm][c]{0.38\linewidth}{\centering\scriptsize TODO: image placeholder}} &
\fbox{\parbox[c][4.0cm][c]{0.38\linewidth}{\centering\scriptsize TODO: image placeholder}} \\[0.5cm]
\fbox{\parbox[c][4.0cm][c]{0.38\linewidth}{\centering\scriptsize TODO: image placeholder}} &
\fbox{\parbox[c][4.0cm][c]{0.38\linewidth}{\centering\scriptsize TODO: image placeholder}}
\end{tabular}
\end{center}

\vfill

\noindent 操作人/审核人：
% #VALUE_ID: LEX-P0010-V0016
% #FIELD_VALUE: 操作人/审核人
% TODO #HANDWRITTEN: illegible signature; handwritten signature is not clearly readable
\fieldvalue{\handwritten{illegible signature}}
\hfill
日期：
% #VALUE_ID: LEX-P0010-V0017
% #FIELD_VALUE: 日期
% #HANDWRITTEN: 2025.8.14
\fieldvalue{\handwritten{2025.8.14}}
% LEXOID_PAGE_COMPLETED: 10/70

\newpage

\begin{center}
\small
铂生物科技（北京）有限公司\\
\textit{PLATINUM LIFE EXCELLENCE BIOTECH}\hfill 文件编码：REC-YZ-SMP-QC-01-001-d\\[2pt]
\textbf{原件请验单}
\end{center}

\begin{center}
\small
\begin{tabular}{|p{2.0cm}|p{4.0cm}|p{2.0cm}|p{4.0cm}|}
\hline
样品名称 &
% #VALUE_ID: LEX-P0011-V0001
% #FIELD_VALUE: 样品名称
% #TODO #HANDWRITTEN: 三线虫疗液（疑似）；字迹模糊
\fieldvalue{\handwritten{三线虫疗液}} &
样品代码 &
% #VALUE_ID: LEX-P0011-V0002
% #FIELD_VALUE: 样品代码
% #TODO #HANDWRITTEN: QK/0001（疑似）；字迹模糊
\fieldvalue{\handwritten{QK/0001}}\\
\hline
样品批号 &
% #VALUE_ID: LEX-P0011-V0003
% #FIELD_VALUE: 样品批号
% #TODO #HANDWRITTEN: A?202?；字迹模糊
\fieldvalue{\handwritten{A?202?}} &
样品规格/数量 &
% #VALUE_ID: LEX-P0011-V0004
% #FIELD_VALUE: 样品规格/数量
% #TODO #HANDWRITTEN: 2瓶，5?；字迹模糊
\fieldvalue{\handwritten{2瓶，5?}}\\
\hline
样品类别 &
\multicolumn{3}{p{10.0cm}|}{
% #VALUE_ID: LEX-P0011-V0005
% #FIELD_VALUE: 生产过程中特检品
生产过程中特检品\fieldvalue{\checkboxfield{checked}}\quad
% #VALUE_ID: LEX-P0011-V0006
% #FIELD_VALUE: 包装包裹品
包装包裹品\fieldvalue{\checkboxfield{unchecked}}\quad
% #VALUE_ID: LEX-P0011-V0007
% #FIELD_VALUE: 成品
成品\fieldvalue{\checkboxfield{unchecked}}\quad
% #VALUE_ID: LEX-P0011-V0008
% #FIELD_VALUE: 其他
其他（\fieldvalue{\checkboxfield{unchecked}}\ ）}\\
\hline
检验类型 &
\multicolumn{3}{p{10.0cm}|}{
% #VALUE_ID: LEX-P0011-V0009
% #FIELD_VALUE: 初检
初检\fieldvalue{\checkboxfield{checked}}\quad
% #VALUE_ID: LEX-P0011-V0010
% #FIELD_VALUE: 复验
复验\fieldvalue{\checkboxfield{unchecked}}\quad
% #VALUE_ID: LEX-P0011-V0011
% #FIELD_VALUE: 退货
退货\fieldvalue{\checkboxfield{unchecked}}\quad
% #VALUE_ID: LEX-P0011-V0012
% #FIELD_VALUE: 稳定性
稳定性\fieldvalue{\checkboxfield{unchecked}}\quad
% #VALUE_ID: LEX-P0011-V0013
% #FIELD_VALUE: 留样
留样\fieldvalue{\checkboxfield{unchecked}}}\\
\hline
请验目的 &
\multicolumn{3}{p{10.0cm}|}{
% #VALUE_ID: LEX-P0011-V0014
% #FIELD_VALUE: 全检
全检\fieldvalue{\checkboxfield{checked}}\quad
% #VALUE_ID: LEX-P0011-V0015
% #FIELD_VALUE: 留样
留样\fieldvalue{\checkboxfield{unchecked}}\quad
% #VALUE_ID: LEX-P0011-V0016
% #FIELD_VALUE: 其他
其他（\fieldvalue{\checkboxfield{unchecked}}\ ）}\\
\hline
请验部门 &
\multicolumn{3}{p{10.0cm}|}{
% #VALUE_ID: LEX-P0011-V0017
% #FIELD_VALUE: 生产部
生产部\fieldvalue{\checkboxfield{checked}}\quad
% #VALUE_ID: LEX-P0011-V0018
% #FIELD_VALUE: 物料管理部
物料管理部\fieldvalue{\checkboxfield{unchecked}}\quad
% #VALUE_ID: LEX-P0011-V0019
% #FIELD_VALUE: 其他
其他（\fieldvalue{\checkboxfield{unchecked}}\ ）}\\
\hline
样品情况 &
\multicolumn{3}{p{10.0cm}|}{
% #VALUE_ID: LEX-P0011-V0020
% #FIELD_VALUE: 随请验单一起提供
随请验单一起提供\fieldvalue{\checkboxfield{checked}}\quad
% #VALUE_ID: LEX-P0011-V0021
% #FIELD_VALUE: 未随请验单一起提供
未随请验单一起提供\fieldvalue{\checkboxfield{unchecked}}}\\
\hline
请验人 &
% #VALUE_ID: LEX-P0011-V0022
% #FIELD_VALUE: 请验人
% #TODO #HANDWRITTEN: 刘某（疑似）；签名难以辨认
\fieldvalue{\handwritten{刘某}} &
请验日期 &
% #VALUE_ID: LEX-P0011-V0023
% #FIELD_VALUE: 请验日期
% #HANDWRITTEN: 2025.07.4
\fieldvalue{\handwritten{2025.07.4}}\\
\hline
收验人 &
% #VALUE_ID: LEX-P0011-V0024
% #FIELD_VALUE: 收验人
% #TODO #HANDWRITTEN: 签名难以辨认；蓝色手写签名
\fieldvalue{\handwritten{签名}} &
收验日期 &
% #VALUE_ID: LEX-P0011-V0025
% #FIELD_VALUE: 收验日期
% #TODO #HANDWRITTEN: 2025.03.14（疑似）；字迹模糊
\fieldvalue{\handwritten{2025.03.14}}\\
\hline
备注 &
\multicolumn{3}{p{10.0cm}|}{
% #VALUE_ID: LEX-P0011-V0026
% #FIELD_VALUE: 备注
% #TODO #HANDWRITTEN: 0h（疑似）；字迹模糊
\fieldvalue{\handwritten{0h}}\qquad
% #VALUE_ID: LEX-P0011-V0027
% #FIELD_VALUE: 备注
% #TODO #HANDWRITTEN: >1L（疑似）；字迹模糊
\fieldvalue{\handwritten{>1L}}}\\
\hline
\end{tabular}
\end{center}
% LEXOID_PAGE_COMPLETED: 11/70

\newpage

\begin{center}
\textbf{文件编码：}
% #VALUE_ID: LEX-P0012-V0001
% #FIELD_VALUE: 文件编码
\fieldvalue{REC-YZ-SMP-QC-01-004-b}
\qquad
版本号：
% #VALUE_ID: LEX-P0012-V0002
% #FIELD_VALUE: 版本号
\fieldvalue{1.4}

\vspace{0.5em}
{\Large\textbf{取样记录}}
\end{center}

% TODO: 页面左上角可见企业标志及企业名称，未使用图像文件，暂以文字结构表示。
\begin{center}
铂立生物科技（北京）有限公司\\
PLATINUM LIFE EXCELLENCE BIOTECH
\end{center}

\small
\renewcommand{\arraystretch}{1.8}
\begin{tabular}{|p{0.15\linewidth}|p{0.35\linewidth}|p{0.15\linewidth}|p{0.25\linewidth}|}
\hline
名称 &
% #VALUE_ID: LEX-P0012-V0003
% #FIELD_VALUE: 名称
\fieldvalue{一级生物反应器细胞液} &
代码 &
% #VALUE_ID: LEX-P0012-V0004
% #FIELD_VALUE: 代码
\fieldvalue{QY0401}
\\ \hline

规格 &
% #VALUE_ID: LEX-P0012-V0005
% #FIELD_VALUE: 规格
\fieldvalue{20ml/袋} &
总数量 &
% #VALUE_ID: LEX-P0012-V0006
% #FIELD_VALUE: 总数量
\fieldvalue{1袋}
\\ \hline

批号 &
% #VALUE_ID: LEX-P0012-V0007
% #FIELD_VALUE: 批号
\fieldvalue{A31Z2202508006} &
进厂批号 &
% #VALUE_ID: LEX-P0012-V0008
% #FIELD_VALUE: 进厂批号
\fieldvalue{/}
\\ \hline

\multicolumn{1}{|l|}{供应单位/厂家} &
\multicolumn{3}{l|}{
% #VALUE_ID: LEX-P0012-V0009
% #FIELD_VALUE: 供应单位/厂家
\fieldvalue{/}}
\\ \hline

取样目的 &
\multicolumn{3}{l|}{
% #VALUE_ID: LEX-P0012-V0010
% #FIELD_VALUE: 过程检验
\fieldvalue{\checkboxfield{checked}} 过程检验
\quad
% #VALUE_ID: LEX-P0012-V0011
% #FIELD_VALUE: 复检
\fieldvalue{\checkboxfield{unchecked}} 复检
\quad
% #VALUE_ID: LEX-P0012-V0012
% #FIELD_VALUE: 稳定性
\fieldvalue{\checkboxfield{unchecked}} 稳定性
}\\[-0.4em]
&\multicolumn{3}{l|}{
% #VALUE_ID: LEX-P0012-V0013
% #FIELD_VALUE: 退货检验
\fieldvalue{\checkboxfield{unchecked}} 退货检验
\quad
% #VALUE_ID: LEX-P0012-V0014
% #FIELD_VALUE: 留样
\fieldvalue{\checkboxfield{unchecked}} 留样
\quad
% #VALUE_ID: LEX-P0012-V0015
% #FIELD_VALUE: 其他
\fieldvalue{\checkboxfield{unchecked}} 其他：\underline{\hspace{3cm}}}
\\ \hline

取样地点 &
\multicolumn{3}{l|}{
% #VALUE_ID: LEX-P0012-V0016
% #FIELD_VALUE: 产品出口
\fieldvalue{\checkboxfield{checked}} 产品出口
\quad
% #VALUE_ID: LEX-P0012-V0017
% #FIELD_VALUE: 物料库
\fieldvalue{\checkboxfield{unchecked}} 物料库
\quad
% #VALUE_ID: LEX-P0012-V0018
% #FIELD_VALUE: 细胞库
\fieldvalue{\checkboxfield{unchecked}} 细胞库
\quad
% #VALUE_ID: LEX-P0012-V0019
% #FIELD_VALUE: 包装间
\fieldvalue{\checkboxfield{unchecked}} 包装间}
\\ \hline

贮存条件 &
\multicolumn{3}{l|}{
% #VALUE_ID: LEX-P0012-V0020
% #FIELD_VALUE: 2--8$^\circ$C
\fieldvalue{\checkboxfield{checked}} 2--8$^\circ$C
\quad
% #VALUE_ID: LEX-P0012-V0021
% #FIELD_VALUE: -20$^\circ$C
\fieldvalue{\checkboxfield{unchecked}} -20$^\circ$C
\quad
% #VALUE_ID: LEX-P0012-V0022
% #FIELD_VALUE: -80$^\circ$C
\fieldvalue{\checkboxfield{unchecked}} -80$^\circ$C
\quad
% #VALUE_ID: LEX-P0012-V0023
% #FIELD_VALUE: 常温
\fieldvalue{\checkboxfield{unchecked}} 常温
\quad
% #VALUE_ID: LEX-P0012-V0024
% #FIELD_VALUE: 其他
\fieldvalue{\checkboxfield{unchecked}} 其他：\underline{\hspace{1.5cm}}}
\\ \hline

盛装容器 &
\multicolumn{3}{l|}{
% #VALUE_ID: LEX-P0012-V0025
% #FIELD_VALUE: 保湿箱
\fieldvalue{\checkboxfield{checked}} 保湿箱
\quad
% #VALUE_ID: LEX-P0012-V0026
% #FIELD_VALUE: 泡沫箱
\fieldvalue{\checkboxfield{unchecked}} 泡沫箱
\quad
% #VALUE_ID: LEX-P0012-V0027
% #FIELD_VALUE: 手提篮
\fieldvalue{\checkboxfield{unchecked}} 手提篮
\quad
% #VALUE_ID: LEX-P0012-V0028
% #FIELD_VALUE: 其他
\fieldvalue{\checkboxfield{unchecked}} 其他：\underline{\hspace{1.5cm}}}
\\ \hline

外观复核 &
\multicolumn{3}{l|}{
% #VALUE_ID: LEX-P0012-V0029
% #FIELD_VALUE: 合格
\fieldvalue{\checkboxfield{checked}} 合格
\quad
% #VALUE_ID: LEX-P0012-V0030
% #FIELD_VALUE: 不合格
\fieldvalue{\checkboxfield{unchecked}} 不合格}
\\ \hline

取样量 &
% #VALUE_ID: LEX-P0012-V0031
% #FIELD_VALUE: 取样量
\fieldvalue{1袋} &
检验量 &
% #VALUE_ID: LEX-P0012-V0032
% #FIELD_VALUE: 检验量
\fieldvalue{1袋}
\\ \hline

\multicolumn{1}{|l|}{} &
\multicolumn{1}{c|}{} &
\multicolumn{1}{l|}{留样量}
&
% #VALUE_ID: LEX-P0012-V0033
% #FIELD_VALUE: 留样量
\fieldvalue{/}
\\ \hline

取样人/日期 &
\multicolumn{1}{c|}{
% #VALUE_ID: LEX-P0012-V0034
% #FIELD_VALUE: 取样人
% #TODO #HANDWRITTEN: 签名（难辨）; 蓝色手写签名无法可靠辨认
\fieldvalue{\handwritten{签名（难辨）}}}
&
\multicolumn{1}{l|}{}
&
% #VALUE_ID: LEX-P0012-V0035
% #FIELD_VALUE: 取样日期
% #HANDWRITTEN: 2025.04.12
\fieldvalue{\handwritten{2025.04.12}}
\\ \hline

复核人/日期 &
\multicolumn{1}{c|}{
% #VALUE_ID: LEX-P0012-V0036
% #FIELD_VALUE: 复核人
% #TODO #HANDWRITTEN: 签名（难辨）; 蓝色手写签名无法可靠辨认
\fieldvalue{\handwritten{签名（难辨）}}}
&
\multicolumn{1}{l|}{}
&
% #VALUE_ID: LEX-P0012-V0037
% #FIELD_VALUE: 复核日期
% #HANDWRITTEN: 2025.04.12
\fieldvalue{\handwritten{2025.04.12}}
\\ \hline

备注 &
\multicolumn{3}{l|}{
% #VALUE_ID: LEX-P0012-V0038
% #FIELD_VALUE: 备注
% #HANDWRITTEN: 0h
\fieldvalue{\handwritten{0h}}
\qquad
% #VALUE_ID: LEX-P0012-V0039
% #FIELD_VALUE: 备注
% #TODO #HANDWRITTEN: 2.1L; 手写单位和数字略有模糊
\fieldvalue{\handwritten{2.1L}}}
\\ \hline
\end{tabular}

\vfill
\begin{center}
第 1 页 共 1 页
\end{center}
% LEXOID_PAGE_COMPLETED: 12/70

\newpage

\begin{center}
\textbf{生物反应器细胞液检验操作记录}\\
文件编码：REC-YZ-SOP-QC-99-006-a
\end{center}

\subsection*{葡萄糖浓度检验}

\begin{tabular}{|p{0.14\linewidth}|p{0.30\linewidth}|p{0.14\linewidth}|p{0.30\linewidth}|}
\hline
样品名称 &
% #VALUE_ID: LEX-P0013-V0001
% #FIELD_VALUE: 样品名称
\fieldvalue{二级生物反应器细胞液} &
样品规格 &
% #VALUE_ID: LEX-P0013-V0002
% #FIELD_VALUE: 样品规格
% #HANDWRITTEN: 20
\fieldvalue{\handwritten{20}}\ \ensuremath{\mathrm{ml}}/袋
\\ \hline
样品批号 &
% #VALUE_ID: LEX-P0013-V0003
% #FIELD_VALUE: 样品批号
% #TODO #HANDWRITTEN: 21?1?906; handwriting unclear
\fieldvalue{\handwritten{21?1?906}} &
取样日期 &
% #VALUE_ID: LEX-P0013-V0004
% #FIELD_VALUE: 取样日期
% #HANDWRITTEN: 2020年03月14日
\fieldvalue{\handwritten{2020年03月14日}}
\\ \hline
\end{tabular}

\section*{1\quad 仪器设备、试剂耗材}

\subsection*{1.1\quad 仪器设备}

\begin{tabular}{|p{0.15\linewidth}|p{0.27\linewidth}|p{0.15\linewidth}|p{0.20\linewidth}|p{0.18\linewidth}|}
\hline
名称 & 厂家 & 型号 & 设备编号 & 检查确认 \\ \hline
% #VALUE_ID: LEX-P0013-V0005
% #FIELD_VALUE: 仪器设备名称
\fieldvalue{血糖仪} &
% #VALUE_ID: LEX-P0013-V0006
% #FIELD_VALUE: 仪器设备厂家
\fieldvalue{三诺生物传感股份有限公司} &
% #VALUE_ID: LEX-P0013-V0007
% #FIELD_VALUE: 仪器设备型号
\fieldvalue{安稳+} &
 &
% #VALUE_ID: LEX-P0013-V0008
% #FIELD_VALUE: 检查确认-正常
\fieldvalue{\checkboxfield{checked}} 正常\quad
% #VALUE_ID: LEX-P0013-V0009
% #FIELD_VALUE: 检查确认-异常
\fieldvalue{\checkboxfield{unchecked}} 异常
\\ \hline
% #VALUE_ID: LEX-P0013-V0010
% #FIELD_VALUE: 仪器设备名称
\fieldvalue{移液器} &
% #VALUE_ID: LEX-P0013-V0011
% #FIELD_VALUE: 仪器设备厂家
\fieldvalue{Eppendorf Corporate} &
% #VALUE_ID: LEX-P0013-V0012
% #FIELD_VALUE: 仪器设备型号
\fieldvalue{0.1--2.5\,\ensuremath{\mu\mathrm{l}}} &
% #VALUE_ID: LEX-P0013-V0013
% #FIELD_VALUE: 设备编号
% #TODO #HANDWRITTEN: 073486; handwriting unclear
\fieldvalue{\handwritten{073486}} &
% #VALUE_ID: LEX-P0013-V0014
% #FIELD_VALUE: 检查确认-正常
\fieldvalue{\checkboxfield{checked}} 正常\quad
% #VALUE_ID: LEX-P0013-V0015
% #FIELD_VALUE: 检查确认-异常
\fieldvalue{\checkboxfield{unchecked}} 异常
\\ \hline
\end{tabular}

\subsection*{1.2\quad 试剂耗材}

\begin{tabular}{|p{0.18\linewidth}|p{0.35\linewidth}|p{0.22\linewidth}|p{0.18\linewidth}|}
\hline
名称 & 厂家 & 批号 & 有效期至 \\ \hline
% #VALUE_ID: LEX-P0013-V0016
% #FIELD_VALUE: 试剂耗材名称
\fieldvalue{血糖测试条} &
% #VALUE_ID: LEX-P0013-V0017
% #FIELD_VALUE: 试剂耗材厂家
\fieldvalue{三诺生物传感股份有限公司} &
% #VALUE_ID: LEX-P0013-V0018
% #FIELD_VALUE: 批号
% #TODO #HANDWRITTEN: 442?4706; handwriting unclear
\fieldvalue{\handwritten{442?4706}} &
% #VALUE_ID: LEX-P0013-V0019
% #FIELD_VALUE: 有效期至
% #TODO #HANDWRITTEN: 2021-? -?8; handwriting unclear
\fieldvalue{\handwritten{2021-? -?8}}
\\ \hline
\end{tabular}

\section*{2\quad 操作步骤}

\begin{tabular}{|p{0.58\linewidth}|p{0.34\linewidth}|}
\hline
\multicolumn{2}{|c|}{操作过程} \\ \hline
取出测试条，取样品自然沉降后的上清液 0.6\,\ensuremath{\mu\mathrm{l}} 加入测试条电极，应重，血糖仪自动发出“嘀”提示音，吸入样品完成，血糖仪自动倒计时，开始测试； &
取上清液：
% #VALUE_ID: LEX-P0013-V0020
% #FIELD_VALUE: 取上清液体积
% #HANDWRITTEN: 0.6
\fieldvalue{\handwritten{0.6}}\ \ensuremath{\mu\mathrm{l}}
\\ \hline
重复以上步骤 2 次，共计数 3 次。 &
每个样品分别共计数
% #VALUE_ID: LEX-P0013-V0021
% #FIELD_VALUE: 样品计数次数
% #HANDWRITTEN: 3
\fieldvalue{\handwritten{3}}
次。
\\ \hline
\multicolumn{2}{|c|}{检测结果} \\ \hline
\multicolumn{2}{|c|}{葡萄糖浓度（mmol/L）} \\ \hline
\multicolumn{2}{|c|}{\begin{tabular}{cccc}
第一次 & 第二次 & 第三次 & 平均值（保留一位小数） \\ \hline
% #VALUE_ID: LEX-P0013-V0022
% #FIELD_VALUE: 第一次检测结果
% #HANDWRITTEN: 24.0
\fieldvalue{\handwritten{24.0}} &
% #VALUE_ID: LEX-P0013-V0023
% #FIELD_VALUE: 第二次检测结果
% #HANDWRITTEN: 21.6
\fieldvalue{\handwritten{21.6}} &
% #VALUE_ID: LEX-P0013-V0024
% #FIELD_VALUE: 第三次检测结果
% #HANDWRITTEN: 22.4
\fieldvalue{\handwritten{22.4}} &
% #VALUE_ID: LEX-P0013-V0025
% #FIELD_VALUE: 平均值
% #HANDWRITTEN: 21.0
\fieldvalue{\handwritten{21.0}}
\end{tabular}} \\ \hline
备注 &
% #VALUE_ID: LEX-P0013-V0026
% #FIELD_VALUE: 备注
% #HANDWRITTEN: /
\fieldvalue{\handwritten{/}}
\\ \hline
实验人/日期 &
% #VALUE_ID: LEX-P0013-V0027
% #FIELD_VALUE: 实验人
% #TODO #HANDWRITTEN: 徐?; handwriting unclear
\fieldvalue{\handwritten{徐?}}\quad
% #VALUE_ID: LEX-P0013-V0028
% #FIELD_VALUE: 实验日期
% #HANDWRITTEN: 2020.3.14
\fieldvalue{\handwritten{2020.3.14}}
\\ \hline
复核人/日期 &
% #VALUE_ID: LEX-P0013-V0029
% #FIELD_VALUE: 复核人
% #TODO #HANDWRITTEN: 李?; handwriting unclear
\fieldvalue{\handwritten{李?}}\quad
% #VALUE_ID: LEX-P0013-V0030
% #FIELD_VALUE: 复核日期
% #HANDWRITTEN: 2020.3.16
\fieldvalue{\handwritten{2020.3.16}}
\\ \hline
\end{tabular}

% TODO: The bottom footer indicates this is page 1 of 5; continuation content is not visible on this PDF page.
% LEXOID_PAGE_COMPLETED: 13/70

\newpage

\begin{center}
\textbf{生物反应器细胞液检验记录}
\end{center}

\noindent
文件编码：REC-YZ-SOR-QC-99-006-a

\begin{center}
\begin{tabular}{|p{2.2cm}|p{5.2cm}|p{2.2cm}|p{4.0cm}|}
\hline
\multicolumn{4}{|c|}{密度与浑浊检验}\\
\hline
样品名称 &
% #VALUE_ID: LEX-P0014-V0001
% #FIELD_VALUE: 样品名称
\fieldvalue{二级生物反应器细胞液} &
样品规格 &
% #VALUE_ID: LEX-P0014-V0002
% #FIELD_VALUE: 样品规格
% #HANDWRITTEN: 20 ml/袋
\fieldvalue{\handwritten{20 ml/袋}}\\
\hline
样品批号 &
% #VALUE_ID: LEX-P0014-V0003
% #FIELD_VALUE: 样品批号
% TODO #HANDWRITTEN: [批号]; handwriting is difficult to read
\fieldvalue{\handwritten{[批号]}} &
取样日期 &
% #VALUE_ID: LEX-P0014-V0004
% #FIELD_VALUE: 取样日期
% #HANDWRITTEN: 2024年08月04日
\fieldvalue{\handwritten{2024年08月04日}}\\
\hline
\end{tabular}
\end{center}

\section*{1\quad 仪器设备、试剂耗材}

\subsection*{1.1\quad 仪器设备}

\begin{center}
\small
\begin{tabular}{|p{2.3cm}|p{2.2cm}|p{2.6cm}|p{2.0cm}|p{2.3cm}|p{1.2cm}|p{1.2cm}|}
\hline
名称 & 厂家 & 型号 & 设备编号 & 校准有效期至 & \multicolumn{2}{c|}{检查确认}\\
\cline{6-7}
&&&&&正常&异常\\
\hline
% #VALUE_ID: LEX-P0014-V0005
% #FIELD_VALUE: 名称
\fieldvalue{细胞计数} &
% #VALUE_ID: LEX-P0014-V0006
% #FIELD_VALUE: 厂家
\fieldvalue{Nexcelom} &
% #VALUE_ID: LEX-P0014-V0007
% #FIELD_VALUE: 型号
\fieldvalue{Cellometer K2} &
% #VALUE_ID: LEX-P0014-V0008
% #FIELD_VALUE: 设备编号
\fieldvalue{1410905} &
% #VALUE_ID: LEX-P0014-V0009
% #FIELD_VALUE: 校准有效期至
% #HANDWRITTEN: 2024.11.10
\fieldvalue{\handwritten{2024.11.10}} &
% #VALUE_ID: LEX-P0014-V0010
% #FIELD_VALUE: 正常
% #HANDWRITTEN: ✓
\fieldvalue{\handwritten{\checkboxfield{checked}}} &
% #VALUE_ID: LEX-P0014-V0011
% #FIELD_VALUE: 异常
\fieldvalue{\checkboxfield{unchecked}}\\
\hline
% #VALUE_ID: LEX-P0014-V0012
% #FIELD_VALUE: 名称
\fieldvalue{高速冷冻离心机} &
% #VALUE_ID: LEX-P0014-V0013
% #FIELD_VALUE: 厂家
\fieldvalue{Thermo} &
% #VALUE_ID: LEX-P0014-V0014
% #FIELD_VALUE: 型号
\fieldvalue{ST16R} &
% #VALUE_ID: LEX-P0014-V0015
% #FIELD_VALUE: 设备编号
\fieldvalue{1410810} &
% #VALUE_ID: LEX-P0014-V0016
% #FIELD_VALUE: 校准有效期至
% #HANDWRITTEN: 2024.8.21
\fieldvalue{\handwritten{2024.8.21}} &
% #VALUE_ID: LEX-P0014-V0017
% #FIELD_VALUE: 正常
% #HANDWRITTEN: ✓
\fieldvalue{\handwritten{\checkboxfield{checked}}} &
% #VALUE_ID: LEX-P0014-V0018
% #FIELD_VALUE: 异常
\fieldvalue{\checkboxfield{unchecked}}\\
\hline
% #VALUE_ID: LEX-P0014-V0019
% #FIELD_VALUE: 名称
\fieldvalue{电热恒温水槽} &
% #VALUE_ID: LEX-P0014-V0020
% #FIELD_VALUE: 厂家
\fieldvalue{上海一恒} &
% #VALUE_ID: LEX-P0014-V0021
% #FIELD_VALUE: 型号
\fieldvalue{DK-8AXX} &
% #VALUE_ID: LEX-P0014-V0022
% #FIELD_VALUE: 设备编号
\fieldvalue{1411001} &
% #VALUE_ID: LEX-P0014-V0023
% #FIELD_VALUE: 校准有效期至
% #HANDWRITTEN: 2024.10.22
\fieldvalue{\handwritten{2024.10.22}} &
% #VALUE_ID: LEX-P0014-V0024
% #FIELD_VALUE: 正常
% #HANDWRITTEN: ✓
\fieldvalue{\handwritten{\checkboxfield{checked}}} &
% #VALUE_ID: LEX-P0014-V0025
% #FIELD_VALUE: 异常
\fieldvalue{\checkboxfield{unchecked}}\\
\hline
\end{tabular}
\end{center}

\subsection*{1.2\quad 试剂耗材}

\begin{center}
\small
\begin{tabular}{|p{4.0cm}|p{3.5cm}|p{3.8cm}|p{3.3cm}|}
\hline
名称 & 厂家 & 批号 & 有效期至\\
\hline
% #VALUE_ID: LEX-P0014-V0026
% #FIELD_VALUE: 名称
\fieldvalue{AO/PI染液} &
% #VALUE_ID: LEX-P0014-V0027
% #FIELD_VALUE: 厂家
\fieldvalue{Nexcelom} &
% #VALUE_ID: LEX-P0014-V0028
% #FIELD_VALUE: 批号
% #HANDWRITTEN: 737?47
\fieldvalue{\handwritten{737?47}} &
% #VALUE_ID: LEX-P0014-V0029
% #FIELD_VALUE: 有效期至
% #HANDWRITTEN: 2026.01.22
\fieldvalue{\handwritten{2026.01.22}}\\
\hline
% #VALUE_ID: LEX-P0014-V0030
% #FIELD_VALUE: 名称
\fieldvalue{0.9\%氯化钠注射液} &
% #VALUE_ID: LEX-P0014-V0031
% #FIELD_VALUE: 厂家
\fieldvalue{石家庄四药} &
% #VALUE_ID: LEX-P0014-V0032
% #FIELD_VALUE: 批号
% #HANDWRITTEN: 24/05/3?8
\fieldvalue{\handwritten{24/05/3?8}} &
% #VALUE_ID: LEX-P0014-V0033
% #FIELD_VALUE: 有效期至
% #HANDWRITTEN: 2026/09/06
\fieldvalue{\handwritten{2026/09/06}}\\
\hline
% #VALUE_ID: LEX-P0014-V0034
% #FIELD_VALUE: 名称
\fieldvalue{20\%人血白蛋白} &
% #VALUE_ID: LEX-P0014-V0035
% #FIELD_VALUE: 厂家
\fieldvalue{Baxter AG} &
% #VALUE_ID: LEX-P0014-V0036
% #FIELD_VALUE: 批号
% #HANDWRITTEN: 4/18?4?4
\fieldvalue{\handwritten{4/18?4?4}} &
% #VALUE_ID: LEX-P0014-V0037
% #FIELD_VALUE: 有效期至
% #HANDWRITTEN: 2024.03.31
\fieldvalue{\handwritten{2024.03.31}}\\
\hline
\end{tabular}
\end{center}

\section*{2\quad 操作步骤：}

\begin{center}
\begin{tabular}{|p{8.0cm}|p{5.3cm}|}
\hline
\multicolumn{2}{|c|}{溶液配制}\\
\hline
1\%人白蛋白化细胞液配制：取5 ml 20\%人白蛋白加入到95 ml 的0.9\%氯化钠注射液中，混匀即得。 &
20\%人血白蛋白
% #VALUE_ID: LEX-P0014-V0038
% #FIELD_VALUE: 20\%人血白蛋白用量
% #HANDWRITTEN: 1
\fieldvalue{\handwritten{1}}\ ml\\
\cline{2-2}
&
0.9\%氯化钠注射液
% #VALUE_ID: LEX-P0014-V0039
% #FIELD_VALUE: 0.9\%氯化钠注射液用量
% #HANDWRITTEN: 99
\fieldvalue{\handwritten{99}}\ ml\\
\hline
\multicolumn{2}{|c|}{
\begin{tabular}{p{3.5cm}p{5.0cm}p{4.0cm}}
名称 & 配制批号 & 有效期至\\
% #VALUE_ID: LEX-P0014-V0040
% #FIELD_VALUE: 名称
\fieldvalue{5$\times$裂解液} &
% #VALUE_ID: LEX-P0014-V0041
% #FIELD_VALUE: 配制批号
% #HANDWRITTEN: S?07?7?0
\fieldvalue{\handwritten{S?07?7?0}} &
% #VALUE_ID: LEX-P0014-V0042
% #FIELD_VALUE: 有效期至
% #HANDWRITTEN: 2024.12.21
\fieldvalue{\handwritten{2024.12.21}}
\end{tabular}}\\
\hline
\end{tabular}
\end{center}

\noindent
按照50 ml离心管的刻度读取样品体积，往样品中加入对应体积的5$\times$裂解液（每4 ml样品加入1 ml 5$\times$裂解液），按照“只进不舍”原则，如26 ml样品，加入7 ml 5$\times$裂解液，充分混匀后倒入至3.3的样品中，置于37℃电热恒温水槽中裂解20 min后，取出离心管摇晃混匀，使样品充分裂解；再置于37℃电热恒温水槽中孵育10 min得到细胞悬液。

\begin{center}
\begin{tabular}{|p{8.0cm}|p{5.3cm}|}
\hline
样品体积：
% #VALUE_ID: LEX-P0014-V0043
% #FIELD_VALUE: 样品体积
% #HANDWRITTEN: 24
\fieldvalue{\handwritten{24}}\ ml &
裂解液：
% #VALUE_ID: LEX-P0014-V0044
% #FIELD_VALUE: 裂解液用量
% #HANDWRITTEN: 7
\fieldvalue{\handwritten{7}}\ ml\\
\hline
电热恒温水槽
% #VALUE_ID: LEX-P0014-V0045
% #FIELD_VALUE: 电热恒温水槽温度
% #HANDWRITTEN: 37
\fieldvalue{\handwritten{37}}~$^\circ$C &
裂解：
% #VALUE_ID: LEX-P0014-V0046
% #FIELD_VALUE: 裂解时间
% #HANDWRITTEN: 20
\fieldvalue{\handwritten{20}}\ min\\
\hline
\end{tabular}
\end{center}

\noindent
取出裂解好的细胞悬液300$\times g$离心5分钟，按照读取样品体积加入对应体积的1\%人白蛋白化细胞液（每1.2 ml样品体积加入0.1 ml 1\%人白蛋白化细胞液，按照“只进不舍”原则，如26 ml）。

\begin{center}
\begin{tabular}{|p{8.0cm}|p{5.3cm}|}
\hline
取出裂解好的细胞悬液300$\times g$离心5分钟，按照读取样品体积加入对应体积的1\%人白蛋白化细胞液 &
% #VALUE_ID: LEX-P0014-V0047
% #FIELD_VALUE: 离心力
% #HANDWRITTEN: 300
\fieldvalue{\handwritten{300}}\ $\times g$离心
% #VALUE_ID: LEX-P0014-V0048
% #FIELD_VALUE: 离心时间
% #HANDWRITTEN: 5
\fieldvalue{\handwritten{5}}\ 分钟\\
\cline{2-2}
&
加入1\%人白蛋白化细胞液总体积：
% #VALUE_ID: LEX-P0014-V0049
% #FIELD_VALUE: 1\%人白蛋白化细胞液总体积
% #HANDWRITTEN: 1
\fieldvalue{\handwritten{1}}\ ml\\
\cline{2-2}
&
吸取
% #VALUE_ID: LEX-P0014-V0050
% #FIELD_VALUE: 吸取体积
% #HANDWRITTEN: 20
\fieldvalue{\handwritten{20}}\ $\mu$l\\
\hline
\end{tabular}
\end{center}

\begin{center}
第2页 共5页
\end{center}
% LEXOID_PAGE_COMPLETED: 14/70

\newpage

\noindent
\textbf{文件编号：} REC-YZ-SOP-QC-99-006-a

\vspace{0.5em}

\noindent
加入 2.2 ml 1\% 人白细胞稀释液，反复吹打至细胞悬液均匀。\\
撕去计数板上下两层保护膜，将细胞悬液混匀，从中吸取 20~$\mu$l\\
与 20~$\mu$l AO/PI 染料混合均匀。\\
吸取
% #VALUE_ID: LEX-P0015-V0001
% #FIELD_VALUE: 吸取体积
% #HANDWRITTEN: 20
\fieldvalue{\handwritten{20}}~$\mu$l 计数。\\
吸取 20~$\mu$l 细胞悬液与 AO/PI 染料的混合液，注入到计数板的一个检测室中。\\
将加入混合液的一端插入进样口。\\
选择 Assay 类型，使用 AO/PI 染料进行双荧光存活率检测；在主界面左侧 ID 栏中输入样本批号。\\
点击主界面左下角“Preview”预览细胞图片。\\
调整仪器焦距直至观察到荧光通道下的细胞中心透亮，轮廓清晰。\\
点击“count”按钮进行计数并自动拍摄图片。\\
重复以上步骤 2 次，共计数 3 次。拍摄好数据。\\
数据保存路径：

\noindent
% #VALUE_ID: LEX-P0015-V0002
% #FIELD_VALUE: 共计数次数
% #HANDWRITTEN: 3
\fieldvalue{\handwritten{3}} 次。

\vspace{0.5em}

\noindent
\begin{tabular}{|p{0.95\linewidth}|}
\hline
\centering 检测结果\tabularnewline
\hline
\end{tabular}

\noindent
检测结果：
% #VALUE_ID: LEX-P0015-V0003
% #FIELD_VALUE: 一级反应
\fieldvalue{\checkboxfield{checked}} 一级反应，
% #VALUE_ID: LEX-P0015-V0004
% #FIELD_VALUE: 二级反应
\fieldvalue{\checkboxfield{unchecked}} 二级反应，
% #VALUE_ID: LEX-P0015-V0005
% #FIELD_VALUE: 三级反应
\fieldvalue{\checkboxfield{unchecked}} 三级反应

\vspace{0.3em}

\noindent
\small
\begin{tabular}{|p{0.09\linewidth}|p{0.18\linewidth}|p{0.18\linewidth}|p{0.14\linewidth}|p{0.13\linewidth}|p{0.13\linewidth}|p{0.14\linewidth}|}
\hline
检测次数 & 活细胞密度（个细胞/ml） & 样品密度（个细胞/ml） & 活细胞密度 CV 值（\%） & 细胞活率（\%） & 样品活率（\%） & 细胞活率 CV 值（\%）\\
\hline
第一次 &
% #VALUE_ID: LEX-P0015-V0006
% #FIELD_VALUE: 第一次活细胞密度
% #HANDWRITTEN: 2.78$\times 10^5$
\fieldvalue{\handwritten{$2.78\times10^5$}} &
&
&
% #VALUE_ID: LEX-P0015-V0007
% #FIELD_VALUE: 第一次细胞活率
% #HANDWRITTEN: 89.7
\fieldvalue{\handwritten{89.7}} &
&
\\
\hline
第二次 &
% #VALUE_ID: LEX-P0015-V0008
% #FIELD_VALUE: 第二次活细胞密度
% #HANDWRITTEN: 2.72$\times 10^5$
\fieldvalue{\handwritten{$2.72\times10^5$}} &
% #VALUE_ID: LEX-P0015-V0009
% #FIELD_VALUE: 样品密度
% #HANDWRITTEN: 2.18$\times 10^4$
\fieldvalue{\handwritten{$2.18\times10^4$}} &
% #VALUE_ID: LEX-P0015-V0010
% #FIELD_VALUE: 活细胞密度 CV 值
% #HANDWRITTEN: 9
\fieldvalue{\handwritten{9}} &
% #VALUE_ID: LEX-P0015-V0011
% #FIELD_VALUE: 第二次细胞活率
% #HANDWRITTEN: 87.7
\fieldvalue{\handwritten{87.7}} &
% #VALUE_ID: LEX-P0015-V0012
% #FIELD_VALUE: 样品活率
% #HANDWRITTEN: 88
\fieldvalue{\handwritten{88}} &
% #VALUE_ID: LEX-P0015-V0013
% #FIELD_VALUE: 细胞活率 CV 值
% #HANDWRITTEN: 2
\fieldvalue{\handwritten{2}}\\
\hline
第三次 &
% #VALUE_ID: LEX-P0015-V0014
% #FIELD_VALUE: 第三次活细胞密度
% #HANDWRITTEN: 2.3$\times 10^5$
\fieldvalue{\handwritten{$2.3\times10^5$}} &
&
&
% #VALUE_ID: LEX-P0015-V0015
% #FIELD_VALUE: 第三次细胞活率
% #HANDWRITTEN: 87.1
\fieldvalue{\handwritten{87.1}} &
&
\\
\hline
\end{tabular}
\normalsize

\vspace{0.4em}

\noindent
\begin{tabular}{|p{0.28\linewidth}|p{0.21\linewidth}|p{0.28\linewidth}|p{0.18\linewidth}|}
\hline
发游离细胞体积（ml） &
% #VALUE_ID: LEX-P0015-V0016
% #FIELD_VALUE: 发游离细胞体积
% #HANDWRITTEN: 1/10
\fieldvalue{\handwritten{1/10}} &
细胞数量（个细胞/瓶） &
% #VALUE_ID: LEX-P0015-V0017
% #FIELD_VALUE: 细胞数量
% #HANDWRITTEN: 4.58
\fieldvalue{\handwritten{4.58}}\\
\hline
\end{tabular}

\vspace{0.3em}

\noindent
计算要求：样品密度为 3 次活细胞密度结果平均值乘以 1\% 人白细胞稀释液最终体积，再除以生物反应器细胞液体积；细胞数量为样品密度乘以发游离体积。

\vspace{0.3em}

\noindent
\begin{tabular}{|p{0.16\linewidth}|p{0.80\linewidth}|}
\hline
备注 &
% #VALUE_ID: LEX-P0015-V0018
% #FIELD_VALUE: 备注
% #HANDWRITTEN: /
\fieldvalue{\handwritten{/}}\\
\hline
实验人/日期 &
% #VALUE_ID: LEX-P0015-V0019
% #FIELD_VALUE: 实验人
% #TODO #HANDWRITTEN: illegible; handwritten signature
\fieldvalue{\handwritten{[字迹不清]}}\quad
% #VALUE_ID: LEX-P0015-V0020
% #FIELD_VALUE: 实验日期
% #HANDWRITTEN: 2025.08.14
\fieldvalue{\handwritten{2025.08.14}}\\
\hline
复核人/日期 &
% #VALUE_ID: LEX-P0015-V0021
% #FIELD_VALUE: 复核人
% #TODO #HANDWRITTEN: illegible; handwritten signature
\fieldvalue{\handwritten{[字迹不清]}}\quad
% #VALUE_ID: LEX-P0015-V0022
% #FIELD_VALUE: 复核日期
% #TODO #HANDWRITTEN: 2025.08.40; final digit unclear
\fieldvalue{\handwritten{2025.08.40}}\\
\hline
\end{tabular}

% TODO: The upper-right instruction block contains additional handwritten markings around the AO/PI volume and counting instruction; only the clearly legible editable values have been transcribed above.
% LEXOID_PAGE_COMPLETED: 15/70

\newpage

\noindent\textbf{文件编码：}%
% #VALUE_ID: LEX-P0016-V0001
% #FIELD_VALUE: 文件编码
\fieldvalue{REC-YZ-SOR-QC-99-0006-a}
\hfill
\textbf{版本号：}%
% #VALUE_ID: LEX-P0016-V0002
% #FIELD_VALUE: 版本号
\fieldvalue{1/4}

\begin{center}
\textbf{生物反应器细胞液检测操作记录}
\end{center}

\begin{center}
\begin{tabular}{|p{0.14\linewidth}|p{0.36\linewidth}|p{0.14\linewidth}|p{0.25\linewidth}|}
\hline
\multicolumn{4}{|c|}{细胞生长状态检验}\\
\hline
样品名称 &
% #VALUE_ID: LEX-P0016-V0003
% #FIELD_VALUE: 样品名称
\fieldvalue{缓生物反应器细胞液} &
样品规格 & ml/袋\\
\hline
样品批号 & & 取样日期 & 年\quad 月\quad 日\\
\hline
\end{tabular}
\end{center}

\section*{1.\quad 仪器设备及试剂：}

\subsection*{1.1.\quad 仪器设备：}

\begin{center}
\begin{tabular}{|p{0.22\linewidth}|p{0.17\linewidth}|p{0.22\linewidth}|p{0.20\linewidth}|p{0.19\linewidth}|}
\hline
仪器名称 & 厂家 & 型号 & 设备编号 & 检查确认\\
\hline
% #VALUE_ID: LEX-P0016-V0004
% #FIELD_VALUE: 仪器名称
\fieldvalue{荧光显微镜} &
% #VALUE_ID: LEX-P0016-V0005
% #FIELD_VALUE: 仪器厂家
\fieldvalue{保卡} &
% #VALUE_ID: LEX-P0016-V0006
% #FIELD_VALUE: 仪器型号
\fieldvalue{Dmi1 LED} &
% #VALUE_ID: LEX-P0016-V0007
% #FIELD_VALUE: 设备编号
\fieldvalue{1431404} &
正常\,
% #VALUE_ID: LEX-P0016-V0008
% #FIELD_VALUE: 检查确认--正常
\fieldvalue{\checkboxfield{unchecked}}
\quad
异常\,
% #VALUE_ID: LEX-P0016-V0009
% #FIELD_VALUE: 检查确认--异常
\fieldvalue{\checkboxfield{unchecked}}\\
\hline
\end{tabular}
\end{center}

\subsection*{1.2.\quad 试剂耗材}

\begin{center}
\begin{tabular}{|p{0.25\linewidth}|p{0.29\linewidth}|p{0.22\linewidth}|p{0.20\linewidth}|p{0.18\linewidth}|}
\hline
名称 & 厂家 & 货号 & 批号 & 有效期至\\
\hline
% #VALUE_ID: LEX-P0016-V0010
% #FIELD_VALUE: 试剂名称
\fieldvalue{细胞活力检测试剂盒} &
% #VALUE_ID: LEX-P0016-V0011
% #FIELD_VALUE: 试剂厂家
\fieldvalue{江苏凯基生物技术股份有限公司} &
% #VALUE_ID: LEX-P0016-V0012
% #FIELD_VALUE: 货号
\fieldvalue{KGA9501-1000} & & \\
\hline
\end{tabular}
\end{center}

% #VALUE_ID: LEX-P0016-V0013
% TODO #HANDWRITTEN: [蓝色斜向字迹不清]; 页面上的蓝色手写批注无法辨认
\handwritten{[蓝色斜向字迹不清]}

\section*{2.\quad 操作步骤：}

\begin{center}
\begin{tabular}{|p{0.55\linewidth}|p{0.40\linewidth}|}
\hline
\multicolumn{2}{|c|}{溶液配制}\\
\hline
20$\times$的荧光染液配制：取\underline{\hspace{1.2cm}}$\mu$l 组分A 和
\underline{\hspace{1.2cm}}$\mu$l 组分B 混合均匀配成20$\times$的荧光染液。 &
\\
\hline
取混匀后样品200$\mu$l/孔至96孔板中，取2孔。 &
样品\underline{\hspace{1.2cm}}$\mu$l/孔\\
& 每个样品取\underline{\hspace{1.2cm}}孔\\
\hline
自然沉降后，吸弃上清，注意不要吸到沉淀样品。 &
% #VALUE_ID: LEX-P0016-V0014
% #FIELD_VALUE: 吸弃上清
\fieldvalue{\checkboxfield{unchecked}} 吸弃上清\\
\hline
取20$\times$的荧光染液加入到96孔板中，200$\mu$l/孔，避光孵育
30--60min。 &
20$\times$的荧光染液\underline{\hspace{1.2cm}}$\mu$l/孔\\
& 避光孵育\underline{\hspace{0.7cm}}$\sim$\underline{\hspace{0.7cm}}\\
\hline
染色结束后，吸弃上清，取PBS加入到以96孔板中，200$\mu$l/孔，吸弃上清，注意不要吸到沉淀样品，清洗3次，之后每孔加入100$\mu$l PBS进行拍照； &
加入PBS\underline{\hspace{1.2cm}}$\mu$l/孔\\
& 清洗\underline{\hspace{1.2cm}}次\\
\hline
将荧光显微镜调至4$\times$物镜，软件中调节荧光相应参数，调节到适合的光源和像距，拍摄4张清晰图像。 &
\\
\hline
选一张清晰图像，打印并粘贴在实验记录上。 &
\\
\hline
\end{tabular}
\end{center}

\begin{tabular}{|p{0.25\linewidth}|p{0.70\linewidth}|}
\hline
数据保存路径 & \\
\hline
\multicolumn{2}{|c|}{检测结果}\\
\hline
\end{tabular}

\begin{center}
第4页\quad 共5页
\end{center}
% LEXOID_PAGE_COMPLETED: 16/70

\newpage

\noindent
\begin{minipage}{\linewidth}
\small
\begin{flushleft}
\textsf{铂生物科技（上海）有限公司}\\
\textsf{PLATINUM LIFE EXCELLENCE BIOTECH}
\end{flushleft}

\begin{center}
\begin{tabular}{|c|}
\hline
铂生物\\
% #VALUE_ID: LEX-P0017-V0001
% #FIELD_VALUE: 文件版本状态
\fieldvalue{原件}\\
\hline
\end{tabular}
\end{center}

\noindent
\textbf{文件编码：}
% #VALUE_ID: LEX-P0017-V0002
% #FIELD_VALUE: 文件编码
\fieldvalue{REC-YZ-SOP-QC-99-006-}

\hfill
\begin{tabular}{|c|}
\hline
铂生物\\
% #VALUE_ID: LEX-P0017-V0003
% #FIELD_VALUE: 更改实施编号
\fieldvalue{更改实施:4}\\
\hline
\end{tabular}

\vspace{0.2cm}

\noindent
\begin{tabular}{|p{0.92\linewidth}|}
\hline
\vspace{20.3cm}\\
\hline
\begin{tabular}{|p{0.16\linewidth}|p{0.76\linewidth}|}
\hline
备注 & \\
\hline
实验人/日期 & \\
\hline
复核人/日期 & \\
\hline
\end{tabular}\\
\hline
\end{tabular}

\vspace{-18.5cm}
\begin{flushright}
% #VALUE_ID: LEX-P0017-V0004
% #TODO #HANDWRITTEN: 检验 2024.01.11; handwriting is rotated and partially illegible
% #FIELD_VALUE: 斜线手写标记
\fieldvalue{\handwritten{检验 2024.01.11}}
\end{flushright}

\vspace{18.5cm}

\begin{center}
第 5 页 共 5 页
\end{center}
\end{minipage}
% LEXOID_PAGE_COMPLETED: 17/70

\newpage

\noindent\textbf{Assay:} 
% #VALUE_ID: LEX-P0018-V0001
% #FIELD_VALUE: Assay
\fieldvalue{MSC-1}

\medskip
\noindent\textbf{Cell Type F1:} 
% #VALUE_ID: LEX-P0018-V0002
% #FIELD_VALUE: Cell Type F1
\fieldvalue{MSC (AO)}

\noindent\textbf{Cell Type F2:} 
% #VALUE_ID: LEX-P0018-V0003
% #FIELD_VALUE: Cell Type F2
\fieldvalue{MSC (PI)}

\medskip
\noindent\textbf{Sample ID:} 
% #VALUE_ID: LEX-P0018-V0004
% #FIELD_VALUE: Sample ID
\fieldvalue{QC-JS-QY0401-A31Z2202508006-20250814-A01-0H-1}

\noindent\textbf{Dilution:} 
% #VALUE_ID: LEX-P0018-V0005
% #FIELD_VALUE: Dilution
\fieldvalue{2.00}

\medskip
\hrule
\medskip

\noindent\textbf{Results:}

\medskip
\noindent
\begin{tabular}{@{}p{0.29\linewidth}p{0.34\linewidth}p{0.29\linewidth}@{}}
\textbf{Count} & \textbf{Concentration} & \textbf{Mean Diameter} \\
\underline{\hspace{0.23\linewidth}} &
\underline{\hspace{0.28\linewidth}} &
\underline{\hspace{0.23\linewidth}} \\[-2pt]

Total:
% #VALUE_ID: LEX-P0018-V0006
% #FIELD_VALUE: Total count
\fieldvalue{88}
&
% #VALUE_ID: LEX-P0018-V0007
% #FIELD_VALUE: Total concentration
\fieldvalue{$3.10\times10^5$ cells/mL}
&
% #VALUE_ID: LEX-P0018-V0008
% #FIELD_VALUE: Total mean diameter
\fieldvalue{13.7 microns}
\\

Live:
% #VALUE_ID: LEX-P0018-V0009
% #FIELD_VALUE: Live count
\fieldvalue{79}
&
% #VALUE_ID: LEX-P0018-V0010
% #FIELD_VALUE: Live concentration
\fieldvalue{$2.78\times10^5$ cells/mL}
&
% #VALUE_ID: LEX-P0018-V0011
% #FIELD_VALUE: Live mean diameter
\fieldvalue{13.9 microns}
\\

Dead:
% #VALUE_ID: LEX-P0018-V0012
% #FIELD_VALUE: Dead count
\fieldvalue{9}
&
% #VALUE_ID: LEX-P0018-V0013
% #FIELD_VALUE: Dead concentration
\fieldvalue{$3.20\times10^4$ cells/mL}
&
% #VALUE_ID: LEX-P0018-V0014
% #FIELD_VALUE: Dead mean diameter
\fieldvalue{12.5 microns}
\\

Viability:
% #VALUE_ID: LEX-P0018-V0015
% #FIELD_VALUE: Viability
\fieldvalue{89.7\%}
&
&
\end{tabular}

\medskip

\noindent
\begin{minipage}[t]{0.47\linewidth}
\centering
% TODO: Insert the upper-left microscopy image visible on this page.
\fbox{\rule{0pt}{0.22\textheight}\rule{0.95\linewidth}{0pt}}
\end{minipage}
\hfill
\begin{minipage}[t]{0.47\linewidth}
\centering
% TODO: Insert the upper-right microscopy image visible on this page.
\fbox{\rule{0pt}{0.22\textheight}\rule{0.95\linewidth}{0pt}}
\end{minipage}

\medskip

\noindent
\begin{minipage}[t]{0.47\linewidth}
\centering
% TODO: Insert the lower-left microscopy image visible on this page.
\fbox{\rule{0pt}{0.22\textheight}\rule{0.95\linewidth}{0pt}}
\end{minipage}
\hfill
\begin{minipage}[t]{0.47\linewidth}
\centering
% TODO: Insert the lower-right microscopy image visible on this page.
\fbox{\rule{0pt}{0.22\textheight}\rule{0.95\linewidth}{0pt}}
\end{minipage}

\vfill

\noindent 操作人/审核人：
% #VALUE_ID: LEX-P0018-V0016
% #FIELD_VALUE: 操作人/审核人
% TODO #HANDWRITTEN: illegible; handwritten signature is not clearly readable
\fieldvalue{\handwritten{illegible}}
\hfill
日期：
% #VALUE_ID: LEX-P0018-V0017
% #FIELD_VALUE: 日期
% #HANDWRITTEN: 2025-8-14
\fieldvalue{\handwritten{2025-8-14}}
% LEXOID_PAGE_COMPLETED: 18/70

\newpage

\noindent
Assay: %
% #VALUE_ID: LEX-P0019-V0001
% #FIELD_VALUE: Assay
\fieldvalue{MSC-1}

\medskip
\noindent
Cell Type F1: %
% #VALUE_ID: LEX-P0019-V0002
% #FIELD_VALUE: Cell Type F1
\fieldvalue{MSC (AO)}

\noindent
Cell Type F2: %
% #VALUE_ID: LEX-P0019-V0003
% #FIELD_VALUE: Cell Type F2
\fieldvalue{MSC (PI)}

\medskip
\noindent
Sample ID: %
% #VALUE_ID: LEX-P0019-V0004
% #FIELD_VALUE: Sample ID
\fieldvalue{QC-JS-QY0401-A31Z2202508006-20250814-A01-0H-2}

\noindent
Dilution: %
% #VALUE_ID: LEX-P0019-V0005
% #FIELD_VALUE: Dilution
\fieldvalue{2.00}

\medskip
\noindent
\rule{\linewidth}{0.4pt}

\noindent
Results:

\medskip
\noindent
\begin{tabular}{@{}p{0.28\linewidth}p{0.34\linewidth}p{0.28\linewidth}@{}}
Count & Concentration & Mean Diameter \\
\texttt{============} & \texttt{=============} & \texttt{==============} \\[2pt]
Total:
% #VALUE_ID: LEX-P0019-V0006
% #FIELD_VALUE: Total count
\fieldvalue{88}
&
% #VALUE_ID: LEX-P0019-V0007
% #FIELD_VALUE: Total concentration
\fieldvalue{$3.11\times10^{5}\ \mathrm{cells/mL}$}
&
% #VALUE_ID: LEX-P0019-V0008
% #FIELD_VALUE: Total mean diameter
\fieldvalue{14.1 micron}
\\
Live:
% #VALUE_ID: LEX-P0019-V0009
% #FIELD_VALUE: Live count
\fieldvalue{77}
&
% #VALUE_ID: LEX-P0019-V0010
% #FIELD_VALUE: Live concentration
\fieldvalue{$2.72\times10^{5}\ \mathrm{cells/mL}$}
&
% #VALUE_ID: LEX-P0019-V0011
% #FIELD_VALUE: Live mean diameter
\fieldvalue{14.1 microns}
\\
Dead:
% #VALUE_ID: LEX-P0019-V0012
% #FIELD_VALUE: Dead count
\fieldvalue{11}
&
% #VALUE_ID: LEX-P0019-V0013
% #FIELD_VALUE: Dead concentration
\fieldvalue{$3.90\times10^{4}\ \mathrm{cells/mL}$}
&
% #VALUE_ID: LEX-P0019-V0014
% #FIELD_VALUE: Dead mean diameter
\fieldvalue{13.7 micron}
\\
Viability:
% #VALUE_ID: LEX-P0019-V0015
% #FIELD_VALUE: Viability
\fieldvalue{87.5\%}
&
&
\end{tabular}

\medskip

\noindent
\begin{tabular}{@{}c@{\hspace{0.04\linewidth}}c@{}}
\begin{minipage}[t]{0.42\linewidth}
\centering
% TODO: Insert the visible upper-left microscopy image; no filename is available.
\fbox{\rule{0pt}{0.25\linewidth}\rule{0.9\linewidth}{0pt}}
\end{minipage}
&
\begin{minipage}[t]{0.42\linewidth}
\centering
% TODO: Insert the visible upper-right dark microscopy image; no filename is available.
\fbox{\rule{0pt}{0.25\linewidth}\rule{0.9\linewidth}{0pt}}
\end{minipage}
\\[1em]
\begin{minipage}[t]{0.42\linewidth}
\centering
% TODO: Insert the visible lower-left microscopy image; no filename is available.
\fbox{\rule{0pt}{0.25\linewidth}\rule{0.9\linewidth}{0pt}}
\end{minipage}
&
\begin{minipage}[t]{0.42\linewidth}
\centering
% TODO: Insert the visible lower-right dark microscopy image; no filename is available.
\fbox{\rule{0pt}{0.25\linewidth}\rule{0.9\linewidth}{0pt}}
\end{minipage}
\end{tabular}

\vfill

\noindent
操作人/审核人：
% #VALUE_ID: LEX-P0019-V0016
% #FIELD_VALUE: 操作人/审核人
% #TODO #HANDWRITTEN: illegible; cursive handwritten signature
\fieldvalue{\handwritten{illegible}}
\hfill
日期：
% #VALUE_ID: LEX-P0019-V0017
% #FIELD_VALUE: 日期
% #TODO #HANDWRITTEN: 2024.3.4; final digit is somewhat unclear
\fieldvalue{\handwritten{2024.3.4}}
% LEXOID_PAGE_COMPLETED: 19/70

\newpage

\noindent
\textbf{Assay:}
% #VALUE_ID: LEX-P0020-V0001
% #FIELD_VALUE: Assay
\fieldvalue{MSC-1}

\vspace{0.6em}

\noindent
\textbf{Cell Type F1:}
% #VALUE_ID: LEX-P0020-V0002
% #FIELD_VALUE: Cell Type F1
\fieldvalue{MSC (AO)}

\noindent
\textbf{Cell Type F2:}
% #VALUE_ID: LEX-P0020-V0003
% #FIELD_VALUE: Cell Type F2
\fieldvalue{MSC (PI)}

\vspace{0.6em}

\noindent
\textbf{Sample ID:}
% #VALUE_ID: LEX-P0020-V0004
% #FIELD_VALUE: Sample ID
\fieldvalue{QC-JS-QY0401-A31Z2202508006-20250814-A01-0H-3}

\noindent
\textbf{Dilution:}
% #VALUE_ID: LEX-P0020-V0005
% #FIELD_VALUE: Dilution
\fieldvalue{2.00}

\vspace{0.8em}
\noindent\rule{\linewidth}{0.4pt}

\noindent\textbf{Results:}

\vspace{0.4em}

\renewcommand{\arraystretch}{1.15}
\begin{tabular}{@{}p{0.30\linewidth}p{0.34\linewidth}p{0.30\linewidth}@{}}
\textbf{Count} & \textbf{Concentration} & \textbf{Mean Diameter} \\
\texttt{============} & \texttt{=============} & \texttt{==============} \\[0.2em]

Total:
% #VALUE_ID: LEX-P0020-V0006
% #FIELD_VALUE: Count -- Total
\fieldvalue{77}
&
% #VALUE_ID: LEX-P0020-V0007
% #FIELD_VALUE: Concentration -- Total
\fieldvalue{$2.71\times10^{5}$ cells/mL}
&
% #VALUE_ID: LEX-P0020-V0008
% #FIELD_VALUE: Mean Diameter -- Total
\fieldvalue{13.8 microns}
\\

Live:
% #VALUE_ID: LEX-P0020-V0009
% #FIELD_VALUE: Count -- Live
\fieldvalue{67}
&
% #VALUE_ID: LEX-P0020-V0010
% #FIELD_VALUE: Concentration -- Live
\fieldvalue{$2.36\times10^{5}$ cells/mL}
&
% #VALUE_ID: LEX-P0020-V0011
% #FIELD_VALUE: Mean Diameter -- Live
\fieldvalue{13.9 microns}
\\

Dead:
% #VALUE_ID: LEX-P0020-V0012
% #FIELD_VALUE: Count -- Dead
\fieldvalue{10}
&
% #VALUE_ID: LEX-P0020-V0013
% #FIELD_VALUE: Concentration -- Dead
\fieldvalue{$3.50\times10^{4}$ cells/mL}
&
% #VALUE_ID: LEX-P0020-V0014
% #FIELD_VALUE: Mean Diameter -- Dead
\fieldvalue{13.2 microns}
\\

Viability:
% #VALUE_ID: LEX-P0020-V0015
% #FIELD_VALUE: Viability
\fieldvalue{87.1\%}
&
&
\end{tabular}

\vspace{1.4em}

\noindent
% TODO: Image visible on page; source filename unavailable.
\fbox{\parbox[c][4.1cm][c]{0.43\linewidth}{\centering
\textit{Image placeholder}}}
\hfill
% TODO: Image visible on page; source filename unavailable.
\fbox{\parbox[c][4.1cm][c]{0.43\linewidth}{\centering
\textit{Image placeholder}}}

\vspace{0.8em}

\noindent
% TODO: Image visible on page; source filename unavailable.
\fbox{\parbox[c][4.1cm][c]{0.43\linewidth}{\centering
\textit{Image placeholder}}}
\hfill
% TODO: Image visible on page; source filename unavailable.
\fbox{\parbox[c][4.1cm][c]{0.43\linewidth}{\centering
\textit{Image placeholder}}}

\vfill

\noindent
操作人/审核人：
% #VALUE_ID: LEX-P0020-V0016
% #FIELD_VALUE: 操作人/审核人
% #TODO #HANDWRITTEN: illegible; cursive handwritten signature
\fieldvalue{\handwritten{[illegible]}}

\hfill
日期：
% #VALUE_ID: LEX-P0020-V0017
% #FIELD_VALUE: 日期
% #HANDWRITTEN: 2024.8.14
\fieldvalue{\handwritten{2024.8.14}}
% LEXOID_PAGE_COMPLETED: 20/70

\newpage

\begin{center}
\textbf{铂金生物技术（北京）有限公司}\\
\textit{PLATINUM LIFE SCIENCE BIOTECH}\\[2pt]
\textbf{原件请验单}
\end{center}

\noindent
文件编码：REC-YZ-SMP-QC-01-001-d\hfill 版本号：1.4

\vspace{4pt}

\begin{center}
\small
\begin{tabular}{|p{2.2cm}|p{5.0cm}|p{2.2cm}|p{5.0cm}|}
\hline
样品名称 &
% #VALUE_ID: LEX-P0021-V0001
% #FIELD_VALUE: 样品名称
% #HANDWRITTEN: 二氯甲烷色谱纯
\fieldvalue{\handwritten{二氯甲烷色谱纯}} &
样品代码 &
% #VALUE_ID: LEX-P0021-V0002
% #FIELD_VALUE: 样品代码
% TODO #HANDWRITTEN: Q?Y?02; handwriting is unclear
\fieldvalue{\handwritten{Q?Y?02}} \\
\hline
样品批号 &
% #VALUE_ID: LEX-P0021-V0003
% #FIELD_VALUE: 样品批号
% TODO #HANDWRITTEN: A?20?; handwriting is unclear
\fieldvalue{\handwritten{A?20?}} &
样品规格/数量 &
% #VALUE_ID: LEX-P0021-V0004
% #FIELD_VALUE: 样品规格/数量
% TODO #HANDWRITTEN: approximately 21 mL/bottle, 1 bottle; handwriting is unclear
\fieldvalue{\handwritten{约21\,mL/瓶，1瓶}} \\
\hline
样品类别 &
% #VALUE_ID: LEX-P0021-V0005
% #FIELD_VALUE: 生产过程检验品
生产过程检验品\ \fieldvalue{\checkboxfield{checked}}\quad
% #VALUE_ID: LEX-P0021-V0006
% #FIELD_VALUE: 代包装产品
代包装产品\ \fieldvalue{\checkboxfield{unchecked}}\quad
% #VALUE_ID: LEX-P0021-V0007
% #FIELD_VALUE: 成品
成品\ \fieldvalue{\checkboxfield{unchecked}}\quad
% #VALUE_ID: LEX-P0021-V0008
% #FIELD_VALUE: 其他
其他（\fieldvalue{\checkboxfield{unchecked}} \hspace{1cm}） &
\multicolumn{2}{l|}{} \\
\hline
检验类型 &
% #VALUE_ID: LEX-P0021-V0009
% #FIELD_VALUE: 初验
初验\ \fieldvalue{\checkboxfield{checked}}\quad
% #VALUE_ID: LEX-P0021-V0010
% #FIELD_VALUE: 复验
复验\ \fieldvalue{\checkboxfield{unchecked}}\quad
% #VALUE_ID: LEX-P0021-V0011
% #FIELD_VALUE: 退货
退货\ \fieldvalue{\checkboxfield{unchecked}}\quad
% #VALUE_ID: LEX-P0021-V0012
% #FIELD_VALUE: 稳定性
稳定性\ \fieldvalue{\checkboxfield{unchecked}}\quad
% #VALUE_ID: LEX-P0021-V0013
% #FIELD_VALUE: 留样
留样\ \fieldvalue{\checkboxfield{unchecked}} &
\multicolumn{2}{l|}{} \\
\hline
请验目的 &
% #VALUE_ID: LEX-P0021-V0014
% #FIELD_VALUE: 全检
全检\ \fieldvalue{\checkboxfield{checked}}\quad
% #VALUE_ID: LEX-P0021-V0015
% #FIELD_VALUE: 留样
留样\ \fieldvalue{\checkboxfield{unchecked}}\quad
% #VALUE_ID: LEX-P0021-V0016
% #FIELD_VALUE: 其他
其他（\fieldvalue{\checkboxfield{unchecked}} \hspace{1cm}） &
\multicolumn{2}{l|}{} \\
\hline
请验部门 &
% #VALUE_ID: LEX-P0021-V0017
% #FIELD_VALUE: 生产部
生产部\ \fieldvalue{\checkboxfield{checked}}\quad
% #VALUE_ID: LEX-P0021-V0018
% #FIELD_VALUE: 物料管理部
物料管理部\ \fieldvalue{\checkboxfield{unchecked}}\quad
% #VALUE_ID: LEX-P0021-V0019
% #FIELD_VALUE: 其他
其他（\fieldvalue{\checkboxfield{unchecked}} \hspace{1cm}） &
\multicolumn{2}{l|}{} \\
\hline
样品情况 &
% #VALUE_ID: LEX-P0021-V0020
% #FIELD_VALUE: 随箱请验单一起提供
随箱请验单一起提供\ \fieldvalue{\checkboxfield{checked}}\quad
% #VALUE_ID: LEX-P0021-V0021
% #FIELD_VALUE: 未随箱请验单一起提供
未随箱请验单一起提供\ \fieldvalue{\checkboxfield{unchecked}} &
\multicolumn{2}{l|}{} \\
\hline
请验人 &
% #VALUE_ID: LEX-P0021-V0022
% #FIELD_VALUE: 请验人
% TODO #HANDWRITTEN: 李春玲; handwriting is unclear
\fieldvalue{\handwritten{李春玲}} &
请验日期 &
% #VALUE_ID: LEX-P0021-V0023
% #FIELD_VALUE: 请验日期
% #HANDWRITTEN: 2024.08.9
\fieldvalue{\handwritten{2024.08.9}} \\
\hline
收验人 &
% #VALUE_ID: LEX-P0021-V0024
% #FIELD_VALUE: 收验人
% TODO #HANDWRITTEN: 赵长?; handwriting is unclear
\fieldvalue{\handwritten{赵长?}} &
收验日期 &
% #VALUE_ID: LEX-P0021-V0025
% #FIELD_VALUE: 收验日期
% #HANDWRITTEN: 2024.08.19
\fieldvalue{\handwritten{2024.08.19}} \\
\hline
备注 &
% #VALUE_ID: LEX-P0021-V0026
% #FIELD_VALUE: 备注
% #HANDWRITTEN: 48h，13.6L
\fieldvalue{\handwritten{48h，13.6L}} &
\multicolumn{2}{l|}{} \\
\hline
\end{tabular}
\end{center}
% LEXOID_PAGE_COMPLETED: 21/70

\newpage

% TODO: 页面左上角包含公司标志，原图未提供可用图像文件。
\begin{center}
\textbf{取样记录}
\end{center}

\noindent
\begin{tabular}{|p{0.13\linewidth}|p{0.22\linewidth}|p{0.13\linewidth}|p{0.20\linewidth}|p{0.13\linewidth}|p{0.17\linewidth}|}
\hline
名称 &
% #VALUE_ID: LEX-P0022-V0001
% #FIELD_VALUE: 名称
\fieldvalue{二级生物反应器细胞液} &
代码 &
% #VALUE_ID: LEX-P0022-V0002
% #FIELD_VALUE: 代码
\fieldvalue{QY0402} &
\multicolumn{2}{c|}{} \\
\hline
规格 &
% #VALUE_ID: LEX-P0022-V0003
% #FIELD_VALUE: 规格
\fieldvalue{20ml/袋} &
总数量 &
% #VALUE_ID: LEX-P0022-V0004
% #FIELD_VALUE: 总数量
\fieldvalue{1袋} &
\multicolumn{2}{c|}{} \\
\hline
批号 &
% #VALUE_ID: LEX-P0022-V0005
% #FIELD_VALUE: 批号
\fieldvalue{A31Z2202508006} &
进厂批号 &
% #VALUE_ID: LEX-P0022-V0006
% #FIELD_VALUE: 进厂批号
\fieldvalue{/} &
\multicolumn{2}{c|}{} \\
\hline
供应单位/厂家 &
\multicolumn{5}{c|}{
% #VALUE_ID: LEX-P0022-V0007
% #FIELD_VALUE: 供应单位/厂家
\fieldvalue{/}} \\
\hline
取样目的 &
\multicolumn{5}{c|}{
% #VALUE_ID: LEX-P0022-V0008
% #FIELD_VALUE: 取样目的--正常检验
\fieldvalue{\checkboxfield{checked}} 正常检验
\quad
% #VALUE_ID: LEX-P0022-V0009
% #FIELD_VALUE: 取样目的--复检
\fieldvalue{\checkboxfield{unchecked}} 复检
\quad
% #VALUE_ID: LEX-P0022-V0010
% #FIELD_VALUE: 取样目的--稳定性
\fieldvalue{\checkboxfield{unchecked}} 稳定性
\newline
% #VALUE_ID: LEX-P0022-V0011
% #FIELD_VALUE: 取样目的--追加检验
\fieldvalue{\checkboxfield{unchecked}} 追加检验
\quad
% #VALUE_ID: LEX-P0022-V0012
% #FIELD_VALUE: 取样目的--留样
\fieldvalue{\checkboxfield{unchecked}} 留样
\quad
% #VALUE_ID: LEX-P0022-V0013
% #FIELD_VALUE: 取样目的--其他
\fieldvalue{\checkboxfield{unchecked}} 其他：\underline{\hspace{2cm}}} \\
\hline
取样地点 &
\multicolumn{5}{c|}{
% #VALUE_ID: LEX-P0022-V0014
% #FIELD_VALUE: 取样地点--产品出厂
\fieldvalue{\checkboxfield{checked}} 产品出厂
\quad
% #VALUE_ID: LEX-P0022-V0015
% #FIELD_VALUE: 取样地点--物料库
\fieldvalue{\checkboxfield{unchecked}} 物料库
\quad
% #VALUE_ID: LEX-P0022-V0016
% #FIELD_VALUE: 取样地点--细胞库
\fieldvalue{\checkboxfield{unchecked}} 细胞库
\quad
% #VALUE_ID: LEX-P0022-V0017
% #FIELD_VALUE: 取样地点--包装间
\fieldvalue{\checkboxfield{unchecked}} 包装间} \\
\hline
贮存条件 &
\multicolumn{5}{c|}{
% #VALUE_ID: LEX-P0022-V0018
% #FIELD_VALUE: 贮存条件--2-8℃
\fieldvalue{\checkboxfield{checked}} 2--8℃
\quad
% #VALUE_ID: LEX-P0022-V0019
% #FIELD_VALUE: 贮存条件--20℃
\fieldvalue{\checkboxfield{unchecked}} -20℃
\quad
% #VALUE_ID: LEX-P0022-V0020
% #FIELD_VALUE: 贮存条件--80℃
\fieldvalue{\checkboxfield{unchecked}} -80℃
\quad
% #VALUE_ID: LEX-P0022-V0021
% #FIELD_VALUE: 贮存条件--常温
\fieldvalue{\checkboxfield{unchecked}} 常温
\quad
% #VALUE_ID: LEX-P0022-V0022
% #FIELD_VALUE: 贮存条件--其他
\fieldvalue{\checkboxfield{unchecked}} 其他：\underline{\hspace{1.5cm}}} \\
\hline
盛装容器 &
\multicolumn{5}{c|}{
% #VALUE_ID: LEX-P0022-V0023
% #FIELD_VALUE: 盛装容器--保温箱
\fieldvalue{\checkboxfield{checked}} 保温箱
\quad
% #VALUE_ID: LEX-P0022-V0024
% #FIELD_VALUE: 盛装容器--泡沫箱
\fieldvalue{\checkboxfield{unchecked}} 泡沫箱
\quad
% #VALUE_ID: LEX-P0022-V0025
% #FIELD_VALUE: 盛装容器--手提篮
\fieldvalue{\checkboxfield{unchecked}} 手提篮
\quad
% #VALUE_ID: LEX-P0022-V0026
% #FIELD_VALUE: 盛装容器--其他
\fieldvalue{\checkboxfield{unchecked}} 其他：\underline{\hspace{1.5cm}}} \\
\hline
外观复核 &
\multicolumn{5}{c|}{
% #VALUE_ID: LEX-P0022-V0027
% #FIELD_VALUE: 外观复核--合格
\fieldvalue{\checkboxfield{checked}} 合格
\quad
% #VALUE_ID: LEX-P0022-V0028
% #FIELD_VALUE: 外观复核--不合格
\fieldvalue{\checkboxfield{unchecked}} 不合格} \\
\hline
取样量 &
% #VALUE_ID: LEX-P0022-V0029
% #FIELD_VALUE: 取样量
\fieldvalue{1袋} &
检验量 &
% #VALUE_ID: LEX-P0022-V0030
% #FIELD_VALUE: 检验量
\fieldvalue{1袋} &
留样量 &
% #VALUE_ID: LEX-P0022-V0031
% #FIELD_VALUE: 留样量
\fieldvalue{/} \\
\hline
取样人/日期 &
\multicolumn{2}{c|}{
% #VALUE_ID: LEX-P0022-V0032
% #FIELD_VALUE: 取样人
% #TODO #HANDWRITTEN: 杨??; 手写姓名辨识不清
\fieldvalue{\handwritten{杨??}}
\quad
% #VALUE_ID: LEX-P0022-V0033
% #FIELD_VALUE: 取样日期
% #HANDWRITTEN: 2025.8.9
\fieldvalue{\handwritten{2025.8.9}}} &
复核人/日期 &
\multicolumn{2}{c|}{
% #VALUE_ID: LEX-P0022-V0034
% #FIELD_VALUE: 复核人
% #TODO #HANDWRITTEN: 李??; 手写姓名辨识不清
\fieldvalue{\handwritten{李??}}
\quad
% #VALUE_ID: LEX-P0022-V0035
% #FIELD_VALUE: 复核日期
% #HANDWRITTEN: 2025.8.17
\fieldvalue{\handwritten{2025.8.17}}} \\
\hline
备注 &
\multicolumn{5}{c|}{
% #VALUE_ID: LEX-P0022-V0036
% #FIELD_VALUE: 备注
% #HANDWRITTEN: 48H
\fieldvalue{\handwritten{48H}}
\quad
% #VALUE_ID: LEX-P0022-V0037
% #FIELD_VALUE: 备注
% #TODO #HANDWRITTEN: 13?L; 手写单位及数值辨识不清
\fieldvalue{\handwritten{13?L}}} \\
\hline
\end{tabular}
% LEXOID_PAGE_COMPLETED: 22/70

\newpage

\begin{center}
{\large 生物反应器细胞液检测操作记录}
\end{center}

\noindent
文件编号：REC-YZ-SOP-QC-99-006-a

\begin{center}
\begin{tabular}{|p{2.5cm}|p{5.0cm}|p{2.5cm}|p{4.0cm}|}
\hline
\multicolumn{4}{|c|}{葡萄糖浓度检验} \\
\hline
样品名称 &
% #VALUE_ID: LEX-P0023-V0001
% #FIELD_VALUE: 样品名称
\fieldvalue{二线生物反应器细胞液} &
样品规格 &
% #VALUE_ID: LEX-P0023-V0002
% #FIELD_VALUE: 样品规格
% #HANDWRITTEN: 20 ml/袋
\fieldvalue{\handwritten{20 ml/袋}} \\
\hline
样品批号 &
% #VALUE_ID: LEX-P0023-V0003
% #FIELD_VALUE: 样品批号
% #TODO #HANDWRITTEN: 211?..., handwriting unclear
\fieldvalue{\handwritten{211?...}} &
取样日期 &
% #VALUE_ID: LEX-P0023-V0004
% #FIELD_VALUE: 取样日期
% #HANDWRITTEN: 2025年03月19日
\fieldvalue{\handwritten{2025年03月19日}} \\
\hline
\end{tabular}
\end{center}

\section*{1\quad 仪器设备、试剂耗材}

\subsection*{1.1\quad 仪器设备}

\begin{center}
\begin{tabular}{|p{2.0cm}|p{4.0cm}|p{2.0cm}|p{3.0cm}|p{3.5cm}|}
\hline
名称 & 厂家 & 型号 & 设备编号 & 检查确认 \\
\hline
血糖仪 &
% #VALUE_ID: LEX-P0023-V0005
% #FIELD_VALUE: 血糖仪厂家
\fieldvalue{三诺生物传感股份有限公司} &
% #VALUE_ID: LEX-P0023-V0006
% #FIELD_VALUE: 血糖仪型号
\fieldvalue{安稳+} &
 &
正常\quad
% #VALUE_ID: LEX-P0023-V0007
% #FIELD_VALUE: 血糖仪检查确认--正常
\fieldvalue{\checkboxfield{checked}}
\quad 异常\quad
% #VALUE_ID: LEX-P0023-V0008
% #FIELD_VALUE: 血糖仪检查确认--异常
\fieldvalue{\checkboxfield{unchecked}} \\
\hline
移液器 &
% #VALUE_ID: LEX-P0023-V0009
% #FIELD_VALUE: 移液器厂家
\fieldvalue{Eppendorf Corporation} &
% #VALUE_ID: LEX-P0023-V0010
% #FIELD_VALUE: 移液器量程
\fieldvalue{0.1--2.5\,\textmu l} &
% #VALUE_ID: LEX-P0023-V0011
% #FIELD_VALUE: 移液器设备编号
% #TODO #HANDWRITTEN: B?7?4?6L; handwriting unclear
\fieldvalue{\handwritten{B?7?4?6L}} &
正常\quad
% #VALUE_ID: LEX-P0023-V0012
% #FIELD_VALUE: 移液器检查确认--正常
\fieldvalue{\checkboxfield{checked}}
\quad 异常\quad
% #VALUE_ID: LEX-P0023-V0013
% #FIELD_VALUE: 移液器检查确认--异常
\fieldvalue{\checkboxfield{unchecked}} \\
\hline
\end{tabular}
\end{center}

\subsection*{1.2\quad 试剂耗材}

\begin{center}
\begin{tabular}{|p{3.0cm}|p{5.5cm}|p{3.0cm}|p{3.5cm}|}
\hline
名称 & 厂家 & 批号 & 有效期至 \\
\hline
血糖测试条 &
% #VALUE_ID: LEX-P0023-V0014
% #FIELD_VALUE: 血糖测试条厂家
\fieldvalue{三诺生物传感股份有限公司} &
% #VALUE_ID: LEX-P0023-V0015
% #FIELD_VALUE: 血糖测试条批号
% #TODO #HANDWRITTEN: 444?4?0-1; handwriting unclear
\fieldvalue{\handwritten{444?4?0-1}} &
% #VALUE_ID: LEX-P0023-V0016
% #FIELD_VALUE: 血糖测试条有效期
% #HANDWRITTEN: 2027.01.08
\fieldvalue{\handwritten{2027.01.08}} \\
\hline
\end{tabular}
\end{center}

\section*{2\quad 操作步骤：}

\begin{center}
\begin{tabular}{|p{7.2cm}|p{6.8cm}|}
\hline
\multicolumn{2}{|c|}{操作过程} \\
\hline
取出测试条，取样品自然沉降后的上清液 $0.6\,\mu l$ 加入测试条电极反应区，血糖仪自动发出“嘀”提示音，吸入样品完成，血糖仪自动进行测试，开始测试。 &
取上清液：
% #VALUE_ID: LEX-P0023-V0017
% #FIELD_VALUE: 取上清液体积
% #HANDWRITTEN: 0.6
\fieldvalue{\handwritten{0.6}}\,\ensuremath{\mu l}
\\
\hline
重复以上步骤 $2$ 次，共计数 $3$ 次。 &
每个样品分别共计数
% #VALUE_ID: LEX-P0023-V0018
% #FIELD_VALUE: 每个样品分别共计数
% #HANDWRITTEN: 3
\fieldvalue{\handwritten{3}} 次。 \\
\hline
\multicolumn{2}{|c|}{检测结果} \\
\hline
\multicolumn{2}{|c|}{葡萄糖浓度（mmol/L）} \\
\hline
第一次 & 第二次 \quad\quad 第三次 \quad\quad 平均值（保留一位小数） \\
\hline
% #VALUE_ID: LEX-P0023-V0019
% #FIELD_VALUE: 第一次葡萄糖浓度
% #HANDWRITTEN: 19.7
\fieldvalue{\handwritten{19.7}} &
% #VALUE_ID: LEX-P0023-V0020
% #FIELD_VALUE: 第二次葡萄糖浓度
% #HANDWRITTEN: 20.1
\fieldvalue{\handwritten{20.1}}
\quad\quad
% #VALUE_ID: LEX-P0023-V0021
% #FIELD_VALUE: 第三次葡萄糖浓度
% #HANDWRITTEN: 20.2
\fieldvalue{\handwritten{20.2}}
\quad\quad
% #VALUE_ID: LEX-P0023-V0022
% #FIELD_VALUE: 平均葡萄糖浓度
% #HANDWRITTEN: 20.0
\fieldvalue{\handwritten{20.0}} \\
\hline
备注 &
% #VALUE_ID: LEX-P0023-V0023
% #FIELD_VALUE: 备注
% #HANDWRITTEN: ✓
\fieldvalue{\handwritten{\checkmark}} \\
\hline
实验人/日期 &
% #VALUE_ID: LEX-P0023-V0024
% #FIELD_VALUE: 实验人
% #TODO #HANDWRITTEN: 郭?; handwriting unclear
\fieldvalue{\handwritten{郭?}}
\quad
% #VALUE_ID: LEX-P0023-V0025
% #FIELD_VALUE: 实验日期
% #HANDWRITTEN: 2025.03.19
\fieldvalue{\handwritten{2025.03.19}} \\
\hline
复核人/日期 &
% #VALUE_ID: LEX-P0023-V0026
% #FIELD_VALUE: 复核人
% #TODO #HANDWRITTEN: 张?; handwriting unclear
\fieldvalue{\handwritten{张?}}
\quad
% #VALUE_ID: LEX-P0023-V0027
% #FIELD_VALUE: 复核日期
% #HANDWRITTEN: 2025.03.19
\fieldvalue{\handwritten{2025.03.19}} \\
\hline
\end{tabular}
\end{center}

\begin{center}
第 1 页 共 5 页
\end{center}
% LEXOID_PAGE_COMPLETED: 23/70

\newpage

\begin{center}
\textbf{生物反应器细胞液检测记录}
\end{center}

\noindent
文件编号：REC-YZ-SOP-QC-99-006-a

\begin{center}
\textbf{密度与活率检验}
\end{center}

\begin{tabular}{|p{0.16\linewidth}|p{0.34\linewidth}|p{0.16\linewidth}|p{0.24\linewidth}|}
\hline
样品名称 &
% #VALUE_ID: LEX-P0024-V0001
% #FIELD_VALUE: 样品名称
\fieldvalue{\handwritten{二}}级生物反应器细胞液 &
样品规格 &
% #VALUE_ID: LEX-P0024-V0002
% #FIELD_VALUE: 样品规格
\fieldvalue{\handwritten{20}} ml/袋
\\ \hline
样品批号 &
% #VALUE_ID: LEX-P0024-V0003
% #FIELD_VALUE: 样品批号
% #TODO #HANDWRITTEN: A?230?/300?6; handwriting unclear
\fieldvalue{\handwritten{A?230?/300?6}} &
取样日期 &
% #VALUE_ID: LEX-P0024-V0004
% #FIELD_VALUE: 取样日期
\fieldvalue{\handwritten{2024年5月?日}}
\\ \hline
\end{tabular}

\section*{1\quad 仪器设备、试剂耗材}

\subsection*{1.1\quad 仪器设备}

\begin{tabular}{|p{0.18\linewidth}|p{0.17\linewidth}|p{0.18\linewidth}|p{0.15\linewidth}|p{0.17\linewidth}|p{0.18\linewidth}|}
\hline
名称 & 厂家 & 型号 & 设备编号 & 校准有效期至 & 检查确认 \\ \hline
细胞计数仪 & Nexcelom & Cellometer K2 & 1410905 &
% #VALUE_ID: LEX-P0024-V0005
% #FIELD_VALUE: 细胞计数仪校准有效期至
\fieldvalue{\handwritten{2024.11.19}} &
正常\ 
% #VALUE_ID: LEX-P0024-V0006
% #FIELD_VALUE: 检查确认--正常
% #HANDWRITTEN: ✓
\fieldvalue{\handwritten{\checkboxfield{checked}}}
\quad 异常\ 
% #VALUE_ID: LEX-P0024-V0007
% #FIELD_VALUE: 检查确认--异常
\fieldvalue{\checkboxfield{unchecked}}
\\ \hline
高速冷冻离心机 & Thermo & ST16R & 1410810 &
% #VALUE_ID: LEX-P0024-V0008
% #FIELD_VALUE: 高速冷冻离心机校准有效期至
% #TODO #HANDWRITTEN: 2024.10.19?; final digits unclear
\fieldvalue{\handwritten{2024.10.19?}} &
正常\ 
% #VALUE_ID: LEX-P0024-V0009
% #FIELD_VALUE: 检查确认--正常
% #HANDWRITTEN: ✓
\fieldvalue{\handwritten{\checkboxfield{checked}}}
\quad 异常\ 
% #VALUE_ID: LEX-P0024-V0010
% #FIELD_VALUE: 检查确认--异常
\fieldvalue{\checkboxfield{unchecked}}
\\ \hline
电热恒温水槽 & 上海一恒 & DK-8AXX & 1411001 &
% #VALUE_ID: LEX-P0024-V0011
% #FIELD_VALUE: 电热恒温水槽校准有效期至
\fieldvalue{\handwritten{2024.07.11}} &
正常\ 
% #VALUE_ID: LEX-P0024-V0012
% #FIELD_VALUE: 检查确认--正常
% #HANDWRITTEN: ✓
\fieldvalue{\handwritten{\checkboxfield{checked}}}
\quad 异常\ 
% #VALUE_ID: LEX-P0024-V0013
% #FIELD_VALUE: 检查确认--异常
\fieldvalue{\checkboxfield{unchecked}}
\\ \hline
\end{tabular}

\subsection*{1.2\quad 试剂耗材}

\begin{tabular}{|p{0.22\linewidth}|p{0.25\linewidth}|p{0.27\linewidth}|p{0.22\linewidth}|}
\hline
名称 & 厂家 & 批号 & 有效期至 \\ \hline
AO/PI 染液 & Nexcelom &
% #VALUE_ID: LEX-P0024-V0014
% #FIELD_VALUE: AO/PI染液批号
% #TODO #HANDWRITTEN: 3?4?9?; handwriting unclear
\fieldvalue{\handwritten{3?4?9?}} &
% #VALUE_ID: LEX-P0024-V0015
% #FIELD_VALUE: AO/PI染液有效期至
\fieldvalue{\handwritten{2024.?.??}}
\\ \hline
0.9\%氯化钠注射液 & 石家庄四药 &
% #VALUE_ID: LEX-P0024-V0016
% #FIELD_VALUE: 0.9%氯化钠注射液批号
% #TODO #HANDWRITTEN: 24?5?730?; handwriting unclear
\fieldvalue{\handwritten{24?5?730?}} &
% #VALUE_ID: LEX-P0024-V0017
% #FIELD_VALUE: 0.9%氯化钠注射液有效期至
\fieldvalue{\handwritten{2024.10.06}}
\\ \hline
20\%人白蛋白 & Baxter AG &
% #VALUE_ID: LEX-P0024-V0018
% #FIELD_VALUE: 20%人白蛋白批号
% #TODO #HANDWRITTEN: unreadable; handwriting unclear
\fieldvalue{\handwritten{???}} &
% #VALUE_ID: LEX-P0024-V0019
% #FIELD_VALUE: 20%人白蛋白有效期至
\fieldvalue{\handwritten{2024.?.31}}
\\ \hline
\end{tabular}

\section*{2\quad 操作步骤}

\subsection*{溶液准备}

\noindent
1\%人血白蛋白溶液配制：取5 ml 20\%人血白蛋白加入到95 ml
的0.9\%氯化钠注射液中，混匀即得。

\begin{tabular}{|p{0.25\linewidth}|p{0.38\linewidth}|p{0.30\linewidth}|}
\hline
名称 & 配制批号 & 有效期至 \\ \hline
5$\times$裂解液 &
% #VALUE_ID: LEX-P0024-V0020
% #FIELD_VALUE: 5倍裂解液配制批号
% #TODO #HANDWRITTEN: 5?0?d?750?; handwriting unclear
\fieldvalue{\handwritten{5?0?d?750?}} &
% #VALUE_ID: LEX-P0024-V0021
% #FIELD_VALUE: 5倍裂解液有效期至
\fieldvalue{\handwritten{2024.?.??}}
\\ \hline
\end{tabular}

\subsection*{操作过程}

\begin{tabular}{|p{0.57\linewidth}|p{0.37\linewidth}|}
\hline
\textbf{操作步骤} & \textbf{样品记录} \\ \hline
按照50 ml离心管的刻度读取样品体积，在取样袋中加入对应体积的5$\times$裂解液（每4 ml样品加入1 ml 5$\times$裂解液），按照“只进不舍”原则，如26 ml样品，加入7 ml 5$\times$裂解液，充分混匀后倒入至2.3的样品中，置于37℃电热恒温水槽中裂解20 min后，取出离心管轻混匀，使样品充分裂解；再置于37℃电热恒温水槽中裂解10 min得到细胞悬液。 &
样品体积：
% #VALUE_ID: LEX-P0024-V0022
% #FIELD_VALUE: 样品体积
\fieldvalue{\handwritten{32}} ml

裂解液：
% #VALUE_ID: LEX-P0024-V0023
% #FIELD_VALUE: 裂解液用量
\fieldvalue{\handwritten{6}} ml

电热恒温水槽：37℃

裂解时间：
% #VALUE_ID: LEX-P0024-V0024
% #FIELD_VALUE: 裂解时间
% #TODO #HANDWRITTEN: 2?4; handwriting unclear
\fieldvalue{\handwritten{2?4}} \quad
% #VALUE_ID: LEX-P0024-V0025
% #FIELD_VALUE: 裂解记录时间
\fieldvalue{\handwritten{11:54}}
\\ \hline
取出裂解好的细胞悬液300$\times g$离心5分钟，按照读取样品体积加入对应体积的1\%人白蛋白稀释液（每1.2 ml样品体积加入0.1 ml 1\%人白蛋白稀释液），按照“只进不舍”原则，如26 ml样品，加入2.2 ml 1\%人白蛋白稀释液。 &
% #VALUE_ID: LEX-P0024-V0026
% #FIELD_VALUE: 离心力
\fieldvalue{\handwritten{300}}$\times g$离心
% #VALUE_ID: LEX-P0024-V0027
% #FIELD_VALUE: 离心时间
\fieldvalue{\handwritten{5}} 分钟

加入1\%人白蛋白稀释液总体积：
% #VALUE_ID: LEX-P0024-V0028
% #FIELD_VALUE: 1%人白蛋白稀释液总体积
\fieldvalue{\handwritten{2.7}} ml

吸取
% #VALUE_ID: LEX-P0024-V0029
% #FIELD_VALUE: 吸取体积
% #TODO #HANDWRITTEN: 20?; handwriting unclear
\fieldvalue{\handwritten{20?}} $\mu$l
\\ \hline
\end{tabular}

\begin{center}
第 2 页 共 5 页
\end{center}
% LEXOID_PAGE_COMPLETED: 24/70

\newpage

\begin{flushleft}
\textbf{细胞生物学检测（北京）有限公司}\\
文件编号：RECZ-SY-OP-CC-99-006-a
\end{flushleft}

\noindent
加人 $2.2\mathrm{ml}$ 的人白蛋白稀释液，反复吹打至单细胞悬液。\\
撕去计数板上下两层保护膜，将细胞悬液混匀，从中吸取 $20\mu l$ 与 $20\mu l$ AO/PI 染料混合均匀。\\
吸取 $20\mu l$ 细胞悬液与 AO/PI 染料的混合液，注入到计数板的一 个检测室中。\\
将加入混合液的一端插入进板口。\\
选择 Assay 类型，使用 AO/PI 染料进行双荧光存活率检测；在主界面左侧 ID 栏中输入样本批号。\\
点击主界面左下角“Preview”预览细胞图片。\\
调整仪器焦距直至观察到荧光更清晰下的细胞中心透亮，轮廓清晰。\\
点击“count”按钮进行计数并自动拍摄图片。\\
重复以上步骤 $2$ 次，共计数 $3$ 次。拍摄数据结果。

\begin{flushright}
\begin{tabular}{@{}l@{}}
吸取
% #VALUE_ID: LEX-P0025-V0001
% #FIELD_VALUE: AO/PI染料混合液吸取体积
% #HANDWRITTEN: 20
\fieldvalue{\handwritten{20}}$\mu l$ AO/PI 染料混合均匀。\\
吸取
% #VALUE_ID: LEX-P0025-V0002
% #FIELD_VALUE: 细胞悬液吸取体积
% #HANDWRITTEN: 20
\fieldvalue{\handwritten{20}}$\mu l$ 计数。\\
每个样品共计数
% #VALUE_ID: LEX-P0025-V0003
% #FIELD_VALUE: 样品计数次数
% #HANDWRITTEN: 2
\fieldvalue{\handwritten{2}}次。
\end{tabular}
\end{flushright}

\noindent
数据保存路径：
% #VALUE_ID: LEX-P0025-V0007
% #FIELD_VALUE: 数据保存路径
% TODO #HANDWRITTEN: 4/19...；手写路径部分字迹不清
\fieldvalue{\handwritten{4/19\ldots}}

\medskip

\begin{center}
\textbf{检测结果}
\end{center}

\noindent
检测结果：
% #VALUE_ID: LEX-P0025-V0004
% #FIELD_VALUE: 一级反应峰
\fieldvalue{\checkboxfield{unchecked}} 一级反应峰；
% #VALUE_ID: LEX-P0025-V0005
% #FIELD_VALUE: 二级反应峰
% #HANDWRITTEN: ✓
\fieldvalue{\handwritten{\checkboxfield{checked}}} 二级反应峰；
% #VALUE_ID: LEX-P0025-V0006
% #FIELD_VALUE: 三级反应峰
\fieldvalue{\checkboxfield{unchecked}} 三级反应峰。

\begin{center}
\small
\begin{tabular}{|p{1.3cm}|p{2.3cm}|p{2.6cm}|p{1.9cm}|p{1.7cm}|p{1.7cm}|p{1.9cm}|}
\hline
检测次数 &
活细胞密度（个细胞/ml） &
样品密度（个细胞/ml） &
细胞密度 CV 值（\%） &
细胞活率（\%） &
样品活率（\%） &
细胞活率 CV 值（\%）\\
\hline
第一次 &
% #VALUE_ID: LEX-P0025-V0008
% #FIELD_VALUE: 第一次活细胞密度
% #HANDWRITTEN: 2.78×10^6
\fieldvalue{\handwritten{$2.78\times10^{6}$}} &
\multirow{3}{*}{
% #VALUE_ID: LEX-P0025-V0011
% #FIELD_VALUE: 样品密度
% #HANDWRITTEN: 2.76×10^6
\fieldvalue{\handwritten{$2.76\times10^{6}$}}
} &
\multirow{3}{*}{
% #VALUE_ID: LEX-P0025-V0012
% #FIELD_VALUE: 细胞密度CV值
% #HANDWRITTEN: 9
\fieldvalue{\handwritten{9}}
} &
% #VALUE_ID: LEX-P0025-V0013
% #FIELD_VALUE: 第一次细胞活率
% #HANDWRITTEN: 81/12
\fieldvalue{\handwritten{81/12}} &
\multirow{3}{*}{
% #VALUE_ID: LEX-P0025-V0016
% #FIELD_VALUE: 样品活率
% #HANDWRITTEN: 78
\fieldvalue{\handwritten{78}}
} &
\multirow{3}{*}{
% #VALUE_ID: LEX-P0025-V0017
% #FIELD_VALUE: 细胞活率CV值
% #HANDWRITTEN: 2
\fieldvalue{\handwritten{2}}
}\\
\hline
第二次 &
% #VALUE_ID: LEX-P0025-V0009
% #FIELD_VALUE: 第二次活细胞密度
% #HANDWRITTEN: 2.46×10^6
\fieldvalue{\handwritten{$2.46\times10^{6}$}} &
&
&
% #VALUE_ID: LEX-P0025-V0014
% #FIELD_VALUE: 第二次细胞活率
% #HANDWRITTEN: 83/12
\fieldvalue{\handwritten{83/12}} &
&\\
\hline
第三次 &
% #VALUE_ID: LEX-P0025-V0010
% #FIELD_VALUE: 第三次活细胞密度
% #HANDWRITTEN: 2.47×10^6
\fieldvalue{\handwritten{$2.47\times10^{6}$}} &
&
&
% #VALUE_ID: LEX-P0025-V0015
% #FIELD_VALUE: 第三次细胞活率
% #HANDWRITTEN: 94/?
\fieldvalue{\handwritten{94/?}} &
&\\
\hline
发螺细胞总体积（ml） &
% #VALUE_ID: LEX-P0025-V0018
% #FIELD_VALUE: 发螺细胞总体积
% #HANDWRITTEN: 138.90
\fieldvalue{\handwritten{138.90}} &
细胞数量（个细胞/缸） &
% #VALUE_ID: LEX-P0025-V0019
% #FIELD_VALUE: 细胞数量
% #HANDWRITTEN: 2.??×10^?
\fieldvalue{\handwritten{$2.??\times10^{?}$}} &
\multicolumn{3}{l|}{}\\
\hline
\end{tabular}
\end{center}

\noindent
计算要求：样品密度为 $3$ 次细胞密度结果平均值乘以 $1\%$ 白蛋白化稀释液最终体积，再除以生物反应器细胞液体积；细胞数量为样品密度乘以发螺体积。

\begin{center}
\begin{tabular}{|p{2cm}|p{12cm}|}
\hline
备注 &
% #VALUE_ID: LEX-P0025-V0020
% #FIELD_VALUE: 备注
% #HANDWRITTEN: /
\fieldvalue{\handwritten{/}}\\
\hline
实验人/日期 &
% #VALUE_ID: LEX-P0025-V0021
% #FIELD_VALUE: 实验人
% TODO #HANDWRITTEN: 刘?；姓名字迹不清
\fieldvalue{\handwritten{刘?}}
\quad
% #VALUE_ID: LEX-P0025-V0022
% #FIELD_VALUE: 实验日期
% #HANDWRITTEN: 2025.08.19
\fieldvalue{\handwritten{2025.08.19}}\\
\hline
复核人/日期 &
% #VALUE_ID: LEX-P0025-V0023
% #FIELD_VALUE: 复核人
% TODO #HANDWRITTEN: 邓?；姓名字迹不清
\fieldvalue{\handwritten{邓?}}
\quad
% #VALUE_ID: LEX-P0025-V0024
% #FIELD_VALUE: 复核日期
% #HANDWRITTEN: 2025.08.19
\fieldvalue{\handwritten{2025.08.19}}\\
\hline
\end{tabular}
\end{center}
% LEXOID_PAGE_COMPLETED: 25/70

\newpage

\begin{center}
{\large \textbf{生物反应器细胞液检测操作记录}}
\end{center}

\begin{center}
\begin{tabular}{|p{2.2cm}|p{5.0cm}|p{2.0cm}|p{4.0cm}|}
\hline
\multicolumn{4}{c}{细胞生长状态检验}\\
\hline
样品名称 &
% #VALUE_ID: LEX-P0026-V0001
% #FIELD_VALUE: 样品名称
% #HANDWRITTEN: 二
\fieldvalue{\handwritten{二}级生物反应器细胞液} &
样品规格 &
% #VALUE_ID: LEX-P0026-V0002
% #FIELD_VALUE: 样品规格
% #HANDWRITTEN: 0
约\fieldvalue{\handwritten{0}}\ ml/袋\\
\hline
样品批号 &
% #VALUE_ID: LEX-P0026-V0003
% #FIELD_VALUE: 样品批号
% #TODO #HANDWRITTEN: [难以辨认]; 蓝色手写批号
\fieldvalue{\handwritten{[难以辨认]}} &
取样日期 &
% #VALUE_ID: LEX-P0026-V0004
% #FIELD_VALUE: 取样日期
% #TODO #HANDWRITTEN: 2024年03月[难以辨认]日; 日期末位不清晰
\fieldvalue{\handwritten{2024年03月[难以辨认]日}}\\
\hline
\end{tabular}
\end{center}

\section{仪器设备及试剂}

\subsection{仪器设备}

\begin{center}
\begin{tabular}{|p{3.0cm}|p{2.5cm}|p{3.0cm}|p{3.0cm}|p{3.2cm}|}
\hline
仪器名称 & 厂家 & 型号 & 设备编号 & 检查确认\\
\hline
% #VALUE_ID: LEX-P0026-V0005
% #FIELD_VALUE: 仪器名称
\fieldvalue{荧光显微镜} &
% #VALUE_ID: LEX-P0026-V0006
% #FIELD_VALUE: 厂家
\fieldvalue{徕卡} &
% #VALUE_ID: LEX-P0026-V0007
% #FIELD_VALUE: 型号
\fieldvalue{Dmi1 LED} &
% #VALUE_ID: LEX-P0026-V0008
% #FIELD_VALUE: 设备编号
\fieldvalue{1431404} &
正常\ 
% #VALUE_ID: LEX-P0026-V0009
% #FIELD_VALUE: 正常
\fieldvalue{\checkboxfield{checked}}
\quad 异常\ 
% #VALUE_ID: LEX-P0026-V0010
% #FIELD_VALUE: 异常
\fieldvalue{\checkboxfield{unchecked}}\\
\hline
\end{tabular}
\end{center}

\subsection{试剂耗材}

\begin{center}
\begin{tabular}{|p{4.0cm}|p{4.5cm}|p{3.0cm}|p{3.0cm}|p{3.0cm}|}
\hline
名称 & 厂家 & 货号 & 批号 & 有效期至\\
\hline
% #VALUE_ID: LEX-P0026-V0011
% #FIELD_VALUE: 名称
\fieldvalue{细胞活力检测试剂盒} &
% #VALUE_ID: LEX-P0026-V0012
% #FIELD_VALUE: 厂家
\fieldvalue{江苏凯基生物技术股份有限公司} &
% #VALUE_ID: LEX-P0026-V0013
% #FIELD_VALUE: 货号
\fieldvalue{KGA9501-1000} &
% #VALUE_ID: LEX-P0026-V0014
% #FIELD_VALUE: 批号
% #TODO #HANDWRITTEN: 202?+?40; 手写批号辨识不完整
\fieldvalue{\handwritten{202?+?40}} &
% #VALUE_ID: LEX-P0026-V0015
% #FIELD_VALUE: 有效期至
% #TODO #HANDWRITTEN: 202?10.29; 手写日期部分不清晰
\fieldvalue{\handwritten{202?10.29}}\\
\hline
\end{tabular}
\end{center}

\section{操作步骤}

\begin{center}
\begin{tabular}{|p{8.0cm}|p{8.0cm}|}
\hline
\multicolumn{2}{c}{溶液配制}\\
\hline
20$\times$的荧光染液配制：取
% #VALUE_ID: LEX-P0026-V0016
% #FIELD_VALUE: 组分A用量
% #HANDWRITTEN: 0.2
\fieldvalue{\handwritten{0.2}}\ $\mu$l 组分 A 和
% #VALUE_ID: LEX-P0026-V0017
% #FIELD_VALUE: 组分B用量
% #HANDWRITTEN: 0.2
\fieldvalue{\handwritten{0.2}}\ $\mu$l 组分 B 与
% #VALUE_ID: LEX-P0026-V0018
% #FIELD_VALUE: PBS用量
% #HANDWRITTEN: 99.6
\fieldvalue{\handwritten{99.6}}\ $\mu$l PBS 混合均匀配成20$\times$的荧光染液。 &
\\
\hline
\multicolumn{2}{c}{操作过程}\\
\hline
取混匀后样品200$\mu$l/孔至96孔板中，取2孔； &
样品
% #VALUE_ID: LEX-P0026-V0019
% #FIELD_VALUE: 样品用量
% #HANDWRITTEN: 20
\fieldvalue{\handwritten{20}}\ $\mu$l/孔\\
& 每个样品取
% #VALUE_ID: LEX-P0026-V0020
% #FIELD_VALUE: 每个样品取孔数
% #HANDWRITTEN: 2
\fieldvalue{\handwritten{2}}\ 孔\\
\hline
自然沉降后，吸弃上清，注意不要吸到沉淀样品。 &
% #VALUE_ID: LEX-P0026-V0021
% #FIELD_VALUE: 吸弃上清
\fieldvalue{\checkboxfield{checked}}\ 吸弃上清\\
\hline
取20$\times$的荧光染液加至96孔板中，200$\mu$l/孔，避光孵育30--60 min； &
20$\times$的荧光染液
% #VALUE_ID: LEX-P0022-V0022
% #FIELD_VALUE: 荧光染液用量
% #HANDWRITTEN: 200
\fieldvalue{\handwritten{200}}\ $\mu$l/孔\\
& 避光孵育
% #VALUE_ID: LEX-P0026-V0023
% #FIELD_VALUE: 避光孵育开始时间
% #HANDWRITTEN: 11:22
\fieldvalue{\handwritten{11:22}}\ \textasciitilde\
% #VALUE_ID: LEX-P0026-V0024
% #FIELD_VALUE: 避光孵育结束时间
% #TODO #HANDWRITTEN: [难以辨认]; 手写结束时间不清晰
\fieldvalue{\handwritten{[难以辨认]}}\\
\hline
染色结束后，吸弃上清，取PBS加入到96孔板中，200$\mu$l/孔，吸弃上清，注意不要吸到沉淀样品，清洗3次，之后每孔加入100$\mu$l PBS进行拍照。 &
加入PBS
% #VALUE_ID: LEX-P0026-V0025
% #FIELD_VALUE: PBS加入量
% #HANDWRITTEN: 200
\fieldvalue{\handwritten{200}}\ $\mu$l/孔\\
& 清洗
% #VALUE_ID: LEX-P0026-V0026
% #FIELD_VALUE: 清洗次数
% #HANDWRITTEN: 3
\fieldvalue{\handwritten{3}}\ 次\\
\hline
将荧光显微镜调至4$\times$物镜，软件中调节芥相应物镜参数，调节到适合的光源和视距，拍摄4张清晰图像。 &
\\
\hline
选一张清晰图像，打印并粘贴在实验记录上。 &
\\
\hline
\multicolumn{2}{l|}{数据保存路径：%
% #VALUE_ID: LEX-P0026-V0027
% #FIELD_VALUE: 数据保存路径
% #TODO #HANDWRITTEN: [难以辨认]; 蓝色手写保存路径不清晰
\fieldvalue{\handwritten{[难以辨认]}}}\\
\hline
\multicolumn{2}{c}{检测结果}\\
\hline
\end{tabular}
\end{center}

% TODO: “检测结果”区域在本页仅显示表头，后续内容位于后续页面。
% LEXOID_PAGE_COMPLETED: 26/70

\newpage

\noindent
\begin{minipage}{\linewidth}
\small
\begin{tabular}{@{}p{0.28\linewidth}p{0.32\linewidth}p{0.36\linewidth}@{}}
华卓生物科技（北京）有限公司 &
\centering
\begin{tabular}{c}
实验生物\\
% #VALUE_ID: LEX-P0027-V0001
% #FIELD_VALUE: 文件状态
\fieldvalue{原件}
\end{tabular}
&
\raggedleft
\begin{tabular}{c}
实验生物\\
版本号：%
% #VALUE_ID: LEX-P0027-V0002
% #FIELD_VALUE: 版本号
\fieldvalue{1.4}
\end{tabular}
\end{tabular}

\vspace{1mm}

\hfill 文件编码：REC-YZ-SOP-QC-99-006-a
\end{minipage}

\vspace{2mm}

\noindent
\fbox{%
\begin{minipage}[t][0.765\textheight][c]{0.96\linewidth}
\centering
% TODO: Insert the visible microscopy image here.
\fbox{\rule{0pt}{3.0cm}\rule{4.8cm}{0pt}}

\vspace{1mm}
% #VALUE_ID: LEX-P0027-V0003
% #TODO #HANDWRITTEN: 字迹不清; handwritten blue annotation below the image
\handwritten{字迹不清}
\end{minipage}%
}

\vspace{1mm}

\noindent
\begin{tabular}{|p{0.16\linewidth}|p{0.78\linewidth}|}
\hline
\centering 备注 &
% #VALUE_ID: LEX-P0027-V0004
% #FIELD_VALUE: 备注
% #HANDWRITTEN: /
\fieldvalue{\handwritten{/}}\\
\hline
\centering 实验人/日期 &
% #VALUE_ID: LEX-P0027-V0005
% #FIELD_VALUE: 实验人
% #TODO #HANDWRITTEN: 姓名不清; handwritten signature is illegible
\fieldvalue{\handwritten{姓名不清}}\quad
% #VALUE_ID: LEX-P0027-V0006
% #FIELD_VALUE: 实验日期
% #HANDWRITTEN: 2024.08.4
\fieldvalue{\handwritten{2024.08.4}}\\
\hline
\centering 复核人/日期 &
% #VALUE_ID: LEX-P0027-V0007
% #FIELD_VALUE: 复核人
% #TODO #HANDWRITTEN: 姓名不清; handwritten signature is illegible
\fieldvalue{\handwritten{姓名不清}}\quad
% #VALUE_ID: LEX-P0027-V0008
% #FIELD_VALUE: 复核日期
% #HANDWRITTEN: 2024.08.9
\fieldvalue{\handwritten{2024.08.9}}\\
\hline
\end{tabular}

\begin{center}
\scriptsize 第 5 页 共 5 页
\end{center}
% LEXOID_PAGE_COMPLETED: 27/70

\newpage

\noindent Assay: 
% #VALUE_ID: LEX-P0028-V0001
% #FIELD_VALUE: Assay
\fieldvalue{MSC-1}

\vspace{0.5em}

\noindent Cell Type F1: 
% #VALUE_ID: LEX-P0028-V0002
% #FIELD_VALUE: Cell Type F1
\fieldvalue{MSC (AO)}

\noindent Cell Type F2: 
% #VALUE_ID: LEX-P0028-V0003
% #FIELD_VALUE: Cell Type F2
\fieldvalue{MSC (PI)}

\vspace{0.5em}

\noindent Sample ID: 
% #VALUE_ID: LEX-P0028-V0004
% #FIELD_VALUE: Sample ID
\fieldvalue{QC-JS-QY0402-A31Z2202508006-20250819-48H-1}

\noindent Dilution: 
% #VALUE_ID: LEX-P0028-V0005
% #FIELD_VALUE: Dilution
\fieldvalue{2.00}

\vspace{0.5em}
\noindent\rule{\linewidth}{0.4pt}

\noindent Results:

\vspace{0.5em}

\begin{tabular}{@{}p{0.28\linewidth}p{0.34\linewidth}p{0.30\linewidth}@{}}
Count & Concentration & Mean Diameter \\
\texttt{===========} & \texttt{===========} & \texttt{=============} \\[0.2em]
Total:
% #VALUE_ID: LEX-P0028-V0006
% #FIELD_VALUE: Total count
\fieldvalue{1014}
&
% #VALUE_ID: LEX-P0028-V0007
% #FIELD_VALUE: Total concentration
\fieldvalue{$3.60\times10^{6}$ cells/mL}
&
% #VALUE_ID: LEX-P0028-V0008
% #FIELD_VALUE: Total mean diameter
\fieldvalue{15.4 micron}
\\
Live:
% #VALUE_ID: LEX-P0028-V0009
% #FIELD_VALUE: Live count
\fieldvalue{821}
&
% #VALUE_ID: LEX-P0028-V0010
% #FIELD_VALUE: Live concentration
\fieldvalue{$2.92\times10^{6}$ cells/mL}
&
% #VALUE_ID: LEX-P0028-V0011
% #FIELD_VALUE: Live mean diameter
\fieldvalue{16.2 micron}
\\
Dead:
% #VALUE_ID: LEX-P0028-V0012
% #FIELD_VALUE: Dead count
\fieldvalue{193}
&
% #VALUE_ID: LEX-P0028-V0013
% #FIELD_VALUE: Dead concentration
\fieldvalue{$6.76\times10^{5}$ cells/mL}
&
% #VALUE_ID: LEX-P0028-V0014
% #FIELD_VALUE: Dead mean diameter
\fieldvalue{11.8 micron}
\\
Viability:
% #VALUE_ID: LEX-P0028-V0015
% #FIELD_VALUE: Viability
\fieldvalue{81.2\%}
&
&
\end{tabular}

\vspace{1.5em}

\noindent
\begin{minipage}[t]{0.47\linewidth}
% TODO: Image visible on page; no source filename is available.
\fbox{\rule{0pt}{0.25\textheight}\rule{0.95\linewidth}{0pt}}

\vspace{0.8em}

% TODO: Image visible on page; no source filename is available.
\fbox{\rule{0pt}{0.25\textheight}\rule{0.95\linewidth}{0pt}}
\end{minipage}
\hfill
\begin{minipage}[t]{0.47\linewidth}
% TODO: Image visible on page; no source filename is available.
\fbox{\rule{0pt}{0.25\textheight}\rule{0.95\linewidth}{0pt}}

\vspace{0.8em}

% TODO: Image visible on page; no source filename is available.
\fbox{\rule{0pt}{0.25\textheight}\rule{0.95\linewidth}{0pt}}
\end{minipage}

\vfill

\noindent
\begin{tabular}{@{}p{0.48\linewidth}p{0.42\linewidth}@{}}
操作人/审核人：
% #VALUE_ID: LEX-P0028-V0016
% #FIELD_VALUE: 操作人/审核人
% TODO #HANDWRITTEN: illegible signature; handwritten blue signature
\fieldvalue{\handwritten{[illegible signature]}}
&
日期：
% #VALUE_ID: LEX-P0028-V0017
% #FIELD_VALUE: 日期
% #HANDWRITTEN: 2024.8.19
\fieldvalue{\handwritten{2024.8.19}}
\end{tabular}
% LEXOID_PAGE_COMPLETED: 28/70

\newpage

\noindent\textbf{Assay:} 
% #VALUE_ID: LEX-P0029-V0001
% #FIELD_VALUE: Assay
\fieldvalue{MSC-1}

\medskip
\noindent\textbf{Cell Type F1:} 
% #VALUE_ID: LEX-P0029-V0002
% #FIELD_VALUE: Cell Type F1
\fieldvalue{MSC (AO)}

\noindent\textbf{Cell Type F2:} 
% #VALUE_ID: LEX-P0029-V0003
% #FIELD_VALUE: Cell Type F2
\fieldvalue{MSC (PI)}

\medskip
\noindent\textbf{Sample ID:} 
% #VALUE_ID: LEX-P0029-V0004
% #FIELD_VALUE: Sample ID
\fieldvalue{QC-JS-QY0402-A31Z2202508006-20250819-48H-2}

\noindent\textbf{Dilution:} 
% #VALUE_ID: LEX-P0029-V0005
% #FIELD_VALUE: Dilution
\fieldvalue{2.00}

\medskip
\noindent\rule{\linewidth}{0.4pt}

\noindent\textbf{Results:}

\medskip
\noindent
\begin{tabular}{@{}p{0.29\linewidth}p{0.34\linewidth}p{0.30\linewidth}@{}}
\textbf{Count} & \textbf{Concentration} & \textbf{Mean Diameter} \\
\texttt{=============} & \texttt{=============} & \texttt{================} \\[2pt]
Total:
% #VALUE_ID: LEX-P0029-V0006
% #FIELD_VALUE: Total count
\fieldvalue{902}
&
% #VALUE_ID: LEX-P0029-V0010
% #FIELD_VALUE: Total concentration
\fieldvalue{$3.19\times 10^{6}$ cells/mL}
&
% #VALUE_ID: LEX-P0029-V0013
% #FIELD_VALUE: Total mean diameter
\fieldvalue{14.8 micron}
\\
Live:
% #VALUE_ID: LEX-P0029-V0007
% #FIELD_VALUE: Live count
\fieldvalue{750}
&
% #VALUE_ID: LEX-P0029-V0011
% #FIELD_VALUE: Live concentration
\fieldvalue{$2.66\times 10^{6}$ cells/mL}
&
% #VALUE_ID: LEX-P0029-V0014
% #FIELD_VALUE: Live mean diameter
\fieldvalue{15.5 microns}
\\
Dead:
% #VALUE_ID: LEX-P0029-V0008
% #FIELD_VALUE: Dead count
\fieldvalue{152}
&
% #VALUE_ID: LEX-P0029-V0012
% #FIELD_VALUE: Dead concentration
\fieldvalue{$5.32\times 10^{5}$ cells/mL}
&
% #VALUE_ID: LEX-P0029-V0015
% #FIELD_VALUE: Dead mean diameter
\fieldvalue{11.8 microns}
\\
Viability:
% #VALUE_ID: LEX-P0029-V0009
% #FIELD_VALUE: Viability
\fieldvalue{83.3\%}
&
&
\end{tabular}

\medskip

\begin{center}
\begin{tabular}{@{}c@{\hspace{1.2cm}}c@{}}
\fbox{\parbox[c][4.0cm][c]{0.38\linewidth}{\centering\% TODO: visible image; no filename available}}
&
\fbox{\parbox[c][4.0cm][c]{0.38\linewidth}{\centering\% TODO: visible image; no filename available}}
\\[0.5cm]
\fbox{\parbox[c][4.0cm][c]{0.38\linewidth}{\centering\% TODO: visible image; no filename available}}
&
\fbox{\parbox[c][4.0cm][c]{0.38\linewidth}{\centering\% TODO: visible image; no filename available}}
\end{tabular}
\end{center}

\vfill

\noindent 操作人/审核人：
% #VALUE_ID: LEX-P0029-V0016
% #FIELD_VALUE: 操作人/审核人
% #TODO #HANDWRITTEN: illegible signature; handwritten signature is not clearly readable
\fieldvalue{\handwritten{[signature]}}

\hfill 日期：
% #VALUE_ID: LEX-P0029-V0017
% #FIELD_VALUE: 日期
% #HANDWRITTEN: 2025.8.19
\fieldvalue{\handwritten{2025.8.19}}
% LEXOID_PAGE_COMPLETED: 29/70

\newpage

Assay: %
% #VALUE_ID: LEX-P0030-V0001
% #FIELD_VALUE: Assay
\fieldvalue{MSC-1}

\medskip

Cell Type F1: %
% #VALUE_ID: LEX-P0030-V0002
% #FIELD_VALUE: Cell Type F1
\fieldvalue{MSC (AO)}

Cell Type F2: %
% #VALUE_ID: LEX-P0030-V0003
% #FIELD_VALUE: Cell Type F2
\fieldvalue{MSC (PI)}

\medskip

Sample ID: %
% #VALUE_ID: LEX-P0030-V0004
% #FIELD_VALUE: Sample ID
\fieldvalue{QC-JS-QY0402-A31222025080600-20250819-48H-3}

Duration: %
% #VALUE_ID: LEX-P0030-V0005
% #FIELD_VALUE: Duration
\fieldvalue{2.00}

\medskip
\hrule
\medskip

Results:

\begin{center}
\begin{tabular}{@{}p{0.29\linewidth}p{0.34\linewidth}p{0.29\linewidth}@{}}
Count & Concentration & Mean Diameter \\
\underline{\hspace{0.75in}} &
\underline{\hspace{1.05in}} &
\underline{\hspace{1.05in}} \\[2pt]
% #VALUE_ID: LEX-P0030-V0006
% #FIELD_VALUE: Count -- Total
\fieldvalue{Total: 829} &
% #VALUE_ID: LEX-P0030-V0007
% #FIELD_VALUE: Concentration -- Total
\fieldvalue{2.94$\times$10$^{6}$ cells/mL} &
% #VALUE_ID: LEX-P0030-V0008
% #FIELD_VALUE: Mean Diameter -- Total
\fieldvalue{15.3 micron} \\

% #VALUE_ID: LEX-P0030-V0009
% #FIELD_VALUE: Count -- Live
\fieldvalue{Live: 695} &
% #VALUE_ID: LEX-P0030-V0010
% #FIELD_VALUE: Concentration -- Live
\fieldvalue{2.47$\times$10$^{6}$ cells/mL} &
% #VALUE_ID: LEX-P0030-V0011
% #FIELD_VALUE: Mean Diameter -- Live
\fieldvalue{15.9 microns} \\

% #VALUE_ID: LEX-P0030-V0012
% #FIELD_VALUE: Count -- Dead
\fieldvalue{Dead: 134} &
% #VALUE_ID: LEX-P0030-V0013
% #FIELD_VALUE: Concentration -- Dead
\fieldvalue{4.69$\times$10$^{5}$ cells/mL} &
% #VALUE_ID: LEX-P0030-V0014
% #FIELD_VALUE: Mean Diameter -- Dead
\fieldvalue{11.9 microns} \\

% #VALUE_ID: LEX-P0030-V0015
% #FIELD_VALUE: Count -- Viability
\fieldvalue{Viability: 84.0\%} &
& 
\end{tabular}
\end{center}

\medskip

\begin{center}
\begin{tabular}{@{}cc@{}}
% TODO: Upper-left microscopy image visible on this page; source filename unavailable.
\fbox{\parbox[c][1.65in][c]{2.35in}{\centering Image placeholder}} &
% TODO: Upper-right fluorescence image visible on this page; source filename unavailable.
\fbox{\parbox[c][1.65in][c]{2.35in}{\centering Image placeholder}} \\[0.25in]
% TODO: Lower-left microscopy image visible on this page; source filename unavailable.
\fbox{\parbox[c][1.65in][c]{2.35in}{\centering Image placeholder}} &
% TODO: Lower-right fluorescence image visible on this page; source filename unavailable.
\fbox{\parbox[c][1.65in][c]{2.35in}{\centering Image placeholder}}
\end{tabular}
\end{center}

\medskip

操作人/审核人：%
% #VALUE_ID: LEX-P0030-V0016
% #FIELD_VALUE: 操作人/审核人
% #TODO #HANDWRITTEN: illegible; handwritten signature
\fieldvalue{\handwritten{illegible}}

\hfill

日期：%
% #VALUE_ID: LEX-P0030-V0017
% #FIELD_VALUE: 日期
% #HANDWRITTEN: 2025.8.19
\fieldvalue{\handwritten{2025.8.19}}
% LEXOID_PAGE_COMPLETED: 30/70

\newpage

\begin{center}
\textbf{原料请验单}
\end{center}

\noindent
文件编码：REC-YZ-SMP-QC-01-001

\begin{center}
\small
\begin{tabular}{|p{0.16\linewidth}|p{0.34\linewidth}|p{0.20\linewidth}|p{0.25\linewidth}|}
\hline
样品名称 &
% #VALUE_ID: LEX-P0031-V0001
% #FIELD_VALUE: 样品名称
% #TODO #HANDWRITTEN: 二氯甲烷（色谱级，99.8\%）; 部分字迹不清
\fieldvalue{\handwritten{二氯甲烷（色谱级，99.8\%）}} &
样品代号 &
% #VALUE_ID: LEX-P0031-V0002
% #FIELD_VALUE: 样品代号
% #TODO #HANDWRITTEN: QY0402; 字迹部分不清
\fieldvalue{\handwritten{QY0402}}
\\
\hline
样品批号 &
% #VALUE_ID: LEX-P0031-V0003
% #FIELD_VALUE: 样品批号
% #TODO #HANDWRITTEN: A20250802; 字迹部分不清
\fieldvalue{\handwritten{A20250802}} &
样品规格/数量 &
% #VALUE_ID: LEX-P0031-V0004
% #FIELD_VALUE: 样品规格/数量
% #HANDWRITTEN: 2ml/瓶$\times$1
\fieldvalue{\handwritten{2ml/瓶$\times$1}}
\\
\hline
样品类别 &
% #VALUE_ID: LEX-P0031-V0005
% #FIELD_VALUE: 生产过程中检验样品
% #HANDWRITTEN: checked
\fieldvalue{\handwritten{\checkboxfield{checked}}} 生产过程中检验样品
&
\multicolumn{2}{l|}{
% #VALUE_ID: LEX-P0031-V0006
% #FIELD_VALUE: 待包装产品
\fieldvalue{\checkboxfield{unchecked}} 待包装产品\quad
% #VALUE_ID: LEX-P0031-V0007
% #FIELD_VALUE: 成品
\fieldvalue{\checkboxfield{unchecked}} 成品
}
\\
\cline{2-4}
&
% #VALUE_ID: LEX-P0031-V0008
% #FIELD_VALUE: 其他
\fieldvalue{\checkboxfield{unchecked}} 其他（\hspace{2cm}）
& & \\
\hline
检验类型 &
% #VALUE_ID: LEX-P0031-V0009
% #FIELD_VALUE: 初验
% #HANDWRITTEN: checked
\fieldvalue{\handwritten{\checkboxfield{checked}}} 初验\quad
% #VALUE_ID: LEX-P0031-V0010
% #FIELD_VALUE: 复验
\fieldvalue{\checkboxfield{unchecked}} 复验\quad
% #VALUE_ID: LEX-P0031-V0011
% #FIELD_VALUE: 退货
\fieldvalue{\checkboxfield{unchecked}} 退货\quad
% #VALUE_ID: LEX-P0031-V0012
% #FIELD_VALUE: 稳定性
\fieldvalue{\checkboxfield{unchecked}} 稳定性\quad
% #VALUE_ID: LEX-P0031-V0013
% #FIELD_VALUE: 留样
\fieldvalue{\checkboxfield{unchecked}} 留样
& \multicolumn{2}{l|}{} \\
\hline
请验目的 &
% #VALUE_ID: LEX-P0031-V0014
% #FIELD_VALUE: 全检
% #HANDWRITTEN: checked
\fieldvalue{\handwritten{\checkboxfield{checked}}} 全检\quad
% #VALUE_ID: LEX-P0031-V0015
% #FIELD_VALUE: 留样
\fieldvalue{\checkboxfield{unchecked}} 留样\quad
% #VALUE_ID: LEX-P0031-V0016
% #FIELD_VALUE: 其他
\fieldvalue{\checkboxfield{unchecked}} 其他（\hspace{2cm}）
& \multicolumn{2}{l|}{} \\
\hline
请验部门 &
% #VALUE_ID: LEX-P0031-V0017
% #FIELD_VALUE: 生产部
% #HANDWRITTEN: checked
\fieldvalue{\handwritten{\checkboxfield{checked}}} 生产部\quad
% #VALUE_ID: LEX-P0031-V0018
% #FIELD_VALUE: 物料管理部
\fieldvalue{\checkboxfield{unchecked}} 物料管理部\quad
% #VALUE_ID: LEX-P0031-V0019
% #FIELD_VALUE: 其他
\fieldvalue{\checkboxfield{unchecked}} 其他（\hspace{2cm}）
& \multicolumn{2}{l|}{} \\
\hline
样品情况 &
% #VALUE_ID: LEX-P0031-V0020
% #FIELD_VALUE: 随请验单一起提供
\fieldvalue{\checkboxfield{unchecked}} 随请验单一起提供\quad
% #VALUE_ID: LEX-P0031-V0021
% #FIELD_VALUE: 未随请验单一起提供
\fieldvalue{\checkboxfield{unchecked}} 未随请验单一起提供
& \multicolumn{2}{l|}{} \\
\hline
请验人 &
% #VALUE_ID: LEX-P0031-V0022
% #FIELD_VALUE: 请验人
% #TODO #HANDWRITTEN: 赵永明; 姓名笔迹不清
\fieldvalue{\handwritten{赵永明}} &
请验日期 &
% #VALUE_ID: LEX-P0031-V0023
% #FIELD_VALUE: 请验日期
% #HANDWRITTEN: 2025.08.18
\fieldvalue{\handwritten{2025.08.18}}
\\
\hline
收验人 &
% #VALUE_ID: LEX-P0031-V0024
% #FIELD_VALUE: 收验人
% #TODO #HANDWRITTEN: illegible; 签名难以辨认
\fieldvalue{\handwritten{签名}} &
收验日期 &
% #VALUE_ID: LEX-P0031-V0025
% #FIELD_VALUE: 收验日期
% #HANDWRITTEN: 2025.08.18
\fieldvalue{\handwritten{2025.08.18}}
\\
\hline
备注 &
% #VALUE_ID: LEX-P0031-V0026
% #FIELD_VALUE: 备注
% #TODO #HANDWRITTEN: 24h 9L; 字迹及单位部分不清
\fieldvalue{\handwritten{24h\quad 9L}}
& \multicolumn{2}{l|}{} \\
\hline
\end{tabular}
\end{center}
% LEXOID_PAGE_COMPLETED: 31/70

\newpage

\section*{取样记录}

\begin{center}
\small
\begin{tabular}{|p{0.14\linewidth}|p{0.32\linewidth}|p{0.14\linewidth}|p{0.30\linewidth}|}
\hline
名称 &
% #VALUE_ID: LEX-P0032-V0001
% #FIELD_VALUE: 名称
\fieldvalue{二级生物反应器细胞液} &
代码 &
% #VALUE_ID: LEX-P0032-V0002
% #FIELD_VALUE: 代码
\fieldvalue{QY0402}
\\ \hline
规格 &
% #VALUE_ID: LEX-P0032-V0003
% #FIELD_VALUE: 规格
\fieldvalue{20ml/袋} &
总数量 &
% #VALUE_ID: LEX-P0032-V0004
% #FIELD_VALUE: 总数量
\fieldvalue{1袋}
\\ \hline
批号 &
% #VALUE_ID: LEX-P0032-V0005
% #FIELD_VALUE: 批号
\fieldvalue{A31Z220508006} &
进厂批号 &
% #VALUE_ID: LEX-P0032-V0006
% #FIELD_VALUE: 进厂批号
\fieldvalue{/}
\\ \hline
\multicolumn{1}{|l|}{供货单位/厂家} &
\multicolumn{3}{l|}{
% #VALUE_ID: LEX-P0032-V0007
% #FIELD_VALUE: 供货单位/厂家
\fieldvalue{/}}
\\ \hline
取样目的 &
\multicolumn{3}{l|}{
% #VALUE_ID: LEX-P0032-V0008
% #FIELD_VALUE: 正常检验
\fieldvalue{\checkboxfield{checked}} 正常检验\quad
% #VALUE_ID: LEX-P0032-V0009
% #FIELD_VALUE: 复检
\fieldvalue{\checkboxfield{unchecked}} 复检\quad
% #VALUE_ID: LEX-P0032-V0010
% #FIELD_VALUE: 稳定性
\fieldvalue{\checkboxfield{unchecked}} 稳定性
\newline
\quad
% #VALUE_ID: LEX-P0032-V0011
% #FIELD_VALUE: 退货检验
\fieldvalue{\checkboxfield{unchecked}} 退货检验\quad
% #VALUE_ID: LEX-P0032-V0012
% #FIELD_VALUE: 留样
\fieldvalue{\checkboxfield{unchecked}} 留样\quad
% #VALUE_ID: LEX-P0032-V0013
% #FIELD_VALUE: 其他
\fieldvalue{\checkboxfield{unchecked}} 其他：\underline{\hspace{2cm}}}
\\ \hline
取样地点 &
\multicolumn{3}{l|}{
% #VALUE_ID: LEX-P0032-V0014
% #FIELD_VALUE: 产品出品
\fieldvalue{\checkboxfield{checked}} 产品出品\quad
% #VALUE_ID: LEX-P0032-V0015
% #FIELD_VALUE: 物料库
\fieldvalue{\checkboxfield{unchecked}} 物料库\quad
% #VALUE_ID: LEX-P0032-V0016
% #FIELD_VALUE: 细胞库
\fieldvalue{\checkboxfield{unchecked}} 细胞库\quad
% #VALUE_ID: LEX-P0032-V0017
% #FIELD_VALUE: 包装间
\fieldvalue{\checkboxfield{unchecked}} 包装间}
\\ \hline
贮存条件 &
\multicolumn{3}{l|}{
% #VALUE_ID: LEX-P0032-V0018
% #FIELD_VALUE: 2--8$^\circ$C
\fieldvalue{\checkboxfield{checked}} 2--8$^\circ$C\quad
% #VALUE_ID: LEX-P0032-V0019
% #FIELD_VALUE: -20$^\circ$C
\fieldvalue{\checkboxfield{unchecked}} -20$^\circ$C\quad
% #VALUE_ID: LEX-P0032-V0020
% #FIELD_VALUE: -80$^\circ$C
\fieldvalue{\checkboxfield{unchecked}} -80$^\circ$C\quad
% #VALUE_ID: LEX-P0032-V0021
% #FIELD_VALUE: 常温
\fieldvalue{\checkboxfield{unchecked}} 常温\quad
% #VALUE_ID: LEX-P0032-V0022
% #FIELD_VALUE: 其他
\fieldvalue{\checkboxfield{unchecked}} 其他：\underline{\hspace{1.5cm}}}
\\ \hline
盛装容器 &
\multicolumn{3}{l|}{
% #VALUE_ID: LEX-P0032-V0023
% #FIELD_VALUE: 保温箱
\fieldvalue{\checkboxfield{checked}} 保温箱\quad
% #VALUE_ID: LEX-P0032-V0024
% #FIELD_VALUE: 泡沫箱
\fieldvalue{\checkboxfield{unchecked}} 泡沫箱\quad
% #VALUE_ID: LEX-P0032-V0025
% #FIELD_VALUE: 手提篮
\fieldvalue{\checkboxfield{unchecked}} 手提篮\quad
% #VALUE_ID: LEX-P0032-V0026
% #FIELD_VALUE: 其他
\fieldvalue{\checkboxfield{unchecked}} 其他：\underline{\hspace{1.5cm}}}
\\ \hline
外观复核 &
\multicolumn{3}{l|}{
% #VALUE_ID: LEX-P0032-V0027
% #FIELD_VALUE: 合格
\fieldvalue{\checkboxfield{checked}} 合格\quad
% #VALUE_ID: LEX-P0032-V0028
% #FIELD_VALUE: 不合格
\fieldvalue{\checkboxfield{unchecked}} 不合格}
\\ \hline
取样量 &
\multicolumn{1}{c|}{
% #VALUE_ID: LEX-P0032-V0029
% #FIELD_VALUE: 取样量
\fieldvalue{1袋}} &
检验量 &
\multicolumn{1}{c|}{
% #VALUE_ID: LEX-P0032-V0030
% #FIELD_VALUE: 检验量
\fieldvalue{1袋}} \\
\hline
留样量 &
\multicolumn{1}{c|}{
% #VALUE_ID: LEX-P0032-V0031
% #FIELD_VALUE: 留样量
\fieldvalue{/}} &
\multicolumn{2}{c|}{} \\
\hline
取样人/日期 &
\multicolumn{1}{c|}{
% #VALUE_ID: LEX-P0032-V0032
% #FIELD_VALUE: 取样人
% #TODO #HANDWRITTEN: 签名不清; 蓝色草写字迹难以辨认
\fieldvalue{\handwritten{签名不清}} \quad
% #VALUE_ID: LEX-P0032-V0033
% #FIELD_VALUE: 取样日期
% #HANDWRITTEN: 2025.8.18
\fieldvalue{\handwritten{2025.8.18}}} &
复核人/日期 &
% #VALUE_ID: LEX-P0032-V0034
% #FIELD_VALUE: 复核人
% #TODO #HANDWRITTEN: 签名不清; 蓝色草写字迹难以辨认
\fieldvalue{\handwritten{签名不清}} \quad
% #VALUE_ID: LEX-P0032-V0035
% #FIELD_VALUE: 复核日期
% #HANDWRITTEN: 2025.08.08
\fieldvalue{\handwritten{2025.08.08}}
\\ \hline
备注 &
\multicolumn{3}{l|}{
% #VALUE_ID: LEX-P0032-V0036
% #FIELD_VALUE: 备注
% #TODO #HANDWRITTEN: 字迹不清，疑似“7M 72”；蓝色手写内容难以辨认
\fieldvalue{\handwritten{7M 72}}}
\\ \hline
\end{tabular}
\end{center}

\begin{center}
第 1 页 共 1 页
\end{center}
% LEXOID_PAGE_COMPLETED: 32/70

\newpage

\begin{center}
\textbf{生物反应器细胞液检测记录}
\end{center}

\noindent
文 件 编 码：REC-YZ-SOP-QC-99-006-a

\begin{center}
\begin{tabular}{|p{0.18\linewidth}|p{0.30\linewidth}|p{0.18\linewidth}|p{0.24\linewidth}|}
\hline
\multicolumn{2}{c|}{宿主生物} & \multicolumn{2}{c|}{宿主生物}\\
\hline
制备状态 &
% #VALUE_ID: LEX-P0033-V0001
% #FIELD_VALUE: 制备状态
\fieldvalue{原件}
& 版本号 & 1.4\\
\hline
样品名称 &
% #VALUE_ID: LEX-P0033-V0002
% #FIELD_VALUE: 样品名称
\fieldvalue{二级生物反应器细胞液}
& 样品规格 &
% #VALUE_ID: LEX-P0033-V0003
% #FIELD_VALUE: 样品规格
\fieldvalue{20}\,ml/袋\\
\hline
样品批号 &
% #VALUE_ID: LEX-P0033-V0004
% #FIELD_VALUE: 样品批号
% #HANDWRITTEN: A12?W?80-1
\fieldvalue{\handwritten{A12?W?80-1}}
& 取样日期 &
% #VALUE_ID: LEX-P0033-V0005
% #FIELD_VALUE: 取样日期
% #HANDWRITTEN: 2024年08月18日
\fieldvalue{\handwritten{2024年08月18日}}\\
\hline
\end{tabular}
\end{center}

\noindent\textbf{1\quad 仪器设备、试剂耗材：}

\noindent\textbf{1.1\quad 仪器设备}

\begin{center}
\small
\begin{tabular}{|p{0.16\linewidth}|p{0.28\linewidth}|p{0.16\linewidth}|p{0.19\linewidth}|p{0.19\linewidth}|}
\hline
名称 & 厂家 & 型号 & 设备编号 & 检查确认\\
\hline
% #VALUE_ID: LEX-P0033-V0006
% #FIELD_VALUE: 仪器名称
\fieldvalue{血糖仪}
&
% #VALUE_ID: LEX-P0033-V0007
% #FIELD_VALUE: 厂家
\fieldvalue{三诺生物传感股份有限公司}
&
% #VALUE_ID: LEX-P0033-V0008
% #FIELD_VALUE: 型号
\fieldvalue{安稳+}
&
% #VALUE_ID: LEX-P0033-V0009
% #FIELD_VALUE: 设备编号
% #HANDWRITTEN: /
\fieldvalue{\handwritten{/}}
&
正常\ 
% #VALUE_ID: LEX-P0033-V0010
% #FIELD_VALUE: 正常
\fieldvalue{\checkboxfield{checked}}
\quad 异常\ 
% #VALUE_ID: LEX-P0033-V0011
% #FIELD_VALUE: 异常
\fieldvalue{\checkboxfield{unchecked}}
\\
\hline
% #VALUE_ID: LEX-P0033-V0012
% #FIELD_VALUE: 仪器名称
\fieldvalue{移液器}
&
% #VALUE_ID: LEX-P0033-V0013
% #FIELD_VALUE: 厂家
\fieldvalue{Eppendorf Corporate}
&
% #VALUE_ID: LEX-P0033-V0014
% #FIELD_VALUE: 型号
\fieldvalue{0.1--2.5\,$\mu$l}
&
% #VALUE_ID: LEX-P0033-V0015
% #FIELD_VALUE: 设备编号
% #HANDWRITTEN: 127?436
\fieldvalue{\handwritten{127?436}}
&
正常\ 
% #VALUE_ID: LEX-P0033-V0016
% #FIELD_VALUE: 正常
\fieldvalue{\checkboxfield{checked}}
\quad 异常\ 
% #VALUE_ID: LEX-P0033-V0017
% #FIELD_VALUE: 异常
\fieldvalue{\checkboxfield{unchecked}}
\\
\hline
\end{tabular}
\end{center}

\noindent\textbf{1.2\quad 试剂耗材}

\begin{center}
\small
\begin{tabular}{|p{0.20\linewidth}|p{0.34\linewidth}|p{0.22\linewidth}|p{0.20\linewidth}|}
\hline
名称 & 厂家 & 批号 & 有效期至\\
\hline
% #VALUE_ID: LEX-P0033-V0018
% #FIELD_VALUE: 耗材名称
\fieldvalue{血糖测试条}
&
% #VALUE_ID: LEX-P0033-V0019
% #FIELD_VALUE: 厂家
\fieldvalue{三诺生物传感股份有限公司}
&
% #VALUE_ID: LEX-P0033-V0020
% #FIELD_VALUE: 批号
% #HANDWRITTEN: 444?9?30
\fieldvalue{\handwritten{444?9?30}}
&
% #VALUE_ID: LEX-P0033-V0021
% #FIELD_VALUE: 有效期至
% #HANDWRITTEN: 2027.3.18
\fieldvalue{\handwritten{2027.3.18}}
\\
\hline
\end{tabular}
\end{center}

\noindent\textbf{2\quad 操作步骤：}

\begin{center}
\small
\begin{tabular}{|p{0.40\linewidth}|p{0.52\linewidth}|}
\hline
\multicolumn{2}{c|}{操作过程}\\
\hline
取出测试条，取样品自然沉降后的上清液
0.6$\mu$l加入测试条电极，血糖仪自动发出“嘀”提示音，吸入样品完成，血糖仪自动进入测试时，开始测试； &
取上清液：
% #VALUE_ID: LEX-P0033-V0022
% #FIELD_VALUE: 取上清液体积
\fieldvalue{0.1}\,$\mu$l\\
\hline
重复以上步骤2次，共计数3次。 &
每个样品分别共计数
% #VALUE_ID: LEX-P0033-V0023
% #FIELD_VALUE: 样品计数次数
\fieldvalue{3}
次。\\
\hline
\multicolumn{2}{c|}{检测结果}\\
\hline
\multicolumn{2}{c|}{葡萄糖浓度（mmol/L）}\\
\hline
\multicolumn{2}{c|}{\begin{tabular}{cccc}
第一次 & 第二次 & 第三次 & 平均值（保留一位小数）\\
\hline
% #VALUE_ID: LEX-P0033-V0024
% #FIELD_VALUE: 第一次检测结果
% #HANDWRITTEN: 18.3
\fieldvalue{\handwritten{18.3}}
&
% #VALUE_ID: LEX-P0033-V0025
% #FIELD_VALUE: 第二次检测结果
% #HANDWRITTEN: 18.7
\fieldvalue{\handwritten{18.7}}
&
% #VALUE_ID: LEX-P0033-V0026
% #FIELD_VALUE: 第三次检测结果
% #HANDWRITTEN: 19.2
\fieldvalue{\handwritten{19.2}}
&
% #VALUE_ID: LEX-P0033-V0027
% #FIELD_VALUE: 平均值
% #HANDWRITTEN: 18.9
\fieldvalue{\handwritten{18.9}}
\end{tabular}}\\
\hline
\multicolumn{2}{|l|}{备注：\space
% #VALUE_ID: LEX-P0033-V0028
% #FIELD_VALUE: 备注
% #HANDWRITTEN: /
\fieldvalue{\handwritten{/}}}\\
\hline
\multicolumn{2}{|l|}{实验人/日期：\space
% #VALUE_ID: LEX-P0033-V0029
% #FIELD_VALUE: 实验人
% #TODO #HANDWRITTEN: 姓名不清；蓝色手写字迹难以辨认
\fieldvalue{\handwritten{姓名不清}}
\quad
% #VALUE_ID: LEX-P0033-V0030
% #FIELD_VALUE: 实验日期
% #HANDWRITTEN: 2024.08.18
\fieldvalue{\handwritten{2024.08.18}}}\\
\hline
\multicolumn{2}{|l|}{复核人/日期：\space
% #VALUE_ID: LEX-P0033-V0031
% #FIELD_VALUE: 复核人
% #TODO #HANDWRITTEN: 姓名不清；蓝色手写字迹难以辨认
\fieldvalue{\handwritten{姓名不清}}
\quad
% #VALUE_ID: LEX-P0033-V0032
% #FIELD_VALUE: 复核日期
% #HANDWRITTEN: 2024.08.18
\fieldvalue{\handwritten{2024.08.18}}}\\
\hline
\end{tabular}
\end{center}

\noindent\hfill 第 1 页 共 5 页
% LEXOID_PAGE_COMPLETED: 33/70

\newpage

\begin{center}
\textbf{生物反应器细胞液检测记录}
\end{center}

\noindent 文件编码：REC-YZ-SOP-QC-099-006-a

\begin{tabular}{|p{0.16\linewidth}|p{0.34\linewidth}|p{0.16\linewidth}|p{0.24\linewidth}|}
\hline
\multicolumn{4}{c|}{密度与语密检验} \\ \hline
样品名称 & 一级生物反应器细胞液 & 样品规格 & ml/袋 \\ \hline
样品批号 &
% #VALUE_ID: LEX-P0034-V0001
% #FIELD_VALUE: 样品批号
% #HANDWRITTEN: 21?2?0?2?6
\fieldvalue{\handwritten{21?2?0?2?6}} &
取样日期 &
% #VALUE_ID: LEX-P0034-V0002
% #FIELD_VALUE: 取样日期
% #HANDWRITTEN: 2024年8月12日
\fieldvalue{\handwritten{2024年8月12日}} \\ \hline
\end{tabular}

\section{仪器设备、试剂耗材}

\subsection{仪器设备}

\begin{tabular}{|p{0.18\linewidth}|p{0.18\linewidth}|p{0.16\linewidth}|p{0.18\linewidth}|p{0.16\linewidth}|l|}
\hline
名称 & 厂家 & 型号 & 设备编号 & 校准有效期至 & 检查确认 \\ \hline
细胞计数仪 & Nexcelom & Cellometer K2 & 1410905 &
% #VALUE_ID: LEX-P0034-V0003
% #FIELD_VALUE: 校准有效期至
% #HANDWRITTEN: 2025.11.1
\fieldvalue{\handwritten{2025.11.1}} &
% #VALUE_ID: LEX-P0034-V0004
% #FIELD_VALUE: 正常
% #HANDWRITTEN: ✓
\fieldvalue{\handwritten{\checkboxfield{checked}}}正常\quad
% #VALUE_ID: LEX-P0034-V0005
% #FIELD_VALUE: 异常
\fieldvalue{\checkboxfield{unchecked}}异常 \\ \hline
高速冷冻离心机 & Thermo & ST16R & 1410810 &
% #VALUE_ID: LEX-P0034-V0006
% #FIELD_VALUE: 校准有效期至
% #HANDWRITTEN: 2025.7.1
\fieldvalue{\handwritten{2025.7.1}} &
% #VALUE_ID: LEX-P0034-V0007
% #FIELD_VALUE: 正常
% #HANDWRITTEN: ✓
\fieldvalue{\handwritten{\checkboxfield{checked}}}正常\quad
% #VALUE_ID: LEX-P0034-V0008
% #FIELD_VALUE: 异常
\fieldvalue{\checkboxfield{unchecked}}异常 \\ \hline
电热恒温水槽 & 上海一恒 & DK-8AXX & 1411001 &
% #VALUE_ID: LEX-P0034-V0009
% #FIELD_VALUE: 校准有效期至
% #HANDWRITTEN: 2024.2.20
\fieldvalue{\handwritten{2024.2.20}} &
% #VALUE_ID: LEX-P0034-V0010
% #FIELD_VALUE: 正常
% #HANDWRITTEN: ✓
\fieldvalue{\handwritten{\checkboxfield{checked}}}正常\quad
% #VALUE_ID: LEX-P0034-V0011
% #FIELD_VALUE: 异常
\fieldvalue{\checkboxfield{unchecked}}异常 \\ \hline
\end{tabular}

\subsection{试剂耗材}

\begin{tabular}{|p{0.20\linewidth}|p{0.34\linewidth}|p{0.22\linewidth}|l|}
\hline
名称 & 厂家 & 批号 & 有效期至 \\ \hline
AO/PI 染液 & Nexcelom &
% #VALUE_ID: LEX-P0034-V0012
% #FIELD_VALUE: 批号
% #HANDWRITTEN: 33449?
\fieldvalue{\handwritten{33449?}} &
% #VALUE_ID: LEX-P0034-V0013
% #FIELD_VALUE: 有效期至
% #HANDWRITTEN: 2025.7.21
\fieldvalue{\handwritten{2025.7.21}} \\ \hline
0.9\%氯化钠注射液 & 石家庄四药 &
% #VALUE_ID: LEX-P0034-V0014
% #FIELD_VALUE: 批号
% #HANDWRITTEN: 20241212
\fieldvalue{\handwritten{20241212}} &
% #VALUE_ID: LEX-P0034-V0015
% #FIELD_VALUE: 有效期至
% #HANDWRITTEN: 2026.09.06
\fieldvalue{\handwritten{2026.09.06}} \\ \hline
20\%人血白蛋白 & Baxter AG &
% #VALUE_ID: LEX-P0034-V0016
% #FIELD_VALUE: 批号
% #HANDWRITTEN: 21?4?4?0
\fieldvalue{\handwritten{21?4?4?0}} &
% #VALUE_ID: LEX-P0034-V0017
% #FIELD_VALUE: 有效期至
% #HANDWRITTEN: 2027.08.31
\fieldvalue{\handwritten{2027.08.31}} \\ \hline
\end{tabular}

\section{操作步骤}

\begin{tabular}{|p{0.56\linewidth}|p{0.36\linewidth}|}
\hline
\multicolumn{2}{c|}{溶液配准} \\ \hline
1\%人白蛋白溶液配液：取5 ml 20\%人血白蛋白加入到95 ml的0.9\%氯化钠注射液中，混匀即得。 &
20\%人血白蛋白
% #VALUE_ID: LEX-P0034-V0018
% #FIELD_VALUE: 20\%人血白蛋白加入量
% #HANDWRITTEN: 1
\fieldvalue{\handwritten{1}} ml \\ \cline{2-2}
&
0.9\%氯化钠注射液
% #VALUE_ID: LEX-P0034-V0019
% #FIELD_VALUE: 0.9\%氯化钠注射液加入量
% #HANDWRITTEN: 9
\fieldvalue{\handwritten{9}} ml \\ \hline
\multicolumn{2}{c|}{} \\[-6pt]
\multicolumn{2}{c|}{\begin{tabular}{ccc}
名称 & 配制批号 & 有效期至 \\ \hline
5$\times$裂解液 &
% #VALUE_ID: LEX-P0034-V0020
% #FIELD_VALUE: 配制批号
% #HANDWRITTEN: 202407280
\fieldvalue{\handwritten{202407280}} &
% #VALUE_ID: LEX-P0034-V0021
% #FIELD_VALUE: 有效期至
% #HANDWRITTEN: 2027.10.27
\fieldvalue{\handwritten{2027.10.27}}
\end{tabular}} \\ \hline
\end{tabular}

\noindent\textbf{操作过程}

\begin{tabular}{|p{0.56\linewidth}|p{0.36\linewidth}|}
\hline
按照50 ml离心管的刻度读取样品体积，往取样管中加入对应体积的5$\times$裂解液（每4 ml样品加入1 ml 5$\times$裂解液，按照“只进不舍”原则，如26 ml样品，加入7 ml 5$\times$裂解液），充分混匀后倒入至2.3的样品中，置于37℃电热恒温水槽中裂解20分钟后，取出离心管振荡混匀，使样品充分裂解，再置于37℃电热恒温水槽中裂解10min，得到细胞悬液。 &
样品体积：
% #VALUE_ID: LEX-P0034-V0022
% #FIELD_VALUE: 样品体积
% #HANDWRITTEN: 3.6
\fieldvalue{\handwritten{3.6}} ml \\ \cline{2-2}
&
裂解液：
% #VALUE_ID: LEX-P0034-V0023
% #FIELD_VALUE: 裂解液加入量
% #HANDWRITTEN: 1
\fieldvalue{\handwritten{1}} ml \\ \cline{2-2}
&
电热恒温水槽
% #VALUE_ID: LEX-P0034-V0024
% #FIELD_VALUE: 电热恒温水槽温度
% #HANDWRITTEN: 37
\fieldvalue{\handwritten{37}}℃ \\ \cline{2-2}
&
裂解：
% #VALUE_ID: LEX-P0034-V0025
% #FIELD_VALUE: 裂解时间
% #HANDWRITTEN: 11分钟
\fieldvalue{\handwritten{11分钟}} \\ \hline
取出裂解好的细胞悬液300$\times g$离心5分钟，按照读取样品体积加入对应体积的1\%人白蛋白溶液重悬（每1.2 ml样品体积加入0.1 ml 1\%人白蛋白溶液，按照“只进不舍”原则，如26 ml样品体积，加入2 ml 1\%人白蛋白溶液）。 &
% #VALUE_ID: LEX-P0034-V0026
% #FIELD_VALUE: 离心速度
% #HANDWRITTEN: 300
\fieldvalue{\handwritten{300}}$\times g$离心
% #VALUE_ID: LEX-P0034-V0027
% #FIELD_VALUE: 离心时间
% #HANDWRITTEN: 4
\fieldvalue{\handwritten{4}}分钟 \\ \cline{2-2}
&
加入1\%人白蛋白溶液重悬体积：
% #VALUE_ID: LEX-P0034-V0028
% #FIELD_VALUE: 1\%人白蛋白溶液重悬体积
% #HANDWRITTEN: 3
\fieldvalue{\handwritten{3}} ml \\ \cline{2-2}
&
吸取
% #VALUE_ID: LEX-P0034-V0029
% #FIELD_VALUE: 吸取量
% #HANDWRITTEN: 2
\fieldvalue{\handwritten{2}} $\mu$l \\ \hline
\end{tabular}

\begin{center}
第 2 页 共 5 页
\end{center}
% LEXOID_PAGE_COMPLETED: 34/70

\newpage

\noindent
\begin{minipage}{\linewidth}
\small
\begin{flushright}
文件编号：REC-ZY-SOP-QC-99-006-a\\
\textcolor{red}{版本号：1.4}
\end{flushright}

\noindent
\begin{tabular}{p{0.57\linewidth}|p{0.38\linewidth}}
加入 2.2ml 1\%人白蛋白消化液，反复吹打至单细胞悬液。 & 与
% #VALUE_ID: LEX-P0035-V0001
% #FIELD_VALUE: 与AO/PI染料混合的体积
% #TODO #HANDWRITTEN: 20；手写数字略模糊
\fieldvalue{\handwritten{20}}$\mu$l AO/PI染料混合均匀。\\
撕去计数板上下两层保护膜，将细胞悬液混匀，从中吸取 20$\mu$l 与 20$\mu$l AO/PI 染料混合均匀。 & 吸取
% #VALUE_ID: LEX-P0035-V0002
% #FIELD_VALUE: 吸取的混合液体积
% #HANDWRITTEN: 20
\fieldvalue{\handwritten{20}}$\mu$l 计数。\\
吸取 20$\mu$l 细胞悬液与 AO/PI 染料的混合液，注入到计数板的一个检测室中。 & 每个样品总共计数
% #VALUE_ID: LEX-P0035-V0003
% #FIELD_VALUE: 每个样品总共计数次数
% #HANDWRITTEN: 2
\fieldvalue{\handwritten{2}}次。\\
将加入混合液的一端插入进液口。 & \\
选择 Assay 类型，使用 AO/PI 染料进行双荧光存活率检测； & \\
在主界面左侧 ID 栏中输入样本批号。 & \\
点击主界面左下角“Preview”预览细胞图片。 & \\
调整仪器焦距直至观察到荧光通道下的细胞中心透亮，轮廓清晰。 & \\
点击“count”按钮进行计数并自动拍摄图片。 & \\
重复以上步骤 2 次，共计数 3 次。打印数据结果。 & \\
\end{tabular}

\vspace{2pt}
\noindent
数据保存路径：
% #VALUE_ID: LEX-P0035-V0004
% #FIELD_VALUE: 数据保存路径
% #TODO #HANDWRITTEN: 14/??/??-??、L2014-??、??；手写路径难以辨认
\fieldvalue{\handwritten{14/??/??-??、L2014-??、??}}

\vspace{2pt}
\noindent
\begin{tabular}{|p{0.13\linewidth}|p{0.12\linewidth}|p{0.12\linewidth}|p{0.12\linewidth}|p{0.12\linewidth}|p{0.12\linewidth}|p{0.12\linewidth}|}
\hline
\multicolumn{7}{|c|}{检测结果}\\
\hline
\multicolumn{7}{|c|}{检测结果：\ 
% #VALUE_ID: LEX-P0035-V0005
% #FIELD_VALUE: 一级反应
\fieldvalue{\checkboxfield{unchecked}} 一级反应、
% #VALUE_ID: LEX-P0035-V0006
% #FIELD_VALUE: 二级反应
\fieldvalue{\checkboxfield{unchecked}} 二级反应、
% #VALUE_ID: LEX-P0035-V0007
% #FIELD_VALUE: 三级反应
\fieldvalue{\checkboxfield{unchecked}} 三级反应}\\
\hline
检测次 数 &
活细胞密度(个细胞/ml) &
样品密度(个细胞/ml) &
细胞密度 CV值（\%）&
细胞活率（\%）&
样品活率（\%）&
细胞活率 CV值（\%）\\
\hline
第一次 &
% #VALUE_ID: LEX-P0035-V0008
% #FIELD_VALUE: 第一次活细胞密度
% #TODO #HANDWRITTEN: 1.64$\times 10^6$；指数及末位数字略模糊
\fieldvalue{\handwritten{$1.64\times10^6$}} &
&
&
% #VALUE_ID: LEX-P0035-V0013
% #FIELD_VALUE: 第一次细胞活率
% #HANDWRITTEN: 85.0
\fieldvalue{\handwritten{85.0}} &
& \\
\hline
第二次 &
% #VALUE_ID: LEX-P0035-V0009
% #FIELD_VALUE: 第二次活细胞密度
% #TODO #HANDWRITTEN: 7.12$\times 10^6$；手写数字略模糊
\fieldvalue{\handwritten{$7.12\times10^6$}} &
\multirow{2}{*}{
% #VALUE_ID: LEX-P0035-V0011
% #FIELD_VALUE: 样品密度
% #TODO #HANDWRITTEN: 1.21$\times10^7$；手写指数略模糊
\fieldvalue{\handwritten{$1.21\times10^7$}}} &
\multirow{2}{*}{
% #VALUE_ID: LEX-P0035-V0012
% #FIELD_VALUE: 细胞密度CV值
% #HANDWRITTEN: 14
\fieldvalue{\handwritten{14}}} &
% #VALUE_ID: LEX-P0035-V0014
% #FIELD_VALUE: 第二次细胞活率
% #TODO #HANDWRITTEN: 81.0；末位数字略模糊
\fieldvalue{\handwritten{81.0}} &
\multirow{2}{*}{
% #VALUE_ID: LEX-P0035-V0016
% #FIELD_VALUE: 样品活率
% #HANDWRITTEN: 89
\fieldvalue{\handwritten{89}}} &
\multirow{2}{*}{
% #VALUE_ID: LEX-P0035-V0017
% #FIELD_VALUE: 细胞活率CV值
% #HANDWRITTEN: 2
\fieldvalue{\handwritten{2}}}\\
\cline{1-2}\cline{5-5}
第三次 &
% #VALUE_ID: LEX-P0035-V0010
% #FIELD_VALUE: 第三次活细胞密度
% #TODO #HANDWRITTEN: 6.14$\times 10^6$；手写数字略模糊
\fieldvalue{\handwritten{$6.14\times10^6$}} &
&
&
% #VALUE_ID: LEX-P0035-V0015
% #FIELD_VALUE: 第三次细胞活率
% #TODO #HANDWRITTEN: 89.0；手写末位数字略模糊
\fieldvalue{\handwritten{89.0}} &
& \\
\hline
发酵细胞体积（ml）&
% #VALUE_ID: LEX-P0035-V0018
% #FIELD_VALUE: 发酵细胞体积
% #TODO #HANDWRITTEN: 9.09；手写数字略模糊
\fieldvalue{\handwritten{9.09}} &
\multicolumn{2}{c|}{细胞数量（个细胞/罐）}&
% #VALUE_ID: LEX-P0035-V0019
% #FIELD_VALUE: 细胞数量
% #TODO #HANDWRITTEN: 7.4$\times10^8$；单位及指数按表头识别
\fieldvalue{\handwritten{$7.4\times10^8$}} &
\multicolumn{2}{c|}{}\\
\hline
\multicolumn{7}{|p{0.97\linewidth}|}{
计算要求：样品密度为 3 次活细胞密度结果取平均值乘以 1\%人白蛋白化消毒液最终体积，再除以生物反应器细胞液体积；细胞数量为样品密度乘以发酵罐体积。
}\\
\hline
备注&
\multicolumn{6}{l|}{
% #VALUE_ID: LEX-P0035-V0020
% #FIELD_VALUE: 备注
% #HANDWRITTEN: /
\fieldvalue{\handwritten{/}}}\\
\hline
实验人/日期&
\multicolumn{6}{l|}{
% #VALUE_ID: LEX-P0035-V0021
% #FIELD_VALUE: 实验人
% #TODO #HANDWRITTEN: 彭??；姓名手写难以辨认
\fieldvalue{\handwritten{彭??}}
\quad
% #VALUE_ID: LEX-P0035-V0022
% #FIELD_VALUE: 实验日期
% #HANDWRITTEN: 2024.08.18
\fieldvalue{\handwritten{2024.08.18}}}\\
\hline
复核人/日期&
\multicolumn{6}{l|}{
% #VALUE_ID: LEX-P0035-V0023
% #FIELD_VALUE: 复核人
% #TODO #HANDWRITTEN: 石??；姓名手写难以辨认
\fieldvalue{\handwritten{石??}}
\quad
% #VALUE_ID: LEX-P0035-V0024
% #FIELD_VALUE: 复核日期
% #HANDWRITTEN: 2025.08.18
\fieldvalue{\handwritten{2025.08.18}}}\\
\hline
\end{tabular}

\end{minipage}

\begin{center}
第 3 页 共 5 页
\end{center}
% LEXOID_PAGE_COMPLETED: 35/70

\newpage

\begin{center}
\textbf{生物反应器细胞液检测操作记录}
\end{center}

\noindent 文件编码：REC-YZ-SOP-QC-099-006-a

\begin{center}
\begin{tabular}{|p{0.14\linewidth}|p{0.38\linewidth}|p{0.14\linewidth}|p{0.25\linewidth}|}
\hline
\multicolumn{4}{|c|}{细胞生长状态检验}\\
\hline
样品名称 &
% #VALUE_ID: LEX-P0036-V0001
% #FIELD_VALUE: 样品名称
\fieldvalue{凝生物反应器细胞液} &
样品规格 &
% #VALUE_ID: LEX-P0036-V0002
% #FIELD_VALUE: 样品规格
\fieldvalue{20 ml/袋}\\
\hline
样品批号 &
% #VALUE_ID: LEX-P0036-V0003
% #FIELD_VALUE: 样品批号
% TODO #HANDWRITTEN: [批号难以辨认]
\fieldvalue{\handwritten{[批号难以辨认]}} &
取样日期 &
% #VALUE_ID: LEX-P0036-V0004
% #FIELD_VALUE: 取样日期
% #HANDWRITTEN: 2024年03月18日
\fieldvalue{\handwritten{2024年03月18日}}\\
\hline
\end{tabular}
\end{center}

\textbf{1.\quad 仪器设备及试剂：}

\textbf{1.1\quad 仪器设备：}

\begin{center}
\begin{tabular}{|p{0.20\linewidth}|p{0.17\linewidth}|p{0.17\linewidth}|p{0.19\linewidth}|p{0.19\linewidth}|}
\hline
仪器名称 & 厂家 & 型号 & 设备编号 & 检查确认\\
\hline
% #VALUE_ID: LEX-P0036-V0005
% #FIELD_VALUE: 仪器名称
\fieldvalue{荧光显微镜} &
% #VALUE_ID: LEX-P0036-V0006
% #FIELD_VALUE: 厂家
\fieldvalue{徕卡} &
% #VALUE_ID: LEX-P0036-V0007
% #FIELD_VALUE: 型号
\fieldvalue{Dmil LED} &
% #VALUE_ID: LEX-P0036-V0008
% #FIELD_VALUE: 设备编号
\fieldvalue{1431404} &
% #VALUE_ID: LEX-P0036-V0009
% #FIELD_VALUE: 正常
\fieldvalue{\checkboxfield{checked}}正常\quad
% #VALUE_ID: LEX-P0036-V0010
% #FIELD_VALUE: 异常
\fieldvalue{\checkboxfield{unchecked}}异常\\
\hline
\end{tabular}
\end{center}

\textbf{1.2\quad 试剂耗材}

\begin{center}
\begin{tabular}{|p{0.25\linewidth}|p{0.27\linewidth}|p{0.18\linewidth}|p{0.15\linewidth}|p{0.15\linewidth}|}
\hline
名称 & 厂家 & 货号 & 批号 & 有效期至\\
\hline
% #VALUE_ID: LEX-P0036-V0011
% #FIELD_VALUE: 名称
\fieldvalue{细胞活力检测试剂盒} &
% #VALUE_ID: LEX-P0036-V0012
% #FIELD_VALUE: 厂家
\fieldvalue{江苏凯基生物技术股份有限公司} &
% #VALUE_ID: LEX-P0036-V0013
% #FIELD_VALUE: 货号
\fieldvalue{KGA9501-1000} &
% #VALUE_ID: LEX-P0036-V0014
% #FIELD_VALUE: 批号
% TODO #HANDWRITTEN: 20250410; handwriting partially unclear
\fieldvalue{\handwritten{20250410}} &
% #VALUE_ID: LEX-P0036-V0015
% #FIELD_VALUE: 有效期至
% TODO #HANDWRITTEN: 20251009; handwriting partially unclear
\fieldvalue{\handwritten{20251009}}\\
\hline
\end{tabular}
\end{center}

\textbf{2.\quad 操作步骤：}

\begin{center}
\begin{tabular}{|p{0.94\linewidth}|}
\hline
\multicolumn{1}{|c|}{溶液配制}\\
\hline
20$\times$的荧光染液配制：取
% #VALUE_ID: LEX-P0036-V0016
% #FIELD_VALUE: 组分A体积
% #HANDWRITTEN: 2.2
\fieldvalue{\handwritten{2.2}}~$\mu$l 组分A 和
% #VALUE_ID: LEX-P0036-V0017
% #FIELD_VALUE: 组分B体积
% #HANDWRITTEN: 20.2
\fieldvalue{\handwritten{20.2}}~$\mu$l 组分B 与
% #VALUE_ID: LEX-P0036-V0018
% #FIELD_VALUE: PBS体积
% #HANDWRITTEN: 39.6
\fieldvalue{\handwritten{39.6}}~$\mu$l PBS 混合均匀配成20$\times$的荧光染液。\\
\hline
\multicolumn{1}{|c|}{操作过程}\\
\hline
取复染后样品200$\mu$l/孔至96孔板中，取
% #VALUE_ID: LEX-P0036-V0019
% #FIELD_VALUE: 取孔数
% #HANDWRITTEN: 2
\fieldvalue{\handwritten{2}}孔；\hspace{1em}样品
% #VALUE_ID: LEX-P0036-V0020
% #FIELD_VALUE: 样品体积
% #HANDWRITTEN: 20
\fieldvalue{\handwritten{20}}~$\mu$l/孔\\
\hspace{1em}每个样品取
% #VALUE_ID: LEX-P0036-V0021
% #FIELD_VALUE: 每个样品取孔数
% #HANDWRITTEN: 2
\fieldvalue{\handwritten{2}}孔\\
\hline
自然沉降后，吸弃上清，注意不要吸到沉淀样品；
% #VALUE_ID: LEX-P0036-V0022
% #FIELD_VALUE: 吸弃上清
\fieldvalue{\checkboxfield{checked}}吸弃上清\\
\hline
取20$\times$的荧光染液加入到96孔板中，200$\mu$l/孔，避光孵育30--60min；
\hspace{1em}20$\times$的荧光染液
% #VALUE_ID: LEX-P0036-V0023
% #FIELD_VALUE: 荧光染液加入体积
% #HANDWRITTEN: 20
\fieldvalue{\handwritten{20}}~$\mu$l/孔\\
\hspace{1em}避光孵育
% #VALUE_ID: LEX-P0036-V0024
% #FIELD_VALUE: 避光孵育起始时间
% #HANDWRITTEN: 11
\fieldvalue{\handwritten{11}}~$\text{min}\sim$
% #VALUE_ID: LEX-P0036-V0025
% #FIELD_VALUE: 避光孵育结束时间
% #HANDWRITTEN: 14
\fieldvalue{\handwritten{14}}\\
\hline
染色结束后，吸弃上清，取PBS加入到96孔板中，200$\mu$l/孔，吸弃上清，注意不要吸到沉淀样品，清洗3次，之后每孔加入100$\mu$l PBS进行拍照。
\\[-2.2ex]
\hspace{1em}加入PBS
% #VALUE_ID: LEX-P0036-V0026
% #FIELD_VALUE: PBS加入体积
% #HANDWRITTEN: 200
\fieldvalue{\handwritten{200}}~$\mu$l/孔；清洗
% #VALUE_ID: LEX-P0036-V0027
% #FIELD_VALUE: 清洗次数
% #HANDWRITTEN: 2
\fieldvalue{\handwritten{2}}次\\
\hline
将荧光显微镜调至4$\times$镜，软件中调节至相应物镜参数，调节到适合的光源和曝光，拍摄4张清晰图像。\\
进一步一张清晰图，打印并粘贴在实验记录上。\\
\hline
数据保存路径：%
% #VALUE_ID: LEX-P0036-V0028
% #FIELD_VALUE: 数据保存路径
% TODO #HANDWRITTEN: [路径内容难以辨认]
\fieldvalue{\handwritten{[路径内容难以辨认]}}\\
\hline
\multicolumn{1}{|c|}{检测结果}\\
\hline
\end{tabular}
\end{center}

\begin{center}
第4页\quad 共5页
\end{center}
% LEXOID_PAGE_COMPLETED: 36/70

\newpage

\noindent
\begin{minipage}{\linewidth}
\small
\raggedright
铂生卓越生物科技（北京）有限公司\\
\textit{PLATINUM LIFE SCIENCE BIOTECH}\\[2pt]
\hfill 文件编码：REC-YZ-SOP-QC-99-006-a
\end{minipage}

\vspace{0.3cm}

\noindent
\begin{tabular}{|p{0.94\linewidth}|}
\hline
\vspace{15.5cm}

\begin{center}
% TODO: Embedded microscopy image visible on the page.
\fbox{\parbox[c][3.0cm][c]{6.8cm}{\centering
\textit{显微图像}
}}
\end{center}

\begin{center}
% #VALUE_ID: LEX-P0037-V0001
% #HANDWRITTEN: R?...
\handwritten{R?...}
\hspace{1.5cm}
% #VALUE_ID: LEX-P0037-V0002
% #HANDWRITTEN: M?1108
\handwritten{M?1108}
\end{center}
\\[0.2cm]
\hline
\end{tabular}

\vspace{-0.02cm}

\noindent
\begin{tabular}{|p{1.7cm}|p{10.9cm}|}
\hline
备注 &
% #VALUE_ID: LEX-P0037-V0003
% #FIELD_VALUE: 备注
% #HANDWRITTEN: /
\fieldvalue{\handwritten{/}}\\
\hline
实验人/日期 &
% #VALUE_ID: LEX-P0037-V0004
% #FIELD_VALUE: 实验人/日期
% #HANDWRITTEN: 2025.8.18
\fieldvalue{\handwritten{2025.8.18}}\\
\hline
复核人/日期 &
% #VALUE_ID: LEX-P0037-V0005
% #FIELD_VALUE: 复核人/日期
% #HANDWRITTEN: 2025.08.18
\fieldvalue{\handwritten{2025.08.18}}\\
\hline
\end{tabular}

\vfill

\begin{center}
\small 第 5 页 共 5 页
\end{center}
% LEXOID_PAGE_COMPLETED: 37/70

\newpage

Assay: \fieldvalue{MSC-1}

Cell Type F1: \fieldvalue{MSC (AO)}

Cell Type F2: \fieldvalue{MSC (PI)}

Sample ID: \fieldvalue{QC-JS-QY0402-A31Z220250806-20250818-B01-24H-1}

Dilution: \fieldvalue{2.00}

\noindent\rule{\linewidth}{0.4pt}

Results:

\medskip

\begin{tabular}{@{}p{0.29\linewidth}p{0.34\linewidth}p{0.30\linewidth}@{}}
\textbf{Count} & \textbf{Concentration} & \textbf{Mean Diameter} \\
\underline{\hspace{0.23\linewidth}} &
\underline{\hspace{0.25\linewidth}} &
\underline{\hspace{0.22\linewidth}} \\[2pt]
Total: \fieldvalue{210} &
\fieldvalue{$7.47 \times 10^{5}$ cells/mL} &
\fieldvalue{16.4 micron} \\
Live: \fieldvalue{187} &
\fieldvalue{$6.65 \times 10^{5}$ cells/mL} &
\fieldvalue{16.5 microns} \\
Dead: \fieldvalue{23} &
\fieldvalue{$8.20 \times 10^{4}$ cells/mL} &
\fieldvalue{15.6 micron} \\
Viability: \fieldvalue{89.0\%} & & \\
\end{tabular}

\bigskip

\begin{center}
\begin{tabular}{@{}cc@{}}
\fbox{\parbox[c][0.22\textheight][c]{0.39\linewidth}{\centering\% TODO: Insert upper-left image from PDF.}} &
\fbox{\parbox[c][0.22\textheight][c]{0.39\linewidth}{\centering\% TODO: Insert upper-right image from PDF.}} \\[1em]
\fbox{\parbox[c][0.22\textheight][c]{0.39\linewidth}{\centering\% TODO: Insert lower-left image from PDF.}} &
\fbox{\parbox[c][0.22\textheight][c]{0.39\linewidth}{\centering\% TODO: Insert lower-right image from PDF.}}
\end{tabular}
\end{center}

\bigskip

操作人/审核人：
% #VALUE_ID: LEX-P0038-V0018
% #FIELD_VALUE: 操作人/审核人
% #TODO #HANDWRITTEN: illegible signature; cursive handwriting
\fieldvalue{\handwritten{[illegible signature]}}

\hfill 日期：
% #VALUE_ID: LEX-P0038-V0019
% #FIELD_VALUE: 日期
% #HANDWRITTEN: 2024.8.18
\fieldvalue{\handwritten{2024.8.18}}
% LEXOID_PAGE_COMPLETED: 38/70

\newpage

\noindent Assay: MSC-1

\medskip
\noindent Cell Type F1: MSC (AO)\\
\noindent Cell Type F2: MSC (PI)

\medskip
\noindent Sample ID: QC-JS-QY0402-A31Z220250806-20250818-B01-24H-2\\
\noindent Dilution: 2.00

\medskip
\hrule
\medskip

\noindent Results:

\medskip

\begin{tabular}{p{0.29\linewidth}p{0.34\linewidth}p{0.29\linewidth}}
\textbf{Count} & \textbf{Concentration} & \textbf{Mean Diameter} \\
\texttt{=============} & \texttt{=============} & \texttt{=============} \\
Total: 256 & $9.10\times10^{5}$ cells/mL & 16.5 micron \\
Live: 234 & $8.32\times10^{5}$ cells/mL & 16.7 microns \\
Dead: 22 & $7.80\times10^{4}$ cells/mL & 15.0 microns \\
Viability: 91.4\% & & \\
\end{tabular}

\medskip

\begin{center}
\begin{tabular}{@{}p{0.43\linewidth}@{\hspace{0.04\linewidth}}p{0.43\linewidth}@{}}
% TODO: Image visible on the source page; no source filename is available.
\fbox{\rule{0pt}{0.27\textheight}\rule{0.42\linewidth}{0pt}} &
% TODO: Image visible on the source page; no source filename is available.
\fbox{\rule{0pt}{0.27\textheight}\rule{0.42\linewidth}{0pt}} \\[1em]
% TODO: Image visible on the source page; no source filename is available.
\fbox{\rule{0pt}{0.27\textheight}\rule{0.42\linewidth}{0pt}} &
% TODO: Image visible on the source page; no source filename is available.
\fbox{\rule{0pt}{0.27\textheight}\rule{0.42\linewidth}{0pt}}
\end{tabular}
\end{center}

\vfill

\noindent 操作人/审核人：
% #VALUE_ID: LEX-P0039-V0001
% #FIELD_VALUE: 操作人/审核人
% #TODO #HANDWRITTEN: illegible signature; stylized handwritten name
\fieldvalue{\handwritten{illegible signature}}
\hfill
日期：
% #VALUE_ID: LEX-P0039-V0002
% #FIELD_VALUE: 日期
% #HANDWRITTEN: 2024.8.18
\fieldvalue{\handwritten{2024.8.18}}
% LEXOID_PAGE_COMPLETED: 39/70

\newpage

\noindent
\textbf{Assay:}
% #VALUE_ID: LEX-P0040-V0001
% #FIELD_VALUE: Assay
\fieldvalue{MSC-1}

\vspace{0.5em}
\noindent
\textbf{Cell Type F1:}
% #VALUE_ID: LEX-P0040-V0002
% #FIELD_VALUE: Cell Type F1
\fieldvalue{MSC (AO)}

\vspace{0.25em}
\noindent
\textbf{Cell Type F2:}
% #VALUE_ID: LEX-P0040-V0003
% #FIELD_VALUE: Cell Type F2
\fieldvalue{MSC (PI)}

\vspace{0.5em}
\noindent
\textbf{Sample ID:}
% #VALUE_ID: LEX-P0040-V0004
% #FIELD_VALUE: Sample ID
\fieldvalue{\texttt{QC-JS-QY0402-A31Z2202508006-20250818-B01-24H-3}}

\vspace{0.25em}
\noindent
\textbf{Dilution:}
% #VALUE_ID: LEX-P0040-V0005
% #FIELD_VALUE: Dilution
\fieldvalue{2.00}

\vspace{1em}
\noindent\rule{\linewidth}{0.4pt}

\noindent Results:

\vspace{0.5em}
\small
\begin{tabular}{@{}p{0.29\linewidth}p{0.35\linewidth}p{0.29\linewidth}@{}}
\textbf{Count} & \textbf{Concentration} & \textbf{Mean Diameter} \\[-0.25em]
\texttt{=============} & \texttt{=============} & \texttt{=============} \\[0.25em]

Total:
% #VALUE_ID: LEX-P0040-V0006
% #FIELD_VALUE: Total count
\fieldvalue{213}
&
% #VALUE_ID: LEX-P0040-V0010
% #FIELD_VALUE: Total concentration
\fieldvalue{$7.57 \times 10^{5}$ cells/mL}
&
% #VALUE_ID: LEX-P0040-V0013
% #FIELD_VALUE: Total mean diameter
\fieldvalue{16.4 micron}
\\

Live:
% #VALUE_ID: LEX-P0040-V0007
% #FIELD_VALUE: Live count
\fieldvalue{187}
&
% #VALUE_ID: LEX-P0040-V0011
% #FIELD_VALUE: Live concentration
\fieldvalue{$6.66 \times 10^{5}$ cells/mL}
&
% #VALUE_ID: LEX-P0040-V0014
% #FIELD_VALUE: Live mean diameter
\fieldvalue{16.9 microns}
\\

Dead:
% #VALUE_ID: LEX-P0040-V0008
% #FIELD_VALUE: Dead count
\fieldvalue{26}
&
% #VALUE_ID: LEX-P0040-V0012
% #FIELD_VALUE: Dead concentration
\fieldvalue{$9.10 \times 10^{4}$ cells/mL}
&
% #VALUE_ID: LEX-P0040-V0015
% #FIELD_VALUE: Dead mean diameter
\fieldvalue{13.0 microns}
\\

Viability:
% #VALUE_ID: LEX-P0040-V0009
% #FIELD_VALUE: Viability
\fieldvalue{88.0\%}
&
&
\end{tabular}
\normalsize

\vspace{1.5em}

\begin{center}
\begin{tabular}{@{}c@{\hspace{1.5em}}c@{}}
% TODO: Image visible on page; source filename unavailable.
\fbox{\rule{0pt}{0.27\textheight}\rule{0.43\linewidth}{0pt}}
&
% TODO: Image visible on page; source filename unavailable.
\fbox{\rule{0pt}{0.27\textheight}\rule{0.43\linewidth}{0pt}}
\\[1em]
% TODO: Image visible on page; source filename unavailable.
\fbox{\rule{0pt}{0.27\textheight}\rule{0.43\linewidth}{0pt}}
&
% TODO: Image visible on page; source filename unavailable.
\fbox{\rule{0pt}{0.27\textheight}\rule{0.43\linewidth}{0pt}}
\end{tabular}
\end{center}

\vfill

\noindent
操作人/审核人：
% #VALUE_ID: LEX-P0040-V0016
% #FIELD_VALUE: 操作人/审核人
% #TODO #HANDWRITTEN: illegible signature; handwritten signature is not clearly legible
\fieldvalue{\handwritten{illegible signature}}
\hfill
日期：
% #VALUE_ID: LEX-P0040-V0017
% #FIELD_VALUE: 日期
% #TODO #HANDWRITTEN: 2025/9/18; handwritten date is partially unclear
\fieldvalue{\handwritten{2025/9/18}}
% LEXOID_PAGE_COMPLETED: 40/70

\newpage

\begin{center}
\textbf{铂金生物医药科技（北京）有限公司}\\
PLATINUM LIFE EXCELLENCE BIOTECH

\vspace{0.5em}
\textbf{原件请验单}
\end{center}

\noindent 文件编码：REC-YZ-SMP-QC-01-001 \hfill 版本号：1.4

\begin{center}
\textcolor{red}{受控文件}
\end{center}

\small
\renewcommand{\arraystretch}{1.45}
\begin{tabular}{|p{0.14\linewidth}|p{0.29\linewidth}|p{0.14\linewidth}|p{0.29\linewidth}|}
\hline
样品名称 &
% #VALUE_ID: LEX-P0041-V0001
% #FIELD_VALUE: 样品名称
% #TODO #HANDWRITTEN: 二级梯度胶囊；字迹较模糊
\fieldvalue{\handwritten{二级梯度胶囊}} &
样品代码 &
% #VALUE_ID: LEX-P0041-V0002
% #FIELD_VALUE: 样品代码
% #TODO #HANDWRITTEN: QX0402；首字母及末位字迹较模糊
\fieldvalue{\handwritten{QX0402}}
\\ \hline

样品批号 &
% #VALUE_ID: LEX-P0041-V0003
% #FIELD_VALUE: 样品批号
% #TODO #HANDWRITTEN: A3122252006；部分字符难以辨认
\fieldvalue{\handwritten{A3122252006}} &
样品规格/数量 &
% #VALUE_ID: LEX-P0041-V0004
% #FIELD_VALUE: 样品规格/数量
% #TODO #HANDWRITTEN: 42\,ml/数/1袋；字迹较模糊
\fieldvalue{\handwritten{42\,ml/数/1袋}}
\\ \hline

\multirow{2}{*}{样品类别} &
% #VALUE_ID: LEX-P0041-V0005
% #FIELD_VALUE: 生产过程中检验样品
\fieldvalue{\checkboxfield{checked}} 生产过程中检验样品
&
% #VALUE_ID: LEX-P0041-V0006
% #FIELD_VALUE: 特色包装产品
\fieldvalue{\checkboxfield{unchecked}} 特色包装产品
&
\\
&
% #VALUE_ID: LEX-P0041-V0007
% #FIELD_VALUE: 成品
\fieldvalue{\checkboxfield{unchecked}} 成品
&
% #VALUE_ID: LEX-P0041-V0008
% #FIELD_VALUE: 其他
\fieldvalue{\checkboxfield{unchecked}} 其他（\hspace{2em}）
&
\\ \hline

检验类型 &
% #VALUE_ID: LEX-P0041-V0009
% #FIELD_VALUE: 初验
\fieldvalue{\checkboxfield{checked}} 初验
&
% #VALUE_ID: LEX-P0041-V0010
% #FIELD_VALUE: 复验
\fieldvalue{\checkboxfield{unchecked}} 复验
\quad
% #VALUE_ID: LEX-P0041-V0011
% #FIELD_VALUE: 退货
\fieldvalue{\checkboxfield{unchecked}} 退货
\quad
% #VALUE_ID: LEX-P0041-V0012
% #FIELD_VALUE: 稳定性
\fieldvalue{\checkboxfield{unchecked}} 稳定性
\quad
% #VALUE_ID: LEX-P0041-V0013
% #FIELD_VALUE: 留样
\fieldvalue{\checkboxfield{unchecked}} 留样
&
\\ \hline

请验目的 &
% #VALUE_ID: LEX-P0041-V0014
% #FIELD_VALUE: 全检
\fieldvalue{\checkboxfield{checked}} 全检
\quad
% #VALUE_ID: LEX-P0041-V0015
% #FIELD_VALUE: 留样
\fieldvalue{\checkboxfield{unchecked}} 留样
\quad
% #VALUE_ID: LEX-P0041-V0016
% #FIELD_VALUE: 其他
\fieldvalue{\checkboxfield{unchecked}} 其他（\hspace{2em}）
&
& 
\\ \hline

请验部门 &
% #VALUE_ID: LEX-P0041-V0017
% #FIELD_VALUE: 生产部
\fieldvalue{\checkboxfield{checked}} 生产部
\quad
% #VALUE_ID: LEX-P0041-V0018
% #FIELD_VALUE: 物料管理部
\fieldvalue{\checkboxfield{unchecked}} 物料管理部
\quad
% #VALUE_ID: LEX-P0041-V0019
% #FIELD_VALUE: 其他
\fieldvalue{\checkboxfield{unchecked}} 其他（\hspace{2em}）
&
&
\\ \hline

样品情况 &
% #VALUE_ID: LEX-P0041-V0020
% #FIELD_VALUE: 随请验单一起提供
\fieldvalue{\checkboxfield{checked}} 随请验单一起提供
\quad
% #VALUE_ID: LEX-P0041-V0021
% #FIELD_VALUE: 未随请验单一起提供
\fieldvalue{\checkboxfield{unchecked}} 未随请验单一起提供
&
&
\\ \hline

请验人 &
% #VALUE_ID: LEX-P0041-V0022
% #FIELD_VALUE: 请验人
% #TODO #HANDWRITTEN: 唐春玲；字迹较模糊
\fieldvalue{\handwritten{唐春玲}} &
请验日期 &
% #VALUE_ID: LEX-P0041-V0023
% #FIELD_VALUE: 请验日期
% #HANDWRITTEN: 2025.08.7
\fieldvalue{\handwritten{2025.08.7}}
\\ \hline

收验人 &
% #VALUE_ID: LEX-P0041-V0024
% #FIELD_VALUE: 收验人
% #TODO #HANDWRITTEN: 胡梦芳；字迹较模糊
\fieldvalue{\handwritten{胡梦芳}} &
收验日期 &
% #VALUE_ID: LEX-P0041-V0025
% #FIELD_VALUE: 收验日期
% #HANDWRITTEN: 2025.08.7
\fieldvalue{\handwritten{2025.08.7}}
\\ \hline

备注 &
% #VALUE_ID: LEX-P0041-V0026
% #FIELD_VALUE: 备注
% #TODO #HANDWRITTEN: 0h 9ml；单位及数字字迹较模糊
\fieldvalue{\handwritten{0h\quad 9\,ml}} &
&
\\ \hline
\end{tabular}
\normalsize
% LEXOID_PAGE_COMPLETED: 41/70

\newpage

\begin{center}
\textbf{取样记录}

\vspace{0.5em}

\begin{tabular}{|p{0.14\linewidth}|p{0.36\linewidth}|p{0.14\linewidth}|p{0.26\linewidth}|}
\hline
名称 &
% #VALUE_ID: LEX-P0042-V0001
% #FIELD_VALUE: 名称
% #TODO #HANDWRITTEN: 冻干粉针剂……；字迹部分难以辨认
\fieldvalue{\handwritten{冻干粉针剂……}} &
代码 &
% #VALUE_ID: LEX-P0042-V0002
% #FIELD_VALUE: 代码
% #TODO #HANDWRITTEN: Q7?00?；字迹难以辨认
\fieldvalue{\handwritten{Q7?00?}} \\
\hline
规格 &
% #VALUE_ID: LEX-P0042-V0003
% #FIELD_VALUE: 规格
% #TODO #HANDWRITTEN: 20 2ml/盒；部分字迹不清
\fieldvalue{\handwritten{20 2ml/盒}} &
总数量 &
% #VALUE_ID: LEX-P0042-V0004
% #FIELD_VALUE: 总数量
% #HANDWRITTEN: 1
\fieldvalue{\handwritten{1}} \\
\hline
批号 &
% #VALUE_ID: LEX-P0042-V0005
% #FIELD_VALUE: 批号
% #TODO #HANDWRITTEN: A?3?m?；字迹难以辨认
\fieldvalue{\handwritten{A?3?m?}} &
进厂批号 &
% #VALUE_ID: LEX-P0042-V0006
% #FIELD_VALUE: 进厂批号
% #HANDWRITTEN: /
\fieldvalue{\handwritten{/}} \\
\hline
供货单位/厂家 &
\multicolumn{3}{p{0.76\linewidth}|}{
% #VALUE_ID: LEX-P0042-V0007
% #FIELD_VALUE: 供货单位/厂家
% #HANDWRITTEN: /
\fieldvalue{\handwritten{/}}
} \\
\hline
\multicolumn{4}{|p{0.90\linewidth}|}{
取样目的\quad
正常检验
% #VALUE_ID: LEX-P0042-V0008
% #FIELD_VALUE: 取样目的－正常检验
% #HANDWRITTEN: ✓
\fieldvalue{\handwritten{\checkboxfield{checked}}}
\quad
复检
% #VALUE_ID: LEX-P0042-V0009
% #FIELD_VALUE: 取样目的－复检
\fieldvalue{\checkboxfield{unchecked}}
\quad
稳定性
% #VALUE_ID: LEX-P0042-V0010
% #FIELD_VALUE: 取样目的－稳定性
\fieldvalue{\checkboxfield{unchecked}}

\par\noindent
\hspace*{4.7em}
日常检验
% #VALUE_ID: LEX-P0042-V0011
% #FIELD_VALUE: 取样目的－日常检验
\fieldvalue{\checkboxfield{unchecked}}
\quad
留样
% #VALUE_ID: LEX-P0042-V0012
% #FIELD_VALUE: 取样目的－留样
\fieldvalue{\checkboxfield{unchecked}}
\quad
其他：
% #VALUE_ID: LEX-P0042-V0013
% #FIELD_VALUE: 取样目的－其他
\fieldvalue{\checkboxfield{unchecked}}
} \\
\hline
\multicolumn{4}{|p{0.90\linewidth}|}{
取样地点\quad
产品出口
% #VALUE_ID: LEX-P0042-V0014
% #FIELD_VALUE: 取样地点－产品出口
% #HANDWRITTEN: ✓
\fieldvalue{\handwritten{\checkboxfield{checked}}}
\quad
物料库
% #VALUE_ID: LEX-P0042-V0015
% #FIELD_VALUE: 取样地点－物料库
\fieldvalue{\checkboxfield{unchecked}}
\quad
细胞库
% #VALUE_ID: LEX-P0042-V0016
% #FIELD_VALUE: 取样地点－细胞库
\fieldvalue{\checkboxfield{unchecked}}
\quad
包装间
% #VALUE_ID: LEX-P0042-V0017
% #FIELD_VALUE: 取样地点－包装间
\fieldvalue{\checkboxfield{unchecked}}
} \\
\hline
\multicolumn{4}{|p{0.90\linewidth}|}{
贮存条件\quad
2--8$^\circ$C
% #VALUE_ID: LEX-P0042-V0018
% #FIELD_VALUE: 贮存条件－2--8$^\circ$C
% #HANDWRITTEN: ✓
\fieldvalue{\handwritten{\checkboxfield{checked}}}
\quad
$-20^\circ$C
% #VALUE_ID: LEX-P0042-V0019
% #FIELD_VALUE: 贮存条件－$-20^\circ$C
\fieldvalue{\checkboxfield{unchecked}}
\quad
$-80^\circ$C
% #VALUE_ID: LEX-P0042-V0020
% #FIELD_VALUE: 贮存条件－$-80^\circ$C
\fieldvalue{\checkboxfield{unchecked}}
\quad
常温
% #VALUE_ID: LEX-P0042-V0021
% #FIELD_VALUE: 贮存条件－常温
\fieldvalue{\checkboxfield{unchecked}}
\quad
其他：
% #VALUE_ID: LEX-P0042-V0022
% #FIELD_VALUE: 贮存条件－其他
\fieldvalue{\checkboxfield{unchecked}}
\underline{\hspace{1.5em}}
} \\
\hline
\multicolumn{4}{|p{0.90\linewidth}|}{
盛装容器\quad
保湿箱
% #VALUE_ID: LEX-P0042-V0023
% #FIELD_VALUE: 盛装容器－保湿箱
% #HANDWRITTEN: ✓
\fieldvalue{\handwritten{\checkboxfield{checked}}}
\quad
泡沫箱
% #VALUE_ID: LEX-P0042-V0024
% #FIELD_VALUE: 盛装容器－泡沫箱
\fieldvalue{\checkboxfield{unchecked}}
\quad
手提篮
% #VALUE_ID: LEX-P0042-V0025
% #FIELD_VALUE: 盛装容器－手提篮
\fieldvalue{\checkboxfield{unchecked}}
\quad
其他：
% #VALUE_ID: LEX-P0042-V0026
% #FIELD_VALUE: 盛装容器－其他
\fieldvalue{\checkboxfield{unchecked}}
\underline{\hspace{1.5em}}
} \\
\hline
\multicolumn{4}{|p{0.90\linewidth}|}{
外观复核\quad
合格
% #VALUE_ID: LEX-P0042-V0027
% #FIELD_VALUE: 外观复核－合格
% #HANDWRITTEN: ✓
\fieldvalue{\handwritten{\checkboxfield{checked}}}
\quad
不合格
% #VALUE_ID: LEX-P0042-V0028
% #FIELD_VALUE: 外观复核－不合格
\fieldvalue{\checkboxfield{unchecked}}
} \\
\hline
取样量 &
% #VALUE_ID: LEX-P0042-V0029
% #FIELD_VALUE: 取样量
% #HANDWRITTEN: 1
\fieldvalue{\handwritten{1}} &
检验量 &
% #VALUE_ID: LEX-P0042-V0030
% #FIELD_VALUE: 检验量
% #HANDWRITTEN: 1
\fieldvalue{\handwritten{1}} \\
\hline
留样量 &
% #VALUE_ID: LEX-P0042-V0031
% #FIELD_VALUE: 留样量
% #HANDWRITTEN: /
\fieldvalue{\handwritten{/}} &
\multicolumn{2}{p{0.40\linewidth}|}{} \\
\hline
取样人/日期 &
% #VALUE_ID: LEX-P0042-V0032
% #FIELD_VALUE: 取样人
% #TODO #HANDWRITTEN: 杨黎生；姓名字迹不清
\fieldvalue{\handwritten{杨黎生}} \quad
% #VALUE_ID: LEX-P0042-V0033
% #FIELD_VALUE: 取样日期
% #HANDWRITTEN: 2025.08.11
\fieldvalue{\handwritten{2025.08.11}} &
复核人/日期 &
% #VALUE_ID: LEX-P0042-V0034
% #FIELD_VALUE: 复核人
% #HANDWRITTEN: 孙志明
\fieldvalue{\handwritten{孙志明}} \quad
% #VALUE_ID: LEX-P0042-V0035
% #FIELD_VALUE: 复核日期
% #HANDWRITTEN: 2025.08.11
\fieldvalue{\handwritten{2025.08.11}} \\
\hline
备注 &
\multicolumn{3}{p{0.76\linewidth}|}{
% #VALUE_ID: LEX-P0042-V0036
% #FIELD_VALUE: 备注
% #TODO #HANDWRITTEN: 02 3.0；字迹难以辨认
\fieldvalue{\handwritten{02 3.0}}
} \\
\hline
\end{tabular}
\end{center}
% LEXOID_PAGE_COMPLETED: 42/70

\newpage

\begin{center}
\textbf{生物反应器细胞液检验记录}
\end{center}

\noindent
文件编码：REC-YZ-SOP-QC-99-006-a

\begin{center}
\begin{tabular}{|p{2.4cm}|p{5.0cm}|p{2.4cm}|p{5.0cm}|}
\hline
\multicolumn{4}{|c|}{葡萄糖浓度检验}\\
\hline
样品名称 &
% #VALUE_ID: LEX-P0043-V0001
% #FIELD_VALUE: 样品名称
\fieldvalue{二级生物反应器细胞液} &
样品规格 &
% #VALUE_ID: LEX-P0043-V0002
% #FIELD_VALUE: 样品规格
% #HANDWRITTEN: 约50 ml/袋
\fieldvalue{\handwritten{约50 ml/袋}}\\
\hline
样品批号 &
% #VALUE_ID: LEX-P0043-V0003
% #FIELD_VALUE: 样品批号
% TODO #HANDWRITTEN: A?2?20250808; handwriting is partially illegible
\fieldvalue{\handwritten{A?2?20250808}} &
取样日期 &
% #VALUE_ID: LEX-P0043-V0004
% #FIELD_VALUE: 取样日期
% #HANDWRITTEN: 2026年08月21日
\fieldvalue{\handwritten{2026年08月21日}}\\
\hline
\end{tabular}
\end{center}

\section{仪器设备、试剂耗材}

\subsection{仪器设备}

\begin{center}
\begin{tabular}{|p{2.1cm}|p{4.4cm}|p{2.0cm}|p{2.8cm}|p{3.5cm}|}
\hline
名称 & 厂家 & 型号 & 设备编号 & 检查确认\\
\hline
% #VALUE_ID: LEX-P0043-V0005
% #FIELD_VALUE: 仪器名称
\fieldvalue{血糖仪} &
% #VALUE_ID: LEX-P0043-V0006
% #FIELD_VALUE: 血糖仪厂家
\fieldvalue{三诺生物传感股份有限公司} &
% #VALUE_ID: LEX-P0043-V0007
% #FIELD_VALUE: 血糖仪型号
\fieldvalue{安稳+} &
 &
正常
% #VALUE_ID: LEX-P0043-V0008
% #FIELD_VALUE: 血糖仪检查确认--正常
% #HANDWRITTEN: ✓
\fieldvalue{\handwritten{\checkboxfield{checked}}}
\quad 异常
% #VALUE_ID: LEX-P0043-V0009
% #FIELD_VALUE: 血糖仪检查确认--异常
\fieldvalue{\checkboxfield{unchecked}}\\
\hline
% #VALUE_ID: LEX-P0043-V0010
% #FIELD_VALUE: 仪器名称
\fieldvalue{移液器} &
% #VALUE_ID: LEX-P0043-V0011
% #FIELD_VALUE: 移液器厂家
\fieldvalue{Eppendorf Corporate} &
% #VALUE_ID: LEX-P0043-V0012
% #FIELD_VALUE: 移液器型号
\fieldvalue{0.1$\sim$2.5\,\textmu l} &
% #VALUE_ID: LEX-P0043-V0013
% #FIELD_VALUE: 移液器设备编号
% #HANDWRITTEN: 072408
\fieldvalue{\handwritten{072408}} &
正常
% #VALUE_ID: LEX-P0043-V0014
% #FIELD_VALUE: 移液器检查确认--正常
% #HANDWRITTEN: ✓
\fieldvalue{\handwritten{\checkboxfield{checked}}}
\quad 异常
% #VALUE_ID: LEX-P0043-V0015
% #FIELD_VALUE: 移液器检查确认--异常
\fieldvalue{\checkboxfield{unchecked}}\\
\hline
\end{tabular}
\end{center}

\subsection{试剂耗材}

\begin{center}
\begin{tabular}{|p{3.0cm}|p{6.0cm}|p{3.0cm}|p{3.0cm}|}
\hline
名称 & 厂家 & 批号 & 有效期至\\
\hline
% #VALUE_ID: LEX-P0043-V0016
% #FIELD_VALUE: 试剂耗材名称
\fieldvalue{血糖测试条} &
% #VALUE_ID: LEX-P0043-V0017
% #FIELD_VALUE: 试剂耗材厂家
\fieldvalue{三诺生物传感股份有限公司} &
% #VALUE_ID: LEX-P0043-V0018
% #FIELD_VALUE: 试剂耗材批号
% #HANDWRITTEN: 442041-06
\fieldvalue{\handwritten{442041-06}} &
% #VALUE_ID: LEX-P0043-V0019
% #FIELD_VALUE: 试剂耗材有效期至
% #HANDWRITTEN: 2027.01.08
\fieldvalue{\handwritten{2027.01.08}}\\
\hline
\end{tabular}
\end{center}

\section{操作步骤}

\begin{center}
\begin{tabular}{|p{10.7cm}|p{5.0cm}|}
\hline
\multicolumn{2}{|c|}{操作过程}\\
\hline
取出测试条，取样品自然沉降后的上清液0.6\,\textmu l加入测试条电极槽，应确认自动发出“嘀”提示音，吸入样品完成，血糖仪自动进入倒计时，开始测试； &
取上清液：
% #VALUE_ID: LEX-P0043-V0020
% #FIELD_VALUE: 取上清液体积
% #HANDWRITTEN: 0.6
\fieldvalue{\handwritten{0.6}}\,\textmu l\\
\hline
重复以上步骤2次，共计数3次。 &
每个样品分别共计数
% #VALUE_ID: LEX-P0043-V0021
% #FIELD_VALUE: 样品计数次数
% #HANDWRITTEN: 3
\fieldvalue{\handwritten{3}} 次。\\
\hline
\multicolumn{2}{|c|}{检测结果}\\
\hline
\multicolumn{2}{|c|}{葡萄糖浓度（mmol/L）}\\
\hline
\multicolumn{2}{|c|}{\begin{tabular}{cccc}
第一次 & 第二次 & 第三次 & 平均值（保留一位小数）\\
\hline
% #VALUE_ID: LEX-P0043-V0022
% #FIELD_VALUE: 第一次葡萄糖浓度
% #HANDWRITTEN: 21.0
\fieldvalue{\handwritten{21.0}} &
% #VALUE_ID: LEX-P0043-V0023
% #FIELD_VALUE: 第二次葡萄糖浓度
% #HANDWRITTEN: 22.0
\fieldvalue{\handwritten{22.0}} &
% #VALUE_ID: LEX-P0043-V0024
% #FIELD_VALUE: 第三次葡萄糖浓度
% #HANDWRITTEN: 21.0
\fieldvalue{\handwritten{21.0}} &
% #VALUE_ID: LEX-P0043-V0025
% #FIELD_VALUE: 平均葡萄糖浓度
% #HANDWRITTEN: 21.3
\fieldvalue{\handwritten{21.3}}
\end{tabular}}\\
\hline
备注 &
\multicolumn{1}{l|}{
% #VALUE_ID: LEX-P0043-V0026
% #FIELD_VALUE: 备注
% #HANDWRITTEN: -
\fieldvalue{\handwritten{-}}}\\
\hline
\multicolumn{2}{|c|}{\begin{tabular}{llll}
实验人/日期 &
% #VALUE_ID: LEX-P0043-V0027
% #FIELD_VALUE: 实验人
% TODO #HANDWRITTEN: 张媛媛; handwriting is partially unclear
\fieldvalue{\handwritten{张媛媛}}
% #VALUE_ID: LEX-P0043-V0028
% #FIELD_VALUE: 实验日期
% #HANDWRITTEN: 2025.08.17
\fieldvalue{\handwritten{2025.08.17}} &
复核人/日期 &
% #VALUE_ID: LEX-P0043-V0029
% #FIELD_VALUE: 复核人
% #HANDWRITTEN: 杨景宇
\fieldvalue{\handwritten{杨景宇}}
% #VALUE_ID: LEX-P0043-V0030
% #FIELD_VALUE: 复核日期
% #HANDWRITTEN: 2025.08.17
\fieldvalue{\handwritten{2025.08.17}}
\end{tabular}}\\
\hline
\end{tabular}
\end{center}

\noindent
第1页 共5页
% LEXOID_PAGE_COMPLETED: 43/70

\newpage

\begin{center}
\textbf{生物反应器细胞液检验记录}
\end{center}

\noindent 文件编号：REC-YZ-SOP-QC-99-006-a

\begin{center}
\begin{tabular}{|p{2.5cm}|p{5.2cm}|p{2.5cm}|p{3.5cm}|}
\hline
\multicolumn{4}{|c|}{密度与活率检验} \\
\hline
样品名称 &
% #VALUE_ID: LEX-P0044-V0001
% #FIELD_VALUE: 样品名称
% #HANDWRITTEN: 二级生物反应器细胞液
\fieldvalue{\handwritten{二级生物反应器细胞液}} &
样品规格 &
% #VALUE_ID: LEX-P0044-V0002
% #FIELD_VALUE: 样品规格
% #HANDWRITTEN: 约20
\fieldvalue{\handwritten{约20}} ml/袋 \\
\hline
样品批号 &
% #VALUE_ID: LEX-P0044-V0003
% #FIELD_VALUE: 样品批号
% #TODO #HANDWRITTEN: A?122?220?; handwriting unclear
\fieldvalue{\handwritten{A?122?220?}} &
取样日期 &
2025年08月07日 \\
\hline
\end{tabular}
\end{center}

\section{仪器设备、试剂耗材}

\subsection{仪器设备}

\begin{center}
\small
\begin{tabular}{|p{2.5cm}|p{2.4cm}|p{2.5cm}|p{2.1cm}|p{2.1cm}|p{2.7cm}|}
\hline
名称 & 厂家 & 型号 & 设备编号 & 校准有效期至 & 检查确认 \\
\hline
% #VALUE_ID: LEX-P0044-V0004
% #FIELD_VALUE: 名称
\fieldvalue{细胞计数仪} &
% #VALUE_ID: LEX-P0044-V0005
% #FIELD_VALUE: 厂家
\fieldvalue{Nexcelom} &
% #VALUE_ID: LEX-P0044-V0006
% #FIELD_VALUE: 型号
\fieldvalue{Cellometer K2} &
% #VALUE_ID: LEX-P0044-V0007
% #FIELD_VALUE: 设备编号
\fieldvalue{1410905} &
% #VALUE_ID: LEX-P0044-V0008
% #FIELD_VALUE: 校准有效期至
% #HANDWRITTEN: 141905
\fieldvalue{\handwritten{141905}} &
正常\ 
% #VALUE_ID: LEX-P0044-V0009
% #FIELD_VALUE: 检查确认--正常
\fieldvalue{\checkboxfield{checked}}\quad
异常\ 
% #VALUE_ID: LEX-P0044-V0010
% #FIELD_VALUE: 检查确认--异常
\fieldvalue{\checkboxfield{unchecked}} \\
\hline
% #VALUE_ID: LEX-P0044-V0011
% #FIELD_VALUE: 名称
\fieldvalue{高速冷冻离心机} &
% #VALUE_ID: LEX-P0044-V0012
% #FIELD_VALUE: 厂家
\fieldvalue{Thermo} &
% #VALUE_ID: LEX-P0044-V0013
% #FIELD_VALUE: 型号
\fieldvalue{ST16R} &
% #VALUE_ID: LEX-P0044-V0014
% #FIELD_VALUE: 设备编号
\fieldvalue{1410810} &
% #VALUE_ID: LEX-P0044-V0015
% #FIELD_VALUE: 校准有效期至
% #HANDWRITTEN: 141810
\fieldvalue{\handwritten{141810}} &
正常\ 
% #VALUE_ID: LEX-P0044-V0016
% #FIELD_VALUE: 检查确认--正常
\fieldvalue{\checkboxfield{checked}}\quad
异常\ 
% #VALUE_ID: LEX-P0044-V0017
% #FIELD_VALUE: 检查确认--异常
\fieldvalue{\checkboxfield{unchecked}} \\
\hline
% #VALUE_ID: LEX-P0044-V0018
% #FIELD_VALUE: 名称
\fieldvalue{电热恒温水槽} &
% #VALUE_ID: LEX-P0044-V0019
% #FIELD_VALUE: 厂家
\fieldvalue{上海一恒} &
% #VALUE_ID: LEX-P0044-V0020
% #FIELD_VALUE: 型号
\fieldvalue{DK-8AXX} &
% #VALUE_ID: LEX-P0044-V0021
% #FIELD_VALUE: 设备编号
\fieldvalue{1411001} &
% #VALUE_ID: LEX-P0044-V0022
% #FIELD_VALUE: 校准有效期至
% #HANDWRITTEN: 2026.05.22
\fieldvalue{\handwritten{2026.05.22}} &
正常\ 
% #VALUE_ID: LEX-P0044-V0023
% #FIELD_VALUE: 检查确认--正常
\fieldvalue{\checkboxfield{checked}}\quad
异常\ 
% #VALUE_ID: LEX-P0044-V0024
% #FIELD_VALUE: 检查确认--异常
\fieldvalue{\checkboxfield{unchecked}} \\
\hline
\end{tabular}
\end{center}

\subsection{试剂耗材}

\begin{center}
\begin{tabular}{|p{3.5cm}|p{3.5cm}|p{4.5cm}|p{3.2cm}|}
\hline
名称 & 厂家 & 批号 & 有效期至 \\
\hline
% #VALUE_ID: LEX-P0044-V0025
% #FIELD_VALUE: 名称
\fieldvalue{AO/PI 染液} &
% #VALUE_ID: LEX-P0044-V0026
% #FIELD_VALUE: 厂家
\fieldvalue{Nexcelom} &
% #VALUE_ID: LEX-P0044-V0027
% #FIELD_VALUE: 批号
% #HANDWRITTEN: 322?605
\fieldvalue{\handwritten{322?605}} &
% #VALUE_ID: LEX-P0044-V0028
% #FIELD_VALUE: 有效期至
% #HANDWRITTEN: 2026.02.24
\fieldvalue{\handwritten{2026.02.24}} \\
\hline
% #VALUE_ID: LEX-P0044-V0029
% #FIELD_VALUE: 名称
\fieldvalue{0.9\%氯化钠注射液} &
% #VALUE_ID: LEX-P0044-V0030
% #FIELD_VALUE: 厂家
\fieldvalue{石家庄四药} &
% #VALUE_ID: LEX-P0044-V0031
% #FIELD_VALUE: 批号
% #HANDWRITTEN: 241?27?202
\fieldvalue{\handwritten{241?27?202}} &
% #VALUE_ID: LEX-P0044-V0032
% #FIELD_VALUE: 有效期至
% #HANDWRITTEN: 2026.10.06
\fieldvalue{\handwritten{2026.10.06}} \\
\hline
% #VALUE_ID: LEX-P0044-V0033
% #FIELD_VALUE: 名称
\fieldvalue{20\%人血白蛋白} &
% #VALUE_ID: LEX-P0044-V0034
% #FIELD_VALUE: 厂家
\fieldvalue{Baxter AG} &
% #VALUE_ID: LEX-P0044-V0035
% #FIELD_VALUE: 批号
% #HANDWRITTEN: A?25?A?9
\fieldvalue{\handwritten{A?25?A?9}} &
% #VALUE_ID: LEX-P0044-V0036
% #FIELD_VALUE: 有效期至
% #HANDWRITTEN: 2025.08.31
\fieldvalue{\handwritten{2025.08.31}} \\
\hline
\end{tabular}
\end{center}

\section{操作步骤}

\subsection*{溶液准备}

1\%人白化纳溶液配制：取5 ml 20\%人血白蛋白加入到95 ml 的0.9\%氯化钠注射液中，混匀即可。

20\%人血白蛋白
% #VALUE_ID: LEX-P0044-V0037
% #FIELD_VALUE: 20\%人血白蛋白用量
% #HANDWRITTEN: 1
\fieldvalue{\handwritten{1}} ml

0.9\%氯化钠注射液
% #VALUE_ID: LEX-P0044-V0038
% #FIELD_VALUE: 0.9\%氯化钠注射液用量
% #HANDWRITTEN: 9
\fieldvalue{\handwritten{9}} ml

\begin{center}
\begin{tabular}{|p{3.5cm}|p{6.0cm}|p{4.5cm}|}
\hline
名称 & 配制批号 & 有效期至 \\
\hline
% #VALUE_ID: LEX-P0044-V0039
% #FIELD_VALUE: 名称
\fieldvalue{5$\times$裂解液} &
% #VALUE_ID: LEX-P0044-V0040
% #FIELD_VALUE: 配制批号
% #HANDWRITTEN: 500305028-01
\fieldvalue{\handwritten{500305028-01}} &
% #VALUE_ID: LEX-P0044-V0041
% #FIELD_VALUE: 有效期至
% #HANDWRITTEN: 2025.10.27
\fieldvalue{\handwritten{2025.10.27}} \\
\hline
\end{tabular}
\end{center}

\subsection*{操作过程}

按照50 ml离心管的刻度读取样品体积，往样品管中加入对应体积的5$\times$裂解液（每4 ml样品加入1 ml 5$\times$裂解液，按照“只进不舍”原则，如26 ml样品，加入7 ml 5$\times$裂解液），充分混匀后倒入至2.3的样品中，置于37℃电热恒温水槽中裂解20 min后，取出离心管摇匀，使样品充分裂解；再置于37℃电热恒温水槽中裂解10 min得到细胞悬液。

\begin{tabular}{|p{8.3cm}|p{6.0cm}|}
\hline
按照50 ml离心管的刻度读取样品体积，往样品管中加入对应体积的5$\times$裂解液（每4 ml样品加入1 ml 5$\times$裂解液，按照“只进不舍”原则，如26 ml样品，加入7 ml 5$\times$裂解液），充分混匀后倒入至2.3的样品中，置于37℃电热恒温水槽中裂解20 min后，取出离心管摇匀，使样品充分裂解；再置于37℃电热恒温水槽中裂解10 min得到细胞悬液。 &
样品体积：
% #VALUE_ID: LEX-P0044-V0042
% #FIELD_VALUE: 样品体积
% #HANDWRITTEN: 40
\fieldvalue{\handwritten{40}} ml

裂解液
% #VALUE_ID: LEX-P0044-V0043
% #FIELD_VALUE: 裂解液体积
% #HANDWRITTEN: 10
\fieldvalue{\handwritten{10}} ml

电热恒温水槽
% #VALUE_ID: LEX-P0044-V0044
% #FIELD_VALUE: 电热恒温水槽温度
% #HANDWRITTEN: 37
\fieldvalue{\handwritten{37}} ℃

裂解
% #VALUE_ID: LEX-P0044-V0045
% #FIELD_VALUE: 裂解时间
% #TODO #HANDWRITTEN: 14:20--6:00; handwriting unclear
\fieldvalue{\handwritten{14:20--6:00}} \\
\hline
取出裂解好的细胞悬液300$\times$g离心5分钟，按照读取样品体积加入对应体积的1\%人白蛋白溶液重悬（每1.2 ml样品体积加入0.1 ml 1\%人白蛋白溶液，按照“只进不舍”原则，如26 ml吸取2.6 ml）。 &
% #VALUE_ID: LEX-P0044-V0046
% #FIELD_VALUE: 离心条件
% #HANDWRITTEN: 2000×g，5分钟
\fieldvalue{\handwritten{2000$\times$g，5分钟}}

加1\%人白蛋白溶液重悬总体积
% #VALUE_ID: LEX-P0044-V0047
% #FIELD_VALUE: 重悬总体积
% #HANDWRITTEN: 24
\fieldvalue{\handwritten{24}} ml

吸取
% #VALUE_ID: LEX-P0044-V0048
% #FIELD_VALUE: 吸取体积
% #HANDWRITTEN: 20
\fieldvalue{\handwritten{20}} $\mu$l \\
\hline
\end{tabular}

% TODO: The form continues beyond the visible portion of this page.
% LEXOID_PAGE_COMPLETED: 44/70

\newpage

\noindent\textbf{细胞计数记录}

\noindent
编号：REC-YZ-SOP-CQ-99-006-a

\medskip
\noindent
加入 $2.2\,\mathrm{ml}$ 的 AOPI 染料，反复吹打至单细胞悬液。

\noindent
撕去计数板上两层保护膜，将细胞悬液混匀，从中吸取 $20\,\mu\mathrm{l}$

\noindent
与
% #VALUE_ID: LEX-P0045-V0001
% #FIELD_VALUE: AOPI染料混合液体积
% #HANDWRITTEN: 20
\fieldvalue{\handwritten{20}}
$\mu\mathrm{l}$ AOPI 染料混合均匀。

\noindent
吸取
% #VALUE_ID: LEX-P0045-V0002
% #FIELD_VALUE: 计数体积
% #HANDWRITTEN: 20
\fieldvalue{\handwritten{20}}
$\mu\mathrm{l}$ 计数。

\noindent
每个样品均计数
% #VALUE_ID: LEX-P0045-V0003
% #FIELD_VALUE: 样品计数次数
% #HANDWRITTEN: 3
\fieldvalue{\handwritten{3}}
次。

\noindent
将 $20\,\mu\mathrm{l}$ 细胞悬液与 AOPI 染料的混合液，注入到计数板的一个检测室中。

\noindent
选择 Assay 类型，使用 AOPI 染料进行双荧光存活率检测；在主界面左侧 ID 栏中输入样本批号。

\noindent
点击主界面左下角 ``Preview'' 预览细胞图片。

\noindent
调整各器焦距直至观察到荧光通道下的细胞中心透亮、轮廓清晰。

\noindent
点击 ``count'' 按钮进行计数并自动拍摄图片。

\noindent
重复以上步骤 2 次，共计数 3 次。打印数据结果。

\medskip

\noindent
数据保存路径：
% #VALUE_ID: LEX-P0045-V0004
% #FIELD_VALUE: 数据保存路径
% TODO #HANDWRITTEN: M:\...2025...; 路径文字部分难以辨认
\fieldvalue{\handwritten{M:\ldots 2025\ldots}}

\begin{center}
\small
\begin{tabular}{|p{0.18\linewidth}|p{0.72\linewidth}|}
\hline
\multicolumn{2}{|c|}{检测结果}\\
\hline
检测结果 &
% #VALUE_ID: LEX-P0045-V0005
% #FIELD_VALUE: 一级反应
\fieldvalue{\checkboxfield{unchecked}} 一级反应、
% #VALUE_ID: LEX-P0045-V0006
% #FIELD_VALUE: 二级反应
% #HANDWRITTEN: \checkmark
\fieldvalue{\handwritten{\checkboxfield{checked}}} 二级反应、
% #VALUE_ID: LEX-P0045-V0007
% #FIELD_VALUE: 三级反应
\fieldvalue{\checkboxfield{unchecked}} 三级反应
\\
\hline
\end{tabular}

\medskip

\begin{tabular}{|p{0.11\linewidth}|p{0.14\linewidth}|p{0.15\linewidth}|p{0.13\linewidth}|p{0.13\linewidth}|p{0.13\linewidth}|p{0.13\linewidth}|}
\hline
检测次数 &
活细胞密度（个细胞/ml） &
样品密度（个细胞/ml） &
活细胞密度 CV 值（\%） &
细胞活率（\%） &
样品活率（\%） &
细胞活率 CV 值（\%）\\
\hline
第一次 &
% #VALUE_ID: LEX-P0045-V0008
% #FIELD_VALUE: 第一次活细胞密度
% #HANDWRITTEN: 6.42 x 10^5
\fieldvalue{\handwritten{$6.42\times10^{5}$}} &
&
&
% #VALUE_ID: LEX-P0045-V0009
% #FIELD_VALUE: 第一次细胞活率
% #HANDWRITTEN: 98.5
\fieldvalue{\handwritten{98.5}} &
% #VALUE_ID: LEX-P0045-V0010
% #FIELD_VALUE: 第一次样品活率
% #HANDWRITTEN: 97
\fieldvalue{\handwritten{97}} &
\\
\hline
第二次 &
% #VALUE_ID: LEX-P0045-V0011
% #FIELD_VALUE: 第二次活细胞密度
% #HANDWRITTEN: 7.60 x 10^5
\fieldvalue{\handwritten{$7.60\times10^{5}$}} &
% #VALUE_ID: LEX-P0045-V0012
% #FIELD_VALUE: 第二次样品密度
% #HANDWRITTEN: 5.93 x 10^4
\fieldvalue{\handwritten{$5.93\times10^{4}$}} &
% #VALUE_ID: LEX-P0045-V0013
% #FIELD_VALUE: 第二次活细胞密度CV值
% #HANDWRITTEN: 9
\fieldvalue{\handwritten{9}} &
% #VALUE_ID: LEX-P0045-V0014
% #FIELD_VALUE: 第二次细胞活率
% #HANDWRITTEN: 96.4
\fieldvalue{\handwritten{96.4}} &
% #VALUE_ID: LEX-P0045-V0015
% #FIELD_VALUE: 第二次样品活率
% #HANDWRITTEN: 97.2
\fieldvalue{\handwritten{97.2}} &
% #VALUE_ID: LEX-P0045-V0016
% #FIELD_VALUE: 第二次细胞活率CV值
% #HANDWRITTEN: 2
\fieldvalue{\handwritten{2}}\\
\hline
第三次 &
% #VALUE_ID: LEX-P0045-V0017
% #FIELD_VALUE: 第三次活细胞密度
% #HANDWRITTEN: 6.71 x 10^5
\fieldvalue{\handwritten{$6.71\times10^{5}$}} &
&
&
% #VALUE_ID: LEX-P0045-V0018
% #FIELD_VALUE: 第三次细胞活率
% #HANDWRITTEN: 71.1
\fieldvalue{\handwritten{71.1}} &
% #VALUE_ID: LEX-P0045-V0019
% #FIELD_VALUE: 第三次样品活率
% TODO #HANDWRITTEN: 难辨；该单元格中的蓝色手写数值模糊
\fieldvalue{\handwritten{难辨}} &
% #VALUE_ID: LEX-P0045-V0020
% #FIELD_VALUE: 第三次细胞活率CV值
% TODO #HANDWRITTEN: 难辨；该单元格中的蓝色手写数值模糊
\fieldvalue{\handwritten{难辨}}\\
\hline
\end{tabular}

\medskip

\begin{tabular}{|p{0.26\linewidth}|p{0.24\linewidth}|p{0.25\linewidth}|p{0.19\linewidth}|}
\hline
发酵罐细胞体积（ml） &
% #VALUE_ID: LEX-P0045-V0021
% #FIELD_VALUE: 发酵罐细胞体积
% #HANDWRITTEN: 1000
\fieldvalue{\handwritten{1000}} &
细胞数量（个细胞/罐） &
% #VALUE_ID: LEX-P0045-V0022
% #FIELD_VALUE: 细胞数量
% #HANDWRITTEN: 5.34 x 10^8
\fieldvalue{\handwritten{$5.34\times10^{8}$}}\\
\hline
\multicolumn{4}{|p{0.94\linewidth}|}{
计算要求：样品密度为 3 次活细胞密度结果取平均值再乘以 1\% 人白化剂溶液重量体积，再除以生物反应器细胞液体积；细胞数量为样品密度乘以发酵罐体积。
}\\
\hline
\multicolumn{4}{|p{0.94\linewidth}|}{
备注：
% #VALUE_ID: LEX-P0045-V0023
% #FIELD_VALUE: 备注
% #HANDWRITTEN: \checkmark
\fieldvalue{\handwritten{\checkmark}}
}\\
\hline
实验人/日期 &
% #VALUE_ID: LEX-P0045-V0024
% #FIELD_VALUE: 实验人/日期
% TODO #HANDWRITTEN: 石强 2025.08.1；姓名及日期部分难以完全辨认
\fieldvalue{\handwritten{石强\quad 2025.08.1}} &
\multicolumn{2}{l|}{}\\
\hline
复核人/日期 &
% #VALUE_ID: LEX-P0045-V0025
% #FIELD_VALUE: 复核人/日期
% TODO #HANDWRITTEN: 郭? 2025.08.07；姓名部分难以辨认
\fieldvalue{\handwritten{郭?\quad 2025.08.07}} &
\multicolumn{2}{l|}{}\\
\hline
\end{tabular}
\end{center}

\begin{center}
第 3 页 共 5 页
\end{center}
% LEXOID_PAGE_COMPLETED: 45/70

\newpage

\section{生物反应器细胞液检测操作记录}

\begin{center}
\begin{tabular}{|p{0.16\linewidth}|p{0.36\linewidth}|p{0.16\linewidth}|p{0.22\linewidth}|}
\hline
\multicolumn{4}{|c|}{细胞生长状态检验}\\
\hline
样品名称 &
% #VALUE_ID: LEX-P0046-V0001
% #FIELD_VALUE: 样品名称
\fieldvalue{发生物反应器细胞液} &
样品规格 & ml/袋\\
\hline
样品批号 & & 取样日期 & 年\quad 月\quad 日\\
\hline
\end{tabular}
\end{center}

\section{仪器及试剂}

\subsection{仪器设备}

\begin{center}
\begin{tabular}{|p{0.22\linewidth}|p{0.18\linewidth}|p{0.18\linewidth}|p{0.20\linewidth}|p{0.18\linewidth}|}
\hline
仪器名称 & 厂家 & 型号 & 设备编号 & 检查确认\\
\hline
% #VALUE_ID: LEX-P0046-V0002
% #FIELD_VALUE: 仪器名称
\fieldvalue{荧光显微镜} &
% #VALUE_ID: LEX-P0046-V0003
% #FIELD_VALUE: 厂家
\fieldvalue{徕卡} &
% #VALUE_ID: LEX-P0046-V0004
% #FIELD_VALUE: 型号
\fieldvalue{Dmi LED} &
% #VALUE_ID: LEX-P0046-V0005
% #FIELD_VALUE: 设备编号
\fieldvalue{1431404} &
正常\ %
% #VALUE_ID: LEX-P0046-V0006
% #FIELD_VALUE: 正常
\fieldvalue{\checkboxfield{unchecked}}
\quad 异常\ %
% #VALUE_ID: LEX-P0046-V0007
% #FIELD_VALUE: 异常
\fieldvalue{\checkboxfield{unchecked}}\\
\hline
\end{tabular}
\end{center}

\subsection{试剂耗材}

\begin{center}
\begin{tabular}{|p{0.22\linewidth}|p{0.30\linewidth}|p{0.20\linewidth}|p{0.14\linewidth}|p{0.14\linewidth}|}
\hline
名称 & 厂家 & 货号 & 批号 & 有效期至\\
\hline
% #VALUE_ID: LEX-P0046-V0008
% #FIELD_VALUE: 名称
\fieldvalue{细胞活力检测试剂盒} &
% #VALUE_ID: LEX-P0046-V0009
% #FIELD_VALUE: 厂家
\fieldvalue{江苏凯基生物技术股份有限公司} &
% #VALUE_ID: LEX-P0046-V0010
% #FIELD_VALUE: 货号
\fieldvalue{KGA9501-1000} & & \\
\hline
\end{tabular}
\end{center}

\section{操作步骤}

\begin{center}
\small
\begin{tabular}{|p{0.56\linewidth}|p{0.38\linewidth}|}
\hline
\multicolumn{2}{|c|}{溶液配制}\\
\hline
20$\times$的荧光染液配制：取\underline{\hspace{1.2cm}} ml 组分 A 和\underline{\hspace{1.2cm}} ml 组分 B，\underline{\hspace{1.2cm}} ml PBS 混合均匀配成 20$\times$ 的荧光染液。 & \\
\hline
\multicolumn{2}{|c|}{操作过程}\\
\hline
取混匀后样品 200$\mu$l/孔至 96 孔板中，取 2 孔。 &
样品\quad $\mu$l/孔，\\
& 每个样品取\underline{\hspace{1.2cm}}孔\\
\hline
自然沉降后，吸弃上清，注意不要吸到沉淀样品； &
吸弃上清\\
\hline
取 20$\times$ 的荧光染液加入到 96 孔板中，200$\mu$l/孔，避光孵育 30--60 min； &
20$\times$ 的荧光染液\underline{\hspace{1.2cm}} $\mu$l/孔\\
& 避光孵育\underline{\hspace{1.2cm}} $\sim$ \\
\hline
染色结束后，吸弃上清，取 PBS 加入到以上 96 孔板中，200$\mu$l/孔，吸弃上清，注意不要吸到沉淀样品，清洗 3 次，之后每孔加入 100$\mu$l PBS 进行拍照； &
加入 PBS\underline{\hspace{1.2cm}} $\mu$l/孔\\
& 清洗\underline{\hspace{1.2cm}}次\\
\hline
将荧光显微镜调至 4$\times$镜，软件中调节至相应物镜参数，调节到合适的光源和像距，拍摄 4 张清晰图像。 &
\\
\hline
选一张清晰图像，打印并粘贴在实验记录上。 &
\\
\hline
\end{tabular}
\end{center}

% TODO: The page contains a blue diagonal handwritten mark crossing the form.
% #VALUE_ID: LEX-P0046-V0011
% #TODO #HANDWRITTEN: 2023.09.11; diagonal handwriting is partially obscured and difficult to read
% #FIELD_VALUE: handwritten date/annotation
\fieldvalue{\handwritten{2023.09.11}}

\begin{center}
\begin{tabular}{|p{0.20\linewidth}|p{0.74\linewidth}|}
\hline
数据保存路径 & \\
\hline
\multicolumn{2}{|c|}{检测结果}\\
\hline
\end{tabular}
\end{center}
% LEXOID_PAGE_COMPLETED: 46/70

\newpage

\begin{center}
\small
\begin{tabular}{p{0.28\linewidth}p{0.34\linewidth}p{0.28\linewidth}}
\textbf{维亚生物科技（北京）有限公司} &
\centering
\begin{tabular}{|c|}
\hline
留存生物\\
原件\\
\hline
\end{tabular}
&
\raggedleft
\begin{tabular}{|c|}
\hline
留存生物\\
受控文件\\
版本号：1.4\\
\hline
\end{tabular}
\end{tabular}

\vspace{-0.2em}
\raggedright
\textit{PHENALY LIFE SCIENCES BIOTECH}
\hfill 文件编码：REC-YZ-SOP-QC-99-006-a
\end{center}

\vspace{0.3em}

\noindent
\fbox{%
\begin{minipage}[t][0.84\textheight][t]{0.94\linewidth}
\vspace{0.02\textheight}

\begin{center}
\begin{picture}(400,500)
\color{blue}
\put(5,10){\line(1,2){250}}
% TODO: The blue diagonal handwritten annotation is partially illegible.
\put(135,260){\rotatebox{55}{%
% #VALUE_ID: LEX-P0047-V0001
% #TODO #HANDWRITTEN: 张建平 2025.03.11; handwritten name is partially illegible
\handwritten{张建平\ 2025.03.11}}}
\end{picture}
\end{center}

\vfill

\begin{tabular}{|p{0.18\linewidth}|p{0.76\linewidth}|}
\hline
备注 & \\[1.8em]
\hline
实验人/日期 & \\[1.8em]
\hline
复核人/日期 & \\[1.8em]
\hline
\end{tabular}

\end{minipage}%
}

\vspace{0.8em}

\begin{center}
\small 第 5 页 共 5 页
\end{center}
% LEXOID_PAGE_COMPLETED: 47/70

\newpage

\noindent Assay: %
% #VALUE_ID: LEX-P0048-V0001
% #FIELD_VALUE: Assay
\fieldvalue{MSC-1}

\medskip
\noindent Cell Type F1: %
% #VALUE_ID: LEX-P0048-V0002
% #FIELD_VALUE: Cell Type F1
\fieldvalue{MSC (AO)}

\noindent Cell Type F2: %
% #VALUE_ID: LEX-P0048-V0003
% #FIELD_VALUE: Cell Type F2
\fieldvalue{MSC (PI)}

\medskip
\noindent Sample ID: %
% #VALUE_ID: LEX-P0048-V0004
% #FIELD_VALUE: Sample ID
\fieldvalue{QC-JS-QY0402-A31Z2202508006-20250817-B01-0H-1}

\noindent Dilution: %
% #VALUE_ID: LEX-P0048-V0005
% #FIELD_VALUE: Dilution
\fieldvalue{2.00}

\medskip
\noindent\hrulefill

\noindent Results:

\medskip
\begin{tabular}{@{}p{0.28\linewidth}p{0.34\linewidth}p{0.28\linewidth}@{}}
\textbf{Count} & \textbf{Concentration} & \textbf{Mean Diameter} \\
\texttt{============} & \texttt{============} & \texttt{============} \\[2pt]
Total: %
% #VALUE_ID: LEX-P0048-V0006
% #FIELD_VALUE: Total count
\fieldvalue{185}
&
% #VALUE_ID: LEX-P0048-V0007
% #FIELD_VALUE: Total concentration
\fieldvalue{$6.52\times10^{5}$ cells/mL}
&
% #VALUE_ID: LEX-P0048-V0008
% #FIELD_VALUE: Total mean diameter
\fieldvalue{14.0 micron}
\\
Live: %
% #VALUE_ID: LEX-P0048-V0009
% #FIELD_VALUE: Live count
\fieldvalue{182}
&
% #VALUE_ID: LEX-P0048-V0010
% #FIELD_VALUE: Live concentration
\fieldvalue{$6.42\times10^{5}$ cells/mL}
&
% #VALUE_ID: LEX-P0048-V0011
% #FIELD_VALUE: Live mean diameter
\fieldvalue{14.1 micron}
\\
Dead: %
% #VALUE_ID: LEX-P0048-V0012
% #FIELD_VALUE: Dead count
\fieldvalue{3}
&
% #VALUE_ID: LEX-P0048-V0013
% #FIELD_VALUE: Dead concentration
\fieldvalue{$1.00\times10^{4}$ cells/mL}
&
% #VALUE_ID: LEX-P0048-V0014
% #FIELD_VALUE: Dead mean diameter
\fieldvalue{11.1 micron}
\\
Viability: %
% #VALUE_ID: LEX-P0048-V0015
% #FIELD_VALUE: Viability
\fieldvalue{98.5\%}
&
&
\end{tabular}

\medskip

% TODO: Four microscopy images are visible on this page, but no image filenames are available.
\begin{center}
\begin{minipage}[t]{0.43\linewidth}
\centering
\fbox{\parbox[c][4.2cm][c]{0.95\linewidth}{\centering TODO: upper-left microscopy image}}
\end{minipage}
\hfill
\begin{minipage}[t]{0.43\linewidth}
\centering
\fbox{\parbox[c][4.2cm][c]{0.95\linewidth}{\centering TODO: upper-right microscopy image}}
\end{minipage}

\vspace{0.5cm}

\begin{minipage}[t]{0.43\linewidth}
\centering
\fbox{\parbox[c][4.2cm][c]{0.95\linewidth}{\centering TODO: lower-left microscopy image}}
\end{minipage}
\hfill
\begin{minipage}[t]{0.43\linewidth}
\centering
\fbox{\parbox[c][4.2cm][c]{0.95\linewidth}{\centering TODO: lower-right microscopy image}}
\end{minipage}
\end{center}

\vfill

\noindent 操作人/审核人：%
% #VALUE_ID: LEX-P0048-V0016
% #FIELD_VALUE: 操作人/审核人
% TODO #HANDWRITTEN: 石孟男; handwriting transcription uncertain
\fieldvalue{\handwritten{石孟男}}
\hfill
日期：%
% #VALUE_ID: LEX-P0048-V0017
% #FIELD_VALUE: 日期
% #HANDWRITTEN: 2025.08.17
\fieldvalue{\handwritten{2025.08.17}}
% LEXOID_PAGE_COMPLETED: 48/70

\newpage

\noindent
Assay: MSC-1

\medskip

\noindent
Cell Type F1: MSC (AO)\\
Cell Type F2: MSC (PI)

\medskip

\noindent
Sample ID: QC-JS-QY0402-A31Z2202508006-20250817-B01-0H-2\\
Dilution: 2.00

\medskip
\hrule
\medskip

\noindent
Results:

\medskip

\noindent
\begin{tabular}{@{}p{0.29\linewidth}p{0.34\linewidth}p{0.29\linewidth}@{}}
\textbf{Count} & \textbf{Concentration} & \textbf{Mean Diameter} \\
\underline{\hspace{0.20\linewidth}} &
\underline{\hspace{0.24\linewidth}} &
\underline{\hspace{0.21\linewidth}} \\[2pt]
Total: 224 & $7.88 \times 10^{5}$ cells/mL & 13.4 micron \\
Live: 216 & $7.60 \times 10^{5}$ cells/mL & 13.5 microns \\
Dead: 8 & $2.80 \times 10^{4}$ cells/mL & 12.9 microns \\
Viability: 96.4\% & & \\
\end{tabular}

\bigskip

\noindent
\begin{tabular}{@{}p{0.48\linewidth}p{0.48\linewidth}@{}}
\centering
% TODO: Insert the visible upper-left microscopy image; no source filename was provided.
\fbox{\parbox[c][0.22\textheight][c]{0.92\linewidth}{\centering Image placeholder}}
&
\centering
% TODO: Insert the visible upper-right fluorescence image; no source filename was provided.
\fbox{\parbox[c][0.22\textheight][c]{0.92\linewidth}{\centering Image placeholder}}
\tabularnewline[1em]
\centering
% TODO: Insert the visible lower-left microscopy image; no source filename was provided.
\fbox{\parbox[c][0.22\textheight][c]{0.92\linewidth}{\centering Image placeholder}}
&
\centering
% TODO: Insert the visible lower-right dark-field image; no source filename was provided.
\fbox{\parbox[c][0.22\textheight][c]{0.92\linewidth}{\centering Image placeholder}}
\tabularnewline
\end{tabular}

\vfill

\noindent
操作人/审核人：

% #VALUE_ID: LEX-P0049-V0001
% #FIELD_VALUE: 操作人/审核人
% #HANDWRITTEN: 邵逸果
\fieldvalue{\handwritten{邵逸果}}
\hfill
日期：

% #VALUE_ID: LEX-P0049-V0002
% #FIELD_VALUE: 日期
% #TODO #HANDWRITTEN: 202508?7; final digit or separator is unclear
\fieldvalue{\handwritten{202508?7}}
% LEXOID_PAGE_COMPLETED: 49/70

\newpage

\noindent\textbf{Assay:} MSC-1

\vspace{0.8em}

\noindent
\textbf{Cell Type F1:} MSC (AO)\\
\textbf{Cell Type F2:} MSC (PI)

\vspace{0.8em}

\noindent
\textbf{Sample ID:} QC-JS-QY0402-A31Z2202508006-20250817-B01-OH-3\\
\textbf{Dilution:} 2.00

\vspace{0.8em}
\hrule
\vspace{0.3em}

\noindent Results:

\vspace{0.5em}

\noindent
\begin{tabular}{@{}p{0.29\linewidth}p{0.34\linewidth}p{0.29\linewidth}@{}}
\textbf{Count} & \textbf{Concentration} & \textbf{Mean Diameter} \\
\texttt{=============} & \texttt{=============} & \texttt{=============} \\
Total: 202 & $7.12{\times}10^{5}$ cells/mL & 14.0 micron \\
Live: 196 & $6.91{\times}10^{5}$ cells/mL & 14.2 microns \\
Dead: 6 & $2.10{\times}10^{4}$ cells/mL & 10.5 microns \\
Viability: 97.1\% & & \\
\end{tabular}

\vspace{1.2em}

\noindent
\begin{tabular}{@{}p{0.48\linewidth}@{\hspace{0.04\linewidth}}p{0.48\linewidth}@{}}
% TODO: Insert the upper-left microscopy image visible on this page.
\centering\fbox{\parbox[c][0.22\textheight][c]{0.97\linewidth}{\centering Upper-left image placeholder}}
&
% TODO: Insert the upper-right microscopy image visible on this page.
\centering\fbox{\parbox[c][0.22\textheight][c]{0.97\linewidth}{\centering Upper-right image placeholder}}
\tabularnewline[1em]
% TODO: Insert the lower-left microscopy image visible on this page.
\centering\fbox{\parbox[c][0.22\textheight][c]{0.97\linewidth}{\centering Lower-left image placeholder}}
&
% TODO: Insert the lower-right microscopy image visible on this page.
\centering\fbox{\parbox[c][0.22\textheight][c]{0.97\linewidth}{\centering Lower-right image placeholder}}
\tabularnewline
\end{tabular}

\vfill

\noindent
操作人/审核人：

% #VALUE_ID: LEX-P0050-V0001
% #FIELD_VALUE: 操作人/审核人
% #HANDWRITTEN: 石逸昊
\fieldvalue{\handwritten{石逸昊}}
\hfill
日期：

% #VALUE_ID: LEX-P0050-V0002
% #FIELD_VALUE: 日期
% #HANDWRITTEN: 2025.08.17
\fieldvalue{\handwritten{2025.08.17}}
% LEXOID_PAGE_COMPLETED: 50/70

\newpage

\begin{center}
\textbf{铂生生物科技（北京）有限公司}\\
\textit{PLATINUM LIFE EXCELLENCE BIOTECH}\\[2pt]
\begin{tabular}{c}
\textbf{铂生生物}\\
原件\quad \textbf{请验单}
\end{tabular}
\hfill
文件编码：REC-YZ-SMP-QC-01-001-d\qquad 版本号：1.4
\end{center}

\begin{center}
\renewcommand{\arraystretch}{1.45}
\small
\begin{tabular}{|p{2.4cm}|p{4.8cm}|p{2.4cm}|p{4.8cm}|}
\hline
\centering 样品名称 &
% #VALUE_ID: LEX-P0051-V0001
% #FIELD_VALUE: 样品名称
% #TODO #HANDWRITTEN: 二氯（字迹不清，后续文字无法辨认）
\fieldvalue{\handwritten{二氯\ldots}} &
\centering 样品代码 &
% #VALUE_ID: LEX-P0051-V0002
% #FIELD_VALUE: 样品代码
% #TODO #HANDWRITTEN: Q10-02（部分字符不清）
\fieldvalue{\handwritten{Q10-02}}\\
\hline
\centering 样品批号 &
% #VALUE_ID: LEX-P0051-V0003
% #FIELD_VALUE: 样品批号
% #TODO #HANDWRITTEN: A1222?0800（字迹部分不清）
\fieldvalue{\handwritten{A1222?0800}} &
\centering 样品规格/数量 &
% #VALUE_ID: LEX-P0051-V0004
% #FIELD_VALUE: 样品规格/数量
% #TODO #HANDWRITTEN: 50ml小瓶×1瓶（字迹部分不清）
\fieldvalue{\handwritten{50ml小瓶$\times$1瓶}}\\
\hline
\centering 样品类别 &
\multicolumn{3}{l|}{
% #VALUE_ID: LEX-P0051-V0005
% #FIELD_VALUE: 样品类别—生产过程中检验样品
% #HANDWRITTEN: √
\fieldvalue{\handwritten{\checkboxfield{checked}}}\,生产过程中检验样品
\quad
% #VALUE_ID: LEX-P0051-V0006
% #FIELD_VALUE: 样品类别—待包装产品
\fieldvalue{\checkboxfield{unchecked}}\,待包装产品
\quad
% #VALUE_ID: LEX-P0051-V0007
% #FIELD_VALUE: 样品类别—成品
\fieldvalue{\checkboxfield{unchecked}}\,成品}\\
\cline{2-4}
&
\multicolumn{3}{l|}{
% #VALUE_ID: LEX-P0051-V0008
% #FIELD_VALUE: 样品类别—其他
\fieldvalue{\checkboxfield{unchecked}}\,其他（\hspace{3cm}）}\\
\hline
\centering 检验类型 &
\multicolumn{3}{l|}{
% #VALUE_ID: LEX-P0051-V0009
% #FIELD_VALUE: 检验类型—初验
% #HANDWRITTEN: √
\fieldvalue{\handwritten{\checkboxfield{checked}}}\,初验
\quad
% #VALUE_ID: LEX-P0051-V0010
% #FIELD_VALUE: 检验类型—复验
\fieldvalue{\checkboxfield{unchecked}}\,复验
\quad
% #VALUE_ID: LEX-P0051-V0011
% #FIELD_VALUE: 检验类型—退货
\fieldvalue{\checkboxfield{unchecked}}\,退货
\quad
% #VALUE_ID: LEX-P0051-V0012
% #FIELD_VALUE: 检验类型—稳定性
\fieldvalue{\checkboxfield{unchecked}}\,稳定性
\quad
% #VALUE_ID: LEX-P0051-V0013
% #FIELD_VALUE: 检验类型—留样
\fieldvalue{\checkboxfield{unchecked}}\,留样}\\
\hline
\centering 请验目的 &
\multicolumn{3}{l|}{
% #VALUE_ID: LEX-P0051-V0014
% #FIELD_VALUE: 请验目的—全检
% #HANDWRITTEN: √
\fieldvalue{\handwritten{\checkboxfield{checked}}}\,全检
\quad
% #VALUE_ID: LEX-P0051-V0015
% #FIELD_VALUE: 请验目的—留样
\fieldvalue{\checkboxfield{unchecked}}\,留样
\quad
% #VALUE_ID: LEX-P0051-V0016
% #FIELD_VALUE: 请验目的—其他
\fieldvalue{\checkboxfield{unchecked}}\,其他（\hspace{3cm}）}\\
\hline
\centering 请验部门 &
\multicolumn{3}{l|}{
% #VALUE_ID: LEX-P0051-V0017
% #FIELD_VALUE: 请验部门—生产部
% #HANDWRITTEN: √
\fieldvalue{\handwritten{\checkboxfield{checked}}}\,生产部
\quad
% #VALUE_ID: LEX-P0051-V0018
% #FIELD_VALUE: 请验部门—物料管理部
\fieldvalue{\checkboxfield{unchecked}}\,物料管理部
\quad
% #VALUE_ID: LEX-P0051-V0019
% #FIELD_VALUE: 请验部门—其他
\fieldvalue{\checkboxfield{unchecked}}\,其他（\hspace{3cm}）}\\
\hline
\centering 样品情况 &
\multicolumn{3}{l|}{
% #VALUE_ID: LEX-P0051-V0020
% #FIELD_VALUE: 样品情况—随请验单一起提供
% #HANDWRITTEN: √
\fieldvalue{\handwritten{\checkboxfield{checked}}}\,随请验单一起提供
\quad
% #VALUE_ID: LEX-P0051-V0021
% #FIELD_VALUE: 样品情况—未随请验单一起提供
\fieldvalue{\checkboxfield{unchecked}}\,未随请验单一起提供}\\
\hline
\centering 请验人 &
% #VALUE_ID: LEX-P0051-V0022
% #FIELD_VALUE: 请验人
% #TODO #HANDWRITTEN: 江瑶（字迹不清）
\fieldvalue{\handwritten{江瑶}} &
\centering 请验日期 &
% #VALUE_ID: LEX-P0051-V0023
% #FIELD_VALUE: 请验日期
% #HANDWRITTEN: 2025.08.20
\fieldvalue{\handwritten{2025.08.20}}\\
\hline
\centering 收验人 &
% #VALUE_ID: LEX-P0051-V0024
% #FIELD_VALUE: 收验人
% #TODO #HANDWRITTEN: 徐\ldots（字迹不清）
\fieldvalue{\handwritten{徐\ldots}} &
\centering 收验日期 &
% #VALUE_ID: LEX-P0051-V0025
% #FIELD_VALUE: 收验日期
% #HANDWRITTEN: 2025.08.20
\fieldvalue{\handwritten{2025.08.20}}\\
\hline
\centering 备注 &
\multicolumn{3}{l|}{
% #VALUE_ID: LEX-P0051-V0026
% #FIELD_VALUE: 备注
% #TODO #HANDWRITTEN: J?H 13?L（字迹无法完全辨认）
\fieldvalue{\handwritten{J?H\quad 13?L}}}\\
\hline
\end{tabular}
\end{center}
% LEXOID_PAGE_COMPLETED: 51/70

\newpage

\begin{center}
{\small
\begin{tabular}{|p{2.0cm}|p{4.9cm}|p{2.0cm}|p{4.2cm}|}
\hline
\multicolumn{4}{c}{\textbf{取样记录}}\\[-1mm]
\hline
\multicolumn{2}{r}{文件编码：REC-YZ-SMP-QC-01-004-b}\hspace{1cm}
&\multicolumn{2}{l}{版本号：1.4}\\
\hline
名称 &
% #VALUE_ID: LEX-P0052-V0001
% #FIELD_VALUE: 名称
% #TODO #HANDWRITTEN: 冻干重组人……抗体； handwriting partially illegible
\fieldvalue{\handwritten{冻干重组人……抗体}} &
代码 &
% #VALUE_ID: LEX-P0052-V0002
% #FIELD_VALUE: 代码
% #TODO #HANDWRITTEN: Qf602；字符辨认不清
\fieldvalue{\handwritten{Qf602}}\\
\hline
规格 &
% #VALUE_ID: LEX-P0052-V0003
% #FIELD_VALUE: 规格
% #TODO #HANDWRITTEN: 2支 30ml/支；前部字符不清
\fieldvalue{\handwritten{2支 30ml/支}} &
总数量 &
% #VALUE_ID: LEX-P0052-V0004
% #FIELD_VALUE: 总数量
% #HANDWRITTEN: 1
\fieldvalue{\handwritten{1}}\\
\hline
批号 &
% #VALUE_ID: LEX-P0052-V0005
% #FIELD_VALUE: 批号
% #TODO #HANDWRITTEN: A1312200……； handwriting partially illegible
\fieldvalue{\handwritten{A1312200……}} &
进厂批号 &
% #VALUE_ID: LEX-P0052-V0006
% #FIELD_VALUE: 进厂批号
% #TODO #HANDWRITTEN: /； handwritten mark
\fieldvalue{\handwritten{/}}\\
\hline
供货单位/厂家 &
\multicolumn{3}{c|}{
% #VALUE_ID: LEX-P0052-V0007
% #FIELD_VALUE: 供货单位/厂家
\fieldvalue{/}}\\
\hline
取样目的 &
\multicolumn{3}{c|}{
% #VALUE_ID: LEX-P0052-V0008
% #FIELD_VALUE: 正常检验
\fieldvalue{\checkboxfield{checked}} 正常检验\quad
% #VALUE_ID: LEX-P0052-V0009
% #FIELD_VALUE: 复检
\fieldvalue{\checkboxfield{unchecked}} 复检\quad
% #VALUE_ID: LEX-P0052-V0010
% #FIELD_VALUE: 稳定性
\fieldvalue{\checkboxfield{unchecked}} 稳定性\quad
% #VALUE_ID: LEX-P0052-V0011
% #FIELD_VALUE: 退货检验
\fieldvalue{\checkboxfield{unchecked}} 退货检验\quad
% #VALUE_ID: LEX-P0052-V0012
% #FIELD_VALUE: 留样
\fieldvalue{\checkboxfield{unchecked}} 留样\quad
% #VALUE_ID: LEX-P0052-V0013
% #FIELD_VALUE: 其他
\fieldvalue{\checkboxfield{unchecked}} 其他：\underline{\hspace{2cm}}}\\
\hline
取样地点 &
\multicolumn{3}{c|}{
% #VALUE_ID: LEX-P0052-V0014
% #FIELD_VALUE: 产品出口
\fieldvalue{\checkboxfield{checked}} 产品出口\quad
% #VALUE_ID: LEX-P0052-V0015
% #FIELD_VALUE: 物料库
\fieldvalue{\checkboxfield{unchecked}} 物料库\quad
% #VALUE_ID: LEX-P0052-V0016
% #FIELD_VALUE: 干细胞库
\fieldvalue{\checkboxfield{unchecked}} 干细胞库\quad
% #VALUE_ID: LEX-P0052-V0017
% #FIELD_VALUE: 包装间
\fieldvalue{\checkboxfield{unchecked}} 包装间}\\
\hline
贮存条件 &
\multicolumn{3}{c|}{
% #VALUE_ID: LEX-P0052-V0018
% #FIELD_VALUE: 2--8$^\circ$C
\fieldvalue{\checkboxfield{checked}} 2--8$^\circ$C\quad
% #VALUE_ID: LEX-P0052-V0019
% #FIELD_VALUE: $-20^\circ$C
\fieldvalue{\checkboxfield{unchecked}} $-20^\circ$C\quad
% #VALUE_ID: LEX-P0052-V0020
% #FIELD_VALUE: $-80^\circ$C
\fieldvalue{\checkboxfield{unchecked}} $-80^\circ$C\quad
% #VALUE_ID: LEX-P0052-V0021
% #FIELD_VALUE: 常温
\fieldvalue{\checkboxfield{unchecked}} 常温\quad
% #VALUE_ID: LEX-P0052-V0022
% #FIELD_VALUE: 其他
\fieldvalue{\checkboxfield{unchecked}} 其他：\underline{\hspace{1.5cm}}}\\
\hline
盛装容器 &
\multicolumn{3}{c|}{
% #VALUE_ID: LEX-P0052-V0023
% #FIELD_VALUE: 保温箱
\fieldvalue{\checkboxfield{checked}} 保温箱\quad
% #VALUE_ID: LEX-P0052-V0024
% #FIELD_VALUE: 泡沫箱
\fieldvalue{\checkboxfield{unchecked}} 泡沫箱\quad
% #VALUE_ID: LEX-P0052-V0025
% #FIELD_VALUE: 手提篮
\fieldvalue{\checkboxfield{unchecked}} 手提篮\quad
% #VALUE_ID: LEX-P0052-V0026
% #FIELD_VALUE: 其他
\fieldvalue{\checkboxfield{unchecked}} 其他：\underline{\hspace{2cm}}}\\
\hline
外观复核 &
\multicolumn{3}{c|}{
% #VALUE_ID: LEX-P0052-V0027
% #FIELD_VALUE: 合格
\fieldvalue{\checkboxfield{checked}} 合格\quad
% #VALUE_ID: LEX-P0052-V0028
% #FIELD_VALUE: 不合格
\fieldvalue{\checkboxfield{unchecked}} 不合格}\\
\hline
\multicolumn{4}{|c|}{\begin{tabular}{llllll}
取样量 &
% #VALUE_ID: LEX-P0052-V0029
% #FIELD_VALUE: 取样量
% #HANDWRITTEN: 1
\fieldvalue{\handwritten{1}} &
检验量 &
% #VALUE_ID: LEX-P0052-V0030
% #FIELD_VALUE: 检验量
% #HANDWRITTEN: 1
\fieldvalue{\handwritten{1}} &
留样量 &
% #VALUE_ID: LEX-P0052-V0031
% #FIELD_VALUE: 留样量
% #TODO #HANDWRITTEN: /； handwritten mark
\fieldvalue{\handwritten{/}}
\end{tabular}}\\
\hline
取样人/日期 &
\multicolumn{1}{c|}{
% #VALUE_ID: LEX-P0052-V0032
% #FIELD_VALUE: 取样人
% #TODO #HANDWRITTEN: 杨贤？；姓名部分难以辨认
\fieldvalue{\handwritten{杨贤？}}} &
\multicolumn{2}{l|}{
% #VALUE_ID: LEX-P0052-V0033
% #FIELD_VALUE: 取样日期
% #TODO #HANDWRITTEN: 20?5.08.20；年份字符不清
\fieldvalue{\handwritten{20?5.08.20}}}\\
\hline
复核人/日期 &
\multicolumn{1}{c|}{
% #VALUE_ID: LEX-P0052-V0034
% #FIELD_VALUE: 复核人
% #TODO #HANDWRITTEN: 赵弘？；姓名部分难以辨认
\fieldvalue{\handwritten{赵弘？}}} &
\multicolumn{2}{l|}{
% #VALUE_ID: LEX-P0052-V0035
% #FIELD_VALUE: 复核日期
% #TODO #HANDWRITTEN: 20?5.08.20；年份字符不清
\fieldvalue{\handwritten{20?5.08.20}}}\\
\hline
备注 &
\multicolumn{3}{c|}{
% #VALUE_ID: LEX-P0052-V0036
% #FIELD_VALUE: 备注
% #TODO #HANDWRITTEN: 70h 13?L.；末尾字符难以辨认
\fieldvalue{\handwritten{70h\quad 13?L.}}}\\
\hline
\end{tabular}
}
\end{center}

\begin{center}
{\small 第 1 页 共 1 页}
\end{center}
% LEXOID_PAGE_COMPLETED: 52/70

\newpage

\section*{生物反应器细胞液检验操作记录}

\noindent
文件编码：REC-YZ-SOP-QC-99-006-a

\begin{tabular}{|p{2.2cm}|p{5.0cm}|p{2.2cm}|p{4.0cm}|}
\hline
\multicolumn{4}{|c|}{葡萄糖浓度检验} \\ \hline
样品名称 &
% #VALUE_ID: LEX-P0053-V0001
% #FIELD_VALUE: 样品名称
\fieldvalue{二级生物反应器细胞液} &
样品规格 &
% #VALUE_ID: LEX-P0053-V0002
% #FIELD_VALUE: 样品规格
\fieldvalue{\handwritten{20}} ml/袋 \\ \hline
样品批号 &
% #VALUE_ID: LEX-P0053-V0003
% #FIELD_VALUE: 样品批号
% #HANDWRITTEN: 202?4?4?4?06
\fieldvalue{\handwritten{202?4?4?4?06}} &
取样日期 &
% #VALUE_ID: LEX-P0053-V0004
% #FIELD_VALUE: 取样日期
% #HANDWRITTEN: 2025年08月22日
\fieldvalue{\handwritten{2025年08月22日}} \\ \hline
\end{tabular}

\subsection*{1\quad 仪器设备、试剂耗材}

\subsubsection*{1.1\quad 仪器设备}

\begin{tabular}{|p{2.0cm}|p{4.0cm}|p{2.2cm}|p{2.7cm}|p{3.2cm}|}
\hline
名称 & 厂家 & 型号 & 设备编号 & 检查确认 \\ \hline
血糖仪 & 三诺生物传感股份有限公司 & 安稳+ &
% #VALUE_ID: LEX-P0053-V0005
% #FIELD_VALUE: 设备编号
% #TODO #HANDWRITTEN: /; handwritten mark is unclear
\fieldvalue{\handwritten{/}} &
正常\quad
% #VALUE_ID: LEX-P0053-V0006
% #FIELD_VALUE: 正常
\fieldvalue{\checkboxfield{checked}}
\quad 异常\quad
% #VALUE_ID: LEX-P0053-V0007
% #FIELD_VALUE: 异常
\fieldvalue{\checkboxfield{unchecked}} \\ \hline
移液器 & Eppendorf Corporate & 0.1$\sim$2.5\,$\mu$l &
% #VALUE_ID: LEX-P0053-V0008
% #FIELD_VALUE: 设备编号
% #TODO #HANDWRITTEN: 0217486; some digits are unclear
\fieldvalue{\handwritten{0217486}} &
正常\quad
% #VALUE_ID: LEX-P0053-V0009
% #FIELD_VALUE: 正常
\fieldvalue{\checkboxfield{checked}}
\quad 异常\quad
% #VALUE_ID: LEX-P0053-V0010
% #FIELD_VALUE: 异常
\fieldvalue{\checkboxfield{unchecked}} \\ \hline
\end{tabular}

\subsubsection*{1.2\quad 试剂耗材}

\begin{tabular}{|p{2.5cm}|p{5.4cm}|p{3.0cm}|p{3.2cm}|}
\hline
名称 & 厂家 & 批号 & 有效期至 \\ \hline
血糖测试条 & 三诺生物传感股份有限公司 &
% #VALUE_ID: LEX-P0053-V0011
% #FIELD_VALUE: 批号
% #HANDWRITTEN: 444?4?0-6
\fieldvalue{\handwritten{444?4?0-6}} &
% #VALUE_ID: LEX-P0053-V0012
% #FIELD_VALUE: 有效期至
% #HANDWRITTEN: 2027-?1-38
\fieldvalue{\handwritten{2027-?1-38}} \\ \hline
\end{tabular}

\subsection*{2\quad 操作步骤：}

\begin{tabular}{|p{8.7cm}|p{5.4cm}|}
\hline
\multicolumn{2}{|c|}{操作过程} \\ \hline
取出测试条，取样品自然沉降后的上清液 0.6$\mu$l 加入测试条电极反应区，血糖仪自动发出“嘀”提示音，吸入样品完成，血糖仪自动测定。 &
取上清液：
% #VALUE_ID: LEX-P0053-V0013
% #FIELD_VALUE: 取上清液
% #HANDWRITTEN: 0.6
\fieldvalue{\handwritten{0.6}} $\mu$l \\ \hline
入倒计时，开始测试； &
 \\ \hline
重复以上步骤 2 次，共计数 3 次。 &
每个样品分别共计数
% #VALUE_ID: LEX-P0053-V0014
% #FIELD_VALUE: 每个样品分别共计数
% #HANDWRITTEN: 3
\fieldvalue{\handwritten{3}} 次。 \\ \hline
\multicolumn{2}{|c|}{检测结果} \\ \hline
\multicolumn{2}{|c|}{葡萄糖浓度（mmol/L）} \\ \hline
\multicolumn{1}{|c|}{第一次} &
\multicolumn{1}{c|}{第二次} \\ \hline
% #VALUE_ID: LEX-P0053-V0015
% #FIELD_VALUE: 第一次
% #HANDWRITTEN: 14.0
\fieldvalue{\handwritten{14.0}} &
% #VALUE_ID: LEX-P0053-V0016
% #FIELD_VALUE: 第二次
% #HANDWRITTEN: 14.4
\fieldvalue{\handwritten{14.4}} \\ \hline
\multicolumn{1}{|c|}{第三次} &
\multicolumn{1}{c|}{平均值（保留一位小数）} \\ \hline
% #VALUE_ID: LEX-P0053-V0017
% #FIELD_VALUE: 第三次
% #HANDWRITTEN: 14.4
\fieldvalue{\handwritten{14.4}} &
% #VALUE_ID: LEX-P0053-V0018
% #FIELD_VALUE: 平均值（保留一位小数）
% #HANDWRITTEN: 14.3
\fieldvalue{\handwritten{14.3}} \\ \hline
备注 &
% #VALUE_ID: LEX-P0053-V0019
% #FIELD_VALUE: 备注
% #HANDWRITTEN: /
\fieldvalue{\handwritten{/}} \\ \hline
实验人/日期 &
% #VALUE_ID: LEX-P0053-V0020
% #FIELD_VALUE: 实验人
% #TODO #HANDWRITTEN: 程峰; handwriting is unclear
\fieldvalue{\handwritten{程峰}}\quad
% #VALUE_ID: LEX-P0053-V0021
% #FIELD_VALUE: 实验日期
% #HANDWRITTEN: 2025.08.20
\fieldvalue{\handwritten{2025.08.20}} \\ \hline
复核人/日期 &
% #VALUE_ID: LEX-P0053-V0022
% #FIELD_VALUE: 复核人
% #TODO #HANDWRITTEN: 楠泉; handwriting is unclear
\fieldvalue{\handwritten{楠泉}}\quad
% #VALUE_ID: LEX-P0053-V0023
% #FIELD_VALUE: 复核日期
% #HANDWRITTEN: 2025.07.20
\fieldvalue{\handwritten{2025.07.20}} \\ \hline
\end{tabular}

% TODO: The source page footer indicates this is page 1 of a 5-page record; only the content visible on this PDF page is transcribed here.
% LEXOID_PAGE_COMPLETED: 53/70

\newpage

\begin{center}
\textbf{铂生生物（北京）有限公司}\\
\textbf{生物反应器细胞液检测操作记录}\\
文件编码：REC-YZ-SOP-QC-99-006-a
\end{center}

\begin{center}
\begin{tabular}{|p{2.2cm}|p{5.0cm}|p{2.2cm}|p{3.2cm}|}
\hline
\multicolumn{4}{|c|}{密度与活率检测}\\
\hline
样品名称 &
% #VALUE_ID: LEX-P0054-V0001
% #FIELD_VALUE: 样品名称
% #TODO #HANDWRITTEN: 一级生物反应器细胞液; 字迹部分不清
\fieldvalue{\handwritten{一级生物反应器细胞液}} &
样品规格 &
% #VALUE_ID: LEX-P0054-V0002
% #FIELD_VALUE: 样品规格
% #TODO #HANDWRITTEN: 2? ml/袋; 数字不清
\fieldvalue{\handwritten{2? ml/袋}}\\
\hline
样品批号 &
% #VALUE_ID: LEX-P0054-V0003
% #FIELD_VALUE: 样品批号
% #TODO #HANDWRITTEN: A?202?A?; 字迹不清
\fieldvalue{\handwritten{A?202?A?}} &
取样日期 &
% #VALUE_ID: LEX-P0054-V0004
% #FIELD_VALUE: 取样日期
% #TODO #HANDWRITTEN: 202?年?月20日; 部分数字不清
\fieldvalue{\handwritten{202?年?月20日}}\\
\hline
\end{tabular}
\end{center}

\section{仪器设备、试剂耗材}

\subsection{仪器设备}

\begin{center}
\small
\begin{tabular}{|p{2.4cm}|p{2.1cm}|p{2.5cm}|p{2.1cm}|p{2.4cm}|p{2.2cm}|}
\hline
名称 & 厂家 & 型号 & 设备编号 & 校准有效期至 & 检查确认\\
\hline
细胞计数仪 & Nexcelom & Cellometer K2 & 1410905 &
% #VALUE_ID: LEX-P0054-V0005
% #FIELD_VALUE: 校准有效期至
% #HANDWRITTEN: 2025.11.13
\fieldvalue{\handwritten{2025.11.13}} &
正常\ 
% #VALUE_ID: LEX-P0054-V0006
% #FIELD_VALUE: 正常
% #HANDWRITTEN: ✓
\fieldvalue{\handwritten{\checkboxfield{checked}}}\quad
异常\ 
% #VALUE_ID: LEX-P0054-V0007
% #FIELD_VALUE: 异常
\fieldvalue{\checkboxfield{unchecked}}\\
\hline
高速冷冻离心机 & Thermo & ST16R & 1410810 &
% #VALUE_ID: LEX-P0054-V0008
% #FIELD_VALUE: 校准有效期至
% #HANDWRITTEN: 2025.01.17
\fieldvalue{\handwritten{2025.01.17}} &
正常\ 
% #VALUE_ID: LEX-P0054-V0009
% #FIELD_VALUE: 正常
% #HANDWRITTEN: ✓
\fieldvalue{\handwritten{\checkboxfield{checked}}}\quad
异常\ 
% #VALUE_ID: LEX-P0054-V0010
% #FIELD_VALUE: 异常
\fieldvalue{\checkboxfield{unchecked}}\\
\hline
电热恒温水槽 & 上海一恒 & DK-8AXX & 1411001 &
% #VALUE_ID: LEX-P0054-V0011
% #FIELD_VALUE: 校准有效期至
% #HANDWRITTEN: 2025.01.17
\fieldvalue{\handwritten{2025.01.17}} &
正常\ 
% #VALUE_ID: LEX-P0054-V0012
% #FIELD_VALUE: 正常
% #HANDWRITTEN: ✓
\fieldvalue{\handwritten{\checkboxfield{checked}}}\quad
异常\ 
% #VALUE_ID: LEX-P0054-V0013
% #FIELD_VALUE: 异常
\fieldvalue{\checkboxfield{unchecked}}\\
\hline
\end{tabular}
\end{center}

\subsection{试剂耗材}

\begin{center}
\begin{tabular}{|p{3.0cm}|p{3.2cm}|p{4.0cm}|p{3.2cm}|}
\hline
名称 & 厂家 & 批号 & 有效期至\\
\hline
AO/PI 染液 & Nexcelom &
% #VALUE_ID: LEX-P0054-V0014
% #FIELD_VALUE: 批号
% #TODO #HANDWRITTEN: 324?4?9; 字迹不清
\fieldvalue{\handwritten{324?4?9}} &
% #VALUE_ID: LEX-P0054-V0015
% #FIELD_VALUE: 有效期至
% #TODO #HANDWRITTEN: 2025.02?; 字迹部分不清
\fieldvalue{\handwritten{2025.02?}}\\
\hline
0.9\%氯化钠注射液 & 石家庄四药 &
% #VALUE_ID: LEX-P0054-V0016
% #FIELD_VALUE: 批号
% #TODO #HANDWRITTEN: 24?073?32; 字迹不清
\fieldvalue{\handwritten{24?073?32}} &
% #VALUE_ID: LEX-P0054-V0017
% #FIELD_VALUE: 有效期至
% #HANDWRITTEN: 2026.02.06
\fieldvalue{\handwritten{2026.02.06}}\\
\hline
20\%人血白蛋白 & Baxter AG &
% #VALUE_ID: LEX-P0054-V0018
% #FIELD_VALUE: 批号
% #TODO #HANDWRITTEN: 4?1?1?1?0; 字迹不清
\fieldvalue{\handwritten{4?1?1?1?0}} &
% #VALUE_ID: LEX-P0054-V0019
% #FIELD_VALUE: 有效期至
% #HANDWRITTEN: 2025.08.31
\fieldvalue{\handwritten{2025.08.31}}\\
\hline
\end{tabular}
\end{center}

\section{操作步骤}

\subsection*{溶液准备}

\begin{tabular}{|p{7.2cm}|p{6.1cm}|}
\hline
1\%人白蛋白溶液配制：取5 ml 20\%人血白蛋白加入到95 ml的0.9\%氯化钠注射液中，混匀即得。&
20\%人血白蛋白
% #VALUE_ID: LEX-P0054-V0020
% #FIELD_VALUE: 20\%人血白蛋白用量
% #HANDWRITTEN: 1
\fieldvalue{\handwritten{1}}\ ml\\
&
0.9\%氯化钠注射液
% #VALUE_ID: LEX-P0054-V0021
% #FIELD_VALUE: 0.9\%氯化钠注射液用量
% #HANDWRITTEN: 19
\fieldvalue{\handwritten{19}}\ ml\\
\hline
\multicolumn{2}{|c|}{\text{溶液配制记录}}\\
\hline
名称 & 配制批号 \hfill 有效期至\\
\hline
5$\times$裂解液 &
% #VALUE_ID: LEX-P0054-V0022
% #FIELD_VALUE: 配制批号
% #TODO #HANDWRITTEN: 520?473?8?; 字迹不清
\fieldvalue{\handwritten{520?473?8?}}
\hfill
% #VALUE_ID: LEX-P0054-V0023
% #FIELD_VALUE: 有效期至
% #HANDWRITTEN: 2025.10.7
\fieldvalue{\handwritten{2025.10.7}}\\
\hline
\end{tabular}

\subsection*{操作过程}

\begin{tabular}{|p{7.2cm}|p{6.1cm}|}
\hline
按照50 ml离心管的刻度读取样品体积，往样品中加入对应体积的5$\times$裂解液（每4 ml样品加入1 ml 5$\times$裂解液，按照“只进不舍”原则，如26 ml样品，加入7 ml 5$\times$裂解液），充分混匀后倒入23 ml的样品中，置于37℃电热恒温水槽中裂解20 min后，取出离心管摇晃混匀，使样品充分裂解；再置于37℃电热恒温水槽中裂解10 min得到细胞悬液。&
样品体积：
% #VALUE_ID: LEX-P0054-V0024
% #FIELD_VALUE: 样品体积
% #HANDWRITTEN: 36
\fieldvalue{\handwritten{36}}\ ml\\
&
裂解液：
% #VALUE_ID: LEX-P0054-V0025
% #FIELD_VALUE: 裂解液
% #HANDWRITTEN: 9
\fieldvalue{\handwritten{9}}\ ml\\
&
电热恒温水槽：
% #VALUE_ID: LEX-P0054-V0026
% #FIELD_VALUE: 电热恒温水槽温度
% #HANDWRITTEN: 70
\fieldvalue{\handwritten{70}}\ \ensuremath{^\circ\mathrm{C}}\\
&
裂解：
% #VALUE_ID: LEX-P0054-V0027
% #FIELD_VALUE: 裂解时间
% #TODO #HANDWRITTEN: 09:30--10:05; 时间字迹不清
\fieldvalue{\handwritten{09:30--10:05}}\\
\hline
取出裂解好的细胞悬液300$\times g$离心5分钟，按要求读取样品体积加入对应体积的1\%人白蛋白溶液重悬（每1.2 ml样品体积加入0.1 ml 1\%人白蛋白溶液，按照“只进不舍”原则，如26 ml样品加入2.2 ml）。&
% #VALUE_ID: LEX-P0054-V0028
% #FIELD_VALUE: 离心力
% #HANDWRITTEN: 300
\fieldvalue{\handwritten{300}}\ $\times g$离心
% #VALUE_ID: LEX-P0054-V0029
% #FIELD_VALUE: 离心时间
% #HANDWRITTEN: 5
\fieldvalue{\handwritten{5}}\ 分钟\\
&
加入1\%人白蛋白溶液重悬体积：
% #VALUE_ID: LEX-P0054-V0030
% #FIELD_VALUE: 1\%人白蛋白溶液重悬体积
% #HANDWRITTEN: 3
\fieldvalue{\handwritten{3}}\ ml\\
&
吸取
% #VALUE_ID: LEX-P0054-V0031
% #FIELD_VALUE: 吸取体积
% #HANDWRITTEN: 20
\fieldvalue{\handwritten{20}}\ $\mu$l\\
\hline
\end{tabular}

\begin{center}
第 2 页 共 5 页
\end{center}
% LEXOID_PAGE_COMPLETED: 54/70

\newpage

\section*{细胞计数记录}

\noindent
\textbf{文件编码：}REC-YZ-SOP-QC-99-006-a

\medskip
\noindent
加入 $2.2\mathrm{ml}$ 人血气化培养液，反复吹打至单细胞悬液。\\
撕去计数板上下两层保护膜，将细胞悬液混匀，从中吸取 $20\mu\mathrm{l}$ 与 $20\mu\mathrm{l}$ AO/PI 染料混合均匀。\\
吸取 $20\mu\mathrm{l}$ 细胞悬液与 AO/PI 染料的混合液，注入到计数板的一 个检测室中。\\
将加入混合液的一端插入进样口。\\
选择 Assay 类型，使用 AO/PI 染料进行双荧光存活率检测；在主界面左侧 ID 栏中输入样本号。\\
点击主界面左下角“Preview”预览细胞图片。\\
调整各照片距离至观察到荧光通道下的细胞中心透亮、轮廓清晰。\\
点击“count”按钮进行计数并自动拍摄图片。\\
重复以上步骤 $2$ 次，共计数 $3$ 次。打印数据结果。

\medskip
\noindent
每个样品均计数
% #VALUE_ID: LEX-P0055-V0001
% #FIELD_VALUE: 每个样品均计数次数
% #HANDWRITTEN: 2
\fieldvalue{\handwritten{2}}
次。

\bigskip
\noindent
\textbf{数据保存路径：}
% #VALUE_ID: LEX-P0055-V0002
% #FIELD_VALUE: 数据保存路径
% #TODO #HANDWRITTEN: 路径内容难辨; 蓝色手写路径过于模糊
\fieldvalue{\handwritten{[路径内容难辨]}}

\medskip
\noindent
检测结果：
% #VALUE_ID: LEX-P0055-V0003
% #FIELD_VALUE: 一级反应器
\fieldvalue{\checkboxfield{unchecked}} 一级反应器；
% #VALUE_ID: LEX-P0055-V0004
% #FIELD_VALUE: 二级反应器
% #HANDWRITTEN: ✓
\fieldvalue{\handwritten{\checkboxfield{checked}}} 二级反应器；
% #VALUE_ID: LEX-P0055-V0005
% #FIELD_VALUE: 三级反应器
\fieldvalue{\checkboxfield{unchecked}} 三级反应器。

\medskip
\begin{center}
\begin{tabular}{|p{0.10\linewidth}|p{0.18\linewidth}|p{0.17\linewidth}|p{0.12\linewidth}|p{0.12\linewidth}|p{0.12\linewidth}|p{0.12\linewidth}|}
\hline
\multicolumn{7}{|c|}{检测结果} \\ \hline
检测次数 &
活细胞浓度（个细胞/ml） &
样品密度（个细胞/ml） &
活细胞密度 CV 值（\%） &
细胞污染率（\%） &
样品活率（\%） &
细胞活率 CV 值（\%） \\ \hline

第一次 &
% #VALUE_ID: LEX-P0055-V0006
% #FIELD_VALUE: 第一次活细胞浓度
% #HANDWRITTEN: 3.23×10^6
\fieldvalue{\handwritten{$3.23\times10^{6}$}} &
&
&
% #VALUE_ID: LEX-P0055-V0007
% #FIELD_VALUE: 第一次细胞污染率
% #HANDWRITTEN: 95.5
\fieldvalue{\handwritten{95.5}} &
& \\ \hline

第二次 &
% #VALUE_ID: LEX-P0055-V0008
% #FIELD_VALUE: 第二次活细胞浓度
% #HANDWRITTEN: 2.3×10^6
\fieldvalue{\handwritten{$2.3\times10^{6}$}} &
% #VALUE_ID: LEX-P0055-V0009
% #FIELD_VALUE: 样品密度
% #HANDWRITTEN: 2.23×10^6
\fieldvalue{\handwritten{$2.23\times10^{6}$}} &
% #VALUE_ID: LEX-P0055-V0010
% #FIELD_VALUE: 活细胞密度CV值
% #HANDWRITTEN: 2
\fieldvalue{\handwritten{2}} &
% #VALUE_ID: LEX-P0055-V0011
% #FIELD_VALUE: 第二次细胞污染率
% #HANDWRITTEN: 96.1
\fieldvalue{\handwritten{96.1}} &
% #VALUE_ID: LEX-P0055-V0012
% #FIELD_VALUE: 第二次样品活率
% #HANDWRITTEN: 95
\fieldvalue{\handwritten{95}} &
% #VALUE_ID: LEX-P0055-V0013
% #FIELD_VALUE: 细胞活率CV值
% #HANDWRITTEN: 1
\fieldvalue{\handwritten{1}} \\ \hline

第三次 &
% #VALUE_ID: LEX-P0055-V0014
% #FIELD_VALUE: 第三次活细胞浓度
% #HANDWRITTEN: 3.33×10^6
\fieldvalue{\handwritten{$3.33\times10^{6}$}} &
&
&
% #VALUE_ID: LEX-P0055-V0015
% #FIELD_VALUE: 第三次细胞污染率
% #HANDWRITTEN: 94.7
\fieldvalue{\handwritten{94.7}} &
& \\ \hline
\end{tabular}
\end{center}

\noindent
\begin{tabular}{|p{0.25\linewidth}|p{0.25\linewidth}|p{0.25\linewidth}|p{0.20\linewidth}|}
\hline
发酵罐细胞体积（ml） &
% #VALUE_ID: LEX-P0055-V0016
% #FIELD_VALUE: 发酵罐细胞体积
% #HANDWRITTEN: 3700
\fieldvalue{\handwritten{3700}} &
细胞数量（个细胞/罐） &
% #VALUE_ID: LEX-P0055-V0017
% #FIELD_VALUE: 细胞数量
% #HANDWRITTEN: 3.7×10^10
\fieldvalue{\handwritten{$3.7\times10^{10}$}} \\ \hline
\multicolumn{4}{|p{0.95\linewidth}|}{
计算要求：样品密度为 $3$ 次活细胞密度测定结果取平均值，以 $1\%$ 人白细胞溶液悬液总体积，再除以生物反应器细胞液体积，细胞数量为样品密度乘以发酵罐体积。
} \\ \hline
\end{tabular}

\medskip
\noindent
\begin{tabular}{|p{0.15\linewidth}|p{0.80\linewidth}|}
\hline
备注 &
% #VALUE_ID: LEX-P0055-V0018
% #FIELD_VALUE: 备注
% #HANDWRITTEN: /
\fieldvalue{\handwritten{/}} \\ \hline
实验人员/日期 &
% #VALUE_ID: LEX-P0055-V0019
% #FIELD_VALUE: 实验人员
% #TODO #HANDWRITTEN: 姓名难辨; 蓝色手写姓名不清晰
\fieldvalue{\handwritten{[姓名难辨]}}\quad
% #VALUE_ID: LEX-P0055-V0020
% #FIELD_VALUE: 实验日期
% #HANDWRITTEN: 2025.09.20
\fieldvalue{\handwritten{2025.09.20}} \\ \hline
复核人/日期 &
% #VALUE_ID: LEX-P0055-V0021
% #FIELD_VALUE: 复核人
% #TODO #HANDWRITTEN: 姓名难辨; 蓝色手写姓名不清晰
\fieldvalue{\handwritten{[姓名难辨]}}\quad
% #VALUE_ID: LEX-P0055-V0022
% #FIELD_VALUE: 复核日期
% #HANDWRITTEN: 2025.09.20
\fieldvalue{\handwritten{2025.09.20}} \\ \hline
\end{tabular}

\begin{center}
第 3 页 共 5 页
\end{center}
% LEXOID_PAGE_COMPLETED: 55/70

\newpage

\begin{center}
\textbf{生物反应器细胞液检验记录}\\
文件编码：REC-YZ-SOP-QC-99-006-a
\end{center}

\begin{center}
\begin{tabular}{|p{0.16\linewidth}|p{0.30\linewidth}|p{0.16\linewidth}|p{0.26\linewidth}|}
\hline
\multicolumn{4}{|c|}{细胞生长状态检验}\\
\hline
样品名称 &
% #VALUE_ID: LEX-P0056-V0001
% #FIELD_VALUE: 样品名称
% #HANDWRITTEN: 二级生物反应器细胞液
\fieldvalue{\handwritten{二级生物反应器细胞液}} &
样品规格 &
% #VALUE_ID: LEX-P0056-V0002
% #FIELD_VALUE: 样品规格
% #HANDWRITTEN: 2?
\fieldvalue{\handwritten{2?}}\,ml/袋\\
\hline
样品批号 &
% #VALUE_ID: LEX-P0056-V0003
% #FIELD_VALUE: 样品批号
% TODO #HANDWRITTEN: A?202?3800; handwriting partially illegible
\fieldvalue{\handwritten{A?202?3800}} &
取样日期 &
% #VALUE_ID: LEX-P0056-V0004
% #FIELD_VALUE: 取样日期
% TODO #HANDWRITTEN: 2024年8月?日; handwriting partially illegible
\fieldvalue{\handwritten{2024年8月?日}}\\
\hline
\end{tabular}
\end{center}

\section*{1.\ 仪器设备及试剂}

\subsection*{1.1\ 仪器设备}

\begin{center}
\begin{tabular}{|p{0.20\linewidth}|p{0.16\linewidth}|p{0.18\linewidth}|p{0.20\linewidth}|p{0.18\linewidth}|}
\hline
仪器名称 & 厂家 & 型号 & 设备编号 & 检查确认\\
\hline
% #VALUE_ID: LEX-P0056-V0005
% #FIELD_VALUE: 仪器名称
\fieldvalue{荧光显微镜} &
% #VALUE_ID: LEX-P0056-V0006
% #FIELD_VALUE: 厂家
\fieldvalue{徕卡} &
% #VALUE_ID: LEX-P0056-V0007
% #FIELD_VALUE: 型号
\fieldvalue{Dmi LED} &
% #VALUE_ID: LEX-P0056-V0008
% #FIELD_VALUE: 设备编号
\fieldvalue{1431404} &
% #VALUE_ID: LEX-P0056-V0009
% #FIELD_VALUE: 正常
\fieldvalue{\checkboxfield{checked}}\,正常\quad
% #VALUE_ID: LEX-P0056-V0010
% #FIELD_VALUE: 异常
\fieldvalue{\checkboxfield{unchecked}}\,异常\\
\hline
\end{tabular}
\end{center}

\subsection*{1.2\ 试剂耗材}

\begin{center}
\begin{tabular}{|p{0.25\linewidth}|p{0.25\linewidth}|p{0.19\linewidth}|p{0.15\linewidth}|p{0.16\linewidth}|}
\hline
名称 & 厂家 & 货号 & 批号 & 有效期至\\
\hline
% #VALUE_ID: LEX-P0056-V0011
% #FIELD_VALUE: 名称
\fieldvalue{细胞活力检测试剂盒} &
% #VALUE_ID: LEX-P0056-V0012
% #FIELD_VALUE: 厂家
\fieldvalue{江苏凯基生物技术股份有限公司} &
% #VALUE_ID: LEX-P0056-V0013
% #FIELD_VALUE: 货号
\fieldvalue{KGA9501-1000} &
% #VALUE_ID: LEX-P0056-V0014
% #FIELD_VALUE: 批号
% TODO #HANDWRITTEN: 202?4010; handwriting partially illegible
\fieldvalue{\handwritten{202?4010}} &
% #VALUE_ID: LEX-P0056-V0015
% #FIELD_VALUE: 有效期至
% TODO #HANDWRITTEN: 2025.10.?; handwriting partially illegible
\fieldvalue{\handwritten{2025.10.?}}\\
\hline
\end{tabular}
\end{center}

\section*{2.\ 操作步骤}

\begin{center}
\begin{tabular}{|p{0.52\linewidth}|p{0.42\linewidth}|}
\hline
\multicolumn{2}{|c|}{溶液配制}\\
\hline
20$\times$的荧光染液配制：取
% #VALUE_ID: LEX-P0056-V0016
% #FIELD_VALUE: 组分A体积
% TODO #HANDWRITTEN: 0.2; handwriting partially illegible
\fieldvalue{\handwritten{0.2}}\,\textmu l
组分A和
% #VALUE_ID: LEX-P0056-V0017
% #FIELD_VALUE: 组分B体积
% TODO #HANDWRITTEN: 0.2; handwriting partially illegible
\fieldvalue{\handwritten{0.2}}\,\textmu l
组分B与
% #VALUE_ID: LEX-P0056-V0018
% #FIELD_VALUE: PBS体积
% TODO #HANDWRITTEN: 5?; handwriting partially illegible
\fieldvalue{\handwritten{5?}}\,\textmu l
PBS混合均匀配成20$\times$的荧光染液。 &
% #VALUE_ID: LEX-P0056-V0019
% #FIELD_VALUE: 配制后总体积
% TODO #HANDWRITTEN: 580; handwriting partially illegible
\fieldvalue{\handwritten{580}}\,\textmu l PBS\\
\hline
\multicolumn{2}{|c|}{操作过程}\\
\hline
取混合后样品200\,\textmu l/孔至96孔板中，取2孔； &
样品
% #VALUE_ID: LEX-P0056-V0020
% #FIELD_VALUE: 样品用量
% #HANDWRITTEN: 200
\fieldvalue{\handwritten{200}}\,\textmu l/孔\\
& 每个样品取
% #VALUE_ID: LEX-P0056-V0021
% #FIELD_VALUE: 每个样品取孔数
% #HANDWRITTEN: 2
\fieldvalue{\handwritten{2}}\,孔\\
\hline
自然沉降后，吸弃上清，注意不要吸到沉淀样品； &
% #VALUE_ID: LEX-P0056-V0022
% #FIELD_VALUE: 吸弃上清
\fieldvalue{\checkboxfield{checked}}\,吸弃上清\\
\hline
取20$\times$的荧光染液加入到96孔板中，200\,\textmu l/孔，避光孵育30--60\,min。 &
20$\times$的荧光染液
% #VALUE_ID: LEX-P0056-V0023
% #FIELD_VALUE: 荧光染液用量
% #HANDWRITTEN: 200
\fieldvalue{\handwritten{200}}\,\textmu l/孔\\
& 避光孵育
% #VALUE_ID: LEX-P0056-V0024
% #FIELD_VALUE: 避光孵育开始时间
% #HANDWRITTEN: 09:30
\fieldvalue{\handwritten{09:30}}\,--
% #VALUE_ID: LEX-P0056-V0025
% #FIELD_VALUE: 避光孵育结束时间
% #HANDWRITTEN: 10:00
\fieldvalue{\handwritten{10:00}}\\
\hline
染色结束后，吸弃上清，取PBS加入到96孔板中，200\,\textmu l/孔，吸弃上清，注意不要吸到沉淀样品，清洗3次，之后每孔加入100\,\textmu l PBS进行拍照。 &
加入PBS
% #VALUE_ID: LEX-P0056-V0026
% #FIELD_VALUE: PBS用量
% #HANDWRITTEN: 200
\fieldvalue{\handwritten{200}}\,\textmu l/孔\\
& 清洗
% #VALUE_ID: LEX-P0056-V0027
% #FIELD_VALUE: 清洗次数
% #HANDWRITTEN: 3
\fieldvalue{\handwritten{3}}\,次\\
\hline
将荧光显微镜调至4$\times$物镜，软件中调节节相应物镜参数，调节到合适的光源和像距，拍摄4张清晰图像。 &
\\
\hline
进一步清晰图像，打印并粘贴在实验记录上。 &
\\
\hline
数据保存路径 &
% #VALUE_ID: LEX-P0056-V0028
% #FIELD_VALUE: 数据保存路径
% TODO #HANDWRITTEN: 127\textbackslash 4电脑\textbackslash 2蓄--原图?--2024; handwriting partially illegible
\fieldvalue{\handwritten{127\textbackslash 4电脑\textbackslash 2蓄--原图?--2024}}\\
\hline
\multicolumn{2}{|c|}{检测结果}\\
\hline
\end{tabular}
\end{center}

\begin{center}
第4页 共5页
\end{center}
% LEXOID_PAGE_COMPLETED: 56/70

\newpage

\noindent
\begin{minipage}{\linewidth}
\small
\textbf{铂金唯实生物科技（北京）有限公司}\hfill
文件编码号：REC-YZ-SOP-QC-99-006-a

\medskip
\centering
\begin{tabular}{|c|}
\hline
原件\\
\hline
\end{tabular}
\hfill
\begin{tabular}{|c|}
\hline
副本号/副本\\
\hline
\end{tabular}
\end{minipage}

\vspace{2cm}

\begin{center}
% TODO: 此处为显微图像；原 PDF 未提供可用的图像文件名。
\fbox{\rule{0pt}{3.4cm}\rule{9.2cm}{0pt}}

% #VALUE_ID: LEX-P0057-V0001
% #TODO #HANDWRITTEN: illegible; blue handwritten annotation beside the image
\handwritten{[illegible handwritten annotation]}

% #VALUE_ID: LEX-P0057-V0002
% #FIELD_VALUE: 图像旁日期注记
% #HANDWRITTEN: 2025.08.20
\fieldvalue{\handwritten{2025.08.20}}
\end{center}

\vfill

\noindent
\begin{tabular}{|p{2.0cm}|p{12.0cm}|}
\hline
备注 &
% #VALUE_ID: LEX-P0057-V0003
% #FIELD_VALUE: 备注
% #HANDWRITTEN: /
\fieldvalue{\handwritten{/}}\\
\hline
实验人/日期 &
% #VALUE_ID: LEX-P0057-V0004
% #FIELD_VALUE: 实验人
% #TODO #HANDWRITTEN: illegible; cursive Chinese name
\fieldvalue{\handwritten{[illegible Chinese name]}}
\quad
% #VALUE_ID: LEX-P0057-V0005
% #FIELD_VALUE: 实验日期
% #HANDWRITTEN: 2025.08.20
\fieldvalue{\handwritten{2025.08.20}}\\
\hline
复核人/日期 &
% #VALUE_ID: LEX-P0057-V0006
% #FIELD_VALUE: 复核人
% #TODO #HANDWRITTEN: illegible; cursive Chinese name
\fieldvalue{\handwritten{[illegible Chinese name]}}
\quad
% #VALUE_ID: LEX-P0057-V0007
% #FIELD_VALUE: 复核日期
% #HANDWRITTEN: 2025.08.20
\fieldvalue{\handwritten{2025.08.20}}\\
\hline
\end{tabular}

\begin{center}
\small 第 5 页 共 5 页
\end{center}
% LEXOID_PAGE_COMPLETED: 57/70

\newpage

\noindent
\textbf{Assay:}
% #VALUE_ID: LEX-P0058-V0001
% #FIELD_VALUE: Assay
\fieldvalue{MSC-1}

\vspace{0.5em}

\noindent
\textbf{Cell Type F1:}
% #VALUE_ID: LEX-P0058-V0002
% #FIELD_VALUE: Cell Type F1
\fieldvalue{MSC (AO)}

\noindent
\textbf{Cell Type F2:}
% #VALUE_ID: LEX-P0058-V0003
% #FIELD_VALUE: Cell Type F2
\fieldvalue{MSC (PI)}

\vspace{0.5em}

\noindent
\textbf{Sample ID:}
% #VALUE_ID: LEX-P0058-V0004
% #FIELD_VALUE: Sample ID
\fieldvalue{QC-JS-QY0402-A31Z2202508006-20250820-70H-1}

\noindent
\textbf{Dilution:}
% #VALUE_ID: LEX-P0058-V0005
% #FIELD_VALUE: Dilution
\fieldvalue{2.00}

\vspace{0.5em}
\hrule
\vspace{0.3em}

\noindent
Results:

\vspace{0.5em}

\noindent
\begin{tabular}{@{}p{0.29\linewidth}p{0.34\linewidth}p{0.29\linewidth}@{}}
\textbf{Count} & \textbf{Concentration} & \textbf{Mean Diameter} \\
\texttt{==========} & \texttt{=============} & \texttt{=============} \\[0.2em]

Total:
% #VALUE_ID: LEX-P0058-V0006
% #FIELD_VALUE: Total Count
\fieldvalue{959}
&
% #VALUE_ID: LEX-P0058-V0007
% #FIELD_VALUE: Total Concentration
\fieldvalue{\(3.38\times10^{6}\) cells/mL}
&
% #VALUE_ID: LEX-P0058-V0008
% #FIELD_VALUE: Total Mean Diameter
\fieldvalue{14.0 microns}
\\

Live:
% #VALUE_ID: LEX-P0058-V0009
% #FIELD_VALUE: Live Count
\fieldvalue{916}
&
% #VALUE_ID: LEX-P0058-V0010
% #FIELD_VALUE: Live Concentration
\fieldvalue{\(3.23\times10^{6}\) cells/mL}
&
% #VALUE_ID: LEX-P0058-V0011
% #FIELD_VALUE: Live Mean Diameter
\fieldvalue{14.0 microns}
\\

Dead:
% #VALUE_ID: LEX-P0058-V0012
% #FIELD_VALUE: Dead Count
\fieldvalue{43}
&
% #VALUE_ID: LEX-P0058-V0013
% #FIELD_VALUE: Dead Concentration
\fieldvalue{\(1.51\times10^{5}\) cells/mL}
&
% #VALUE_ID: LEX-P0058-V0014
% #FIELD_VALUE: Dead Mean Diameter
\fieldvalue{13.0 microns}
\\

Viability:
% #VALUE_ID: LEX-P0058-V0015
% #FIELD_VALUE: Viability
\fieldvalue{95.5\%}
&
&
\end{tabular}

\vspace{1.5em}

\noindent
\begin{tabular}{@{}cc@{}}
\fbox{\parbox[c][4.0cm][c]{0.39\linewidth}{\centering TODO: image placeholder}} &
\fbox{\parbox[c][4.0cm][c]{0.39\linewidth}{\centering TODO: image placeholder}} \\[0.8em]
\fbox{\parbox[c][4.0cm][c]{0.39\linewidth}{\centering TODO: image placeholder}} &
\fbox{\parbox[c][4.0cm][c]{0.39\linewidth}{\centering TODO: image placeholder}}
\end{tabular}

\vfill

\noindent
操作人/审核人：
% #VALUE_ID: LEX-P0058-V0016
% #FIELD_VALUE: 操作人/审核人
% TODO #HANDWRITTEN: 疑似“强波”；签名字迹不清
\fieldvalue{\handwritten{强波}}
\hfill
日期：
% #VALUE_ID: LEX-P0058-V0017
% #FIELD_VALUE: 日期
% #HANDWRITTEN: 2025.08.20
\fieldvalue{\handwritten{2025.08.20}}
% LEXOID_PAGE_COMPLETED: 58/70

\newpage

\noindent
% #VALUE_ID: LEX-P0059-V0001
% #FIELD_VALUE: Assay
Assay: \fieldvalue{MSC-1}

\medskip
\noindent
% #VALUE_ID: LEX-P0059-V0002
% #FIELD_VALUE: Cell Type F1
Cell Type F1: \fieldvalue{MSC (AO)}

\noindent
% #VALUE_ID: LEX-P0059-V0003
% #FIELD_VALUE: Cell Type F2
Cell Type F2: \fieldvalue{MSC (PI)}

\medskip
\noindent
% #VALUE_ID: LEX-P0059-V0004
% #FIELD_VALUE: Sample ID
Sample ID: \fieldvalue{-QC-JS-QY0402-A31Z2202508006-20250820-70H-2\_20250820\_101639-2}

\noindent
% #VALUE_ID: LEX-P0059-V0005
% #FIELD_VALUE: Dilution
Dilution: \fieldvalue{2.00}

\medskip
\noindent\rule{\linewidth}{0.4pt}

\noindent Results:

\medskip
\noindent
\begin{tabular}{@{}p{0.29\linewidth}p{0.34\linewidth}p{0.29\linewidth}@{}}
\textbf{Count} & \textbf{Concentration} & \textbf{Mean Diameter} \\
\underline{\hspace{0.21\linewidth}} &
\underline{\hspace{0.24\linewidth}} &
\underline{\hspace{0.22\linewidth}} \\[2pt]

Total:
% #VALUE_ID: LEX-P0059-V0006
% #FIELD_VALUE: Count — Total
\fieldvalue{960}
&
% #VALUE_ID: LEX-P0059-V0007
% #FIELD_VALUE: Concentration — Total
\fieldvalue{$3.38\times10^{6}$ cells/mL}
&
% #VALUE_ID: LEX-P0059-V0008
% #FIELD_VALUE: Mean Diameter — Total
\fieldvalue{13.3 micron}
\\

Live:
% #VALUE_ID: LEX-P0059-V0009
% #FIELD_VALUE: Count — Live
\fieldvalue{923}
&
% #VALUE_ID: LEX-P0059-V0010
% #FIELD_VALUE: Concentration — Live
\fieldvalue{$3.25\times10^{6}$ cells/mL}
&
% #VALUE_ID: LEX-P0059-V0011
% #FIELD_VALUE: Mean Diameter — Live
\fieldvalue{13.3 microns}
\\

Dead:
% #VALUE_ID: LEX-P0059-V0012
% #FIELD_VALUE: Count — Dead
\fieldvalue{37}
&
% #VALUE_ID: LEX-P0059-V0013
% #FIELD_VALUE: Concentration — Dead
\fieldvalue{$1.30\times10^{6}$ cells/mL}
&
% #VALUE_ID: LEX-P0059-V0014
% #FIELD_VALUE: Mean Diameter — Dead
\fieldvalue{12.5 microns}
\\

Viability:
% #VALUE_ID: LEX-P0059-V0015
% #FIELD_VALUE: Viability
\fieldvalue{96.2\%}
&
&
\end{tabular}

\medskip

% TODO: Four microscopy images are visible in a 2-by-2 arrangement; no image filenames are available.
\begin{center}
\begin{tabular}{@{}p{0.46\linewidth}@{\hspace{0.04\linewidth}}p{0.46\linewidth}@{}}
\centering
\fbox{\parbox[c][4.0cm][c]{0.95\linewidth}{\centering TODO: microscopy image}}
&
\centering
\fbox{\parbox[c][4.0cm][c]{0.95\linewidth}{\centering TODO: microscopy image}}
\tabularnewline[0.5cm]
\centering
\fbox{\parbox[c][4.0cm][c]{0.95\linewidth}{\centering TODO: microscopy image}}
&
\centering
\fbox{\parbox[c][4.0cm][c]{0.95\linewidth}{\centering TODO: microscopy image}}
\tabularnewline
\end{tabular}
\end{center}

\vfill

\noindent
操作人/审核人：
% #VALUE_ID: LEX-P0059-V0016
% #FIELD_VALUE: 操作人/审核人
% TODO #HANDWRITTEN: illegible Chinese signature; handwriting is unclear
\fieldvalue{\handwritten{签名}}

\hfill
日期：
% #VALUE_ID: LEX-P0059-V0017
% #FIELD_VALUE: 日期
% TODO #HANDWRITTEN: 2025-08-20; handwritten date is partially unclear
\fieldvalue{\handwritten{2025-08-20}}
% LEXOID_PAGE_COMPLETED: 59/70

\newpage

\noindent\textbf{Assay:} 
% #VALUE_ID: LEX-P0060-V0001
% #FIELD_VALUE: Assay
\fieldvalue{MSC-1}

\medskip

\noindent\textbf{Cell Type F1:} 
% #VALUE_ID: LEX-P0060-V0002
% #FIELD_VALUE: Cell Type F1
\fieldvalue{MSC (AO)}

\noindent\textbf{Cell Type F2:} 
% #VALUE_ID: LEX-P0060-V0003
% #FIELD_VALUE: Cell Type F2
\fieldvalue{MSC (PI)}

\medskip

\noindent\textbf{Sample ID:} 
% #VALUE_ID: LEX-P0060-V0004
% #FIELD_VALUE: Sample ID
\fieldvalue{-QC-JS-QY0402-A31Z2202508006-20250820-70H-3\_20250820\_101841-2}

\noindent\textbf{Dilution:} 
% #VALUE_ID: LEX-P0060-V0005
% #FIELD_VALUE: Dilution
\fieldvalue{2.00}

\medskip
\hrule
\medskip

\noindent Results:

\medskip

\noindent
\begin{tabular}{@{}p{0.28\linewidth}p{0.34\linewidth}p{0.28\linewidth}@{}}
\textbf{Count} & \textbf{Concentration} & \textbf{Mean Diameter} \\[-1pt]
\underline{\hspace{0.22\linewidth}} &
\underline{\hspace{0.27\linewidth}} &
\underline{\hspace{0.22\linewidth}} \\[3pt]
Total:
% #VALUE_ID: LEX-P0060-V0006
% #FIELD_VALUE: Total count
\fieldvalue{1000}
&
% #VALUE_ID: LEX-P0060-V0007
% #FIELD_VALUE: Total concentration
\fieldvalue{$3.52\times10^{6}$ cells/mL}
&
% #VALUE_ID: LEX-P0060-V0008
% #FIELD_VALUE: Total mean diameter
\fieldvalue{13.2 micron}
\\
Live:
% #VALUE_ID: LEX-P0060-V0009
% #FIELD_VALUE: Live count
\fieldvalue{947}
&
% #VALUE_ID: LEX-P0060-V0010
% #FIELD_VALUE: Live concentration
\fieldvalue{$3.33\times10^{6}$ cells/mL}
&
% #VALUE_ID: LEX-P0060-V0011
% #FIELD_VALUE: Live mean diameter
\fieldvalue{13.2 microns}
\\
Dead:
% #VALUE_ID: LEX-P0060-V0012
% #FIELD_VALUE: Dead count
\fieldvalue{53}
&
% #VALUE_ID: LEX-P0060-V0013
% #FIELD_VALUE: Dead concentration
\fieldvalue{$1.86\times10^{5}$ cells/mL}
&
% #VALUE_ID: LEX-P0060-V0014
% #FIELD_VALUE: Dead mean diameter
\fieldvalue{12.9 microns}
\\
Viability:
% #VALUE_ID: LEX-P0060-V0015
% #FIELD_VALUE: Viability
\fieldvalue{94.7\%}
&
&
\end{tabular}

\bigskip

% TODO: Four microscopy images are visible on this page, but no source filenames are available.
\begin{center}
\begin{tabular}{@{}c@{\hspace{0.04\linewidth}}c@{}}
\fbox{\parbox[c][0.22\textheight][c]{0.40\linewidth}{\centering TODO: microscopy image}} &
\fbox{\parbox[c][0.22\textheight][c]{0.40\linewidth}{\centering TODO: microscopy image}} \\[8pt]
\fbox{\parbox[c][0.22\textheight][c]{0.40\linewidth}{\centering TODO: microscopy image}} &
\fbox{\parbox[c][0.22\textheight][c]{0.40\linewidth}{\centering TODO: microscopy image}}
\end{tabular}
\end{center}

\vfill

\noindent 操作人/审核人：
% #VALUE_ID: LEX-P0060-V0016
% #FIELD_VALUE: 操作人/审核人
% #TODO #HANDWRITTEN: 疑似“陈波”; 签名难以辨认
\fieldvalue{\handwritten{陈波}}
\hfill
日期：
% #VALUE_ID: LEX-P0060-V0017
% #FIELD_VALUE: 日期
% #HANDWRITTEN: 2025-8-20
\fieldvalue{\handwritten{2025-8-20}}
% LEXOID_PAGE_COMPLETED: 60/70

\newpage

\begin{center}
{\small
\textbf{铂金瑞生物科技（北京）有限公司}\\
\textit{PLATINUM LIFE EXCELLENCE BIOCHEMICAL}\\[2pt]
\textbf{请验单}
}
\end{center}

\noindent
检验单编号：REC-YZ-SMP-QC-01-001-d

\begin{center}
\renewcommand{\arraystretch}{1.45}
\small
\begin{tabular}{|p{0.16\linewidth}|p{0.34\linewidth}|p{0.18\linewidth}|p{0.24\linewidth}|}
\hline
样品名称 &
% #VALUE_ID: LEX-P0061-V0001
% #FIELD_VALUE: 样品名称
% #TODO #HANDWRITTEN: 三氯硝基苯; handwritten text is partially unclear
\fieldvalue{\handwritten{三氯硝基苯}} &
样品代码 &
% #VALUE_ID: LEX-P0061-V0002
% #FIELD_VALUE: 样品代码
\fieldvalue{240403}
\\ \hline

样品批号 &
% #VALUE_ID: LEX-P0061-V0003
% #FIELD_VALUE: 样品批号
% #TODO #HANDWRITTEN: 2024.02.06; handwritten batch number is unclear
\fieldvalue{\handwritten{2024.02.06}} &
样品规格/数量 &
% #VALUE_ID: LEX-P0061-V0004
% #FIELD_VALUE: 样品规格/数量
% #TODO #HANDWRITTEN: 20ml/袋×1; handwritten entry is partially unclear
\fieldvalue{\handwritten{20ml/袋×1}}
\\ \hline

样品类别 &
\multicolumn{3}{p{0.76\linewidth}|}{
% #VALUE_ID: LEX-P0061-V0005
% #FIELD_VALUE: 样品类别：生产过程中的检验样品
\fieldvalue{\checkboxfield{checked}} 生产过程中的检验样品\quad
% #VALUE_ID: LEX-P0061-V0006
% #FIELD_VALUE: 样品类别：包装产品
\fieldvalue{\checkboxfield{unchecked}} 包装产品\quad
% #VALUE_ID: LEX-P0061-V0007
% #FIELD_VALUE: 样品类别：成品
\fieldvalue{\checkboxfield{unchecked}} 成品\par\noindent
% #VALUE_ID: LEX-P0061-V0008
% #FIELD_VALUE: 样品类别：其他
\fieldvalue{\checkboxfield{unchecked}} 其他（\hspace{1.5cm}）
}
\\ \hline

检验类型 &
\multicolumn{3}{p{0.76\linewidth}|}{
% #VALUE_ID: LEX-P0061-V0009
% #FIELD_VALUE: 检验类型：初验
\fieldvalue{\checkboxfield{checked}} 初验\quad
% #VALUE_ID: LEX-P0061-V0010
% #FIELD_VALUE: 检验类型：复验
\fieldvalue{\checkboxfield{unchecked}} 复验\quad
% #VALUE_ID: LEX-P0061-V0011
% #FIELD_VALUE: 检验类型：退货
\fieldvalue{\checkboxfield{unchecked}} 退货\quad
% #VALUE_ID: LEX-P0061-V0012
% #FIELD_VALUE: 检验类型：稳定性
\fieldvalue{\checkboxfield{unchecked}} 稳定性\quad
% #VALUE_ID: LEX-P0061-V0013
% #FIELD_VALUE: 检验类型：留样
\fieldvalue{\checkboxfield{unchecked}} 留样
}
\\ \hline

请验目的 &
\multicolumn{3}{p{0.76\linewidth}|}{
% #VALUE_ID: LEX-P0061-V0014
% #FIELD_VALUE: 请验目的：全检
\fieldvalue{\checkboxfield{checked}} 全检\quad
% #VALUE_ID: LEX-P0061-V0015
% #FIELD_VALUE: 请验目的：留样
\fieldvalue{\checkboxfield{unchecked}} 留样\quad
% #VALUE_ID: LEX-P0061-V0016
% #FIELD_VALUE: 请验目的：其他
\fieldvalue{\checkboxfield{unchecked}} 其他（\hspace{1.5cm}）
}
\\ \hline

请验部门 &
\multicolumn{3}{p{0.76\linewidth}|}{
% #VALUE_ID: LEX-P0061-V0017
% #FIELD_VALUE: 请验部门：生产部
\fieldvalue{\checkboxfield{checked}} 生产部\quad
% #VALUE_ID: LEX-P0061-V0018
% #FIELD_VALUE: 请验部门：物料管理部
\fieldvalue{\checkboxfield{unchecked}} 物料管理部\quad
% #VALUE_ID: LEX-P0061-V0019
% #FIELD_VALUE: 请验部门：其他
\fieldvalue{\checkboxfield{unchecked}} 其他（\hspace{1.5cm}）
}
\\ \hline

样品情况 &
\multicolumn{3}{p{0.76\linewidth}|}{
% #VALUE_ID: LEX-P0061-V0020
% #FIELD_VALUE: 样品情况：随请验单一起提供
\fieldvalue{\checkboxfield{checked}} 随请验单一起提供\quad
% #VALUE_ID: LEX-P0061-V0021
% #FIELD_VALUE: 样品情况：未随请验单一起提供
\fieldvalue{\checkboxfield{unchecked}} 未随请验单一起提供
}
\\ \hline

请验人 &
% #VALUE_ID: LEX-P0061-V0022
% #FIELD_VALUE: 请验人
% #TODO #HANDWRITTEN: 崔敏; handwritten name is partially unclear
\fieldvalue{\handwritten{崔敏}} &
请验日期 &
% #VALUE_ID: LEX-P0061-V0023
% #FIELD_VALUE: 请验日期
% #HANDWRITTEN: 2025.08.23
\fieldvalue{\handwritten{2025.08.23}}
\\ \hline

收验人 &
% #VALUE_ID: LEX-P0061-V0024
% #FIELD_VALUE: 收验人
% #TODO #HANDWRITTEN: 姚丹; handwritten name is partially unclear
\fieldvalue{\handwritten{姚丹}} &
收验日期 &
% #VALUE_ID: LEX-P0061-V0025
% #FIELD_VALUE: 收验日期
% #HANDWRITTEN: 2025.8.23
\fieldvalue{\handwritten{2025.8.23}}
\\ \hline

备注 &
\multicolumn{3}{p{0.76\linewidth}|}{
% #VALUE_ID: LEX-P0061-V0026
% #FIELD_VALUE: 备注
% #TODO #HANDWRITTEN: 7.5h; handwritten notation is partially unclear
% #HANDWRITTEN: 7.5h
\fieldvalue{\handwritten{7.5h}}\qquad
% #VALUE_ID: LEX-P0061-V0027
% #FIELD_VALUE: 备注
% #TODO #HANDWRITTEN: 485mg (85.8); parenthetical notation is unclear
\fieldvalue{\handwritten{485mg (85.8)}}
}
\\ \hline
\end{tabular}
\end{center}
% LEXOID_PAGE_COMPLETED: 61/70

\newpage

\begin{center}
\textbf{金华生物制品有限公司}\\
\textit{PLATINUM LIP EXCELLENCE BIOTECH}\\[2mm]
文件编码：REC-YZ-SMP-QC-01-004-b \qquad 版本号：1.4\\[2mm]
\section*{取样记录}
\end{center}

\renewcommand{\arraystretch}{1.6}
\begin{tabular}{|p{0.14\linewidth}|p{0.35\linewidth}|p{0.14\linewidth}|p{0.27\linewidth}|}
\hline
名称 &
% #VALUE_ID: LEX-P0062-V0001
% #FIELD_VALUE: 名称
% #TODO #HANDWRITTEN: 冻干物质控品（部分字迹不清）
\fieldvalue{\handwritten{冻干物质控品}} &
代码 &
% #VALUE_ID: LEX-P0062-V0002
% #FIELD_VALUE: 代码
% #HANDWRITTEN: QC0103
\fieldvalue{\handwritten{QC0103}}
\\
\hline
规格 &
% #VALUE_ID: LEX-P0062-V0003
% #FIELD_VALUE: 规格
% #HANDWRITTEN: 70 ml/袋
\fieldvalue{\handwritten{70 ml/袋}} &
总数量 &
% #VALUE_ID: LEX-P0062-V0004
% #FIELD_VALUE: 总数量
% #HANDWRITTEN: 1袋
\fieldvalue{\handwritten{1袋}}
\\
\hline
批号 &
% #VALUE_ID: LEX-P0062-V0005
% #FIELD_VALUE: 批号
% #TODO #HANDWRITTEN: A?2?20508006；首部字符难以辨认
\fieldvalue{\handwritten{A?2?20508006}} &
进厂批号 &
% #VALUE_ID: LEX-P0062-V0006
% #FIELD_VALUE: 进厂批号
% #HANDWRITTEN: /
\fieldvalue{\handwritten{/}}
\\
\hline
供货单位/厂家 &
\multicolumn{3}{l|}{
% #VALUE_ID: LEX-P0062-V0007
% #FIELD_VALUE: 供货单位/厂家
\fieldvalue{/}
}\\
\hline
取样目的 &
\multicolumn{3}{l|}{
% #VALUE_ID: LEX-P0062-V0008
% #FIELD_VALUE: 正常检验
% #HANDWRITTEN: ✓
\fieldvalue{\handwritten{\checkboxfield{checked}}}\ 正常检验
\quad
% #VALUE_ID: LEX-P0062-V0009
% #FIELD_VALUE: 复检
\fieldvalue{\checkboxfield{unchecked}}\ 复检
\quad
% #VALUE_ID: LEX-P0062-V0010
% #FIELD_VALUE: 稳定性
\fieldvalue{\checkboxfield{unchecked}}\ 稳定性
}\\[-1mm]
&\multicolumn{3}{l|}{
% #VALUE_ID: LEX-P0062-V0011
% #FIELD_VALUE: 退货检验
\fieldvalue{\checkboxfield{unchecked}}\ 退货检验
\quad
% #VALUE_ID: LEX-P0062-V0012
% #FIELD_VALUE: 留样
\fieldvalue{\checkboxfield{unchecked}}\ 留样
\quad
% #VALUE_ID: LEX-P0062-V0013
% #FIELD_VALUE: 其他
\fieldvalue{\checkboxfield{unchecked}}\ 其他：\underline{\hspace{2cm}}
}\\
\hline
取样地点 &
\multicolumn{3}{l|}{
% #VALUE_ID: LEX-P0062-V0014
% #FIELD_VALUE: 产品出口
% #HANDWRITTEN: ✓
\fieldvalue{\handwritten{\checkboxfield{checked}}}\ 产品出口
\quad
% #VALUE_ID: LEX-P0062-V0015
% #FIELD_VALUE: 物料库
\fieldvalue{\checkboxfield{unchecked}}\ 物料库
\quad
% #VALUE_ID: LEX-P0062-V0016
% #FIELD_VALUE: 细胞库
\fieldvalue{\checkboxfield{unchecked}}\ 细胞库
\quad
% #VALUE_ID: LEX-P0062-V0017
% #FIELD_VALUE: 包装间
\fieldvalue{\checkboxfield{unchecked}}\ 包装间
}\\
\hline
贮存条件 &
\multicolumn{3}{l|}{
% #VALUE_ID: LEX-P0062-V0018
% #FIELD_VALUE: 2--8$^\circ$C
% #HANDWRITTEN: ✓
\fieldvalue{\handwritten{\checkboxfield{checked}}}\ 2--8$^\circ$C
\quad
% #VALUE_ID: LEX-P0062-V0019
% #FIELD_VALUE: -20$^\circ$C
\fieldvalue{\checkboxfield{unchecked}}\ -20$^\circ$C
\quad
% #VALUE_ID: LEX-P0062-V0020
% #FIELD_VALUE: -80$^\circ$C
\fieldvalue{\checkboxfield{unchecked}}\ -80$^\circ$C
}\\[-1mm]
&\multicolumn{3}{l|}{
% #VALUE_ID: LEX-P0062-V0021
% #FIELD_VALUE: 常温
\fieldvalue{\checkboxfield{unchecked}}\ 常温
\quad
% #VALUE_ID: LEX-P0062-V0022
% #FIELD_VALUE: 其他
\fieldvalue{\checkboxfield{unchecked}}\ 其他：\underline{\hspace{2cm}}
}\\
\hline
盛装容器 &
\multicolumn{3}{l|}{
% #VALUE_ID: LEX-P0062-V0023
% #FIELD_VALUE: 保温箱
% #HANDWRITTEN: ✓
\fieldvalue{\handwritten{\checkboxfield{checked}}}\ 保温箱
\quad
% #VALUE_ID: LEX-P0062-V0024
% #FIELD_VALUE: 泡沫箱
\fieldvalue{\checkboxfield{unchecked}}\ 泡沫箱
\quad
% #VALUE_ID: LEX-P0062-V0025
% #FIELD_VALUE: 手提篮
\fieldvalue{\checkboxfield{unchecked}}\ 手提篮
\quad
% #VALUE_ID: LEX-P0062-V0026
% #FIELD_VALUE: 其他
\fieldvalue{\checkboxfield{unchecked}}\ 其他：\underline{\hspace{2cm}}
}\\
\hline
外观复核 &
\multicolumn{3}{l|}{
% #VALUE_ID: LEX-P0062-V0027
% #FIELD_VALUE: 合格
% #HANDWRITTEN: ✓
\fieldvalue{\handwritten{\checkboxfield{checked}}}\ 合格
\quad
% #VALUE_ID: LEX-P0062-V0028
% #FIELD_VALUE: 不合格
\fieldvalue{\checkboxfield{unchecked}}\ 不合格
}\\
\hline
取样量 &
% #VALUE_ID: LEX-P0062-V0029
% #FIELD_VALUE: 取样量
% #HANDWRITTEN: 1袋
\fieldvalue{\handwritten{1袋}} &
检验量 &
% #VALUE_ID: LEX-P0062-V0030
% #FIELD_VALUE: 检验量
% #HANDWRITTEN: 1袋
\fieldvalue{\handwritten{1袋}}
\\
\cline{1-4}
留样量 &
\multicolumn{3}{l|}{
% #VALUE_ID: LEX-P0062-V0031
% #FIELD_VALUE: 留样量
% #HANDWRITTEN: /
\fieldvalue{\handwritten{/}}
}\\
\hline
取样人/日期 &
% #VALUE_ID: LEX-P0062-V0032
% #FIELD_VALUE: 取样人/日期
% #TODO #HANDWRITTEN: 姓名及日期，约为“2025.08.15”；姓名字迹难以辨认
\fieldvalue{\handwritten{（姓名不清）\quad 2025.08.15}} &
复核人/日期 &
% #VALUE_ID: LEX-P0062-V0033
% #FIELD_VALUE: 复核人/日期
% #TODO #HANDWRITTEN: 姓名及日期，约为“2025.08.23”；姓名字迹难以辨认
\fieldvalue{\handwritten{（姓名不清）\quad 2025.08.23}}
\\
\hline
备注 &
\multicolumn{3}{l|}{
% #VALUE_ID: LEX-P0062-V0034
% #FIELD_VALUE: 备注
% #TODO #HANDWRITTEN: 7? h；0.85 kg（括号内内容难以辨认）
\fieldvalue{\handwritten{7? h；0.85 kg（内容不清）}}
}\\
\hline
\end{tabular}

\vfill
\begin{center}
第 1 页 共 1 页
\end{center}
% LEXOID_PAGE_COMPLETED: 62/70

\newpage

\begin{center}
\textbf{生物反应器细胞液检验操作记录}
\end{center}

\begin{center}
\begin{tabular}{|p{2.8cm}|p{6.0cm}|p{2.2cm}|p{3.2cm}|}
\hline
\multicolumn{4}{|c|}{密度与活率检验} \\ \hline
样品名称 &
% #VALUE_ID: LEX-P0063-V0001
% #FIELD_VALUE: 样品名称
% #HANDWRITTEN: 三级生物反应器细胞液
\fieldvalue{\handwritten{三级生物反应器细胞液}} &
样品规格 &
% #VALUE_ID: LEX-P0063-V0002
% #FIELD_VALUE: 样品规格
% #HANDWRITTEN: 20
\fieldvalue{\handwritten{20}} ml/袋 \\ \hline
样品批号 &
% #VALUE_ID: LEX-P0063-V0003
% #FIELD_VALUE: 样品批号
% #TODO #HANDWRITTEN: A?120?08?6; 字迹不清
\fieldvalue{\handwritten{A?120?08?6}} &
取样日期 &
% #VALUE_ID: LEX-P0063-V0004
% #FIELD_VALUE: 取样日期
% #TODO #HANDWRITTEN: 202?年08月12日; 年份部分字迹不清
\fieldvalue{\handwritten{202?年08月12日}} \\ \hline
\end{tabular}
\end{center}

\section{仪器设备、试剂耗材}

\subsection{仪器设备}

\begin{center}
\begin{tabular}{|p{2.7cm}|p{2.3cm}|p{2.4cm}|p{2.5cm}|p{3.0cm}|p{2.5cm}|}
\hline
名称 & 厂家 & 型号 & 设备编号 & 校准有效期至 & 检查确认 \\ \hline
细胞计数仪 & Nexcelom & Cellometer K2 & 1410905 &
% #VALUE_ID: LEX-P0063-V0005
% #FIELD_VALUE: 校准有效期至
% #HANDWRITTEN: 2025.1.16
\fieldvalue{\handwritten{2025.1.16}} &
% #VALUE_ID: LEX-P0063-V0006
% #FIELD_VALUE: 正常
\fieldvalue{\checkboxfield{unchecked}} 正常
% #VALUE_ID: LEX-P0063-V0007
% #FIELD_VALUE: 异常
\fieldvalue{\checkboxfield{unchecked}} 异常 \\ \hline
高速冷冻离心机 & Thermo & ST16R & 1410810 &
% #VALUE_ID: LEX-P0063-V0008
% #FIELD_VALUE: 校准有效期至
% #TODO #HANDWRITTEN: 2025.05.09; 日期末位字迹不清
\fieldvalue{\handwritten{2025.05.09}} &
% #VALUE_ID: LEX-P0063-V0009
% #FIELD_VALUE: 正常
\fieldvalue{\checkboxfield{unchecked}} 正常
% #VALUE_ID: LEX-P0063-V0010
% #FIELD_VALUE: 异常
\fieldvalue{\checkboxfield{unchecked}} 异常 \\ \hline
电热恒温水槽 & 上海一恒 & DK-8AXX & 1411001 &
% #VALUE_ID: LEX-P0063-V0011
% #FIELD_VALUE: 校准有效期至
% #TODO #HANDWRITTEN: 2025.05.02; 日期末位字迹不清
\fieldvalue{\handwritten{2025.05.02}} &
% #VALUE_ID: LEX-P0063-V0012
% #FIELD_VALUE: 正常
\fieldvalue{\checkboxfield{unchecked}} 正常
% #VALUE_ID: LEX-P0063-V0013
% #FIELD_VALUE: 异常
\fieldvalue{\checkboxfield{unchecked}} 异常 \\ \hline
\end{tabular}
\end{center}

\subsection{试剂耗材}

\begin{center}
\begin{tabular}{|p{3.2cm}|p{3.7cm}|p{4.0cm}|p{3.2cm}|}
\hline
名称 & 厂家 & 批号 & 有效期至 \\ \hline
AO/PI 染液 & Nexcelom &
% #VALUE_ID: LEX-P0063-V0014
% #FIELD_VALUE: 批号
% #TODO #HANDWRITTEN: 33?0?; 字迹不清
\fieldvalue{\handwritten{33?0?}} &
% #VALUE_ID: LEX-P0063-V0015
% #FIELD_VALUE: 有效期至
% #HANDWRITTEN: 2026.01.02
\fieldvalue{\handwritten{2026.01.02}} \\ \hline
0.9\%氯化钠注射液 & 石家庄四药 &
% #VALUE_ID: LEX-P0063-V0016
% #FIELD_VALUE: 批号
% #TODO #HANDWRITTEN: 24/??/??; 批号字迹不清
\fieldvalue{\handwritten{24/??/??}} &
% #VALUE_ID: LEX-P0063-V0017
% #FIELD_VALUE: 有效期至
% #HANDWRITTEN: 2026.10.16
\fieldvalue{\handwritten{2026.10.16}} \\ \hline
20\%人血白蛋白 & Baxter AG &
% #VALUE_ID: LEX-P0063-V0018
% #FIELD_VALUE: 批号
% #TODO #HANDWRITTEN: A181?844; 字迹不清
\fieldvalue{\handwritten{A181?844}} &
% #VALUE_ID: LEX-P0063-V0019
% #FIELD_VALUE: 有效期至
% #HANDWRITTEN: 2027.02.28
\fieldvalue{\handwritten{2027.02.28}} \\ \hline
\end{tabular}
\end{center}

\section{操作步骤}

\begin{center}
\begin{tabular}{|p{7.4cm}|p{6.6cm}|}
\hline
\multicolumn{2}{|c|}{溶液准备} \\ \hline
1\%人自氯化锂溶液配制：取5 ml 20\%人血白蛋白加入到95 ml的0.9\%氯化钠注射液中，混匀即得。 &
20\%人血白蛋白
% #VALUE_ID: LEX-P0063-V0020
% #FIELD_VALUE: 20\%人血白蛋白用量
% #HANDWRITTEN: 5
\fieldvalue{\handwritten{5}} ml \\[-1mm]
& 0.9\%氯化钠注射液
% #VALUE_ID: LEX-P0063-V0021
% #FIELD_VALUE: 0.9\%氯化钠注射液用量
% #HANDWRITTEN: 95
\fieldvalue{\handwritten{95}} ml \\ \hline
名称 & 配制批号 \hspace{3cm} 有效期至 \\ \hline
5$\times$裂解液 &
% #VALUE_ID: LEX-P0063-V0022
% #FIELD_VALUE: 配制批号
% #HANDWRITTEN: 2025.07.30-01
\fieldvalue{\handwritten{2025.07.30-01}}
\hspace{1.0cm}
% #VALUE_ID: LEX-P0063-V0023
% #FIELD_VALUE: 有效期至
% #HANDWRITTEN: 2025.12.17
\fieldvalue{\handwritten{2025.12.17}} \\ \hline
\multicolumn{2}{|c|}{操作过程} \\ \hline
\multicolumn{2}{|p{7.4cm}|}{按照50 ml离心管的刻度读取样品体积，往样品中加入对应体积的5$\times$裂解液（每4 ml样品加入1 ml 5$\times$裂解液，按照“只进不含”原则，如26 ml样品，加入7 ml 5$\times$裂解液），充分混匀后倒入至2.3的样品中，置于37℃电热恒温水槽中裂解20 min后，取出离心管摇匀，使样品充分裂解；再置于37℃电热恒温水槽中裂解10 min得到细胞悬液。} \\ \hline
样品体积：
% #VALUE_ID: LEX-P0063-V0024
% #FIELD_VALUE: 样品体积
% #HANDWRITTEN: 38
\fieldvalue{\handwritten{38}} ml &
裂解液：
% #VALUE_ID: LEX-P0063-V0025
% #FIELD_VALUE: 裂解液体积
% #HANDWRITTEN: 10
\fieldvalue{\handwritten{10}} ml \\ \hline
电热恒温水槽
% #VALUE_ID: LEX-P0063-V0026
% #FIELD_VALUE: 电热恒温水槽温度
% #HANDWRITTEN: 37
\fieldvalue{\handwritten{37}}\,${}^{\circ}$C &
裂解液
% #VALUE_ID: LEX-P0063-V0027
% #FIELD_VALUE: 裂解时间
% #TODO #HANDWRITTEN: 15:04--15:34; 起始时间字迹不清
\fieldvalue{\handwritten{15:04--15:34}} \\ \hline
\multicolumn{2}{|p{14.0cm}|}{取出裂解好的细胞悬液300$\times g$离心5分钟，按照读取样品体积加入对应体积的1\%人自氯化钠溶液（每1.2 ml样品体积加入0.1 ml 1\%人自氯化钠溶液，按照“只进不含”原则，如26 ml样品，取吸2.6 ml）。} \\ \hline
取出裂解好的细胞悬液
% #VALUE_ID: LEX-P0063-V0028
% #FIELD_VALUE: 离心力
% #HANDWRITTEN: 300
\fieldvalue{\handwritten{300}}$\times g$离心5分钟 &
加入1\%人自氯化钠溶液至总体积：
% #VALUE_ID: LEX-P0063-V0029
% #FIELD_VALUE: 1\%人自氯化钠溶液总体积
% #HANDWRITTEN: 7.2
\fieldvalue{\handwritten{7.2}} ml \\ \hline
取吸
% #VALUE_ID: LEX-P0063-V0030
% #FIELD_VALUE: 取吸体积
% #HANDWRITTEN: 20
\fieldvalue{\handwritten{20}} $\mu$l &
\end{tabular}
\end{center}

\begin{center}
第 2 页 共 5 页
\end{center}
% LEXOID_PAGE_COMPLETED: 63/70

\newpage

\noindent
\textbf{PLATINUM LIFE EXCELLENCE BIOTECH}\hfill
文件编号：REC-YZ-SOP-QC-99-006-a

\medskip

\noindent
\begin{minipage}{0.58\linewidth}
加入 2.2ml 的细胞悬液（含混悬量），反复吹打至单细胞悬液。
  
摇去计数板上下两层保护膜，将细胞悬液混匀，从中吸取 20$\mu$l
与 20$\mu$l AO/PI 染料混合均匀。

吸取 20$\mu$l 细胞悬液与 AO/PI 染料的混合液，注入到计数板的一
个检测室内。

将加入混合液的一端插入进样口。

选择 Assay 类型，使用 AO/PI 染料进行双荧光活存检测，
在主界面左侧 ID 栏中输入样本批号。

点击主界面左下角 ``Preview'' 预览细胞图片。

调整各条亮斑直至观察到荧光通道下的细胞中心透亮，轮廓清
晰。

点击 ``count'' 按钮进行计数并自动拍摄图片。

重复以上步骤 2 次，共计数 3 次。打印钢缆结果。
\end{minipage}
\hfill
\begin{minipage}{0.35\linewidth}
吸取
% #VALUE_ID: LEX-P0064-V0001
% #FIELD_VALUE: 吸取体积
\fieldvalue{\handwritten{20}}$\mu$l 计数。

每个样品均计数
% #VALUE_ID: LEX-P0064-V0002
% #FIELD_VALUE: 样品计数次数
\fieldvalue{\handwritten{2}}次。
\end{minipage}

\medskip

\noindent
数据保存路径\hfill
% #TODO #HANDWRITTEN: 此处为一整行手写保存路径，部分字迹难以辨认
% #VALUE_ID: LEX-P0064-V0003
% #HANDWRITTEN: 体细胞/（疑似）KANAIA[?]细胞检测/2025
\fieldvalue{\handwritten{体细胞/（疑似）KANAIA[?]细胞检测/2025}}

\medskip

\noindent
\begin{tabular}{|p{0.97\linewidth}|}
\hline
\centering 检测结果\tabularnewline
\hline
检测结果：
一级反应
% #VALUE_ID: LEX-P0064-V0004
% #FIELD_VALUE: 一级反应
\fieldvalue{\checkboxfield{unchecked}}
\quad
二级反应
% #VALUE_ID: LEX-P0064-V0005
% #FIELD_VALUE: 二级反应
\fieldvalue{\checkboxfield{unchecked}}
\quad
三级反应
% #VALUE_ID: LEX-P0064-V0006
% #FIELD_VALUE: 三级反应
\fieldvalue{\checkboxfield{checked}}\\
\hline
\end{tabular}

\medskip

\noindent
\begin{tabular}{|p{0.10\linewidth}|p{0.16\linewidth}|p{0.16\linewidth}|p{0.14\linewidth}|p{0.14\linewidth}|p{0.14\linewidth}|p{0.14\linewidth}|}
\hline
检测次数
& 活细胞密度（个细胞/ml）
& 样品密度（个细胞/ml）
& 细胞密度 CV 值（\%）
& 细胞活率（\%）
& 样品活率（\%）
& 细胞活率 CV 值（\%）\\
\hline
第一次
&
% #VALUE_ID: LEX-P0064-V0007
% #FIELD_VALUE: 第一次活细胞密度
% #HANDWRITTEN: 9.40\times 10^{6}
\fieldvalue{\handwritten{$9.40\times 10^{6}$}}
&
&
&
% #VALUE_ID: LEX-P0064-V0008
% #FIELD_VALUE: 第一次细胞活率
% #HANDWRITTEN: 97.3
\fieldvalue{\handwritten{97.3}}
&
&\\
\hline
第二次
&
% #VALUE_ID: LEX-P0064-V0009
% #FIELD_VALUE: 第二次活细胞密度
% #HANDWRITTEN: 10.6\times 10^{6}
\fieldvalue{\handwritten{$10.6\times 10^{6}$}}
&
% #VALUE_ID: LEX-P0064-V0010
% #FIELD_VALUE: 第二次样品密度
% #HANDWRITTEN: 9.04\times 10^{5}
\fieldvalue{\handwritten{$9.04\times 10^{5}$}}
&
% #VALUE_ID: LEX-P0064-V0011
% #FIELD_VALUE: 细胞密度 CV 值
% #HANDWRITTEN: 14
\fieldvalue{\handwritten{14}}
&
% #VALUE_ID: LEX-P0064-V0012
% #FIELD_VALUE: 第二次细胞活率
% #HANDWRITTEN: 96.4
\fieldvalue{\handwritten{96.4}}
&
% #VALUE_ID: LEX-P0064-V0013
% #FIELD_VALUE: 样品活率
% #HANDWRITTEN: 97
\fieldvalue{\handwritten{97}}
&
% #VALUE_ID: LEX-P0064-V0014
% #FIELD_VALUE: 细胞活率 CV 值
% #HANDWRITTEN: 1
\fieldvalue{\handwritten{1}}\\
\hline
第三次
&
% #VALUE_ID: LEX-P0064-V0015
% #FIELD_VALUE: 第三次活细胞密度
% #HANDWRITTEN: 12.7\times 10^{6}
\fieldvalue{\handwritten{$12.7\times 10^{6}$}}
&
&
&
% #VALUE_ID: LEX-P0064-V0016
% #FIELD_VALUE: 第三次细胞活率
% #HANDWRITTEN: 96.7
\fieldvalue{\handwritten{96.7}}
&
&\\
\hline
\end{tabular}

\medskip

\noindent
\begin{tabular}{|p{0.30\linewidth}|p{0.27\linewidth}|p{0.40\linewidth}|}
\hline
发酵罐细胞体积（ml）
&
% #VALUE_ID: LEX-P0064-V0017
% #FIELD_VALUE: 发酵罐细胞体积
% #HANDWRITTEN: 48500
\fieldvalue{\handwritten{48500}}
&
细胞数量（个细胞/罐）
% #VALUE_ID: LEX-P0064-V0018
% #FIELD_VALUE: 细胞数量
% #HANDWRITTEN: 14.30\times 10^{10}
\fieldvalue{\handwritten{$14.30\times 10^{10}$}}\\
\hline
\multicolumn{3}{|p{0.97\linewidth}|}{
计算要求：样品密度为 3 次活细胞密度结果取平均值乘以 1\% 人白血化学溶液悬量总体，再除以生物反应器细胞液体积；细胞数量为样品密度乘以发酵罐体积。
}\\
\hline
备注
&
% #VALUE_ID: LEX-P0064-V0019
% #HANDWRITTEN: /
\fieldvalue{\handwritten{/}}
&\\
\hline
实验人/日期
&
% #TODO #HANDWRITTEN: 张娟（姓名字迹不清）
% #VALUE_ID: LEX-P0064-V0020
% #FIELD_VALUE: 实验人
% #HANDWRITTEN: 张娟
\fieldvalue{\handwritten{张娟}}
&
% #VALUE_ID: LEX-P0064-V0021
% #FIELD_VALUE: 实验日期
% #HANDWRITTEN: 2025.8.23
\fieldvalue{\handwritten{2025.8.23}}\\
\hline
复核人/日期
&
% #TODO #HANDWRITTEN: 复核人姓名难以辨认
% #VALUE_ID: LEX-P0064-V0022
% #FIELD_VALUE: 复核人
% #HANDWRITTEN: （姓名不清）
\fieldvalue{\handwritten{（姓名不清）}}
&
% #TODO #HANDWRITTEN: 复核日期疑似 2025.8.23，末位字迹不清
% #VALUE_ID: LEX-P0064-V0023
% #FIELD_VALUE: 复核日期
% #HANDWRITTEN: 2025.8.23
\fieldvalue{\handwritten{2025.8.23}}\\
\hline
\end{tabular}

\bigskip

\begin{center}
第 3 页 共 5 页
\end{center}
% LEXOID_PAGE_COMPLETED: 64/70

\newpage

Assay: 
% #VALUE_ID: LEX-P0065-V0001
% #FIELD_VALUE: Assay
\fieldvalue{MSC-1}

\vspace{0.5em}

Cell Type F1: 
% #VALUE_ID: LEX-P0065-V0002
% #FIELD_VALUE: Cell Type F1
\fieldvalue{MSC (AO)}

\noindent
Cell Type F2: 
% #VALUE_ID: LEX-P0065-V0003
% #FIELD_VALUE: Cell Type F2
\fieldvalue{MSC (PI)}

\vspace{0.5em}

Sample ID: 
% #VALUE_ID: LEX-P0065-V0004
% #FIELD_VALUE: Sample ID
\fieldvalue{QC-JS-QY0403-A31Z2202508006-20250823-C01-71.8H-1}

\noindent
Dilution: 
% #VALUE_ID: LEX-P0065-V0005
% #FIELD_VALUE: Dilution
\fieldvalue{2.00}

\vspace{0.5em}
\hrule
\vspace{0.3em}

Results:

\vspace{0.3em}

\begin{tabular}{@{}p{0.29\linewidth}p{0.34\linewidth}p{0.29\linewidth}@{}}
\textbf{Count} & \textbf{Concentration} & \textbf{Mean Diameter} \\
\texttt{==============} & \texttt{================} & \texttt{================} \\[0.2em]

Total: &
% #VALUE_ID: LEX-P0065-V0006
% #FIELD_VALUE: Total concentration
\fieldvalue{$9.66\times 10^6$ cells/mL} &
% #VALUE_ID: LEX-P0065-V0007
% #FIELD_VALUE: Total mean diameter
\fieldvalue{9.6 micron} \\

Live: &
% #VALUE_ID: LEX-P0065-V0008
% #FIELD_VALUE: Live concentration
\fieldvalue{$9.40\times 10^6$ cells/mL} &
% #VALUE_ID: LEX-P0065-V0009
% #FIELD_VALUE: Live mean diameter
\fieldvalue{9.5 microns} \\

Dead: &
% #VALUE_ID: LEX-P0065-V0010
% #FIELD_VALUE: Dead concentration
\fieldvalue{$2.56\times 10^5$ cells/mL} &
% #VALUE_ID: LEX-P0065-V0011
% #FIELD_VALUE: Dead mean diameter
\fieldvalue{12.2 microns} \\

Viability: &
&
\\[-0.1em]
% #VALUE_ID: LEX-P0065-V0012
% #FIELD_VALUE: Total count
\fieldvalue{2776} &
&
\\[-1.1em]
% #VALUE_ID: LEX-P0065-V0013
% #FIELD_VALUE: Live count
\fieldvalue{2703} &
&
\\[-1.1em]
% #VALUE_ID: LEX-P0065-V0014
% #FIELD_VALUE: Dead count
\fieldvalue{73} &
&
\\[-1.1em]
% #VALUE_ID: LEX-P0065-V0015
% #FIELD_VALUE: Viability
\fieldvalue{97.3\%} &
&
\end{tabular}

\vspace{1em}

\begin{center}
\begin{tabular}{@{}c@{\hspace{1.5em}}c@{}}
% TODO: Image panel visible on this page; source filename unavailable.
\fbox{\rule{0pt}{3.2cm}\rule{6.0cm}{0pt}} &
% TODO: Image panel visible on this page; source filename unavailable.
\fbox{\rule{0pt}{3.2cm}\rule{6.0cm}{0pt}} \\[1em]
% TODO: Image panel visible on this page; source filename unavailable.
\fbox{\rule{0pt}{3.2cm}\rule{6.0cm}{0pt}} &
% TODO: Image panel visible on this page; source filename unavailable.
\fbox{\rule{0pt}{3.2cm}\rule{6.0cm}{0pt}}
\end{tabular}
\end{center}

\vfill

操作人/审核人：
% #VALUE_ID: LEX-P0065-V0016
% #FIELD_VALUE: 操作人/审核人
% TODO #HANDWRITTEN: illegible Chinese signature; handwritten characters are unclear
\fieldvalue{\handwritten{（签名）}}

\hfill 日期：
% #VALUE_ID: LEX-P0065-V0017
% #FIELD_VALUE: 日期
% #HANDWRITTEN: 2025-08-23
\fieldvalue{\handwritten{2025-08-23}}
% LEXOID_PAGE_COMPLETED: 65/70

\newpage

\noindent\textbf{Assay:}
% #VALUE_ID: LEX-P0066-V0001
% #FIELD_VALUE: Assay
\fieldvalue{MSC-1}

\medskip
\noindent\textbf{Cell Type F1:}
% #VALUE_ID: LEX-P0066-V0002
% #FIELD_VALUE: Cell Type F1
\fieldvalue{MSC (AO)}

\noindent\textbf{Cell Type F2:}
% #VALUE_ID: LEX-P0066-V0003
% #FIELD_VALUE: Cell Type F2
\fieldvalue{MSC (PI)}

\medskip
\noindent\textbf{Sample ID:}
% #VALUE_ID: LEX-P0066-V0004
% #FIELD_VALUE: Sample ID
\fieldvalue{QC-JS-QY0403-A31Z220250806-20250823-C01-71.8H-2}

\noindent\textbf{Dilution:}
% #VALUE_ID: LEX-P0066-V0005
% #FIELD_VALUE: Dilution
\fieldvalue{2.00}

\medskip
\noindent\rule{\linewidth}{0.4pt}

\noindent\textbf{Results:}

\medskip
\noindent
\begin{tabular}{@{}p{0.29\linewidth}p{0.34\linewidth}p{0.29\linewidth}@{}}
\textbf{Count} & \textbf{Concentration} & \textbf{Mean Diameter} \\
\texttt{============} & \texttt{=============} & \texttt{===============} \\[2pt]
Total:
% #VALUE_ID: LEX-P0066-V0006
% #FIELD_VALUE: Total count
\fieldvalue{3161}
&
% #VALUE_ID: LEX-P0066-V0007
% #FIELD_VALUE: Total concentration
\fieldvalue{$1.10\times10^{7}\ \mathrm{cells/mL}$}
&
% #VALUE_ID: LEX-P0066-V0008
% #FIELD_VALUE: Total mean diameter
\fieldvalue{10.1 micron}
\\
Live:
% #VALUE_ID: LEX-P0066-V0009
% #FIELD_VALUE: Live count
\fieldvalue{3047}
&
% #VALUE_ID: LEX-P0066-V0010
% #FIELD_VALUE: Live concentration
\fieldvalue{$1.06\times10^{7}\ \mathrm{cells/mL}$}
&
% #VALUE_ID: LEX-P0066-V0011
% #FIELD_VALUE: Live mean diameter
\fieldvalue{10.0 microns}
\\
Dead:
% #VALUE_ID: LEX-P0066-V0012
% #FIELD_VALUE: Dead count
\fieldvalue{114}
&
% #VALUE_ID: LEX-P0066-V0013
% #FIELD_VALUE: Dead concentration
\fieldvalue{$3.99\times10^{5}\ \mathrm{cells/mL}$}
&
% #VALUE_ID: LEX-P0066-V0014
% #FIELD_VALUE: Dead mean diameter
\fieldvalue{11.9 micron}
\\
Viability:
% #VALUE_ID: LEX-P0066-V0015
% #FIELD_VALUE: Viability
\fieldvalue{96.4\%}
&
&
\end{tabular}

\medskip

\noindent
\begin{tabular}{@{}p{0.46\linewidth}p{0.46\linewidth}@{}}
% TODO: Image visible on this page; no source filename is available.
\fbox{\parbox[c][0.22\linewidth][c]{0.42\linewidth}{\centering Image placeholder}}
&
% TODO: Image visible on this page; no source filename is available.
\fbox{\parbox[c][0.22\linewidth][c]{0.42\linewidth}{\centering Image placeholder}}
\\[8pt]
% TODO: Image visible on this page; no source filename is available.
\fbox{\parbox[c][0.22\linewidth][c]{0.42\linewidth}{\centering Image placeholder}}
&
% TODO: Image visible on this page; no source filename is available.
\fbox{\parbox[c][0.22\linewidth][c]{0.42\linewidth}{\centering Image placeholder}}
\end{tabular}

\vfill

\noindent
操作人/审核人：
% #VALUE_ID: LEX-P0066-V0016
% #FIELD_VALUE: 操作人/审核人
% #TODO #HANDWRITTEN: 姚敏; handwriting is unclear
\fieldvalue{\handwritten{姚敏}}
\hfill
日期：
% #VALUE_ID: LEX-P0066-V0017
% #FIELD_VALUE: 日期
% #HANDWRITTEN: 2025-08-23
\fieldvalue{\handwritten{2025-08-23}}
% LEXOID_PAGE_COMPLETED: 66/70

\newpage

\noindent
% #VALUE_ID: LEX-P0067-V0001
% #FIELD_VALUE: Assay
Assay: \fieldvalue{MSC-1}

\vspace{0.5em}

\noindent
% #VALUE_ID: LEX-P0067-V0002
% #FIELD_VALUE: Cell Type F1
Cell Type F1: \fieldvalue{MSC (AO)}

\noindent
% #VALUE_ID: LEX-P0067-V0003
% #FIELD_VALUE: Cell Type F2
Cell Type F2: \fieldvalue{MSC (PI)}

\vspace{0.5em}

\noindent
% #VALUE_ID: LEX-P0067-V0004
% #FIELD_VALUE: Sample ID
Sample ID: \fieldvalue{\texttt{QC-JS-QY0403-A31Z2202508006-20250823-C01-71.8H-3}}

\noindent
% #VALUE_ID: LEX-P0067-V0005
% #FIELD_VALUE: Dilution
Dilution: \fieldvalue{2.00}

\vspace{0.5em}
\hrule
\vspace{0.5em}

\noindent\textbf{Results:}

\vspace{0.5em}

\small
\begin{tabular}{@{}p{0.29\linewidth}p{0.35\linewidth}p{0.30\linewidth}@{}}
\textbf{Count} & \textbf{Concentration} & \textbf{Mean Diameter} \\
\texttt{=============} & \texttt{=============} & \texttt{=============} \\[0.25em]

Total:
% #VALUE_ID: LEX-P0067-V0006
% #FIELD_VALUE: Count --- Total
\fieldvalue{3613}
&
% #VALUE_ID: LEX-P0067-V0010
% #FIELD_VALUE: Concentration --- Total
\fieldvalue{$1.27\times 10^{7}\ \mathrm{cells/mL}$}
&
% #VALUE_ID: LEX-P0067-V0013
% #FIELD_VALUE: Mean Diameter --- Total
\fieldvalue{12.4 micron}
\\

Live:
% #VALUE_ID: LEX-P0067-V0007
% #FIELD_VALUE: Count --- Live
\fieldvalue{3479}
&
% #VALUE_ID: LEX-P0067-V0011
% #FIELD_VALUE: Concentration --- Live
\fieldvalue{$1.22\times 10^{7}\ \mathrm{cells/mL}$}
&
% #VALUE_ID: LEX-P0067-V0014
% #FIELD_VALUE: Mean Diameter --- Live
\fieldvalue{12.5 microns}
\\

Dead:
% #VALUE_ID: LEX-P0067-V0008
% #FIELD_VALUE: Count --- Dead
\fieldvalue{134}
&
% #VALUE_ID: LEX-P0067-V0012
% #FIELD_VALUE: Concentration --- Dead
\fieldvalue{$4.69\times 10^{5}\ \mathrm{cells/mL}$}
&
% #VALUE_ID: LEX-P0067-V0015
% #FIELD_VALUE: Mean Diameter --- Dead
\fieldvalue{11.9 microns}
\\

Viability:
% #VALUE_ID: LEX-P0067-V0009
% #FIELD_VALUE: Viability
\fieldvalue{96.3\%}
&
&
\end{tabular}
\normalsize

\vspace{1.5em}

\begin{center}
\begin{tabular}{@{}c@{\hspace{1.5em}}c@{}}
% TODO: Image visible on this page; no source filename was provided.
\fbox{\parbox[c][5.1cm][c]{0.39\linewidth}{\centering Image placeholder}}
&
% TODO: Image visible on this page; no source filename was provided.
\fbox{\parbox[c][5.1cm][c]{0.39\linewidth}{\centering Image placeholder}}
\\[1em]
% TODO: Image visible on this page; no source filename was provided.
\fbox{\parbox[c][5.1cm][c]{0.39\linewidth}{\centering Image placeholder}}
&
% TODO: Image visible on this page; no source filename was provided.
\fbox{\parbox[c][5.1cm][c]{0.39\linewidth}{\centering Image placeholder}}
\end{tabular}
\end{center}

\vfill

\noindent
操作人/审核人：
% #VALUE_ID: LEX-P0067-V0016
% #FIELD_VALUE: 操作人/审核人
% #TODO #HANDWRITTEN: 姚虹?; handwritten Chinese name is partially illegible
\fieldvalue{\handwritten{姚虹?}}
\hfill
日期：
% #VALUE_ID: LEX-P0067-V0017
% #FIELD_VALUE: 日期
% #HANDWRITTEN: 25.08.23
\fieldvalue{\handwritten{25.08.23}}
% LEXOID_PAGE_COMPLETED: 67/70

\newpage

\begin{center}
\textbf{生物反应器细胞液检验操作记录}
\end{center}

\noindent
\begin{tabular}{@{}p{0.22\linewidth}p{0.28\linewidth}p{0.20\linewidth}p{0.25\linewidth}@{}}
\multicolumn{4}{l}{文件编码：%
% #VALUE_ID: LEX-P0068-V0001
% #FIELD_VALUE: 文件编码
\fieldvalue{REC-YZ-SOP-QC-99-006-a}}\\
\end{tabular}

\begin{center}
\begin{tabular}{|p{0.18\linewidth}|p{0.30\linewidth}|p{0.18\linewidth}|p{0.25\linewidth}|}
\hline
\multicolumn{4}{|c|}{葡萄糖浓度检验}\\
\hline
样品名称 &
% #VALUE_ID: LEX-P0068-V0002
% #FIELD_VALUE: 样品名称
\fieldvalue{经生物反应器细胞液} &
样品规格 &
% #VALUE_ID: LEX-P0068-V0003
% #FIELD_VALUE: 样品规格
\fieldvalue{ml/袋}\\
\hline
样品批号 & & 取样日期 & 年\quad 月\quad 日\\
\hline
\end{tabular}
\end{center}

\section*{1\quad 仪器设备、试剂耗材}

\subsection*{1.1\quad 仪器设备}

\begin{center}
\begin{tabular}{|p{0.15\linewidth}|p{0.27\linewidth}|p{0.16\linewidth}|p{0.19\linewidth}|p{0.19\linewidth}|}
\hline
名称 & 厂家 & 型号 & 设备编号 & 检查确认\\
\hline
% #VALUE_ID: LEX-P0068-V0004
% #FIELD_VALUE: 仪器名称
\fieldvalue{血糖仪} &
% #VALUE_ID: LEX-P0068-V0005
% #FIELD_VALUE: 厂家
\fieldvalue{三诺生物传感股份有限公司} &
% #VALUE_ID: LEX-P0068-V0006
% #FIELD_VALUE: 型号
\fieldvalue{安捷} &
&
% #VALUE_ID: LEX-P0068-V0007
% #FIELD_VALUE: 正常
\fieldvalue{\checkboxfield{unchecked}} 正常\quad
% #VALUE_ID: LEX-P0068-V0008
% #FIELD_VALUE: 异常
\fieldvalue{\checkboxfield{unchecked}} 异常\\
\hline
% #VALUE_ID: LEX-P0068-V0009
% #FIELD_VALUE: 仪器名称
\fieldvalue{移液器} &
% #VALUE_ID: LEX-P0068-V0010
% #FIELD_VALUE: 厂家
\fieldvalue{Eppendorf Corporate} &
% #VALUE_ID: LEX-P0068-V0011
% #FIELD_VALUE: 型号
\fieldvalue{0.1--2.5\,\textmu l} &
&
% #VALUE_ID: LEX-P0068-V0012
% #FIELD_VALUE: 正常
\fieldvalue{\checkboxfield{unchecked}} 正常\quad
% #VALUE_ID: LEX-P0068-V0013
% #FIELD_VALUE: 异常
\fieldvalue{\checkboxfield{unchecked}} 异常\\
\hline
\end{tabular}
\end{center}

\subsection*{1.2\quad 试剂耗材}

\begin{center}
\begin{tabular}{|p{0.20\linewidth}|p{0.40\linewidth}|p{0.20\linewidth}|p{0.20\linewidth}|}
\hline
名称 & 厂家 & 批号 & 有效期至\\
\hline
% #VALUE_ID: LEX-P0068-V0014
% #FIELD_VALUE: 试剂名称
\fieldvalue{血糖测试条} &
% #VALUE_ID: LEX-P0068-V0015
% #FIELD_VALUE: 厂家
\fieldvalue{三诺生物传感股份有限公司} &
& \\
\hline
\end{tabular}
\end{center}

\section*{2\quad 操作步骤}

\begin{center}
\small
\begin{tabular}{|p{0.57\linewidth}|p{0.38\linewidth}|}
\hline
\multicolumn{2}{|c|}{操作过程}\\
\hline
取出测试条，取样品自然沉降后的上清液 $0.6\,\mu\mathrm{l}$ 加入测试条电极反应区，血糖仪自动发出“嘀”提示音，吸入样品完成，血糖仪自动进入倒计时，开始测试； &
取上清液：\underline{\hspace{2cm}}$\mu\mathrm{l}$\\
\hline
重复以上步骤 2 次，共计数 3 次。 &
每个样品分别共计数\underline{\hspace{1.5cm}}次。\\
\hline
\multicolumn{2}{|c|}{检测结果}\\
\hline
\multicolumn{2}{|c|}{葡萄糖浓度（mmol/L）}\\
\hline
\multicolumn{2}{|c|}{\begin{tabular}{cccc}
第一次 & 第二次 & 第三次 & 平均值（保留一位小数）\\[-1.5ex]
\hline
\multicolumn{1}{|p{0.23\linewidth}|}{} &
\multicolumn{1}{p{0.23\linewidth}|}{} &
\multicolumn{1}{p{0.23\linewidth}|}{} &
\multicolumn{1}{p{0.26\linewidth}|}{}
\end{tabular}}\\
\hline
备注 &
\multicolumn{1}{p{0.70\linewidth}|}{}\\
\hline
实验人/日期 &
\multicolumn{1}{p{0.70\linewidth}|}{}\\
\hline
复核人/日期 &
\multicolumn{1}{p{0.70\linewidth}|}{}\\
\hline
\end{tabular}
\end{center}

% TODO: The form fields on this page are visibly blank and contain no additional filled-in values.
% LEXOID_PAGE_COMPLETED: 68/70

\newpage

\begin{center}
\textbf{生物反应器细胞液检验操作记录}
\end{center}

\begin{center}
文件编码：REC-YZ-SOP-QC-99-006-a
\end{center}

\begin{tabular}{|p{0.14\linewidth}|p{0.36\linewidth}|p{0.14\linewidth}|p{0.28\linewidth}|}
\hline
样品名称 &
% #VALUE_ID: LEX-P0069-V0001
% #FIELD_VALUE: 样品名称
\fieldvalue{一级生物反应器细胞瘤液} &
样品规格 & ml/袋 \\ \hline
样品批号 & & 取样日期 & 年\quad 月\quad 日 \\ \hline
\end{tabular}

\vspace{0.4em}

\textbf{1.\quad 仪器设备及试剂：}

\textbf{1.1.\quad 仪器设备：}

\begin{tabular}{|p{0.22\linewidth}|p{0.18\linewidth}|p{0.18\linewidth}|p{0.20\linewidth}|p{0.16\linewidth}|}
\hline
仪器名称 & 厂家 & 型号 & 设备编号 & 检查确认 \\ \hline
% #VALUE_ID: LEX-P0069-V0002
% #FIELD_VALUE: 仪器名称
\fieldvalue{荧光显微镜} &
% #VALUE_ID: LEX-P0069-V0003
% #FIELD_VALUE: 厂家
\fieldvalue{徕卡} &
% #VALUE_ID: LEX-P0069-V0004
% #FIELD_VALUE: 型号
\fieldvalue{Dmi LED} &
% #VALUE_ID: LEX-P0069-V0005
% #FIELD_VALUE: 设备编号
\fieldvalue{1431404} &
正常\ 
% #VALUE_ID: LEX-P0069-V0006
% #FIELD_VALUE: 正常
\fieldvalue{\checkboxfield{unchecked}}
\quad 异常\ 
% #VALUE_ID: LEX-P0069-V0007
% #FIELD_VALUE: 异常
\fieldvalue{\checkboxfield{unchecked}} \\ \hline
\end{tabular}

\textbf{1.2.\quad 试剂耗材}

\begin{tabular}{|p{0.20\linewidth}|p{0.30\linewidth}|p{0.20\linewidth}|p{0.18\linewidth}|p{0.12\linewidth}|}
\hline
名称 & 厂家 & 货号 & 批号 & 有效期至 \\ \hline
% #VALUE_ID: LEX-P0069-V0008
% #FIELD_VALUE: 名称
\fieldvalue{细胞活力检测试剂} &
% #VALUE_ID: LEX-P0069-V0009
% #FIELD_VALUE: 厂家
\fieldvalue{江苏凯基生物技术股份有限公司} &
% #VALUE_ID: LEX-P0069-V0010
% #FIELD_VALUE: 货号
\fieldvalue{KGA9501-1000} &
& \\ \hline
\end{tabular}

\textbf{2.\quad 操作步骤：}

\begin{tabular}{|p{0.92\linewidth}|}
\hline
\multicolumn{1}{|c|}{溶液配制} \\ \hline
20$\times$的荧光染液配制：取\_\_\_\_ ml 组分A 和\_\_\_\_ ml 组分B 与\_\_\_\_\_\_ ml PBS 混合均匀配成20$\times$的荧光染液。 \\ \hline
\multicolumn{1}{|c|}{操作过程} \\ \hline
取细胞后样品 200$\mu$l/孔至96孔板中，取2孔；
样品\_\_\_\_\_\_ $\mu$l/孔 \\
每个样品\_\_\_\_\_\_孔 \\ \hline
自然沉降后，吸弃上清，注意不要吸到沉淀样品；
% #VALUE_ID: LEX-P0069-V0011
% #FIELD_VALUE: 吸弃上清
\fieldvalue{\checkboxfield{unchecked}}\ 吸弃上清 \\ \hline
取20$\times$的荧光染液加入到96孔板中，200$\mu$l/孔，避光静置30--60mins；
20$\times$的荧光染液\_\_\_\_\_\_ $\mu$l/孔 \\
避光静置\_\_\_\_\_\_ min \\ \hline
染色结束后，吸弃上清，取PBS 加入到以上96孔板中，200$\mu$l/孔，吸弃上清，注意不要吸到沉淀样品，清洗3次，之后每孔加入100$\mu$l PBS 进行拍照。
加入PBS\_\_\_\_\_\_ $\mu$l/孔 \\
清洗\_\_\_\_\_\_次 \\ \hline
将荧光显微镜调至4$\times$物镜，软件中调节节至相应物镜参数，调节到适合的光源和像距，拍摄4张清晰图像。
\\ \hline
统一张清晰图像，打印并粘贴在实验记录本上。
\\ \hline
\multicolumn{1}{|c|}{数据保存路径} \\ \hline
\multicolumn{1}{|c|}{检测结果} \\ \hline
\end{tabular}

\vfill

\begin{center}
第4页 共5页
\end{center}
% LEXOID_PAGE_COMPLETED: 69/70

\newpage

\noindent
\begin{tabular}{@{}p{0.24\linewidth}p{0.18\linewidth}p{0.32\linewidth}p{0.22\linewidth}@{}}
\textit{PLANTIUM LIFE EXCELLENCE BIOTECHEN} &
\begin{center}
\scriptsize 蜀生生物\\[-1pt]
% #VALUE_ID: LEX-P0070-V0001
% #FIELD_VALUE: 文件状态
\fieldvalue{原件}
\end{center}
&
\centering
文件编号：
% #VALUE_ID: LEX-P0070-V0002
% #FIELD_VALUE: 文件编号
\fieldvalue{\texttt{REC-YZ-SOP-QC-99-006-a}}
&
\begin{center}
\scriptsize 蜀生生物\\[-1pt]
版本号：
% #VALUE_ID: LEX-P0070-V0003
% #FIELD_VALUE: 版本号
\fieldvalue{1.4}\\[-1pt]
生效日期：
% #VALUE_ID: LEX-P0070-V0004
% #FIELD_VALUE: 生效日期
\fieldvalue{2025.1.14}
\end{center}
\end{tabular}

\vspace{0.15cm}

\noindent
\begin{tabular}{|p{\linewidth-2\tabcolsep-2\arrayrulewidth}|}
\hline
\vspace{21.8cm}\\
\hline
\begin{tabular}{|p{0.16\linewidth}|p{0.78\linewidth}|}
\hline
备注 & \\[0.38cm]
\hline
实验人/日期 & \\[0.38cm]
\hline
复核人/日期 & \\[0.38cm]
\hline
\end{tabular}
\\
\hline
\end{tabular}

\end{document}
% LEXOID_PAGE_COMPLETED: 70/70
